% Generated mechanically from the frozen duality.tex target.
% Do not edit this generated file; edit the bound locale source instead.

% OK, start here.
%

\setcounter{chapter}{47}
\chapter{概形的对偶理论}



\phantomsection
\label{duality-section-phantom}


\section{引言}
\label{duality-section-introduction}

\noindent
本章研究概形态射的相对对偶理论以及概形上的对偶复形。
参考文献为 \cite{RD}。

\medskip\noindent
Noetherian 环的对偶复形已在《对偶化复形》第
\ref{dualizing-section-dualizing} 节及其后各节中定义并研究。
本章继续研究概形上的对偶复形；见第
\ref{duality-section-dualizing-schemes} 节。

\medskip\noindent
本章的主要篇幅用于研究如下情形中的推前函子之右伴随：
模层的导出范畴，且其上同调层为拟凝聚层。见第
\ref{duality-section-twisted-inverse-image}、
\ref{duality-section-restriction-to-opens}、
\ref{duality-section-base-change-map}、
\ref{duality-section-base-change-II}、
\ref{duality-section-trace}、
\ref{duality-section-compare-with-pullback}、
\ref{duality-section-sections-with-exact-support}、
\ref{duality-section-duality-finite}、
\ref{duality-section-perfect-proper}、
\ref{duality-section-dualizing-Cartier} 和
\ref{duality-section-examples} 节。
这里遵循论文
\cite{Neeman-Grothendieck}、\cite{LN}、
\cite{Lipman-notes} 和 \cite{Neeman-improvement}。

\medskip\noindent
第 \ref{duality-section-upper-shriek}、
\ref{duality-section-upper-shriek-properties} 和
\ref{duality-section-base-change-shriek} 节讨论 Noetherian 概形之间
有限型分离态射的重要而实用的上叹号函子 $f^!$，并最终在总览性的
第 \ref{duality-section-duality} 节汇总。

\medskip\noindent
第 \ref{duality-section-glue} 节说明另一种对偶理论与对偶复形理论：
在一个固定的、带有对偶复形的局部 Noetherian 基上工作
（该节对应 Hartshorne 书中的一条注记）。

\medskip\noindent
其余各节给出若干应用。

\medskip\noindent
下一章把本章内容延伸到代数空间上的对偶理论；见
《空间的对偶理论》第 \ref{spaces-duality-section-introduction} 节。







\section{概形上的对偶复形}
\label{duality-section-dualizing-schemes}

\noindent
局部 Noetherian 概形上的对偶复形被定义为这样一个复形：
仿射局部地，它来自相应环上的对偶复形。这个定义并非完全标准，
但在有限维 Noetherian 概形上与文献中的所有定义一致。

\begin{lemma}
\label{duality-lemma-equivalent-definitions}
设 $X$ 为局部 Noetherian 概形，$K$ 为 $D(\mathcal{O}_X)$ 的对象。
下列条件等价：
\begin{enumerate}
\item 对每个仿射开集 $U = \Spec(A) \subset X$，存在 $A$ 的对偶复形
$\omega_A^\bullet$，使得 $K|_U$ 同构于 $\omega_A^\bullet$ 在函子
$\widetilde{} : D(A) \to D(\mathcal{O}_U)$ 下的像。
\item 存在仿射开覆盖 $X = \bigcup U_i$，$U_i = \Spec(A_i)$，
使得对每个 $i$，均存在 $A_i$ 的对偶复形 $\omega_i^\bullet$，
且 $K|_{U_i}$ 同构于 $\omega_i^\bullet$ 在函子
$\widetilde{} : D(A_i) \to D(\mathcal{O}_{U_i})$ 下的像。
\end{enumerate}
\end{lemma}

\begin{proof}
设 (2) 成立，并令 $U = \Spec(A)$ 为 $X$ 的仿射开集。
条件 (2) 蕴含 $K$ 属于 $D_\QCoh(\mathcal{O}_X)$，因此在 $D(A)$
中存在对象 $\omega_A^\bullet$，其相伴拟凝聚层复形同构于 $K|_U$；
见《概形的导出范畴》，引理
\ref{perfect-lemma-affine-compare-bounded}。只需证明
$\omega_A^\bullet$ 是 $A$ 的对偶复形。

\medskip\noindent
由于 $X = \bigcup U_i$ 是开覆盖，可找到标准开覆盖
$U = D(f_1) \cup \ldots \cup D(f_m)$，使每个 $D(f_j)$
都是某个仿射开集 $U_i$ 中的标准开集；见《概形》，引理
\ref{schemes-lemma-standard-open-two-affines}。设
$D(f_j) = D(g_j)$，其中 $g_j \in A_{i_j}$。则
$A_{f_j} \cong (A_{i_j})_{g_j}$，并且由《概形的导出范畴》，引理
\ref{perfect-lemma-affine-compare-bounded}，在导出范畴中有
$$
(\omega_A^\bullet)_{f_j} \cong (\omega_i^\bullet)_{g_j}
$$
由《对偶化复形》，引理 \ref{dualizing-lemma-dualizing-localize}，
对 $j = 1, \ldots, m$，复形 $(\omega_A^\bullet)_{f_j}$
是 $A_{f_j}$ 上的对偶复形。于是由《对偶化复形》，引理
\ref{dualizing-lemma-dualizing-glue}，$\omega_A^\bullet$ 是对偶复形。
\end{proof}

\begin{definition}
\label{duality-definition-dualizing-scheme}
设 $X$ 为局部 Noetherian 概形。若 $D(\mathcal{O}_X)$ 的对象 $K$
满足引理 \ref{duality-lemma-equivalent-definitions} 中的等价条件，
则称 $K$ 为一个{\it 对偶复形}。
\end{definition}

\noindent
请参见本节开头的说明。

\begin{lemma}
\label{duality-lemma-affine-duality}
设 $A$ 为 Noetherian 环，$X = \Spec(A)$。设 $K, L$ 为 $D(A)$
的对象。若 $K \in D_{\textit{Coh}}(A)$ 且 $L$ 具有有限内射维数，
则在 $D(\mathcal{O}_X)$ 中有
$$
R\SheafHom_{\mathcal{O}_X}(\widetilde{K}, \widetilde{L})
=
\widetilde{R\Hom_A(K, L)}
$$
\end{lemma}

\begin{proof}
可设 $L$ 由内射 $A$-模组成的有限复形 $I^\bullet$ 给出。
对 $I^\bullet$ 的长度作归纳，并利用各构造与可区别三角的相容性，
归约到 $L = I[0]$ 的情形，其中 $I$ 为内射 $A$-模。在此情形下，
《概形的导出范畴》，引理
\ref{perfect-lemma-quasi-coherence-internal-hom} 表明，
$R\SheafHom_{\mathcal{O}_X}(\widetilde{K}, \widetilde{L})$
的第 $n$ 个上同调层是与预层
$$
D(f) \longmapsto \Ext^n_{A_f}(K \otimes_A A_f, I \otimes_A A_f)
$$
相伴的层。由于 $A$ 为 Noetherian 环，$A_f$-模
$I \otimes_A A_f$ 是内射的（《对偶化复形》，引理
\ref{dualizing-lemma-localization-injective-modules}）。因此
\begin{align*}
\Ext^n_{A_f}(K \otimes_A A_f, I \otimes_A A_f)
& =
\Hom_{A_f}(H^{-n}(K \otimes_A A_f), I \otimes_A A_f) \\
& =
\Hom_{A_f}(H^{-n}(K) \otimes_A A_f, I \otimes_A A_f) \\
& =
\Hom_A(H^{-n}(K), I) \otimes_A A_f
\end{align*}
最后一个等式成立，是因为 $H^{-n}(K)$ 是有限 $A$-模；见代数，引理
\ref{algebra-lemma-hom-from-finitely-presented}。这证明了典范映射
$$
\widetilde{R\Hom_A(K, L)}
\longrightarrow
R\SheafHom_{\mathcal{O}_X}(\widetilde{K}, \widetilde{L})
$$
在此情形下为拟同构，证明完成。
\end{proof}

\begin{lemma}
\label{duality-lemma-internal-hom-evaluate-isom}
设 $X$ 为 Noetherian 概形，$K, L, M \in D_\QCoh(\mathcal{O}_X)$。
则《上同调》，引理 \ref{cohomology-lemma-internal-hom-evaluate} 中的映射
$$
R\SheafHom(L, M) \otimes_{\mathcal{O}_X}^\mathbf{L} K
\longrightarrow
R\SheafHom(R\SheafHom(K, L), M)
$$
在下列两种情形中是同构：
\begin{enumerate}
\item $K \in D^-_{\textit{Coh}}(\mathcal{O}_X)$，
$L \in D^+_{\textit{Coh}}(\mathcal{O}_X)$，并且 $M$ 仿射局部地具有
有限内射维数（见证明）；或者
\item $K$ 和 $L$ 属于 $D_{\textit{Coh}}(\mathcal{O}_X)$，
对象 $R\SheafHom(L, M)$ 具有有限 Tor 维数，并且
$L$ 和 $M$ 仿射局部地具有有限内射维数
（特别地，$L$ 和 $M$ 是有界的）。
\end{enumerate}
\end{lemma}

\begin{proof}
先证 (1)。所谓 $M$ 仿射局部地具有有限内射维数，是指 $X$ 有一个
由仿射开集 $U = \Spec(A)$ 构成的开覆盖，使得与 $M|_U$ 对应的
$D(A)$ 中的对象（见《概形的导出范畴》，引理
\ref{perfect-lemma-affine-compare-bounded}）具有有限内射维数
\footnote{此条件不依赖于 Noetherian 概形 $X$ 的仿射开覆盖的选取。
细节从略。}。为证明本引理，可以把 $X$ 替换为 $U$，即可以假设
对某个 Noetherian 环 $A$ 有 $X = \Spec(A)$。注意，由《概形的导出范畴》，
引理 \ref{perfect-lemma-coherent-internal-hom}，
$R\SheafHom(K, L)$ 属于 $D^+_{\textit{Coh}}(\mathcal{O}_X)$。
此外，由引理 \ref{duality-lemma-affine-duality} 及《概形的导出范畴》，引理
\ref{perfect-lemma-quasi-coherence-internal-hom}，箭头左右两端的构造均与函子
$D(A) \to D_\QCoh(\mathcal{O}_X)$ 交换
（严格说来，这里用到了对 $K$、$L$、$M$ 的假设，以及刚刚对
$R\SheafHom(K, L)$ 证明的结论）。最后，由《代数进阶》，引理
\ref{more-algebra-lemma-internal-hom-evaluate-isomorphism} 的第 (2) 部分，
该箭头是同构。

\medskip\noindent
再证 (2)。论证同上。一个小差别是，此时
$R\SheafHom(K, L)$ 属于 $D_{\textit{Coh}}(\mathcal{O}_X)$，因为可以
仿射局部地考察（由引理 \ref{duality-lemma-affine-duality}，这是允许的），
并应用《对偶化复形》，引理
\ref{dualizing-lemma-finite-ext-into-bounded-injective}。
最后由《代数进阶》，引理
\ref{more-algebra-lemma-internal-hom-evaluate-isomorphism-technical} 得出结论。
\end{proof}

\begin{lemma}
\label{duality-lemma-dualizing-schemes}
设 $K$ 是局部 Noetherian 概形 $X$ 上的对偶复形。则 $K$ 是
$D_{\textit{Coh}}(\mathcal{O}_X)$ 的对象，并且
$D = R\SheafHom_{\mathcal{O}_X}(-, K)$ 诱导反等价
$$
D :
D_{\textit{Coh}}(\mathcal{O}_X)
\longrightarrow
D_{\textit{Coh}}(\mathcal{O}_X)
$$
它带有一个典范同构 $\text{id} \to D \circ D$。若 $X$ 拟紧，
则 $D$ 交换 $D^+_{\textit{Coh}}(\mathcal{O}_X)$ 与
$D^-_{\textit{Coh}}(\mathcal{O}_X)$，并诱导反等价
$D^b_{\textit{Coh}}(\mathcal{O}_X) \to D^b_{\textit{Coh}}(\mathcal{O}_X)$。
\end{lemma}

\begin{proof}
设 $U \subset X$ 为仿射开集，记 $U = \Spec(A)$，并设
$\omega_A^\bullet$ 为与 $K|_U$ 对应的 $A$ 的对偶复形，
如引理 \ref{duality-lemma-equivalent-definitions} 中所述。由引理
\ref{duality-lemma-affine-duality}，图表
$$
\xymatrix{
D_{\textit{Coh}}(A) \ar[r] \ar[d]_{R\Hom_A(-, \omega_A^\bullet)} &
D_{\textit{Coh}}(\mathcal{O}_U) \ar[d]^{R\SheafHom_{\mathcal{O}_X}(-, K|_U)} \\
D_{\textit{Coh}}(A) \ar[r] &
D(\mathcal{O}_U)
}
$$
交换。因此 $D$ 将 $D_{\textit{Coh}}(\mathcal{O}_X)$ 映入
$D_{\textit{Coh}}(\mathcal{O}_X)$。此外，对所有 $L$，典范映射
$$
L
\longrightarrow
R\SheafHom_{\mathcal{O}_X}(K, K) \otimes_{\mathcal{O}_X}^\mathbf{L} L
\longrightarrow
R\SheafHom_{\mathcal{O}_X}(R\SheafHom_{\mathcal{O}_X}(L, K), K)
$$
都是同构，其中第二个箭头使用了《上同调》，引理
\ref{cohomology-lemma-internal-hom-evaluate}；这是因为由《对偶化复形》，
引理 \ref{dualizing-lemma-dualizing}，该结论在仿射概形上成立
\footnote{另一种方法是先仿射局部地证明
$R\SheafHom_{\mathcal{O}_X}(K, K) = \mathcal{O}_X$，再用引理
\ref{duality-lemma-internal-hom-evaluate-isom} 的第 (2) 部分说明该映射为同构。}，
而且我们已经看到，在仿射概形上这恢复了代数中的情形。
拟紧情形下函子 $D$ 的有界性陈述，同样由《对偶化复形》，引理
\ref{dualizing-lemma-dualizing} 的相应陈述得到。
\end{proof}

\noindent
设 $X$ 为局部赋环空间。回忆：若 $D(\mathcal{O}_X)$ 的对象 $L$
关于由导出张量积给出的 $D(\mathcal{O}_X)$ 上的对称幺半结构是可逆对象，
则称该对象是{\it 可逆的}。我们在《上同调》，引理
\ref{cohomology-lemma-invertible-derived} 中已经看到，这意味着 $L$ 是完美的，
且存在开覆盖 $X = \bigcup U_i$，使得对某些整数 $n_i$ 有
$L|_{U_i} \cong \mathcal{O}_{U_i}[-n_i]$。在此情形下，函数
$$
x \mapsto n_x,\quad
\text{其中 }n_x\text{ is the unique integer such that }
H^{n_x}(L_x) \not = 0
$$
在 $X$ 上局部常值。特别地，有 $L = \bigoplus H^n(L)[-n]$；
这给出了一个表示 $L$ 的良定义的 $\mathcal{O}_X$-模复形
（其微分全为零）。

\begin{lemma}
\label{duality-lemma-dualizing-unique-schemes}
设 $X$ 为局部 Noetherian 概形。若 $K$ 与 $K'$ 是 $X$ 上的对偶复形，
则对 $D(\mathcal{O}_X)$ 的某个可逆对象 $L$，$K'$ 同构于
$K \otimes_{\mathcal{O}_X}^\mathbf{L} L$。
\end{lemma}

\begin{proof}
令
$$
L = R\SheafHom_{\mathcal{O}_X}(K, K')
$$
这是 $D(\mathcal{O}_X)$ 的可逆对象，因为此结论仿射局部地成立；
见《对偶化复形》，引理 \ref{dualizing-lemma-dualizing-unique} 及其证明。
同理，求值映射 $L \otimes_{\mathcal{O}_X}^\mathbf{L} K \to K'$
是同构。
\end{proof}

\begin{lemma}
\label{duality-lemma-dimension-function-scheme}
设 $X$ 为局部 Noetherian 概形，$\omega_X^\bullet$ 为 $X$ 上的对偶复形。
则 $X$ 是普遍等长链的，并且由
$X \to \mathbf{Z}$ 定义的函数
$$
x \longmapsto \delta(x)\text{ 使得 }
\omega_{X, x}^\bullet[-\delta(x)]
\text{ is a normalized dualizing complex over }
\mathcal{O}_{X, x}
$$
是维数函数。
\end{lemma}

\begin{proof}
由仿射情形的《对偶化复形》，引理
\ref{dualizing-lemma-dimension-function} 及各定义立即得到。
\end{proof}

\begin{lemma}
\label{duality-lemma-sitting-in-degrees}
设 $X$ 为局部 Noetherian 概形，$\omega_X^\bullet$ 为 $X$ 上的对偶复形，
其相伴维数函数为 $\delta$。设 $\mathcal{F}$ 为凝聚 $\mathcal{O}_X$-模。令
$\mathcal{E}^i = \SheafExt^{-i}_{\mathcal{O}_X}(\mathcal{F}, \omega_X^\bullet)$。
则 $\mathcal{E}^i$ 是凝聚 $\mathcal{O}_X$-模，并且对 $x \in X$ 有：
\begin{enumerate}
\item $\mathcal{E}^i_x$ 仅当
$\delta(x) \leq i \leq \delta(x) + \dim(\text{Supp}(\mathcal{F}_x))$ 时可能非零；
\item $\dim(\text{Supp}(\mathcal{E}^{i + \delta(x)}_x)) \leq i$；
\item $\text{depth}(\mathcal{F}_x)$ 是满足
$\mathcal{E}_x^{i + \delta(x)} \not = 0$ 的最小整数 $i \geq 0$；
\item 有
$x \in \text{Supp}(\bigoplus_{j \leq i} \mathcal{E}^j)
\Leftrightarrow
\text{depth}_{\mathcal{O}_{X, x}}(\mathcal{F}_x) + \delta(x) \leq i$。
\end{enumerate}
\end{lemma}

\begin{proof}
由引理 \ref{duality-lemma-dualizing-schemes}，$\mathcal{E}^i$ 是凝聚的。
选取 $x$ 的一个仿射邻域，并应用《概形的导出范畴》，引理
\ref{perfect-lemma-quasi-coherence-internal-hom} 以及《代数进阶》，引理
\ref{more-algebra-lemma-base-change-RHom} 的第 (3) 部分，可得
$$
\mathcal{E}^i_x =
\SheafExt^{-i}_{\mathcal{O}_X}(\mathcal{F}, \omega_X^\bullet)_x =
\Ext^{-i}_{\mathcal{O}_{X, x}}(\mathcal{F}_x,
\omega_{X, x}^\bullet) =
\Ext^{\delta(x) - i}_{\mathcal{O}_{X, x}}(\mathcal{F}_x,
\omega_{X, x}^\bullet[-\delta(x)])
$$
根据引理 \ref{duality-lemma-dimension-function-scheme} 中 $\delta$ 的构造，
第 (1)、(2)、(3) 项便归结为《对偶化复形》，引理
\ref{dualizing-lemma-sitting-in-degrees}。第 (4) 项是第 (3) 项和第 (1) 项的形式推论。
\end{proof}




\section{推前的右伴随}
\label{duality-section-twisted-inverse-image}

\noindent
本节及下一节的参考文献为
\cite{Neeman-Grothendieck}、\cite{LN}、
\cite{Lipman-notes} 和 \cite{Neeman-improvement}。

\medskip\noindent
设 $f : X \to Y$ 为概形态射。本节考察函子
$Rf_* : D_\QCoh(\mathcal{O}_X) \to D_\QCoh(\mathcal{O}_Y)$ 的右伴随。
文献中，当这个函子存在时，有时将其记作 $f^{\times}$。
这种记号并未得到普遍接受，所以这里不采用。我们也不把这样的函子记作
$f^!$，因为对一般态射 $f$，这会与 \cite{RD} 中的记号冲突。

\begin{lemma}
\label{duality-lemma-twisted-inverse-image}
\begin{reference}
这与 \cite[例 4.2]{Neeman-Grothendieck} 几乎相同。
\end{reference}
设 $f : X \to Y$ 为拟分离拟紧概形之间的态射。函子
$Rf_* : D_\QCoh(\mathcal{O}_X) \to D_\QCoh(\mathcal{O}_Y)$ 有右伴随。
\end{lemma}

\begin{proof}
验证《导出范畴》，命题 \ref{derived-proposition-brown} 的假设，即可证明右伴随存在。
首先，范畴 $D_\QCoh(\mathcal{O}_X)$ 有直和；见《概形的导出范畴》，引理
\ref{perfect-lemma-quasi-coherence-direct-sums}。
由《概形的导出范畴》，定理
\ref{perfect-theorem-bondal-van-den-Bergh}，范畴
$D_\QCoh(\mathcal{O}_X)$ 是紧生成的。由于 $X$ 和 $Y$ 均拟紧且拟分离，
$f$ 也如此；见《概形》，引理 \ref{schemes-lemma-compose-after-separated} 和
\ref{schemes-lemma-quasi-compact-permanence}。因此函子 $Rf_*$ 与直和交换；
见《概形的导出范畴》，引理
\ref{perfect-lemma-quasi-coherence-pushforward-direct-sums}。证毕。
\end{proof}

\begin{example}
\label{duality-example-affine-twisted-inverse-image}
设 $A \to B$ 为环同态，$Y = \Spec(A)$、$X = \Spec(B)$，
并设 $f : X \to Y$ 为与 $A \to B$ 对应的态射。通过等价
$D(B) \to D_\QCoh(\mathcal{O}_X)$ 和
$D(A) \to D_\QCoh(\mathcal{O}_Y)$，函子
$Rf_* : D_\QCoh(\mathcal{O}_X) \to D_\QCoh(\mathcal{O}_Y)$
对应于限制函子 $D(B) \to D(A)$。因此，它的右伴随对应于
《对偶化复形》，第 \ref{dualizing-section-trivial} 节中的函子
$K \longmapsto R\Hom(B, K)$。
\end{example}

\begin{example}
\label{duality-example-does-not-preserve-coherent}
若 $f : X \to Y$ 是 Noetherian 概形之间的分离有限型态射，则
$Rf_* : D_\QCoh(\mathcal{O}_X) \to D_\QCoh(\mathcal{O}_Y)$ 的右伴随
并不把 $D_{\textit{Coh}}(\mathcal{O}_Y)$ 映入
$D_{\textit{Coh}}(\mathcal{O}_X)$。事实上，设 $k$ 为域，并考虑态射
$f : \mathbf{A}^1_k \to \Spec(k)$。由例
\ref{duality-example-affine-twisted-inverse-image}，这归结为下述问题：当
$A = k$ 且 $B = k[x]$ 时，$R\Hom(B, -)$ 是否把
$D_{\textit{Coh}}(A)$ 映入 $D_{\textit{Coh}}(B)$。答案是否定的，因为
$$
R\Hom(k[x], k) = \left(\prod\nolimits_{n \geq 0} k\right)[0]
$$
并不是有限 $k[x]$-模。因此 $a(\mathcal{O}_Y)$ 没有凝聚上同调层。
\end{example}

\begin{example}
\label{duality-example-does-not-preserve-bounded-above}
若 $f : X \to Y$ 是 Noetherian 概形之间的固有态射，甚至是有限态射，则
$Rf_* : D_\QCoh(\mathcal{O}_X) \to D_\QCoh(\mathcal{O}_Y)$ 的右伴随
并不把 $D_\QCoh^-(\mathcal{O}_Y)$ 映入
$D_\QCoh^-(\mathcal{O}_X)$。事实上，设 $k$ 为域，$k[\epsilon]$
为 $k$ 上的对偶数，令 $X = \Spec(k)$、$Y = \Spec(k[\epsilon])$。
则对所有 $i \geq 0$，$\Ext^i_{k[\epsilon]}(k, k)$ 均非零。
因此，由例 \ref{duality-example-affine-twisted-inverse-image}，
$a(\mathcal{O}_Y)$ 不是上有界的。
\end{example}

\begin{lemma}
\label{duality-lemma-twisted-inverse-image-bounded-below}
设 $f : X \to Y$ 为拟紧拟分离概形之间的态射。设
$a : D_\QCoh(\mathcal{O}_Y) \to D_\QCoh(\mathcal{O}_X)$ 为引理
\ref{duality-lemma-twisted-inverse-image} 中 $Rf_*$ 的右伴随。则 $a$ 把
$D^+_\QCoh(\mathcal{O}_Y)$ 映入 $D^+_\QCoh(\mathcal{O}_X)$。
事实上，存在整数 $N$，使得 $i \leq c$ 时 $H^i(K) = 0$ 蕴含
$i \leq c - N$ 时 $H^i(a(K)) = 0$。
\end{lemma}

\begin{proof}
由《概形的导出范畴》，引理
\ref{perfect-lemma-quasi-coherence-direct-image}，函子 $Rf_*$
具有有限上同调维数。换言之，存在整数 $N$，使得若 $i \geq c$ 时
$H^i(L) = 0$，则 $i \geq N + c$ 时 $H^i(Rf_*L) = 0$。
设 $K \in D^+_\QCoh(\mathcal{O}_Y)$ 且 $i \leq c$ 时 $H^i(K) = 0$。则
$$
\Hom_{D(\mathcal{O}_X)}(\tau_{\leq c - N}a(K), a(K)) =
\Hom_{D(\mathcal{O}_Y)}(Rf_*\tau_{\leq c - N}a(K), K) = 0
$$
由上述结论，显然这蕴含 $i \leq c - N$ 时 $H^i(a(K)) = 0$。
\end{proof}

\noindent
设 $f : X \to Y$ 为拟分离拟紧概形之间的态射，并以 $a$ 表示
$Rf_* : D_\QCoh(\mathcal{O}_X) \to D_\QCoh(\mathcal{O}_Y)$ 的右伴随。
对每个 $K \in D_\QCoh(\mathcal{O}_Y)$ 和
$L \in D_\QCoh(\mathcal{O}_X)$，得到典范映射
\begin{equation}
\label{duality-equation-sheafy-trace}
Rf_*R\SheafHom_{\mathcal{O}_X}(L, a(K))
\longrightarrow
R\SheafHom_{\mathcal{O}_Y}(Rf_*L, K)
\end{equation}
具体地说，此映射构造为复合
$$
Rf_*R\SheafHom_{\mathcal{O}_X}(L, a(K)) \to
R\SheafHom_{\mathcal{O}_Y}(Rf_*L, Rf_*a(K)) \to
R\SheafHom_{\mathcal{O}_Y}(Rf_*L, K)
$$
其中第一个箭头见《上同调》，注
\ref{cohomology-remark-projection-formula-for-internal-hom}，
第二个箭头则是伴随的余单位 $Rf_*a(K) \to K$。

\begin{lemma}
\label{duality-lemma-iso-on-RSheafHom}
设 $f : X \to Y$ 为拟紧拟分离概形之间的态射，$a$ 为
$Rf_* : D_\QCoh(\mathcal{O}_X) \to D_\QCoh(\mathcal{O}_Y)$ 的右伴随。
设 $L \in D_\QCoh(\mathcal{O}_X)$ 且 $K \in D_\QCoh(\mathcal{O}_Y)$。
对映射 (\ref{duality-equation-sheafy-trace})
$$
Rf_*R\SheafHom_{\mathcal{O}_X}(L, a(K))
\longrightarrow
R\SheafHom_{\mathcal{O}_Y}(Rf_*L, K)
$$
应用《概形的导出范畴》，第 \ref{perfect-section-better-coherator} 节中讨论的函子
$DQ_Y : D(\mathcal{O}_Y) \to D_\QCoh(\mathcal{O}_Y)$ 后，所得映射为同构。
\end{lemma}

\begin{proof}
由《概形的导出范畴》，引理 \ref{perfect-lemma-better-coherator}，
$DQ_Y$ 存在，故此陈述有意义。由于 $DQ_Y$ 是嵌入函子
$D_\QCoh(\mathcal{O}_Y) \to D(\mathcal{O}_Y)$ 的右伴随，
为证明本引理，只需说明对任意 $M \in D_\QCoh(\mathcal{O}_Y)$，
映射 (\ref{duality-equation-sheafy-trace}) 诱导双射
$$
\Hom_Y(M, Rf_*R\SheafHom_{\mathcal{O}_X}(L, a(K)))
\longrightarrow
\Hom_Y(M, R\SheafHom_{\mathcal{O}_Y}(Rf_*L, K))
$$
为此，使用如下等式链：
\begin{align*}
\Hom_Y(M, Rf_*R\SheafHom_{\mathcal{O}_X}(L, a(K)))
& =
\Hom_X(Lf^*M, R\SheafHom_{\mathcal{O}_X}(L, a(K))) \\
& =
\Hom_X(Lf^*M \otimes_{\mathcal{O}_X}^\mathbf{L} L, a(K)) \\
& =
\Hom_Y(Rf_*(Lf^*M \otimes_{\mathcal{O}_X}^\mathbf{L} L), K) \\
& =
\Hom_Y(M \otimes_{\mathcal{O}_Y}^\mathbf{L} Rf_*L, K) \\
& =
\Hom_Y(M, R\SheafHom_{\mathcal{O}_Y}(Rf_*L, K))
\end{align*}
第一个等式由《上同调》，引理 \ref{cohomology-lemma-adjoint}；
第二个等式由《上同调》，引理 \ref{cohomology-lemma-internal-hom}；
第三个等式由 $a$ 的构造；第四个等式由《概形的导出范畴》，引理
\ref{perfect-lemma-cohomology-base-change}（这是关键一步）；
第五个等式由《上同调》，引理 \ref{cohomology-lemma-internal-hom}。
\end{proof}

\begin{example}
\label{duality-example-iso-on-RSheafHom}
如果不应用“拟凝聚化函子” $DQ_Y$，引理 \ref{duality-lemma-iso-on-RSheafHom}
的陈述并不成立。事实上，设 $Y = \Spec(R)$ 且 $X = \mathbf{A}^1_R$，
取 $L = \mathcal{O}_X$ 和 $K = \mathcal{O}_Y$。箭头左端属于
$D_\QCoh(\mathcal{O}_Y)$，但箭头右端同构于
$\prod_{n \geq 0} \mathcal{O}_Y$，后者不是拟凝聚的。
\end{example}

\begin{remark}
\label{duality-remark-iso-on-RSheafHom}
在引理 \ref{duality-lemma-iso-on-RSheafHom} 的情形下，由《概形的导出范畴》，引理
\ref{perfect-lemma-pushforward-better-coherator}，有
$$
DQ_Y(Rf_*R\SheafHom_{\mathcal{O}_X}(L, a(K))) =
Rf_* DQ_X(R\SheafHom_{\mathcal{O}_X}(L, a(K)))
$$
因此，如果 $R\SheafHom_{\mathcal{O}_X}(L, a(K)) \in D_\QCoh(\mathcal{O}_X)$，
便可从箭头左端“删去” $DQ_Y$。另一方面，如果已知
$R\SheafHom_{\mathcal{O}_Y}(Rf_*L, K) \in D_\QCoh(\mathcal{O}_Y)$，
便可从箭头右端“删去” $DQ_Y$。如果两者都成立，就得到
(\ref{duality-equation-sheafy-trace}) 是同构。结合《概形的导出范畴》，引理
\ref{perfect-lemma-quasi-coherence-internal-hom} 可知，在下列情形中
$Rf_*R\SheafHom_{\mathcal{O}_X}(L, a(K)) \to
R\SheafHom_{\mathcal{O}_Y}(Rf_*L, K)$ 是同构：
\begin{enumerate}
\item $L$ 和 $Rf_*L$ 都是完美的；或者
\item $K$ 下有界，且 $L$ 和 $Rf_*L$ 都是伪凝聚的。
\end{enumerate}
对 (2)，用到：若 $K$ 下有界，则 $a(K)$ 也下有界；见引理
\ref{duality-lemma-twisted-inverse-image-bounded-below}。
\end{remark}

\begin{example}
\label{duality-example-iso-on-RSheafHom-noetherian}
设 $f : X \to Y$ 为 Noetherian 概形之间的固有态射，
$L \in D^-_{\textit{Coh}}(X)$ 且 $K \in D^+_{\QCoh}(\mathcal{O}_Y)$。
则映射 $Rf_*R\SheafHom_{\mathcal{O}_X}(L, a(K)) \to
R\SheafHom_{\mathcal{O}_Y}(Rf_*L, K)$ 是同构。事实上，由
《概形的导出范畴》，引理
\ref{perfect-lemma-identify-pseudo-coherent-noetherian} 和
\ref{perfect-lemma-direct-image-coherent}，复形 $L$ 与 $Rf_*L$ 均为伪凝聚的，
故注 \ref{duality-remark-iso-on-RSheafHom} 中的讨论适用。
\end{example}

\begin{lemma}
\label{duality-lemma-iso-global-hom}
设 $f : X \to Y$ 为拟分离拟紧概形之间的态射。对所有
$L \in D_\QCoh(\mathcal{O}_X)$ 和 $K \in D_\QCoh(\mathcal{O}_Y)$，
(\ref{duality-equation-sheafy-trace}) 诱导全局导出 Hom 的同构
$R\Hom_X(L, a(K)) \to R\Hom_Y(Rf_*L, K)$。
\end{lemma}

\begin{proof}
由《上同调》，第 \ref{cohomology-section-global-RHom} 节中的构造，有
$$
R\Hom_X(L, a(K)) =
R\Gamma(X, R\SheafHom_{\mathcal{O}_X}(L, a(K))) =
R\Gamma(Y, Rf_*R\SheafHom_{\mathcal{O}_X}(L, a(K)))
$$
以及
$$
R\Hom_Y(Rf_*L, K) = R\Gamma(Y, R\SheafHom_{\mathcal{O}_Y}(Rf_*L, K))
$$
因此本引理是引理 \ref{duality-lemma-iso-on-RSheafHom} 的推论。具体而言，
$D(\mathcal{O}_Y)$ 中的映射 $E \to E'$ 若诱导同构
$DQ_Y(E) \to DQ_Y(E')$，便诱导拟同构
$R\Gamma(Y, E) \to R\Gamma(Y, E')$。事实上，
$H^i(Y, E) = \Ext^i_Y(\mathcal{O}_Y, E) = \Hom(\mathcal{O}_Y[-i], E) =
\Hom(\mathcal{O}_Y[-i], DQ_Y(E))$，因为 $\mathcal{O}_Y[-i]$ 属于
$D_\QCoh(\mathcal{O}_Y)$，而 $DQ_Y$ 是嵌入函子
$D_\QCoh(\mathcal{O}_Y) \to D(\mathcal{O}_Y)$ 的右伴随。
\end{proof}







\section{推前的右伴随与开集限制}
\label{duality-section-restriction-to-opens}

\noindent
本节研究推前的右伴随在多大程度上与到开子概形的限制交换。
这是一个基变换问题，所以先在较一般的情形下讨论。

\medskip\noindent
我们经常需要判断推前的右伴随是否与基变换交换。为此，考虑拟紧拟分离概形的笛卡尔方块
\begin{equation}
\label{duality-equation-base-change}
\vcenter{
\xymatrix{
X' \ar[r]_{g'} \ar[d]_{f'} & X \ar[d]^f \\
Y' \ar[r]^g & Y
}
}
\end{equation}
以
$$
a  : D_\QCoh(\mathcal{O}_Y) \to D_\QCoh(\mathcal{O}_X)
\quad\text{且}\quad
a'  : D_\QCoh(\mathcal{O}_{Y'}) \to D_\QCoh(\mathcal{O}_{X'})
$$
表示 $Rf_*$ 与 $Rf'_*$ 的右伴随（引理
\ref{duality-lemma-twisted-inverse-image}）。考虑《上同调》，注
\ref{cohomology-remark-base-change} 中的基变换映射。它在具有拟凝聚上同调层的层的导出范畴上诱导函子变换
$$
Lg^* \circ Rf_* \longrightarrow Rf'_* \circ L(g')^*
$$
从而在相反方向诱导右伴随之间的变换
$$
a \circ Rg_* \longleftarrow Rg'_* \circ a'
$$

\begin{lemma}
\label{duality-lemma-flat-precompose-pus}
在图 (\ref{duality-equation-base-change}) 中，假设 $g$ 平坦；更一般地，假设
$f$ 与 $g$ Tor 独立。则 $a \circ Rg_* \leftarrow Rg'_* \circ a'$ 是同构。
\end{lemma}

\begin{proof}
在此情形下，由《概形的导出范畴》，引理
\ref{perfect-lemma-compare-base-change}，对 $D_\QCoh(\mathcal{O}_X)$ 中每个
$K$，基变换映射
$Lg^* \circ Rf_* K \longrightarrow Rf'_* \circ L(g')^*K$ 都是同构。
因此，相应的伴随函子之间的变换也是同构。
\end{proof}

\noindent
设 $f : X \to Y$ 为拟紧拟分离概形之间的态射，$V \subset Y$
为拟紧开子概形，并令 $U = f^{-1}(V)$。这给出如
(\ref{duality-equation-base-change}) 中的笛卡尔方块
$$
\xymatrix{
U \ar[r]_{j'} \ar[d]_{f|_U} & X \ar[d]^f \\
V \ar[r]^j & Y
}
$$
由引理 \ref{duality-lemma-flat-precompose-pus}，映射
$\xi : a \circ Rj_* \leftarrow Rj'_* \circ a'$ 是同构，
其中 $a$ 和 $a'$ 分别是 $Rf_*$ 与 $R(f|_U)_*$ 的右伴随。
于是得到函子 $D_\QCoh(\mathcal{O}_Y) \to D_\QCoh(\mathcal{O}_U)$ 的变换
\begin{equation}
\label{duality-equation-sheafy}
(j')^* \circ a \to
(j')^* \circ a \circ Rj_* \circ j^* \xrightarrow{\xi^{-1}}
(j')^* \circ Rj'_* \circ a' \circ j^* \to a' \circ j^*
\end{equation}
其中第一个箭头来自 $\text{id} \to Rj_* \circ j^*$，最后一个箭头来自同构
$(j')^* \circ Rj'_* \to \text{id}$。特别地，(\ref{duality-equation-sheafy})
在 $K$ 上取值时为同构，当且仅当
$a(K)|_U \to a(Rj_*(K|_V))|_U$ 为同构。

\begin{example}
\label{duality-example-not-supported-on-inverse-image}
存在 Noetherian 概形的有限态射 $f : X \to Y$，使得
(\ref{duality-equation-sheafy}) 在某个
$K \in D_{\textit{Coh}}(\mathcal{O}_Y)$ 上取值时不是同构。
具体地，设 $X = \Spec(B) \to Y = \Spec(A)$，其中
$A = k[x, \epsilon]$，$k$ 是域且 $\epsilon^2 = 0$，并且
$B = k[x] = A/(\epsilon)$。对 $n \in \mathbf{N}$，令
$M_n = A/(\epsilon, x^n)$。注意
$$
\Ext^i_A(B, M_n) = M_n,\quad i \geq 0
$$
因为 $B$ 有自由周期分解 $\ldots \to A \to A \to A$，
其中各映射均由乘以 $\epsilon$ 给出。考虑
$D_{\textit{Coh}}(A)$ 的对象 $K = \bigoplus M_n[n] = \prod M_n[n]$；
这里在 $D(A)$ 中相等，由《导出范畴》，引理
\ref{derived-lemma-direct-sums} 和 \ref{derived-lemma-products}。
由例 \ref{duality-example-affine-twisted-inverse-image}，$a(K)$ 对应于
$R\Hom(B, K)$，而由上述结论
$$
H^0(R\Hom(B, K)) = \Ext^0_A(B, K) =
\prod\nolimits_{n \geq 1} \Ext^n_A(B. M_n) = 
\prod\nolimits_{n \geq 1} M_n
$$
然而，这个模含有不被 $x$ 的任何幂零化的元素，而复形 $K$
的上同调中的每个元素都被 $x$ 的某个幂零化。换言之，取
$V = D(x)$、$U = D(x)$ 和复形 $K$ 时，映射
(\ref{duality-equation-sheafy}) 不可能是同构，因为 $(j')^*(a(K))$ 非零，
而 $a'(j^*K)$ 为零。
\end{example}

\begin{lemma}
\label{duality-lemma-when-sheafy}
设 $f : X \to Y$ 为拟紧拟分离概形之间的态射，$a$ 为
$Rf_* : D_\QCoh(\mathcal{O}_X) \to D_\QCoh(\mathcal{O}_Y)$ 的右伴随。
设 $V \subset Y$ 为拟紧开集，其逆像为 $U \subset X$。
\begin{enumerate}
\item 对每个支撑在 $Y \setminus V$ 上的
$Q \in D_\QCoh^+(\mathcal{O}_Y)$，像 $a(Q)$ 支撑在
$X \setminus U$ 上，当且仅当 (\ref{duality-equation-sheafy}) 在
$D_\QCoh^+(\mathcal{O}_Y)$ 的所有 $K$ 上都是同构。
\item 对每个支撑在 $Y \setminus V$ 上的
$Q \in D_\QCoh(\mathcal{O}_Y)$，像 $a(Q)$ 支撑在
$X \setminus U$ 上，当且仅当 (\ref{duality-equation-sheafy}) 在
$D_\QCoh(\mathcal{O}_Y)$ 的所有 $K$ 上都是同构。
\item 若 $a$ 与直和交换，则 (1) 的等价条件蕴含 (2) 的等价条件。
\end{enumerate}
\end{lemma}

\begin{proof}
先证 (1)。设 $K \in D_\QCoh^+(\mathcal{O}_Y)$，选取可区别三角
$$
K \to Rj_*K|_V \to Q \to K[1]
$$
注意 $Q$ 属于 $D_\QCoh^+(\mathcal{O}_Y)$（《概形的导出范畴》，引理
\ref{perfect-lemma-quasi-coherence-direct-image}），且支撑在
$Y \setminus V$ 上（《概形的导出范畴》，定义
\ref{perfect-definition-supported-on}）。应用 $a$，在 $X$ 上得到可区别三角
$$
a(K) \to a(Rj_*K|_V) \to a(Q) \to a(K)[1]
$$
若 $a(Q)$ 支撑在 $X \setminus U$ 上，则限制到 $U$ 后，映射
$a(K)|_U \to a(Rj_*K|_V)|_U$ 是同构，即 (\ref{duality-equation-sheafy})
在 $K$ 上是同构。逆命题立即成立。

\medskip\noindent
(2) 的证明与 (1) 的证明完全相同。

\medskip\noindent
再证 (3)。假设 (1) 的等价条件成立，令 $T = Y \setminus V$。
用记号 $D_{\QCoh, T}(\mathcal{O}_Y)$ 和
$D_{\QCoh, f^{-1}(T)}(\mathcal{O}_X)$ 分别表示上同调层支撑在
$T$ 和 $f^{-1}(T)$ 上的复形。由于 $a$ 与直和交换，以
$$
\{Q \in D_{\QCoh, T}(\mathcal{O}_Y) \mid
a(Q) \in D_{\QCoh, f^{-1}(T)}(\mathcal{O}_X)\}
$$
为对象的严格满、饱和三角子范畴 $\mathcal{D}$ 在直和下封闭，
因而在导出余极限下也封闭。另一方面，范畴
$D_{\QCoh, T}(\mathcal{O}_Y)$ 由一个完美对象 $E$ 生成；
见《概形的导出范畴》，引理
\ref{perfect-lemma-generator-with-support}。由假设，$E \in \mathcal{D}$。
由《导出范畴》，引理 \ref{derived-lemma-write-as-colimit}，
$D_{\QCoh, T}(\mathcal{O}_Y)$ 的每个对象 $Q$ 都是某个系统
$Q_1 \to Q_2 \to Q_3 \to \ldots$ 的导出余极限，
且各转移映射的锥都是 $E$ 的平移的直和。归纳可知，对所有 $n$，
$Q_n \in \mathcal{D}$，最终 $Q$ 也属于 $\mathcal{D}$。
因此 (2) 的等价条件成立。
\end{proof}

\begin{lemma}
\label{duality-lemma-proper-noetherian}
设 $Y$ 为拟紧拟分离概形，$f : X \to Y$ 为固有态射。若
\footnote{此证明适用于满足下述条件的拟紧拟分离概形态射：
对 $X$ 上每个完美的 $P$，$Rf_*P$ 都是伪凝聚的。由 Kiehl 的一个定理
\cite{Kiehl} 容易推出，当 $f$ 固有且伪凝聚时，此条件成立。
这才是本引理及本章其他若干结果的正确一般性。}
\begin{enumerate}
\item $f$ 平坦且有限表现；或者
\item $Y$ 为 Noetherian 概形，
\end{enumerate}
则对 $Y$ 的所有拟紧开集 $V$，引理 \ref{duality-lemma-when-sheafy}
第 (1) 部分的等价条件均成立。
\end{lemma}

\begin{proof}
设 $Q \in D^+_\QCoh(\mathcal{O}_Y)$ 支撑在 $Y \setminus V$ 上。
反设 $a(Q)$ 不支撑在 $X \setminus U$ 上。则可以找到 $U$ 上的完美复形
$P_U$ 以及非零映射 $P_U \to a(Q)|_U$；这由《概形的导出范畴》，定理
\ref{perfect-theorem-bondal-van-den-Bergh} 得到。再应用《概形的导出范畴》，引理
\ref{perfect-lemma-lift-map-from-perfect-complex-with-support}，
可以假设存在 $X$ 上的完美复形 $P$ 以及映射 $P \to a(Q)$，
其到 $U$ 的限制非零。由 $a$ 的定义，此映射伴随于映射 $Rf_*P \to Q$。

\medskip\noindent
复形 $Rf_*P$ 是伪凝聚的。在情形 (1) 中，这由《概形的导出范畴》，引理
\ref{perfect-lemma-flat-proper-pseudo-coherent-direct-image-general} 得到；
在情形 (2) 中，这由《概形的导出范畴》，引理
\ref{perfect-lemma-direct-image-coherent} 和
\ref{perfect-lemma-identify-pseudo-coherent-noetherian} 得到。
因此可应用《概形的导出范畴》，引理
\ref{perfect-lemma-map-from-pseudo-coherent-to-complex-with-support}，
得到完美复形的映射 $I \to \mathcal{O}_Y$，其到 $V$ 的限制为同构，
并且复合 $I \otimes^\mathbf{L}_{\mathcal{O}_Y} Rf_*P \to Rf_*P \to Q$ 为零。
由《概形的导出范畴》，引理 \ref{perfect-lemma-cohomology-base-change}，有
$I \otimes^\mathbf{L}_{\mathcal{O}_Y} Rf_*P =
Rf_*(Lf^*I \otimes^\mathbf{L}_{\mathcal{O}_X} P)$。
从而复合
$$
Lf^*I \otimes^\mathbf{L}_{\mathcal{O}_X} P \to P \to a(Q)
$$
为零。然而，它到 $U$ 的限制正是映射 $P|_U \to a(Q)|_U$，
而我们假设该映射非零。此矛盾完成证明。
\end{proof}












\section{推前的右伴随与基变换，I}
\label{duality-section-base-change-map}

\noindent
映射 (\ref{duality-equation-sheafy}) 是基变换映射的一种特殊情形。
具体地，假设有拟紧拟分离概形的笛卡尔图
$$
\xymatrix{
X' \ar[r]_{g'} \ar[d]_{f'} & X \ar[d]^f \\
Y' \ar[r]^g & Y
}
$$
即 (\ref{duality-equation-base-change}) 中那样的图。假设 $f$ 与 $g$
{\bf Tor 独立}。于是可考虑由复合
\begin{equation}
\label{duality-equation-base-change-map}
L(g')^* \circ a \to
L(g')^* \circ a \circ Rg_* \circ Lg^* \leftarrow
L(g')^* \circ Rg'_* \circ a' \circ Lg^* \to a' \circ Lg^*
\end{equation}
给出的函子 $D_\QCoh(\mathcal{O}_Y) \to D_\QCoh(\mathcal{O}_{X'})$ 的态射。
第一个箭头来自伴随映射 $\text{id} \to Rg_* Lg^*$，
最后一个箭头来自伴随映射 $L(g')^*Rg'_* \to \text{id}$。
为使中间的箭头可逆，需要 Tor 独立假设；见引理
\ref{duality-lemma-flat-precompose-pus}。也可以利用 $L(g')^*$ 与 $R(g')_*$
的伴随性，把 (\ref{duality-equation-base-change-map}) 看作自然变换
$$
a \to a \circ Rg_* \circ Lg^* \leftarrow Rg'_* \circ a' \circ Lg^*
$$
其中第二个箭头同样可逆。若 $M \in D_\QCoh(\mathcal{O}_X)$ 且
$K \in D_\QCoh(\mathcal{O}_Y)$，则此映射在 Yoneda 函子上由下列映射给出：
\begin{align*}
\Hom_X(M, a(K))
& =
\Hom_Y(Rf_*M, K) \\
& \to
\Hom_Y(Rf_*M, Rg_* Lg^*K) \\
& =
\Hom_{Y'}(Lg^*Rf_*M, Lg^*K) \\
& \leftarrow
\Hom_{Y'}(Rf'_* L(g')^*M, Lg^*K) \\
& =
\Hom_{X'}(L(g')^*M, a'(Lg^*K)) \\
& =
\Hom_X(M, Rg'_*a'(Lg^*K))
\end{align*}
其中向左的箭头由《概形的导出范畴》，引理
\ref{perfect-lemma-compare-base-change} 中的基变换定理可逆；
这使构造更加明确。

\medskip\noindent
本节先证明，该基变换映射在叠置方块时满足若干自然相容性，
正如通常的基变换映射在《上同调》，注
\ref{cohomology-remark-compose-base-change} 和
\ref{cohomology-remark-compose-base-change-horizontal} 中那样。
建议初次阅读时跳过本节余下内容。

\begin{lemma}
\label{duality-lemma-compose-base-change-maps}
考虑拟紧拟分离概形的交换图
$$
\xymatrix{
X' \ar[r]_k \ar[d]_{f'} & X \ar[d]^f \\
Y' \ar[r]^l \ar[d]_{g'} & Y \ar[d]^g \\
Z' \ar[r]^m & Z
}
$$
其中两个方块都是笛卡尔的，并且 $f$ 与 $l$ Tor 独立，
$g$ 与 $m$ 也 Tor 独立。则两个方块的映射
(\ref{duality-equation-base-change-map}) 复合后给出外侧矩形的基变换映射
（精确陈述见证明）。
\end{lemma}

\begin{proof}
由假设可知 $g \circ f$ 与 $m$ Tor 独立（细节从略），故此陈述有意义。
在本证明中，以 $k^*$ 代替 $Lk^*$，以 $f_*$ 代替 $Rf_*$。
令 $a$、$b$、$c$ 分别为引理 \ref{duality-lemma-twisted-inverse-image}
对 $f$、$g$、$g \circ f$ 给出的右伴随，并对带撇版本作类似约定。
对应于上方方块的箭头是复合
$$
\gamma_{top} :
k^* \circ a \to k^* \circ a \circ l_* \circ l^*
\xleftarrow{\xi_{top}} k^* \circ k_* \circ a' \circ l^* \to a' \circ l^*
$$
其中 $\xi_{top} : k_* \circ a' \to a \circ l_*$ 是同构，因而可逆；
它是基变换映射 $l^* \circ f_* \to f'_* \circ k^*$ 的“对偶”箭头。
外侧两个箭头来自典范映射 $1 \to l_* \circ l^*$ 和
$k^* \circ k_* \to 1$。类似地，对第二个方块有
$$
\gamma_{bot} :
l^* \circ b \to l^* \circ b \circ m_* \circ m^*
\xleftarrow{\xi_{bot}} l^* \circ l_* \circ b' \circ m^* \to b' \circ m^*
$$
对外侧矩形则得到
$$
\gamma_{rect} :
k^* \circ c \to k^* \circ c \circ m_* \circ m^*
\xleftarrow{\xi_{rect}} k^* \circ k_* \circ c' \circ m^* \to c' \circ m^*
$$
有 $(g \circ f)_* = g_* \circ f_*$，因而 $c = a \circ b$，
类似地 $c' = a' \circ b'$。本引理断言 $\gamma_{rect}$ 等于复合
$$
k^* \circ c = k^* \circ a \circ b \xrightarrow{\gamma_{top}}
a' \circ l^* \circ b \xrightarrow{\gamma_{bot}}
a' \circ b' \circ m^* = c' \circ m^*
$$
为此，考察下图：
$$
\xymatrix{
& & k^* \circ a \circ b \ar[d] \ar[lldd] \\
& & k^* \circ a \circ l_* \circ l^* \circ b \ar[ld] \\
k^* \circ a \circ b \circ m_* \circ m^* \ar[r] &
k^* \circ a \circ l_* \circ l^* \circ b \circ m_* \circ m^* &
k^* \circ k_* \circ a' \circ l^* \circ b \ar[u]_{\xi_{top}} \ar[d] \ar[ld] \\
& k^*\circ k_* \circ a' \circ l^* \circ b \circ m_* \circ m^*
\ar[u]_{\xi_{top}} \ar[rd] &
a' \circ l^* \circ b \ar[d] \\
k^* \circ k_* \circ a' \circ b' \circ m^* \ar[uu]_{\xi_{rect}} \ar[ddrr] &
k^*\circ k_* \circ a' \circ l^* \circ l_* \circ b' \circ m^*
\ar[u]_{\xi_{bot}} \ar[l] \ar[dr] &
a' \circ l^* \circ b \circ m_* \circ m^* \\
& & a' \circ l^* \circ l_* \circ b' \circ m^* \ar[u]_{\xi_{bot}} \ar[d] \\
& & a' \circ b' \circ m^*
}
$$
沿右侧向下得到上述复合，沿左侧向下得到 $\gamma_{rect}$。
图右侧的所有四边形都交换，这由《范畴》，引理
\ref{categories-lemma-properties-2-cat-cats} 得到；也可更直接地使用
《范畴》，定义 \ref{categories-definition-horizontal-composition} 之前的讨论。
因此，只需说明在将箭头 $\xi_{top}$、$\xi_{bot}$、$\xi_{rect}$
取逆之后，图
$$
\xymatrix{
a \circ l_* \circ l^* \circ b \circ m_* &
a \circ b \circ m_* \ar[l] \\
k_* \circ a' \circ l^* \circ b \circ m_* \ar[u]_{\xi_{top}} & \\
k_* \circ a' \circ l^* \circ l_* \circ b' \ar[u]_{\xi_{bot}} \ar[r] &
k_* \circ a' \circ b' \ar[uu]_{\xi_{rect}}
}
$$
交换；注意，这不同于要求原图交换。然而，由《范畴》，引理
\ref{categories-lemma-properties-2-cat-cats}，图
$$
\xymatrix{
& a \circ l_* \circ l^* \circ b \circ m_* \\
a \circ l_* \circ l^* \circ l_* \circ b'
\ar[ru]^{\xi_{bot}} & &
k_* \circ a' \circ l^* \circ b \circ m_* \ar[ul]_{\xi_{top}} \\
& k_* \circ a' \circ l^* \circ l_* \circ b'
\ar[ul]^{\xi_{top}} \ar[ur]_{\xi_{bot}}
}
$$
交换。又因为图
$$
\vcenter{
\xymatrix{
a \circ l_* \circ l^* \circ b \circ m_* & a \circ b \circ m \ar[l] \\
a \circ l_* \circ l^* \circ l_* \circ b' \ar[u] &
a \circ l_* \circ b' \ar[l] \ar[u]
}
}
\quad\text{且}\quad
\vcenter{
\xymatrix{
a \circ l_* \circ l^* \circ l_* \circ b' \ar[r] & a \circ l_* \circ b' \\
k_* \circ a' \circ l^* \circ l_* \circ b' \ar[u] \ar[r] &
k_* \circ a' \circ b' \ar[u]
}
}
$$
交换（见所引参考），且复合 $l_* \to l_* \circ l^* \circ l_* \to l_*$
是恒等，所以只需证明
$$
k \circ a' \circ b' \xrightarrow{\xi_{bot}} a \circ l_* \circ b
\xrightarrow{\xi_{top}} a \circ b \circ m_*
$$
等于 $\xi_{rect}$（通过识别 $a \circ b = c$ 和
$a' \circ b' = c'$）。这正是《上同调》，注
\ref{cohomology-remark-compose-base-change} 中陈述的对偶，证明完成。
\end{proof}

\begin{lemma}
\label{duality-lemma-compose-base-change-maps-horizontal}
考虑拟紧拟分离概形的交换图
$$
\xymatrix{
X'' \ar[r]_{g'} \ar[d]_{f''} & X' \ar[r]_g \ar[d]_{f'} & X \ar[d]^f \\
Y'' \ar[r]^{h'} & Y' \ar[r]^h & Y
}
$$
其中两个方块都是笛卡尔的，并且 $f$ 与 $h$ Tor 独立，
$f'$ 与 $h'$ 也 Tor 独立。则两个方块的映射
(\ref{duality-equation-base-change-map}) 复合后给出外侧矩形的基变换映射
（精确陈述见证明）。
\end{lemma}

\begin{proof}
由假设可知 $f$ 与 $h \circ h'$ Tor 独立（细节从略），故此陈述有意义。
在本证明中，以 $g^*$ 代替 $Lg^*$，以 $f_*$ 代替 $Rf_*$。
令 $a$、$a'$、$a''$ 分别为引理 \ref{duality-lemma-twisted-inverse-image}
对 $f$、$f'$、$f''$ 给出的右伴随。对应于右侧方块的箭头是复合
$$
\gamma_{right} :
g^* \circ a \to g^* \circ a \circ h_* \circ h^*
\xleftarrow{\xi_{right}} g^* \circ g_* \circ a' \circ h^* \to a' \circ h^*
$$
其中 $\xi_{right} : g_* \circ a' \to a \circ h_*$ 是同构，因而可逆；
它是基变换映射 $h^* \circ f_* \to f'_* \circ g^*$ 的“对偶”箭头。
外侧两个箭头来自典范映射 $1 \to h_* \circ h^*$ 和
$g^* \circ g_* \to 1$。类似地，对左侧方块有
$$
\gamma_{left} :
(g')^* \circ a' \to (g')^* \circ a' \circ (h')_* \circ (h')^*
\xleftarrow{\xi_{left}}
(g')^* \circ (g')_* \circ a'' \circ (h')^* \to a'' \circ (h')^*
$$
对外侧矩形则得到
$$
\gamma_{rect} :
k^* \circ a \to
k^* \circ a \circ m_* \circ m^* \xleftarrow{\xi_{rect}}
k^* \circ k_* \circ a'' \circ m^* \to
a'' \circ m^*
$$
其中 $k = g \circ g'$ 且 $m = h \circ h'$。有
$k^* = (g')^* \circ g^*$ 和 $m^* = (h')^* \circ h^*$。
本引理断言 $\gamma_{rect}$ 等于复合
$$
k^* \circ a =
(g')^* \circ g^* \circ a \xrightarrow{\gamma_{right}}
(g')^* \circ a' \circ h^* \xrightarrow{\gamma_{left}}
a'' \circ (h')^* \circ h^* = a'' \circ m^*
$$
为此，考察下图：
$$
\xymatrix{
& (g')^* \circ g^* \circ a \ar[d] \ar[ddl] \\
& (g')^* \circ g^* \circ a \circ h_* \circ h^* \ar[ld] \\
(g')^* \circ g^* \circ a \circ h_* \circ (h')_* \circ (h')^* \circ h^* &
(g')^* \circ g^* \circ g_* \circ a' \circ h^*
\ar[u]_{\xi_{right}} \ar[d] \ar[ld] \\
(g')^* \circ g^* \circ g_* \circ a' \circ (h')_* \circ (h')^* \circ h^*
\ar[u]_{\xi_{right}} \ar[dr] &
(g')^* \circ a' \circ h^* \ar[d] \\
(g')^* \circ g^* \circ g_* \circ (g')_* \circ a'' \circ (h')^* \circ h^*
\ar[u]_{\xi_{left}} \ar[ddr] \ar[dr] &
(g')^* \circ a' \circ (h')_* \circ (h')^* \circ h^* \\
& (g')^*\circ (g')_* \circ a'' \circ (h')^* \circ h^*
\ar[u]_{\xi_{left}} \ar[d] \\
& a'' \circ (h')^* \circ h^*
}
$$
沿右侧向下得到上述复合，沿左侧向下得到 $\gamma_{rect}$。
图右侧的所有四边形都交换，这由《范畴》，引理
\ref{categories-lemma-properties-2-cat-cats} 得到；也可更直接地使用
《范畴》，定义 \ref{categories-definition-horizontal-composition} 之前的讨论。
因此只需说明
$$
g_* \circ (g')_* \circ a'' \xrightarrow{\xi_{left}}
g_* \circ a' \circ (h')_* \xrightarrow{\xi_{right}}
a \circ h_* \circ (h')_*
$$
等于 $\xi_{rect}$。这正是《上同调》，注
\ref{cohomology-remark-compose-base-change-horizontal} 中陈述的对偶，
证明完成。
\end{proof}

\begin{remark}
\label{duality-remark-going-around}
考虑拟紧拟分离概形的交换图
$$
\xymatrix{
X'' \ar[r]_{k'} \ar[d]_{f''} & X' \ar[r]_k \ar[d]_{f'} & X \ar[d]^f \\
Y'' \ar[r]^{l'} \ar[d]_{g''} & Y' \ar[r]^l \ar[d]_{g'} & Y \ar[d]^g \\
Z'' \ar[r]^{m'} & Z' \ar[r]^m & Z
}
$$
其中所有方块均为笛卡尔的，并且 $(f, l)$、$(g, m)$、
$(f', l')$、$(g', m')$ 都是 Tor 独立的态射对。
令 $a$、$a'$、$a''$、$b$、$b'$、$b''$ 分别为引理
\ref{duality-lemma-twisted-inverse-image} 对 $f$、$f'$、$f''$、$g$、$g'$、$g''$
给出的右伴随。如下把图中的方块标记为 $A$、$B$、$C$、$D$：
$$
\begin{matrix}
A & B \\
C & D
\end{matrix}
$$
则各方块的映射 (\ref{duality-equation-base-change-map}) 为
（这里使用 $k^* = Lk^*$ 等简记）
$$
\begin{matrix}
\gamma_A : (k')^* \circ a' \to a'' \circ (l')^* &
\gamma_B : k^* \circ a \to a' \circ l^* \\
\gamma_C : (l')^* \circ b' \to b'' \circ (m')^* &
\gamma_D : l^* \circ b \to b' \circ m^*
\end{matrix}
$$
对 $2 \times 1$ 和 $1 \times 2$ 的矩形，还有四个基变换映射
$$
\begin{matrix}
\gamma_{A + B} : (k \circ k')^* \circ a \to a'' \circ (l \circ l')^* \\
\gamma_{C + D} : (l \circ l')^* \circ b \to b'' \circ (m \circ m')^* \\
\gamma_{A + C} : (k')^* \circ (a' \circ b') \to (a'' \circ b'') \circ (m')^* \\
\gamma_{B + D} : k^* \circ (a \circ b) \to (a' \circ b') \circ m^*
\end{matrix}
$$
由引理 \ref{duality-lemma-compose-base-change-maps-horizontal}，有
$$
\gamma_{A + B} = \gamma_A \circ \gamma_B, \quad
\gamma_{C + D} = \gamma_C \circ \gamma_D
$$
由引理 \ref{duality-lemma-compose-base-change-maps}，有
$$
\gamma_{A + C} = \gamma_C \circ \gamma_A, \quad
\gamma_{B + D} = \gamma_D \circ \gamma_B
$$
严格地说，按《范畴》，第 \ref{categories-section-formal-cat-cat} 节的记号，
应写成 $\gamma_{A + B} = (\gamma_A \star \text{id}_{l^*}) \circ
(\text{id}_{(k')^*} \star \gamma_B)$，其他情形类似。不过，这里继续沿用
引理 \ref{duality-lemma-compose-base-change-maps} 和
\ref{duality-lemma-compose-base-change-maps-horizontal} 证明中的记号滥用，
省略与恒等变换的 $\star$ 积；只要知道变换的源和目标，便可确定应补入哪些项。
由以上讨论，先验地得到两个变换
$$
(k')^* \circ k^* \circ a \circ b
\longrightarrow
a'' \circ b'' \circ (m')^* \circ m^*
$$
即
$$
\gamma_C \circ \gamma_A \circ \gamma_D \circ \gamma_B =
\gamma_{A + C} \circ \gamma_{B + D}
$$
和
$$
\gamma_C \circ \gamma_D \circ \gamma_A \circ \gamma_B =
\gamma_{C + D} \circ \gamma_{A + B}
$$
本注的要点是指出这两个变换相等。为此，只需证明图
$$
\xymatrix{
(k')^* \circ a' \circ l^* \circ b \ar[r]_{\gamma_D} \ar[d]_{\gamma_A} &
(k')^* \circ a' \circ b' \circ m^* \ar[d]^{\gamma_A} \\
a'' \circ (l')^* \circ l^* \circ b \ar[r]^{\gamma_D} &
a'' \circ (l')^* \circ b' \circ m^*
}
$$
交换。这由《范畴》，引理 \ref{categories-lemma-properties-2-cat-cats}
得到；也可更直接地使用《范畴》，定义
\ref{categories-definition-horizontal-composition} 之前的讨论。
\end{remark}







\section{推前的右伴随与基变换，II}
\label{duality-section-base-change-II}

\noindent
本节证明，第 \ref{duality-section-base-change-map} 节的基变换映射在若干情形下是同构。
首先注意，只要推前的右伴随的形成与开集限制交换，就只需在仿射开集上检验。

\begin{remark}
\label{duality-remark-check-over-affines}
考虑拟紧拟分离概形的笛卡尔图
$$
\xymatrix{
X' \ar[r]_{g'} \ar[d]_{f'} & X \ar[d]^f \\
Y' \ar[r]^g & Y
}
$$
其中 $(g, f)$ Tor 独立。设 $V \subset Y$ 和 $V' \subset Y'$
为满足 $g(V') \subset V$ 的仿射开集。作笛卡尔图
$$
\vcenter{
\xymatrix{
U \ar[r] \ar[d] & X \ar[d] \\
V \ar[r] & Y
}
}
\quad\text{且}\quad
\vcenter{
\xymatrix{
U' \ar[r] \ar[d] & X' \ar[d] \\
V' \ar[r] & Y'
}
}
$$
假设关于 $K$ 和第一个图的 (\ref{duality-equation-sheafy})，以及关于
$Lg^*K$ 和第二个图的 (\ref{duality-equation-sheafy})，都是同构。
则基变换映射 (\ref{duality-equation-base-change-map})
$$
L(g')^*a(K) \longrightarrow a'(Lg^*K)
$$
到 $U'$ 的限制，同构于关于 $K|_V$ 及笛卡尔图
$$
\xymatrix{
U' \ar[r] \ar[d] & U \ar[d] \\
V' \ar[r] & V
}
$$
的基变换映射 (\ref{duality-equation-base-change-map})。这是因为
(\ref{duality-equation-sheafy}) 是基变换映射 (\ref{duality-equation-base-change-map})
的特殊情形，而且水平叠置方块时，基变换映射正确复合；见引理
\ref{duality-lemma-compose-base-change-maps-horizontal}。因此，为检验基变换映射
限制到 $U'$ 后是否为同构，只需处理最后一个图。
\end{remark}

\begin{lemma}
\label{duality-lemma-more-base-change}
在图 (\ref{duality-equation-base-change}) 中，假设
\begin{enumerate}
\item $g : Y' \to Y$ 是仿射概形之间的态射；
\item $f : X \to Y$ 是固有的；并且
\item $f$ 与 $g$ Tor 独立。
\end{enumerate}
则在下列情形中，基变换映射 (\ref{duality-equation-base-change-map}) 诱导同构
$$
L(g')^*a(K) \longrightarrow a'(Lg^*K)
$$
\begin{enumerate}
\item 若 $f$ 平坦且有限表现，则对所有 $K \in D_\QCoh(\mathcal{O}_Y)$；
\item 若 $f$ 完美且 $Y$ 为 Noetherian 概形，则对所有
$K \in D_\QCoh(\mathcal{O}_Y)$；
\item 若 $g$ 具有有限 Tor 维数且 $Y$ 为 Noetherian 概形，则对
$K \in D_\QCoh^+(\mathcal{O}_Y)$。
\end{enumerate}
\end{lemma}

\begin{proof}
记 $Y = \Spec(A)$、$Y' = \Spec(A')$。作为仿射态射的基变换，态射
$g'$ 是仿射的。设 $M$ 为 $D_\QCoh(\mathcal{O}_X)$ 的一个完美生成元；
见《概形的导出范畴》，定理
\ref{perfect-theorem-bondal-van-den-Bergh}。于是 $L(g')^*M$ 是
$D_\QCoh(\mathcal{O}_{X'})$ 的生成元；见《概形的导出范畴》，注
\ref{perfect-remark-pullback-generator}。因此，只需证明
(\ref{duality-equation-base-change-map}) 诱导全局 Hom 复形的同构
\begin{equation}
\label{duality-equation-iso}
R\Hom_{X'}(L(g')^*M, L(g')^*a(K))
\longrightarrow
R\Hom_{X'}(L(g')^*M, a'(Lg^*K))
\end{equation}
见《上同调》，第 \ref{cohomology-section-global-RHom} 节；
因为这将蕴含 $L(g')^*a(K) \to a'(Lg^*K)$ 的锥为零。
证明结构如下：先说明这两个 Hom 复形同构，再在证明最后说明该同构确由
(\ref{duality-equation-iso}) 诱导。

\medskip\noindent
左端。由于 $M$ 是完美的，由《概形的导出范畴》，引理
\ref{perfect-lemma-affine-morphism-and-hom-out-of-perfect}，典范映射
$$
R\Hom_X(M, a(K)) \otimes^\mathbf{L}_A A'
\longrightarrow
R\Hom_{X'}(L(g')^*M, L(g')^*a(K))
$$
是同构。将其与引理 \ref{duality-lemma-iso-global-hom} 的同构
$R\Hom_Y(Rf_*M, K) = R\Hom_X(M, a(K))$ 结合，可知左端等于
$R\Hom_Y(Rf_*M, K) \otimes^\mathbf{L}_A A'$。

\medskip\noindent
右端。先使用引理 \ref{duality-lemma-iso-global-hom} 的同构
$$
R\Hom_{X'}(L(g')^*M, a'(Lg^*K)) = R\Hom_{Y'}(Rf'_*L(g')^*M, Lg^*K)
$$
再由《概形的导出范畴》，引理 \ref{perfect-lemma-compare-base-change}，
基变换映射 $Lg^*Rf_*M \to Rf'_*L(g')^*M$ 是同构，故可将右端改写为
$R\Hom_{Y'}(Lg^*Rf_*M, Lg^*K)$。由于 $Y$、$Y'$ 仿射，且
$K$、$Rf_*M$ 属于 $D_\QCoh(\mathcal{O}_Y)$（《概形的导出范畴》，引理
\ref{perfect-lemma-quasi-coherence-direct-image}），在 $D(A')$ 中有典范映射
$$
\beta :
R\Hom_Y(Rf_*M, K) \otimes^\mathbf{L}_A A'
\longrightarrow
R\Hom_{Y'}(Lg^*Rf_*M, Lg^*K)
$$
这就是《代数进阶》，等式 (\ref{more-algebra-equation-base-change-RHom})
中的箭头；这里使用《概形的导出范畴》，引理
\ref{perfect-lemma-affine-compare-bounded} 和
\ref{perfect-lemma-quasi-coherence-internal-hom} 在代数表述之间来回转换。
\begin{enumerate}
\item 若 $f$ 平坦且有限表现，则由《概形的导出范畴》，引理
\ref{perfect-lemma-flat-proper-perfect-direct-image-general}，复形 $Rf_*M$
在 $Y$ 上完美；再由《代数进阶》，引理
\ref{more-algebra-lemma-base-change-RHom} 的第 (1) 部分，$\beta$ 是同构。
\item 若 $f$ 完美且 $Y$ 为 Noetherian 概形，则由《态射进阶》，引理
\ref{more-morphisms-lemma-perfect-proper-perfect-direct-image}，复形
$Rf_*M$ 在 $Y$ 上完美；于是 $\beta$ 同前为同构。
\item 若 $g$ 具有有限 Tor 维数且 $Y$ 为 Noetherian 概形，则复形
$Rf_*M$ 在 $Y$ 上伪凝聚（《概形的导出范畴》，引理
\ref{perfect-lemma-direct-image-coherent} 和
\ref{perfect-lemma-identify-pseudo-coherent-noetherian}）；由《代数进阶》，
引理 \ref{more-algebra-lemma-base-change-RHom} 的第 (4) 部分，$\beta$ 是同构。
\end{enumerate}
因此得到与上一段相同的答案。

\medskip\noindent
在证明余下部分中，我们说明第二、三段给出的
(\ref{duality-equation-iso}) 左右两端的识别，实际上正由
(\ref{duality-equation-iso}) 给出。为使公式便于处理，记
$(-, -)_X = R\Hom_X(-, -)$，以 $- \otimes A'$ 代替
$- \otimes_A^\mathbf{L} A'$，并简记 $g^* = Lg^*$、$f_* = Rf_*$。
考虑交换图
$$
\xymatrix{
((g')^*M, (g')^*a(K))_{X'} \ar[d] &
(M, a(K))_X \otimes A' \ar[l]^-\alpha \ar[d] &
(f_*M, K)_Y \otimes A' \ar@{=}[l] \ar[d] \\
((g')^*M, (g')^*a(g_*g^*K))_{X'} &
(M, a(g_*g^*K))_X \otimes A' \ar[l]^-\alpha &
(f_*M, g_*g^*K)_Y \otimes A' \ar@{=}[l] \ar@/_4pc/[dd]_{\mu'} \\
((g')^*M, (g')^*g'_*a'(g^*K))_{X'} \ar[u] \ar[d] &
(M, g'_*a'(g^*K))_X \otimes A' \ar[u] \ar[l]^-\alpha \ar[ld]^\mu &
(f_*M, K) \otimes A' \ar[d]^\beta \\
((g')^*M, a'(g^*K))_{X'} &
(f'_*(g')^*M, g^*K)_{Y'} \ar@{=}[l] \ar[r] &
(g^*f_*M, g^*K)_{Y'}
}
$$
标记为 $\alpha$ 的箭头，是《概形的导出范畴》，引理
\ref{perfect-lemma-affine-morphism-and-hom-out-of-perfect}
对四角为 $X', X, Y', Y$ 的图给出的映射。由于水平箭头对各变量函子性，
图的上部交换。中间的竖直箭头来自引理
\ref{duality-lemma-flat-precompose-pus} 的可逆变换
$g'_* \circ a' \to a \circ g_*$，故中间方块交换。沿左侧向下正是
(\ref{duality-equation-iso})。上方的水平箭头给出证明第二段所用的识别；
下方的水平箭头（包括 $\beta$）给出第三段所用的识别。
给定 $E \in D(A)$、$E' \in D(A')$ 及 $D(A)$ 中的
$c : E \to E'$，用 $\mu_c : E \otimes A' \to E'$ 表示由 $c$
以及限制与基变换的伴随性诱导的映射；若 $c$ 明确，就写
$\mu = \mu_c$，即从记号中省略 $c$。图中的 $\mu$ 正是这种映射，
其中 $c$ 由识别
$(M, g'_*a(g^*K))_X = ((g')^*M, a'(g^*K))_{X'}$ 给出；
含 $\mu$ 的三角形由《概形的导出范畴》，注
\ref{perfect-remark-multiplication-map} 交换。

\medskip\noindent
注意图
$$
\xymatrix{
(M, a(g_*g^*K))_X &
(f_*M, g_* g^*K)_Y \ar@{=}[l] &
(g^*f_*M, g^*K)_{Y'} \ar@{=}[l] \\
(M, g'_* a'(g^*K))_X \ar[u] &
((g')^*M, a'(g^*K))_{X'} \ar@{=}[l] &
(f'_*(g')^*M, g^*K)_{Y'} \ar@{=}[l] \ar[u]
}
$$
依变换 $g'_* \circ a' \to a \circ g_*$ 的定义交换。令 $\mu'$
如上对应于识别 $(f_*M, g_*g^*K)_X = (g^*f_*M, g^*K)_{Y'}$，
则六边形也交换。因此，只需证明 $\beta$ 等于
$(f_*M, K)_Y \otimes A' \to (f_*M, g_*g^*K)_X \otimes A'$
与 $\mu'$ 的复合。为此，只需证明所诱导的两个映射
$(f_*M, K)_Y \to (g^*f_*M, g^*K)_{Y'}$ 相同。换言之，
只需证明对所有 $E, K \in D(A)$，图
$$
\xymatrix{
R\Hom_A(E, K) \ar[rr]_{\text{由下列对象诱导： }\beta} \ar[rd] & &
R\Hom_{A'}(E \otimes_A^\mathbf{L} A', K \otimes_A^\mathbf{L} A') \\
& R\Hom_A(E, K \otimes_A^\mathbf{L} A') \ar[ru]
}
$$
交换。而《代数进阶》，第 \ref{more-algebra-section-base-change-RHom} 节
正是如此构造 $\beta$ 的，证明完成。
\end{proof}








\section{推前的右伴随与迹映射}
\label{duality-section-trace}

\noindent
设 $f : X \to Y$ 为拟紧拟分离概形之间的态射，并设
$a : D_\QCoh(\mathcal{O}_Y) \to D_\QCoh(\mathcal{O}_X)$ 为引理
\ref{duality-lemma-twisted-inverse-image} 中的右伴随。由《范畴》，第
\ref{categories-section-adjoint} 节，得到函子变换
$$
\text{Tr}_f : Rf_* \circ a \longrightarrow \text{id}
$$
对 $K \in D_\QCoh(\mathcal{O}_Y)$，相应的映射
$\text{Tr}_{f, K} : Rf_*a(K) \longrightarrow K$ 有时称为{\it 迹映射}。
它具有如下性质：刻画右伴随的双射
$$
\Hom_X(L, a(K)) \longrightarrow \Hom_Y(Rf_*L, K)
$$
对 $L \in D_\QCoh(\mathcal{O}_X)$ 由
$$
\varphi \longmapsto \text{Tr}_{f, K} \circ Rf_*\varphi
$$
给出。映射 (\ref{duality-equation-sheafy-trace})
$$
Rf_*R\SheafHom_{\mathcal{O}_X}(L, a(K))
\longrightarrow
R\SheafHom_{\mathcal{O}_Y}(Rf_*L, K)
$$
由与 $\text{Tr}_{f, K}$ 复合得到。本节将考察的每个迹映射都是这个迹映射的特殊情形。
在讨论若干特殊情形之前，先证明迹映射的形成与基变换交换。

\begin{lemma}[迹映射与基变换]
\label{duality-lemma-trace-map-and-base-change}
假设有图 (\ref{duality-equation-base-change})，其中 $f$ 与 $g$ Tor 独立。
则映射 $1 \star \text{Tr}_f : Lg^* \circ Rf_* \circ a \to Lg^*$ 和
$\text{Tr}_{f'} \star 1 : Rf'_* \circ a' \circ Lg^* \to Lg^*$
通过基变换映射
$\beta : Lg^* \circ Rf_* \to Rf'_* \circ L(g')^*$
（《上同调》，注 \ref{cohomology-remark-base-change}）以及
$\alpha : L(g')^* \circ a \to a' \circ Lg^*$
（\ref{duality-equation-base-change-map}）相一致。更精确地，函子变换的图
$$
\xymatrix{
Lg^* \circ Rf_* \circ a
\ar[d]_{\beta \star 1} \ar[r]_-{1 \star \text{Tr}_f} &
Lg^* \\
Rf'_* \circ L(g')^* \circ a \ar[r]^{1 \star \alpha} &
Rf'_* \circ a' \circ Lg^* \ar[u]_{\text{Tr}_{f'} \star 1}
}
$$
交换。
\end{lemma}

\begin{proof}
本证明中，以 $f_*$ 代替 $Rf_*$，以 $g^*$ 代替 $Lg^*$，
并省略与恒等变换的 $\star$ 积；只要知道变换的源和目标，
便可确定应补入哪些项。回忆，
$\beta : g^* \circ f_* \to f'_* \circ (g')^*$ 是同构，而 $\alpha$
利用同构 $\beta^\vee : g'_* \circ a' \to a \circ g_*$ 定义；
后者是 $\beta$ 的伴随，见引理 \ref{duality-lemma-flat-precompose-pus} 及其证明。
先注意，引理所述图的上方水平箭头等于复合
$$
g^* \circ f_* \circ a \to
g^* \circ f_* \circ a \circ g_* \circ g^* \to
g^* \circ g_* \circ g^* \to g^*
$$
其中第一个箭头是 $(g^*, g_*)$ 的单位，第二个箭头是
$\text{Tr}_f$，第三个箭头是 $(g^*, g_*)$ 的余单位。
这是单位与余单位的复合
$g^* \to g^* \circ g_* \circ g^* \to g^*$ 为恒等这一事实的直接推论。
考虑图
$$
\xymatrix{
& g^* \circ f_* \circ a \ar[ld]_\beta \ar[d] \ar[r]_{\text{Tr}_f} & g^* \\
f'_* \circ (g')^* \circ a \ar[dr] &
g^* \circ f_* \circ a \circ g_* \circ g^* \ar[d]_\beta \ar[ru] &
g^* \circ f_* \circ g'_* \circ a' \circ g^* \ar[l]_{\beta^\vee} \ar[d]_\beta &
f'_* \circ a' \circ g^* \ar[lu]_{\text{Tr}_{f'}} \\
& f'_* \circ (g')^* \circ a \circ g_* \circ g^* &
f'_* \circ (g')^* \circ g'_* \circ a' \circ g^* \ar[ru] \ar[l]_{\beta^\vee}
}
$$
其中两个方块交换，这由《范畴》，引理
\ref{categories-lemma-properties-2-cat-cats} 得到；也可更直接地使用
《范畴》，定义 \ref{categories-definition-horizontal-composition} 之前的讨论。
三角形由上述讨论交换。由《范畴》，引理
\ref{categories-lemma-transformation-between-functors-and-adjoints}，方块
$$
\xymatrix{
g^* \circ f_* \circ g'_* \circ a' \ar[d]_{\beta^\vee} \ar[r]_-\beta &
f'_* \circ (g')^* \circ g'_* \circ a' \ar[d] \\
g^* \circ f_* \circ a \circ g_* \ar[r] &
\text{id}
}
$$
交换，从而大图中的五边形交换。由于 $\beta$ 和 $\beta^\vee$ 均为同构，
并且依定义沿大图外侧的复合等于
$\text{Tr}_f \circ \alpha \circ \beta$，故本引理得证。
\end{proof}

\noindent
设 $f : X \to Y$ 为拟紧拟分离概形之间的态射，并设
$a : D_\QCoh(\mathcal{O}_Y) \to D_\QCoh(\mathcal{O}_X)$ 为引理
\ref{duality-lemma-twisted-inverse-image} 中 $Rf_*$ 的右伴随。由《范畴》，
第 \ref{categories-section-adjoint} 节，得到函子变换
$$
\eta_f : \text{id} \to  a \circ Rf_*
$$
称为该伴随的单位。

\begin{lemma}
\label{duality-lemma-unit-and-base-change}
假设有图 (\ref{duality-equation-base-change})，其中 $f$ 与 $g$ Tor 独立。
则映射 $1 \star \eta_f : L(g')^* \to L(g')^* \circ a \circ Rf_*$ 和
$\eta_{f'} \star 1 : L(g')^* \to a' \circ Rf'_* \circ L(g')^*$
通过基变换映射
$\beta : Lg^* \circ Rf_* \to Rf'_* \circ L(g')^*$
（《上同调》，注 \ref{cohomology-remark-base-change}）以及
$\alpha : L(g')^* \circ a \to a' \circ Lg^*$
（\ref{duality-equation-base-change-map}）相一致。更精确地，函子变换的图
$$
\xymatrix{
L(g')^* \ar[r]_-{1 \star \eta_f} \ar[d]_{\eta_{f'} \star 1} &
L(g')^* \circ a \circ Rf_* \ar[d]^\alpha \\
a' \circ Rf'_* \circ L(g')^* &
a' \circ Lg^* \circ Rf_* \ar[l]_-\beta
}
$$
交换。
\end{lemma}

\begin{proof}
此证明对偶于引理 \ref{duality-lemma-trace-map-and-base-change} 的证明。
本证明中，以 $f_*$ 代替 $Rf_*$，以 $g^*$ 代替 $Lg^*$，
并省略与恒等变换的 $\star$ 积；只要知道变换的源和目标，
便可确定应补入哪些项。回忆，
$\beta : g^* \circ f_* \to f'_* \circ (g')^*$ 是同构，而 $\alpha$
利用同构 $\beta^\vee : g'_* \circ a' \to a \circ g_*$ 定义；
后者是 $\beta$ 的伴随，见引理 \ref{duality-lemma-flat-precompose-pus} 及其证明。
先注意，引理所述图的左侧竖直箭头等于复合
$$
(g')^* \to (g')^* \circ g'_* \circ (g')^* \to
(g')^* \circ g'_* \circ a' \circ f'_* \circ (g')^* \to
a' \circ f'_* \circ (g')^*
$$
其中第一个箭头是 $((g')^*, g'_*)$ 的单位，第二个箭头是
$\eta_{f'}$，第三个箭头是 $((g')^*, g'_*)$ 的余单位。
这是单位与余单位的复合
$(g')^* \to (g')^* \circ (g')_* \circ (g')^* \to (g')^*$
为恒等这一事实的直接推论。考虑图
$$
\xymatrix{
& (g')^* \circ a \circ f_* \ar[r] &
(g')^* \circ a \circ g_* \circ g^* \circ f_*
\ar[ld]_\beta \\
(g')^* \ar[ru]^{\eta_f} \ar[dd]_{\eta_{f'}} \ar[rd] &
(g')^* \circ a \circ g_* \circ f'_* \circ (g')^* &
(g')^* \circ g'_* \circ a' \circ g^* \circ f_*
\ar[u]_{\beta^\vee} \ar[ld]_\beta \ar[d] \\
& (g')^* \circ g'_* \circ a' \circ f'_* \circ (g')^*
\ar[ld] \ar[u]_{\beta^\vee} &
a' \circ g^* \circ f_* \ar[lld]^\beta \\
a' \circ f'_* \circ (g')^*
}
$$
其中两个方块交换，这由《范畴》，引理
\ref{categories-lemma-properties-2-cat-cats} 得到；也可更直接地使用
《范畴》，定义 \ref{categories-definition-horizontal-composition} 之前的讨论。
三角形由上述讨论交换。由《范畴》，引理
\ref{categories-lemma-transformation-between-functors-and-adjoints} 的对偶，方块
$$
\xymatrix{
\text{id} \ar[r] \ar[d] &
g'_* \circ a' \circ g^* \circ f_* \ar[d]^\beta \\
g'_* \circ a' \circ g^* \circ f_* \ar[r]^{\beta^\vee} &
a \circ g_* \circ f'_* \circ (g')^*
}
$$
交换，从而大图中的五边形交换。由于 $\beta$ 和 $\beta^\vee$ 均为同构，
并且依定义沿大图外侧的复合等于
$\beta \circ \alpha \circ \eta_f$，故本引理得证。
\end{proof}

\begin{example}
\label{duality-example-trace-affine}
设 $A \to B$ 为环同态，$Y = \Spec(A)$、$X = \Spec(B)$，
并设 $f : X \to Y$ 为与 $A \to B$ 对应的态射。如例
\ref{duality-example-affine-twisted-inverse-image} 中所见，
$Rf_* : D_\QCoh(\mathcal{O}_X) \to D_\QCoh(\mathcal{O}_Y)$ 的右伴随
把 $D(A) = D_\QCoh(\mathcal{O}_Y)$ 的对象 $K$ 送到
$D(B) = D_\QCoh(\mathcal{O}_X)$ 中的 $R\Hom(B, K)$。迹映射为
$$
\text{Tr}_{f, K} : R\Hom(B, K) \longrightarrow R\Hom(A, K) = K
$$
它由 $A$-模映射 $A \to B$ 诱导。
\end{example}




\section{推前的右伴随与拉回}
\label{duality-section-compare-with-pullback}

\noindent
设 $f : X \to Y$ 为拟紧拟分离概形之间的态射，$a$ 为引理
\ref{duality-lemma-twisted-inverse-image} 中推前的右伴随。对
$K, L \in D_\QCoh(\mathcal{O}_Y)$，有典范映射
$$
Lf^*K \otimes^\mathbf{L}_{\mathcal{O}_X} a(L)
\longrightarrow
a(K \otimes_{\mathcal{O}_Y}^\mathbf{L} L)
$$
具体地，此映射伴随于映射
$$
Rf_*(Lf^*K \otimes^\mathbf{L}_{\mathcal{O}_X} a(L)) =
K \otimes^\mathbf{L}_{\mathcal{O}_Y} Rf_*(a(L))
\longrightarrow
K \otimes^\mathbf{L}_{\mathcal{O}_Y} L
$$
其中等式由《概形的导出范畴》，引理
\ref{perfect-lemma-cohomology-base-change}，而箭头使用迹映射
$Rf_*a(L) \to L$。取 $L = \mathcal{O}_Y$，得到映射
\begin{equation}
\label{duality-equation-compare-with-pullback}
Lf^*K \otimes^\mathbf{L}_{\mathcal{O}_X} a(\mathcal{O}_Y) \longrightarrow a(K)
\end{equation}
它关于 $K$ 函子性，并与可区别三角相容。

\begin{lemma}
\label{duality-lemma-compare-with-pullback-perfect}
设 $f : X \to Y$ 为拟紧拟分离概形之间的态射。对
$K, L \in D_\QCoh(\mathcal{O}_Y)$，若 $K$ 完美，则上面定义的映射
$Lf^*K \otimes^\mathbf{L}_{\mathcal{O}_X} a(L) \to
a(K \otimes_{\mathcal{O}_Y}^\mathbf{L} L)$ 是同构。特别地，若
$K$ 完美，则 (\ref{duality-equation-compare-with-pullback}) 是同构。
\end{lemma}

\begin{proof}
令 $K^\vee$ 为 $K$ 的“对偶”；见《上同调》，引理
\ref{cohomology-lemma-dual-perfect-complex}。对
$M \in D_\QCoh(\mathcal{O}_X)$，有
\begin{align*}
\Hom_{D(\mathcal{O}_Y)}(Rf_*M, K \otimes^\mathbf{L}_{\mathcal{O}_Y} L)
& =
\Hom_{D(\mathcal{O}_Y)}(
Rf_*M \otimes^\mathbf{L}_{\mathcal{O}_Y} K^\vee, L) \\
& =
\Hom_{D(\mathcal{O}_X)}(
M \otimes^\mathbf{L}_{\mathcal{O}_X} Lf^*K^\vee, a(L)) \\
& =
\Hom_{D(\mathcal{O}_X)}(M,
Lf^*K \otimes^\mathbf{L}_{\mathcal{O}_X} a(L))
\end{align*}
第二个等式由 $a$ 的定义和投影公式（《上同调》，引理
\ref{cohomology-lemma-projection-formula-perfect}）得到；也可使用更一般的
《概形的导出范畴》，引理 \ref{perfect-lemma-cohomology-base-change}。
于是由 Yoneda 引理得到结论。
\end{proof}

\begin{lemma}
\label{duality-lemma-restriction-compare-with-pullback}
假设有图 (\ref{duality-equation-base-change})，其中 $f$ 与 $g$ Tor 独立。
设 $K \in D_\QCoh(\mathcal{O}_Y)$。则图
$$
\xymatrix{
L(g')^*(Lf^*K \otimes^\mathbf{L}_{\mathcal{O}_X} a(\mathcal{O}_Y))
\ar[r] \ar[d] & L(g')^*a(K) \ar[d] \\
L(f')^*Lg^*K \otimes_{\mathcal{O}_{X'}}^\mathbf{L} a'(\mathcal{O}_{Y'})
\ar[r] & a'(Lg^*K)
}
$$
交换，其中水平箭头是关于 $K$ 和 $Lg^*K$ 的映射
(\ref{duality-equation-compare-with-pullback})，竖直映射则使用《上同调》，注
\ref{cohomology-remark-base-change} 和
(\ref{duality-equation-base-change-map}) 构造。
\end{lemma}

\begin{proof}
本证明中，以 $f_*$ 代替 $Rf_*$，以 $f^*$ 代替 $Lf^*$ 等，
并以 $\otimes$ 代替 $\otimes^\mathbf{L}_{\mathcal{O}_X}$ 等。
把 (\ref{duality-equation-compare-with-pullback}) 写成复合
\begin{align*}
f^*K \otimes a(\mathcal{O}_Y)
& \to
a(f_*(f^*K \otimes a(\mathcal{O}_Y))) \\
& \leftarrow
a(K \otimes f_*a(\mathcal{O}_K)) \\
& \to
a(K \otimes \mathcal{O}_Y) \\
& \to
a(K)
\end{align*}
其中第一个箭头是单位 $\eta_f$；第二个箭头是将 $a$ 应用于
《上同调》，等式 (\ref{cohomology-equation-projection-formula-map})，
它由《概形的导出范畴》，引理
\ref{perfect-lemma-cohomology-base-change} 为同构；第三个箭头是将
$a$ 应用于 $\text{id}_K \otimes \text{Tr}_f$；第四个箭头是将
$a$ 应用于同构 $K \otimes \mathcal{O}_Y = K$。只需证明这些映射
分别给出引理陈述中的交换方块。对 $\eta_f$ 和 $\text{Tr}_f$，
这分别是引理 \ref{duality-lemma-unit-and-base-change} 和
\ref{duality-lemma-trace-map-and-base-change}。对使用《上同调》，等式
(\ref{cohomology-equation-projection-formula-map}) 的箭头，
这是《上同调》，注 \ref{cohomology-remark-compatible-with-diagram}。
对乘法映射则显然。证明完成。
\end{proof}

\begin{lemma}
\label{duality-lemma-compare-on-open}
设 $f : X \to Y$ 为 Noetherian 概形之间的固有态射。设 $V \subset Y$
为开集，使得 $f^{-1}(V) \to V$ 是同构。则对
$K \in D_\QCoh^+(\mathcal{O}_Y)$，映射
(\ref{duality-equation-compare-with-pullback}) 限制到 $f^{-1}(V)$ 后为同构。
\end{lemma}

\begin{proof}
由引理 \ref{duality-lemma-proper-noetherian}，映射 (\ref{duality-equation-sheafy})
对 $D_\QCoh^+(\mathcal{O}_Y)$ 的对象为同构。因此，由引理
\ref{duality-lemma-restriction-compare-with-pullback}，关于 $K$ 的
(\ref{duality-equation-compare-with-pullback}) 到 $f^{-1}(V)$ 的限制，
就是关于 $K|_V$ 和 $f^{-1}(V) \to V$ 的映射
(\ref{duality-equation-compare-with-pullback})。故只需在 $f$ 为恒等态射时证明该映射为同构；这很清楚。
\end{proof}

\begin{lemma}
\label{duality-lemma-transitivity-compare-with-pullback}
设 $f : X \to Y$ 和 $g : Y \to Z$ 为拟紧拟分离概形的可复合态射，
并令 $h = g \circ f$。设 $a, b, c$ 分别为引理
\ref{duality-lemma-twisted-inverse-image} 对 $f, g, h$ 给出的伴随。
则对任意 $K \in D_\QCoh(\mathcal{O}_Z)$，图
$$
\xymatrix{
Lf^*(Lg^*K \otimes_{\mathcal{O}_Y}^\mathbf{L}
b(\mathcal{O}_Z)) \otimes_{\mathcal{O}_X}^\mathbf{L} a(\mathcal{O}_Y)
\ar@{=}[d] \ar[r] &
a(Lg^*K \otimes_{\mathcal{O}_Y}^\mathbf{L} b(\mathcal{O}_Z)) \ar[r] &
a(b(K)) \ar@{=}[d] \\
Lh^*K \otimes_{\mathcal{O}_X}^\mathbf{L} Lf^*b(\mathcal{O}_Z)
\otimes_{\mathcal{O}_X}^\mathbf{L} a(\mathcal{O}_Y) \ar[r] &
Lh^*K \otimes_{\mathcal{O}_X}^\mathbf{L} c(\mathcal{O}_Z) \ar[r] &
c(K)
}
$$
交换，其中各箭头为 (\ref{duality-equation-compare-with-pullback})，
并使用了 $Lh^* = Lf^* \circ Lg^*$ 和 $c = a \circ b$。
\end{lemma}

\begin{proof}
本证明中，以 $f_*$ 代替 $Rf_*$，以 $f^*$ 代替 $Lf^*$ 等，
并以 $\otimes$ 代替 $\otimes^\mathbf{L}_{\mathcal{O}_X}$ 等。
上方箭头的复合伴随于映射
$$
g_*f_*(f^*(g^*K \otimes b(\mathcal{O}_Z)) \otimes a(\mathcal{O}_Y)) \to K
$$
由《概形的导出范畴》，引理 \ref{perfect-lemma-cohomology-base-change}，
其左端等于 $K \otimes g_*f_*(f^*b(\mathcal{O}_Z) \otimes a(\mathcal{O}_Y))$。
检视各定义可知，该映射来自映射
$$
g_*f_*(f^*b(\mathcal{O}_Z) \otimes a(\mathcal{O}_Y))
\xleftarrow{g_*\epsilon}
g_*(b(\mathcal{O}_Z) \otimes f_*a(\mathcal{O}_Y)) \xrightarrow{g_*\alpha}
g_*(b(\mathcal{O}_Z)) \xrightarrow{\beta} \mathcal{O}_Z
$$
与 $\text{id}_K$ 的张量积。这里 $\epsilon$ 是《概形的导出范畴》，引理
\ref{perfect-lemma-cohomology-base-change} 中的同构，而 $\beta$
来自余单位映射 $g_*b \to \text{id}$。类似地，下方水平箭头的复合
伴随于 $\text{id}_K$ 与下列复合的张量积：
$$
g_*f_*(f^*b(\mathcal{O}_Z) \otimes a(\mathcal{O}_Y)) \xrightarrow{g_*f_*\delta}
g_*f_*(ab(\mathcal{O}_Z)) \xrightarrow{g_*\gamma}
g_*(b(\mathcal{O}_Z)) \xrightarrow{\beta}
\mathcal{O}_Z
$$
其中 $\gamma$ 来自余单位映射 $f_*a \to \text{id}$，而 $\delta$
是以复合
$$
f_*(f^*b(\mathcal{O}_Z) \otimes a(\mathcal{O}_Y))
\xleftarrow{\epsilon}
b(\mathcal{O}_Z) \otimes f_*a(\mathcal{O}_Y) \xrightarrow{\alpha}
b(\mathcal{O}_Z)
$$
为伴随的映射。由伴随函子、伴随映射和余单位的一般性质
（见《范畴》，第 \ref{categories-section-adjoint} 节），
如所需有 $\gamma \circ f_*\delta = \alpha \circ \epsilon^{-1}$。
\end{proof}





\section{闭浸入的推前右伴随}
\label{duality-section-sections-with-exact-support}

\noindent
设 $i : (Z, \mathcal{O}_Z) \to (X, \mathcal{O}_X)$ 为赋环空间的态射，
使得 $i$ 同胚地映到一个闭子集上，并且
$i^\sharp : \mathcal{O}_X \to i_*\mathcal{O}_Z$ 满射
（例如概形的闭浸入）。令 $\mathcal{I} = \Ker(i^\sharp)$。
对 $\mathcal{O}_X$-模层 $\mathcal{F}$，层
$$
\SheafHom_{\mathcal{O}_X}(i_*\mathcal{O}_Z, \mathcal{F})
$$
是被 $\mathcal{I}$ 零化的 $\mathcal{O}_X$-模层。因此，由《模》，引理
\ref{modules-lemma-i-star-equivalence}，存在一个 $\mathcal{O}_Z$-模层，
记作 $\SheafHom(\mathcal{O}_Z, \mathcal{F})$，使得作为
$\mathcal{O}_X$-模有
$$
i_*\SheafHom(\mathcal{O}_Z, \mathcal{F}) =
\SheafHom_{\mathcal{O}_X}(i_*\mathcal{O}_Z, \mathcal{F})
$$
下面详述其含义。

\begin{lemma}
\label{duality-lemma-compute-sheaf-with-exact-support}
采用上述记号。函子 $\SheafHom(\mathcal{O}_Z, -)$ 是函子
$i_* : \textit{Mod}(\mathcal{O}_Z) \to \textit{Mod}(\mathcal{O}_X)$
的右伴随。对开集 $V \subset Z$，有
$$
\Gamma(V, \SheafHom(\mathcal{O}_Z, \mathcal{F})) =
\{s \in \Gamma(U, \mathcal{F}) \mid \mathcal{I}s = 0\}
$$
其中 $U \subset X$ 为与 $Z$ 的交等于 $V$ 的开集。
\end{lemma}

\begin{proof}
设 $\mathcal{G}$ 为 $\mathcal{O}_Z$-模层。则
$$
\Hom_{\mathcal{O}_X}(i_*\mathcal{G}, \mathcal{F}) =
\Hom_{i_*\mathcal{O}_Z}(i_*\mathcal{G},
\SheafHom_{\mathcal{O}_X}(i_*\mathcal{O}_Z, \mathcal{F})) =
\Hom_{\mathcal{O}_Z}(\mathcal{G}, \SheafHom(\mathcal{O}_Z, \mathcal{F}))
$$
第一个等式由《模》，引理 \ref{modules-lemma-adjoint-tensor-restrict}，
第二个等式由 $i_*$ 的全忠实性；见《模》，引理
\ref{modules-lemma-i-star-equivalence}。截面的描述留给读者验证。
\end{proof}

\noindent
函子
$$
\textit{Mod}(\mathcal{O}_X)
\longrightarrow
\textit{Mod}(\mathcal{O}_Z),
\quad
\mathcal{F} \longmapsto \SheafHom(\mathcal{O}_Z, \mathcal{F})
$$
左正合，并有导出延拓
$$
R\SheafHom(\mathcal{O}_Z, -) : D(\mathcal{O}_X) \to D(\mathcal{O}_Z).
$$

\begin{lemma}
\label{duality-lemma-sheaf-with-exact-support-adjoint}
采用上述记号。函子 $R\SheafHom(\mathcal{O}_Z, -)$ 是函子
$Ri_* : D(\mathcal{O}_Z) \to D(\mathcal{O}_X)$ 的右伴随。
\end{lemma}

\begin{proof}
由引理 \ref{duality-lemma-compute-sheaf-with-exact-support}，$i_*$ 与
$\SheafHom(\mathcal{O}_Z, -)$ 是伴随函子，故结论成立；见《导出范畴》，
引理 \ref{derived-lemma-derived-adjoint-functors}。
\end{proof}

\begin{lemma}
\label{duality-lemma-sheaf-with-exact-support-ext}
采用上述记号。对 $D(\mathcal{O}_X)$ 中所有 $K$，在
$D(\mathcal{O}_X)$ 中有
$$
Ri_*R\SheafHom(\mathcal{O}_Z, K) =
R\SheafHom_{\mathcal{O}_X}(i_*\mathcal{O}_Z, K)
$$
\end{lemma}

\begin{proof}
由函子 $R\SheafHom(\mathcal{O}_Z, -)$ 的构造立即得到。
\end{proof}

\begin{lemma}
\label{duality-lemma-sheaf-with-exact-support-internal-home}
采用上述记号。对 $M \in D(\mathcal{O}_Z)$ 及
$D(\mathcal{O}_X)$ 中所有 $K$，在 $D(\mathcal{O}_Z)$ 中有
$$
R\SheafHom_{\mathcal{O}_X}(Ri_*M, K) =
Ri_*R\SheafHom_{\mathcal{O}_Z}(M, R\SheafHom(\mathcal{O}_Z, K))
$$
\end{lemma}

\begin{proof}
这由函子 $R\SheafHom(\mathcal{O}_Z, -)$ 的构造以及下述事实立即得到：
若 $\mathcal{K}^\bullet$ 是 $\mathcal{O}_X$-模的 K-内射复形，
则 $\SheafHom(\mathcal{O}_Z, \mathcal{K}^\bullet)$ 是
$\mathcal{O}_Z$-模的 K-内射复形；见《导出范畴》，引理
\ref{derived-lemma-adjoint-preserve-K-injectives}。
\end{proof}

\begin{lemma}
\label{duality-lemma-sheaf-with-exact-support-quasi-coherent}
设 $i : Z \to X$ 为概形的伪凝聚闭浸入
（若 $X$ 局部 Noetherian，则可取任意闭浸入）。则
\begin{enumerate}
\item $R\SheafHom(\mathcal{O}_Z, -)$ 把 $D^+_\QCoh(\mathcal{O}_X)$
映入 $D^+_\QCoh(\mathcal{O}_Z)$；并且
\item 若 $X = \Spec(A)$ 且 $Z = \Spec(B)$，则图
$$
\xymatrix{
D^+(B) \ar[r] & D_\QCoh^+(\mathcal{O}_Z) \\
D^+(A) \ar[r] \ar[u]^{R\Hom(B, -)} &
D_\QCoh^+(\mathcal{O}_X) \ar[u]_{R\SheafHom(\mathcal{O}_Z, -)}
}
$$
交换。
\end{enumerate}
\end{lemma}

\begin{proof}
先解释括号中的说明。若 $X$ 局部 Noetherian，则由《态射进阶》，引理
\ref{more-morphisms-lemma-Noetherian-pseudo-coherent}，$i$ 是伪凝聚的。

\medskip\noindent
设 $K$ 为 $D^+_\QCoh(\mathcal{O}_X)$ 的对象。为证明 (1)，由《态射》，
引理 \ref{morphisms-lemma-i-star-equivalence}，只需证明对
$H^n(R\SheafHom(\mathcal{O}_Z, K))$ 应用 $i_*$ 后，得到 $X$
上的拟凝聚模。由引理 \ref{duality-lemma-sheaf-with-exact-support-ext}，
这等价于证明 $R\SheafHom_{\mathcal{O}_X}(i_*\mathcal{O}_Z, K)$
属于 $D_\QCoh(\mathcal{O}_X)$。由于 $i$ 伪凝聚，层
$\mathcal{O}_Z$ 是伪凝聚 $\mathcal{O}_X$-模。因此结论由
《概形的导出范畴》，引理
\ref{perfect-lemma-quasi-coherence-internal-hom} 得到。

\medskip\noindent
如 (2) 中假设 $X = \Spec(A)$、$Z = \Spec(B)$。设 $I^\bullet$
为表示 $D^+(A)$ 中对象 $K$ 的下有界内射 $A$-模复形。
则作为 $B$-模复形，有 $R\Hom(B, K) = \Hom_A(B, I^\bullet)$。
选取拟同构
$$
\widetilde{I^\bullet} \longrightarrow \mathcal{I}^\bullet
$$
其中 $\mathcal{I}^\bullet$ 是下有界内射 $\mathcal{O}_X$-模复形。
由引理 \ref{duality-lemma-compute-sheaf-with-exact-support} 中函子
$\SheafHom(\mathcal{O}_Z, -)$ 的描述，存在映射
$$
\Hom_A(B, I^\bullet)
\longrightarrow
\Gamma(Z, \SheafHom(\mathcal{O}_Z, \mathcal{I}^\bullet))
$$
注意 $\SheafHom(\mathcal{O}_Z, \mathcal{I}^\bullet)$ 表示
$R\SheafHom(\mathcal{O}_Z, \widetilde{K})$。应用
$\widetilde{\ }$ 函子的泛性质，在 $D(\mathcal{O}_Z)$ 中得到映射
$$
\widetilde{\Hom_A(B, I^\bullet)}
\longrightarrow
R\SheafHom(\mathcal{O}_Z, \widetilde{K})
$$
可以在应用 $i_*$ 后检验此映射是 $D(\mathcal{O}_Z)$ 中的同构。
而一旦应用 $i_*$，便得到《概形的导出范畴》，引理
\ref{perfect-lemma-quasi-coherence-internal-hom} 中的同构；这里使用了
引理 \ref{duality-lemma-sheaf-with-exact-support-ext} 的识别。
\end{proof}

\begin{lemma}
\label{duality-lemma-sheaf-with-exact-support-coherent}
设 $i : Z \to X$ 为概形的闭浸入，并假设 $X$ 局部 Noetherian。
则 $R\SheafHom(\mathcal{O}_Z, -)$ 把
$D^+_{\textit{Coh}}(\mathcal{O}_X)$ 映入
$D^+_{\textit{Coh}}(\mathcal{O}_Z)$。
\end{lemma}

\begin{proof}
问题关于 $X$ 是局部的，故可假设 $X$ 仿射。记
$X = \Spec(A)$、$Z = \Spec(B)$，其中 $A$ Noetherian，且
$A \to B$ 满射。此时可应用引理
\ref{duality-lemma-sheaf-with-exact-support-quasi-coherent} 把问题转化为代数问题。
相应的代数结论是《对偶化复形》，引理
\ref{dualizing-lemma-exact-support-coherent} 的推论。
\end{proof}

\begin{lemma}
\label{duality-lemma-twisted-inverse-image-closed}
设 $X$ 为拟紧拟分离概形，$i : Z \to X$ 为伪凝聚闭浸入
（若 $X$ 为 Noetherian 概形，则任意闭浸入均为伪凝聚的）。
设 $a : D_\QCoh(\mathcal{O}_X) \to D_\QCoh(\mathcal{O}_Z)$ 为
$Ri_*$ 的右伴随。则对 $K \in D_\QCoh^+(\mathcal{O}_X)$，有函子性同构
$$
a(K) = R\SheafHom(\mathcal{O}_Z, K)
$$
\end{lemma}

\begin{proof}
括号中的陈述由《态射进阶》，引理
\ref{more-morphisms-lemma-Noetherian-pseudo-coherent} 得到。
由引理 \ref{duality-lemma-sheaf-with-exact-support-adjoint}，函子
$R\SheafHom(\mathcal{O}_Z, -)$ 是
$Ri_* : D(\mathcal{O}_Z) \to D(\mathcal{O}_X)$ 的右伴随。
此外，由引理 \ref{duality-lemma-sheaf-with-exact-support-quasi-coherent}
和引理 \ref{duality-lemma-twisted-inverse-image-bounded-below}，
$R\SheafHom(\mathcal{O}_Z, -)$ 与 $a$ 都把
$D_\QCoh^+(\mathcal{O}_X)$ 映入 $D_\QCoh^+(\mathcal{O}_Z)$。
因此，由伴随函子的唯一性得到该同构。
\end{proof}

\begin{example}
\label{duality-example-trace-closed-immersion}
若 $i : Z \to X$ 是 Noetherian 概形的闭浸入，则对
$K \in D_\QCoh^+(\mathcal{O}_X)$，图
$$
\xymatrix{
i_*a(K) \ar[rr]_-{\text{Tr}_{i, K}} \ar@{=}[d] & &
K \ar@{=}[d] \\
i_*R\SheafHom(\mathcal{O}_Z, K) \ar@{=}[r] &
R\SheafHom_{\mathcal{O}_X}(i_*\mathcal{O}_Z, K)
\ar[r] & K
}
$$
交换。这里水平等号由引理 \ref{duality-lemma-sheaf-with-exact-support-ext}，
下方水平箭头由映射 $\mathcal{O}_X \to i_*\mathcal{O}_Z$ 诱导。
图的交换性是引理 \ref{duality-lemma-twisted-inverse-image-closed} 的推论。
\end{example}




\section{闭浸入的推前右伴随与基变换}
\label{duality-section-sections-with-exact-support-base-change}

\noindent
考虑概形的笛卡尔图
$$
\xymatrix{
Z' \ar[r]_{i'} \ar[d]_g & X' \ar[d]^f \\
Z \ar[r]^i & X
}
$$
其中 $i$ 为闭浸入。若 $Z$ 与 $X'$ 在 $X$ 上 Tor 独立，则存在典范基变换映射
\begin{equation}
\label{duality-equation-base-change-exact-support}
Lg^*R\SheafHom(\mathcal{O}_Z, K)
\longrightarrow
R\SheafHom(\mathcal{O}_{Z'}, Lf^*K)
\end{equation}
它在 $D(\mathcal{O}_{Z'})$ 中关于 $D(\mathcal{O}_X)$ 的 $K$ 函子性。
具体地，由引理 \ref{duality-lemma-sheaf-with-exact-support-adjoint} 的伴随性，
这样的箭头等价于 $D(\mathcal{O}_{X'})$ 中的映射
$$
Ri'_*Lg^*R\SheafHom(\mathcal{O}_Z, K)
\longrightarrow
Lf^*K
$$
由 Tor 独立性，有 $Ri'_* \circ Lg^* = Lf^* \circ Ri_*$；
见《概形的导出范畴》，引理
\ref{perfect-lemma-compare-base-change-closed-immersion}。
因此，这又等价于映射
$$
Lf^*Ri_*R\SheafHom(\mathcal{O}_Z, K)
\longrightarrow
Lf^*K
$$
可取 $Lf^*(can)$，其中
$can : Ri_* R\SheafHom(\mathcal{O}_Z, K) \to K$ 是该伴随的余单位。

\begin{lemma}
\label{duality-lemma-check-base-change-is-iso}
在上述情形下，映射 (\ref{duality-equation-base-change-exact-support}) 是同构，
当且仅当《上同调》，注
\ref{cohomology-remark-prepare-fancy-base-change} 中的基变换映射
$$
Lf^*R\SheafHom_{\mathcal{O}_X}(\mathcal{O}_Z, K)
\longrightarrow
R\SheafHom_{\mathcal{O}_{X'}}(\mathcal{O}_{Z'}, Lf^*K)
$$
是同构。
\end{lemma}

\begin{proof}
由所假设的 Tor 独立性，$\mathcal{O}_{Z'} = Lf^*\mathcal{O}_Z$，
故此陈述有意义。由于 $i'_*$ 正合且忠实，只需证明在应用 $Ri'_*$ 后，
映射 (\ref{duality-equation-base-change-exact-support}) 是同构。由所假设的 Tor 独立性及
《概形的导出范畴》，引理
\ref{perfect-lemma-compare-base-change-closed-immersion}，有
$Ri'_* \circ Lg^* = Lf^* \circ Ri_*$，从而得到映射
$$
Lf^*Ri_*R\SheafHom(\mathcal{O}_Z, K)
\longrightarrow
Ri'_*R\SheafHom(\mathcal{O}_{Z'}, Lf^*K)
$$
由引理 \ref{duality-lemma-sheaf-with-exact-support-ext}，其源和目标正如本引理陈述中所示。
这里略去验证它与《上同调》，注
\ref{cohomology-remark-prepare-fancy-base-change} 中构造的是同一个映射。
\end{proof}

\begin{lemma}
\label{duality-lemma-flat-bc-sheaf-with-exact-support}
在上述情形中，假设 $f$ 平坦且 $i$ 伪凝聚。则对
$D^+_\QCoh(\mathcal{O}_X)$ 中的 $K$，
(\ref{duality-equation-base-change-exact-support}) 是同构。
\end{lemma}

\begin{proof}
第一种证明。为证明此映射是同构，可以局部地工作。因此可假设
$X$、$X'$、$Z$、$Z'$ 均为仿射概形，分别对应于环
$A$、$A'$、$B$、$B'$。于是 $B$ 与 $A'$ 在 $A$ 上 Tor 独立。
由引理 \ref{duality-lemma-check-base-change-is-iso}，只需验证
$$
R\Hom_A(B, K) \otimes_A^\mathbf{L} A' =
R\Hom_{A'}(B', K \otimes_A^\mathbf{L} A')
$$
对所有 $K \in D^+(A)$ 在 $D(A')$ 中成立。这里使用《概形的导出范畴》，
引理 \ref{perfect-lemma-quasi-coherence-internal-hom}，以及
$B$（相应地\ $B'$）作为 $A$-模（相应地\ $A'$-模）是伪凝聚的事实，
以便在环和概形两个层面比较导出 Hom。上式由《代数进阶》，引理
\ref{more-algebra-lemma-internal-hom-evaluate-tensor-isomorphism}
第 (3) 部分得到。另见《对偶化复形》，节
\ref{dualizing-section-base-change-trivial-duality} 中的讨论。

\medskip\noindent
第二种证明\footnote{此证明表明，只需假设 $K$ 属于
$D^+(\mathcal{O}_X)$。}。取 $z' \in Z'$，其像为 $z \in Z$。
先证明 (\ref{duality-equation-base-change-exact-support}) 在 $z'$ 处的茎上诱导映射
$$
R\Hom(\mathcal{O}_{Z, z}, K_z)
\otimes_{\mathcal{O}_{Z, x}}^\mathbf{L} \mathcal{O}_{Z', z'}
\longrightarrow
R\Hom(\mathcal{O}_{Z', z'},
K_z \otimes_{\mathcal{O}_{X, z}}^\mathbf{L} \mathcal{O}_{X', z'})
$$
即《对偶化复形》，方程 (\ref{dualizing-equation-base-change}) 中的映射。
事实上，这些映射的构造完全相同。随后应用《对偶化复形》，引理
\ref{dualizing-lemma-flat-bc-surjection}。
\end{proof}

\begin{lemma}
\label{duality-lemma-sheaf-with-exact-support-tensor}
设 $i : Z \to X$ 为概形的伪凝聚闭浸入。设
$M \in D_\QCoh(\mathcal{O}_X)$ 局部具有位于 $[a, \infty)$ 中的
Tor 振幅，并设 $K \in D_\QCoh^+(\mathcal{O}_X)$。则在
$D(\mathcal{O}_Z)$ 中存在典范同构
$$
R\SheafHom(\mathcal{O}_Z, K) \otimes_{\mathcal{O}_Z}^\mathbf{L} Li^*M =
R\SheafHom(\mathcal{O}_Z, K \otimes_{\mathcal{O}_X}^\mathbf{L} M)
$$
\end{lemma}

\begin{proof}
由引理 \ref{duality-lemma-sheaf-with-exact-support-adjoint} 和
\ref{duality-lemma-sheaf-with-exact-support-ext}，从左边到右边的映射等价于映射
$$
Ri_*R\SheafHom(\mathcal{O}_Z, K) \otimes_{\mathcal{O}_X}^\mathbf{L} M
\longrightarrow
K \otimes_{\mathcal{O}_X}^\mathbf{L} M
$$
对此映射，取余单位 $Ri_*R\SheafHom(\mathcal{O}_Z, K) \to K$
与 $\text{id}_M$ 的张量积。为证明在所给假设下此映射是同构，
可用引理 \ref{duality-lemma-sheaf-with-exact-support-quasi-coherent} 将问题转化为代数，
再例如应用《代数进阶》，引理
\ref{more-algebra-lemma-internal-hom-evaluate-tensor-isomorphism} 第 (3) 部分。
也可以不使用引理 \ref{duality-lemma-sheaf-with-exact-support-quasi-coherent}，
而像引理 \ref{duality-lemma-flat-bc-sheaf-with-exact-support} 的第二种证明那样考察茎。
\end{proof}



\section{有限态射的推前右伴随}
\label{duality-section-duality-finite}

\noindent
若 $i : Z \to X$ 为概形的闭浸入，则函子
$i_* : \textit{Mod}(\mathcal{O}_Z) \to \textit{Mod}(\mathcal{O}_X)$
有右伴随 $\SheafHom(\mathcal{O}_Z, -)$，其导出延拓
$R\SheafHom(\mathcal{O}_Z, -)$ 是
$Ri_* : D(\mathcal{O}_Z) \to D(\mathcal{O}_X)$ 的右伴随；见节
\ref{duality-section-sections-with-exact-support}。对于有限态射
$f : Y \to X$，这种策略不可行，因为函子
$f_* : \textit{Mod}(\mathcal{O}_Y) \to \textit{Mod}(\mathcal{O}_X)$
一般并不正合，因而没有右伴随。替代办法是考虑正合函子
$\textit{Mod}(f_*\mathcal{O}_Y) \to \textit{Mod}(\mathcal{O}_X)$，
再考虑相应的右伴随及其导出延拓。

\medskip\noindent
设 $f : Y \to X$ 为概形的仿射态射。对一个
$\mathcal{O}_X$-模层 $\mathcal{F}$，层
$$
\SheafHom_{\mathcal{O}_X}(f_*\mathcal{O}_Y, \mathcal{F})
$$
是一个 $f_*\mathcal{O}_Y$-模层。由此得到函子
$\textit{Mod}(\mathcal{O}_X) \to \textit{Mod}(f_*\mathcal{O}_Y)$，
记作 $\SheafHom(f_*\mathcal{O}_Y, -)$。

\begin{lemma}
\label{duality-lemma-compute-sheafhom-affine}
采用上述记号。函子 $\SheafHom(f_*\mathcal{O}_Y, -)$ 是限制函子
$\textit{Mod}(f_*\mathcal{O}_Y) \to \textit{Mod}(\mathcal{O}_X)$
的右伴随。对仿射开集 $U \subset X$，有
$$
\Gamma(U, \SheafHom(f_*\mathcal{O}_Y, \mathcal{F})) =
\Hom_A(B, \mathcal{F}(U))
$$
其中 $A = \mathcal{O}_X(U)$，且 $B = \mathcal{O}_Y(f^{-1}(U))$。
\end{lemma}

\begin{proof}
伴随性由《模》，引理 \ref{modules-lemma-adjoint-tensor-restrict} 得到。
由于 $f$ 仿射，$f_*\mathcal{O}_Y$ 是由视为 $A$-模的 $B$
所对应的拟凝聚层。因此，$U$ 上截面的描述由《概形》，引理
\ref{schemes-lemma-compare-constructions} 得到。
\end{proof}

\noindent
函子 $\SheafHom(f_*\mathcal{O}_Y, -)$ 左正合。令
$$
R\SheafHom(f_*\mathcal{O}_Y, -) :
D(\mathcal{O}_X)
\longrightarrow
D(f_*\mathcal{O}_Y)
$$
为其导出延拓。

\begin{lemma}
\label{duality-lemma-sheafhom-affine-adjoint}
采用上述记号。函子 $R\SheafHom(f_*\mathcal{O}_Y, -)$ 是函子
$D(f_*\mathcal{O}_Y) \to D(\mathcal{O}_X)$ 的右伴随。
\end{lemma}

\begin{proof}
由引理 \ref{duality-lemma-compute-sheafhom-affine} 和《导出范畴》，引理
\ref{derived-lemma-derived-adjoint-functors} 得到。
\end{proof}

\begin{lemma}
\label{duality-lemma-sheafhom-affine-ext}
采用上述记号。复合函子
$$
D(\mathcal{O}_X) \xrightarrow{R\SheafHom(f_*\mathcal{O}_Y, -)}
D(f_*\mathcal{O}_Y) \to D(\mathcal{O}_X)
$$
即函子 $K \mapsto R\SheafHom_{\mathcal{O}_X}(f_*\mathcal{O}_Y, K)$。
\end{lemma}

\begin{proof}
这由构造立即得到。
\end{proof}

\begin{lemma}
\label{duality-lemma-finite-twisted}
设 $f : Y \to X$ 为概形的有限伪凝聚态射
（Noetherian 概形之间的有限态射是伪凝聚的）。函子
$R\SheafHom(f_*\mathcal{O}_Y, -)$ 把 $D_\QCoh^+(\mathcal{O}_X)$
映入 $D_\QCoh^+(f_*\mathcal{O}_Y)$。若 $X$ 拟紧且拟分离，则图
$$
\xymatrix{
D_\QCoh^+(\mathcal{O}_X) \ar[rr]_a \ar[rd]_{R\SheafHom(f_*\mathcal{O}_Y, -)}
& & D_\QCoh^+(\mathcal{O}_Y) \ar[ld]^\Phi \\
& D_\QCoh^+(f_*\mathcal{O}_Y)
}
$$
交换，其中 $a$ 是引理 \ref{duality-lemma-twisted-inverse-image} 中对应于
$f$ 的右伴随，而 $\Phi$ 是《概形的导出范畴》，引理
\ref{perfect-lemma-affine-morphism-equivalence} 中的等价。
\end{lemma}

\begin{proof}
括号中的说明由《态射进阶》，引理
\ref{more-morphisms-lemma-Noetherian-pseudo-coherent} 得到。
由于 $f$ 伪凝聚，$\mathcal{O}_X$-模 $f_*\mathcal{O}_Y$ 伪凝聚；
见《态射进阶》，引理
\ref{more-morphisms-lemma-finite-pseudo-coherent}。因此
$R\SheafHom(f_*\mathcal{O}_Y, -)$ 把 $D_\QCoh^+(\mathcal{O}_X)$
映入 $D_\QCoh^+(f_*\mathcal{O}_Y)$；见《概形的导出范畴》，引理
\ref{perfect-lemma-quasi-coherence-internal-hom}。
于是 $\Phi \circ a$ 与 $R\SheafHom(f_*\mathcal{O}_Y, -)$ 在
$D_\QCoh^+(\mathcal{O}_X)$ 上相同，因为这两个函子都是限制函子
$D_\QCoh^+(f_*\mathcal{O}_Y) \to D_\QCoh^+(\mathcal{O}_X)$ 的右伴随。
为此使用引理 \ref{duality-lemma-twisted-inverse-image-bounded-below} 和
\ref{duality-lemma-sheafhom-affine-adjoint}。
\end{proof}

\begin{remark}
\label{duality-remark-trace-map-finite}
若 $f : Y \to X$ 是 Noetherian 概形之间的有限态射，则对
$K \in D_\QCoh^+(\mathcal{O}_X)$，图
$$
\xymatrix{
Rf_*a(K) \ar[r]_-{\text{Tr}_{f, K}} \ar@{=}[d] & K \ar@{=}[d] \\
R\SheafHom_{\mathcal{O}_X}(f_*\mathcal{O}_Y, K) \ar[r] & K
}
$$
交换。这由引理 \ref{duality-lemma-finite-twisted} 得到。下方水平箭头由映射
$\mathcal{O}_X \to f_*\mathcal{O}_Y$ 诱导，上方水平箭头是节
\ref{duality-section-trace} 中讨论的迹映射。
\end{remark}






\section{真平坦态射的推前右伴随}
\label{duality-section-proper-flat}

\noindent
对于拟紧拟分离概形之间的真、平坦且有限表示的态射，
推前的右伴随具有若干显著性质。

\begin{lemma}
\label{duality-lemma-proper-flat}
设 $Y$ 为拟紧拟分离概形，$f : X \to Y$ 为真、平坦且有限表示的概形态射。
设 $a$ 为引理 \ref{duality-lemma-twisted-inverse-image} 中
$Rf_* : D_\QCoh(\mathcal{O}_X) \to D_\QCoh(\mathcal{O}_Y)$ 的右伴随。
则 $a$ 与直和交换。
\end{lemma}

\begin{proof}
设 $P$ 为 $D(\mathcal{O}_X)$ 的完美对象。由《概形的导出范畴》，引理
\ref{perfect-lemma-flat-proper-perfect-direct-image-general}，
复形 $Rf_*P$ 在 $Y$ 上完美。设 $K_i$ 为
$D_\QCoh(\mathcal{O}_Y)$ 中的一族对象。于是
\begin{align*}
\Hom_{D(\mathcal{O}_X)}(P, a(\bigoplus K_i))
& =
\Hom_{D(\mathcal{O}_Y)}(Rf_*P, \bigoplus K_i) \\
& =
\bigoplus \Hom_{D(\mathcal{O}_Y)}(Rf_*P, K_i) \\
& =
\bigoplus \Hom_{D(\mathcal{O}_X)}(P, a(K_i))
\end{align*}
因为完美对象是紧对象（《概形的导出范畴》，命题
\ref{perfect-proposition-compact-is-perfect}）。由于
$D_\QCoh(\mathcal{O}_X)$ 有完美生成元（《概形的导出范畴》，定理
\ref{perfect-theorem-bondal-van-den-Bergh}），可知映射
$\bigoplus a(K_i) \to a(\bigoplus K_i)$ 是同构，即 $a$ 与直和交换。
\end{proof}

\begin{lemma}
\label{duality-lemma-proper-flat-relative}
设 $Y$ 为拟紧拟分离概形，$f : X \to Y$ 为真、平坦且有限表示的概形态射。
设 $a$ 为引理 \ref{duality-lemma-twisted-inverse-image} 中
$Rf_* : D_\QCoh(\mathcal{O}_X) \to D_\QCoh(\mathcal{O}_Y)$ 的右伴随。则
\begin{enumerate}
\item 对每个闭集 $T \subset Y$，若 $Q \in D_\QCoh(Y)$ 支撑于 $T$，
则 $a(Q)$ 支撑于 $f^{-1}(T)$；
\item 对每个拟紧开集 $V \subset Y$ 及任意
$K \in D_\QCoh(\mathcal{O}_Y)$，映射 (\ref{duality-equation-sheafy}) 是同构。
\end{enumerate}
\end{lemma}

\begin{proof}
这由引理 \ref{duality-lemma-when-sheafy}、\ref{duality-lemma-proper-noetherian}
和 \ref{duality-lemma-proper-flat} 得到。严格说来，(1) 首先只对补集为拟紧开集的
$T$ 成立；但一般情形也随之成立，因为任意闭子集都是此类闭集的交。
\end{proof}

\begin{lemma}
\label{duality-lemma-compare-with-pullback-flat-proper}
设 $Y$ 为拟紧拟分离概形，$f : X \to Y$ 为真、平坦且有限表示的概形态射。
则对 $D_\QCoh(\mathcal{O}_Y)$ 的每个对象 $K$，映射
(\ref{duality-equation-compare-with-pullback}) 都是同构。
\end{lemma}

\begin{proof}
由引理 \ref{duality-lemma-proper-flat}，$a$ 与直和交换。因此，
$D_\QCoh(\mathcal{O}_Y)$ 中使 (\ref{duality-equation-compare-with-pullback})
成为同构的对象构成 $D_\QCoh(\mathcal{O}_Y)$ 的严格满、饱和三角子范畴，
且该子范畴在取直和下封闭。由于 $D_\QCoh(\mathcal{O}_Y)$ 是由单个完美对象
生成的模范畴（分别见《概形的导出范畴》，定理
\ref{perfect-theorem-DQCoh-is-Ddga} 和定理
\ref{perfect-theorem-bondal-van-den-Bergh}），可如《代数进阶》，注
\ref{more-algebra-remark-P-resolution} 那样论证，说明只需对一个完美对象证明
(\ref{duality-equation-compare-with-pullback}) 是同构。但该结论对完美对象成立；见引理
\ref{duality-lemma-compare-with-pullback-perfect}。
\end{proof}

\noindent
下述引理表明，对有限表示的真平坦态射，基变换映射
(\ref{duality-equation-base-change-map}) 是同构。例
\ref{duality-example-base-change-wrong} 将说明，对完美真态射，这一结论不再总是成立；
此时必须加入 Tor 独立性条件。

\begin{lemma}
\label{duality-lemma-proper-flat-base-change}
设 $g : Y' \to Y$ 为拟紧拟分离概形之间的态射，
$f : X \to Y$ 为有限表示的真平坦态射。则对所有
$K \in D_\QCoh(\mathcal{O}_Y)$，基变换映射
(\ref{duality-equation-base-change-map}) 是同构。
\end{lemma}

\begin{proof}
由引理 \ref{duality-lemma-proper-flat-relative}，函子 $a$ 和 $a'$ 的构造
与到 $Y$ 和 $Y'$ 的开集的限制交换。因此可假设
$Y' \to Y$ 是仿射概形之间的态射；见注
\ref{duality-remark-check-over-affines}。此时结论由引理
\ref{duality-lemma-more-base-change} 得到。
\end{proof}

\begin{remark}
\label{duality-remark-relative-dualizing-complex}
设 $Y$ 为拟紧拟分离概形，$f : X \to Y$ 为有限表示的真平坦态射。
设 $a$ 为引理 \ref{duality-lemma-twisted-inverse-image} 中对应于 $f$ 的伴随。
在此情形下，$\omega_{X/Y}^\bullet = a(\mathcal{O}_Y)$ 有时称为
{\it 相对对偶化复形}。由引理
\ref{duality-lemma-compare-with-pullback-flat-proper}，对
$K \in D_\QCoh(\mathcal{O}_Y)$ 有函子性同构
$a(K) = Lf^*K \otimes_{\mathcal{O}_X}^\mathbf{L} \omega_{X/Y}^\bullet$。
此外，节 \ref{duality-section-trace} 中的迹映射
$$
\text{Tr}_{f, \mathcal{O}_Y} : Rf_*\omega_{X/Y}^\bullet \to \mathcal{O}_Y
$$
诱导 $D_\QCoh(\mathcal{O}_Y)$ 中所有 $K$ 的迹映射。更确切地，图
$$
\xymatrix{
Rf_*a(K) \ar[rrr]_{\text{Tr}_{f, K}} \ar@{=}[d] & & &
K \ar@{=}[d] \\
Rf_*(Lf^*K \otimes_{\mathcal{O}_X}^\mathbf{L} \omega_{X/Y}^\bullet)
\ar@{=}[r] &
K \otimes_{\mathcal{O}_Y}^\mathbf{L} Rf_*\omega_{X/Y}^\bullet
\ar[rr]^-{\text{id}_K \otimes \text{Tr}_{f, \mathcal{O}_Y}} & & K
}
$$
交换，其中下方右侧的等式由《概形的导出范畴》，引理
\ref{perfect-lemma-cohomology-base-change} 给出。若
$g : Y' \to Y$ 是拟紧拟分离概形之间的态射，且
$X' = Y' \times_Y X$，则由引理 \ref{duality-lemma-proper-flat-base-change}，有
$\omega_{X'/Y'}^\bullet = L(g')^*\omega_{X/Y}^\bullet$，其中
$g' : X' \to X$ 为投影；并且由引理
\ref{duality-lemma-trace-map-and-base-change}，对 $f' : X' \to Y'$ 的迹映射
$$
\text{Tr}_{f', \mathcal{O}_{Y'}} :
Rf'_*\omega_{X'/Y'}^\bullet \to \mathcal{O}_{Y'}
$$
是在基变换同构下由 $\text{Tr}_{f, \mathcal{O}_Y}$ 基变换得到的。
\end{remark}

\begin{remark}
\label{duality-remark-relative-dualizing-complex-relative-cup-product}
设 $f : X \to Y$、$\omega^\bullet_{X/Y}$ 和
$\text{Tr}_{f, \mathcal{O}_Y}$ 如注
\ref{duality-remark-relative-dualizing-complex} 中所述。设 $K$ 和 $M$ 属于
$D_\QCoh(\mathcal{O}_X)$，且 $M$ 伪凝聚（例如完美）。假设给定映射
$K \otimes_{\mathcal{O}_X}^\mathbf{L} M \to \omega^\bullet_{X/Y}$，
它通过《上同调》，方程 (\ref{cohomology-equation-internal-hom}) 对应于同构
$K \to R\SheafHom_{\mathcal{O}_X}(M, \omega^\bullet_{X/Y})$。
则相对杯积（《上同调》，注 \ref{cohomology-remark-cup-product}）
$$
Rf_*K \otimes_{\mathcal{O}_Y}^\mathbf{L} Rf_*M
\to
Rf_*(K \otimes_{\mathcal{O}_X}^\mathbf{L} M)
\to
Rf_*\omega^\bullet_{X/Y}
\xrightarrow{\text{Tr}_{f, \mathcal{O}_Y}}
\mathcal{O}_Y
$$
确定一个同构
$Rf_*K \to R\SheafHom_{\mathcal{O}_Y}(Rf_*M, \mathcal{O}_Y)$。
事实上，由于 $\omega^\bullet_{X/Y} = a(\mathcal{O}_Y)$，典范映射
(\ref{duality-equation-sheafy-trace})
$$
Rf_*R\SheafHom_{\mathcal{O}_X}(M, \omega^\bullet_{X/Y}) \to
R\SheafHom_{\mathcal{O}_Y}(Rf_*M, \mathcal{O}_Y)
$$
由引理 \ref{duality-lemma-iso-on-RSheafHom}、注
\ref{duality-remark-iso-on-RSheafHom} 以及 $M$ 与 $Rf_*M$ 伪凝聚这一事实而为同构；
见《概形的导出范畴》，引理
\ref{perfect-lemma-flat-proper-pseudo-coherent-direct-image-general}。
要看出相对杯积诱导的正是此同构，使用《上同调》，注
\ref{cohomology-remark-relative-cup-and-composition} 中图的交换性。
\end{remark}

\begin{lemma}
\label{duality-lemma-properties-relative-dualizing}
设 $Y$ 为拟紧拟分离概形，$f : X \to Y$ 为真、平坦且有限表示的概形态射，
其相对对偶化复形为 $\omega_{X/Y}^\bullet$
（注 \ref{duality-remark-relative-dualizing-complex}）。则
\begin{enumerate}
\item $\omega_{X/Y}^\bullet$ 是 $D(\mathcal{O}_X)$ 中的 $Y$-完美对象；
\item $Rf_*\omega_{X/Y}^\bullet$ 的正次数上同调层均为零；
\item $\mathcal{O}_X \to
R\SheafHom_{\mathcal{O}_X}(\omega_{X/Y}^\bullet, \omega_{X/Y}^\bullet)$
是同构。
\end{enumerate}
\end{lemma}

\begin{proof}
由于 $\omega_{X/Y}^\bullet$ 的构造与基变换交换（见注
\ref{duality-remark-relative-dualizing-complex}），可以并且确实假设 $Y$ 仿射。
对 $D(\mathcal{O}_X)$ 中的完美对象 $E$，有
\begin{align*}
Rf_*(E \otimes_{\mathcal{O}_X}^\mathbf{L} \omega_{X/Y}^\bullet)
& =
Rf_*R\SheafHom_{\mathcal{O}_X}(E^\vee, \omega_{X/Y}^\bullet) \\
& =
R\SheafHom_{\mathcal{O}_Y}(Rf_*E^\vee, \mathcal{O}_Y) \\
& =
(Rf_*E^\vee)^\vee
\end{align*}
第一个等式见《上同调》，引理
\ref{cohomology-lemma-dual-perfect-complex}。第二个等式见引理
\ref{duality-lemma-iso-on-RSheafHom}、注 \ref{duality-remark-iso-on-RSheafHom} 以及
《概形的导出范畴》，引理
\ref{perfect-lemma-flat-proper-perfect-direct-image-general}。
第三个等式是对偶的定义。尤其，这些引文也表明所得对象是
$D(\mathcal{O}_Y)$ 中的完美对象。由《态射进阶》，引理
\ref{more-morphisms-lemma-characterize-relatively-perfect}，可知
$\omega_{X/Y}^\bullet$ 是 $Y$-完美的。这证明了 (1)。

\medskip\noindent
设 $M$ 为 $D_\QCoh(\mathcal{O}_Y)$ 的对象。则
\begin{align*}
\Hom_Y(M, Rf_*\omega_{X/Y}^\bullet) & =
\Hom_X(Lf^*M, \omega_{X/Y}^\bullet) \\
& =
\Hom_Y(Rf_*Lf^*M, \mathcal{O}_Y) \\
& =
\Hom_Y(M \otimes_{\mathcal{O}_Y}^\mathbf{L} Rf_*\mathcal{O}_X, \mathcal{O}_Y)
\end{align*}
第一个等式由《上同调》，引理 \ref{cohomology-lemma-adjoint}。
第二个等式由 $a$ 的构造。第三个等式由《概形的导出范畴》，引理
\ref{perfect-lemma-cohomology-base-change}。回忆
$Rf_*\mathcal{O}_X$ 是 Tor 振幅位于 $[0, N]$ 的完美对象，其中
$N$ 为某个整数；见《概形的导出范畴》，引理
\ref{perfect-lemma-flat-proper-perfect-direct-image-general}。
因此可用位于次数 $[0, N]$ 的有限投射模复形表示
$Rf_*\mathcal{O}_X$（使用《代数进阶》，引理
\ref{more-algebra-lemma-perfect} 以及 $Y$ 仿射这一事实）。
故若对某个 $i > 0$ 有 $M = \mathcal{O}_Y[-i]$，则最后一个群为零。
由于 $Y$ 仿射，可知对 $i > 0$，
$H^i(Rf_*\omega_{X/Y}^\bullet) = 0$。这证明了 (2)。

\medskip\noindent
设 $E$ 为 $D_\QCoh(\mathcal{O}_X)$ 的完美对象。则有
\begin{align*}
\Hom_X(E, R\SheafHom_{\mathcal{O}_X}(\omega_{X/Y}^\bullet, \omega_{X/Y}^\bullet)
& =
\Hom_X(E \otimes_{\mathcal{O}_X}^\mathbf{L} \omega_{X/Y}^\bullet,
\omega_{X/Y}^\bullet) \\
& =
\Hom_Y(Rf_*(E \otimes_{\mathcal{O}_X}^\mathbf{L} \omega_{X/Y}^\bullet),
\mathcal{O}_Y) \\
& =
\Hom_Y(Rf_*(R\SheafHom_{\mathcal{O}_X}(E^\vee, \omega_{X/Y}^\bullet)),
\mathcal{O}_Y) \\
& =
\Hom_Y(R\SheafHom_{\mathcal{O}_Y}(Rf_*E^\vee, \mathcal{O}_Y),
\mathcal{O}_Y) \\
& =
R\Gamma(Y, Rf_*E^\vee) \\
& =
\Hom_X(E, \mathcal{O}_X)
\end{align*}
第一个等式由《上同调》，引理
\ref{cohomology-lemma-internal-hom}。第二个等式是
$\omega_{X/Y}^\bullet$ 的定义。第三个等式来自对偶完美复形
$E^\vee$ 的构造；见《上同调》，引理
\ref{cohomology-lemma-dual-perfect-complex}。第四个等式来自证明第一段中所示的等式
$Rf_*R\SheafHom_{\mathcal{O}_X}(E^\vee, \omega_{X/Y}^\bullet) =
R\SheafHom_{\mathcal{O}_Y}(Rf_*E^\vee, \mathcal{O}_Y)$。
第五个等式由完美复形的双重对偶性（《上同调》，引理
\ref{cohomology-lemma-dual-perfect-complex}）以及 $Rf_*E$ 完美这一事实成立；
后者见《概形的导出范畴》，引理
\ref{perfect-lemma-flat-proper-perfect-direct-image-general}。
最后一个等式是 $f$ 的 Leray 等式。这一串等式本质上由 Yoneda 引理表明 (3)
成立。具体地，由《概形的导出范畴》，引理
\ref{perfect-lemma-quasi-coherence-internal-hom}，对象
$R\SheafHom(\omega_{X/Y}^\bullet, \omega_{X/Y}^\bullet)$ 属于
$D_\QCoh(\mathcal{O}_X)$。在上式中取 $E = \mathcal{O}_X$，得到映射
$\alpha : \mathcal{O}_X \to
R\SheafHom_{\mathcal{O}_X}(\omega_{X/Y}^\bullet, \omega_{X/Y}^\bullet)$，
它对应于 $\text{id}_{\mathcal{O}_X} \in \Hom_X(\mathcal{O}_X, \mathcal{O}_X)$。
由于上述所有同构都关于 $E$ 函子性，可知 $\alpha$ 的锥是
$D_\QCoh(\mathcal{O}_X)$ 的一个对象 $C$，且对所有完美对象 $E$ 都有
$\Hom(E, C) = 0$。由于完美对象生成该范畴（《概形的导出范畴》，定理
\ref{perfect-theorem-bondal-van-den-Bergh}），可知 $\alpha$ 是同构。
\end{proof}

\begin{lemma}[刚性]
\label{duality-lemma-van-den-bergh}
设 $Y$ 为拟紧拟分离概形，$f : X \to Y$ 为有限表示的真平坦态射，
其相对对偶化复形为 $\omega_{X/Y}^\bullet$
（注 \ref{duality-remark-relative-dualizing-complex}）。存在典范同构
\begin{equation}
\label{duality-equation-pre-rigid}
\mathcal{O}_X =
c(L\text{pr}_1^*\omega_{X/Y}^\bullet) =
c(L\text{pr}_2^*\omega_{X/Y}^\bullet)
\end{equation}
以及典范同构
\begin{equation}
\label{duality-equation-rigid}
\omega_{X/Y}^\bullet =
c\left(L\text{pr}_1^*\omega_{X/Y}^\bullet
\otimes_{\mathcal{O}_{X \times_Y X}}^\mathbf{L}
L\text{pr}_2^*\omega_{X/Y}^\bullet\right)
\end{equation}
其中 $c$ 是引理 \ref{duality-lemma-twisted-inverse-image} 中对角态射
$\Delta : X \to X \times_Y X$ 的右伴随。
\end{lemma}

\begin{proof}
设 $a$ 为引理 \ref{duality-lemma-twisted-inverse-image} 中 $Rf_*$ 的右伴随。
考虑笛卡尔方块
$$
\xymatrix{
X \times_Y X \ar[r]_q \ar[d]_p & X \ar[d]_f \\
X \ar[r]^f & Y
}
$$
设 $b$ 为引理 \ref{duality-lemma-twisted-inverse-image} 中对应于 $p$ 的右伴随。则
\begin{align*}
\omega_{X/Y}^\bullet
& =
c(b(\omega_{X/Y}^\bullet)) \\
& =
c(Lp^*\omega_{X/Y}^\bullet
\otimes_{\mathcal{O}_{X \times_Y X}}^\mathbf{L} b(\mathcal{O}_X)) \\
& =
c(Lp^*\omega_{X/Y}^\bullet
\otimes_{\mathcal{O}_{X \times_Y X}}^\mathbf{L}
Lq^*a(\mathcal{O}_Y)) \\
& =
c(Lp^*\omega_{X/Y}^\bullet
\otimes_{\mathcal{O}_{X \times_Y X}}^\mathbf{L}
Lq^*\omega_{X/Y}^\bullet)
\end{align*}
如 (\ref{duality-equation-rigid}) 所示。解释如下：
\begin{enumerate}
\item 第一个等式成立，因为由 $\text{id}_X = p \circ \Delta$ 有
$\text{id} = c \circ b$。
\item 第二个等式由引理
\ref{duality-lemma-compare-with-pullback-flat-proper} 成立。
\item 第三个等式由引理 \ref{duality-lemma-proper-flat-base-change} 以及
$\mathcal{O}_X = Lf^*\mathcal{O}_Y$ 成立。
\item 第四个等式成立，因为 $\omega_{X/Y}^\bullet = a(\mathcal{O}_Y)$。
\end{enumerate}
方程 (\ref{duality-equation-pre-rigid}) 的证明完全相同。
\end{proof}

\begin{remark}
\label{duality-remark-van-den-bergh}
引理 \ref{duality-lemma-van-den-bergh} 表明，我们的相对对偶化复形在与
\cite{vdB-rigid} 所引入概念类似的意义下是{\it 刚性的}。
具体地，(\ref{duality-equation-rigid}) 右侧的函子关于
$\omega_{X/Y}^\bullet$ 是“二次的”，而 (\ref{duality-equation-rigid})
左侧的函子是“线性的”，
因而这在一定程度上“确定”了复形 $\omega_{X/Y}^\bullet$。
对偶理论还有一种使用“刚性”相对对偶化复形的处理方式；例如见
\cite{Neeman-rigid}、\cite{Yekutieli-rigid} 和 \cite{Yekutieli-Zhang}。
我们将在节 \ref{duality-section-relative-dualizing-complexes} 中回到这一点。
\end{remark}





\section{完美真态射的推前右伴随}
\label{duality-section-perfect-proper}

\noindent
本节恰当的一般性应当是考虑拟紧拟分离概形之间的完美真态射；见
\cite{LN}。

\begin{lemma}
\label{duality-lemma-proper-flat-noetherian}
设 $f : X \to Y$ 为 Noetherian 概形之间的完美真态射。
设 $a$ 为引理 \ref{duality-lemma-twisted-inverse-image} 中
$Rf_* : D_\QCoh(\mathcal{O}_X) \to D_\QCoh(\mathcal{O}_Y)$ 的右伴随。
则 $a$ 与直和交换。
\end{lemma}

\begin{proof}
设 $P$ 为 $D(\mathcal{O}_X)$ 的完美对象。由《态射进阶》，引理
\ref{more-morphisms-lemma-perfect-proper-perfect-direct-image}，
复形 $Rf_*P$ 在 $Y$ 上完美。设 $K_i$ 为
$D_\QCoh(\mathcal{O}_Y)$ 中的一族对象。则
\begin{align*}
\Hom_{D(\mathcal{O}_X)}(P, a(\bigoplus K_i))
& =
\Hom_{D(\mathcal{O}_Y)}(Rf_*P, \bigoplus K_i) \\
& =
\bigoplus \Hom_{D(\mathcal{O}_Y)}(Rf_*P, K_i) \\
& =
\bigoplus \Hom_{D(\mathcal{O}_X)}(P, a(K_i))
\end{align*}
因为完美对象是紧对象（《概形的导出范畴》，命题
\ref{perfect-proposition-compact-is-perfect}）。由于
$D_\QCoh(\mathcal{O}_X)$ 有完美生成元（《概形的导出范畴》，定理
\ref{perfect-theorem-bondal-van-den-Bergh}），可知映射
$\bigoplus a(K_i) \to a(\bigoplus K_i)$ 是同构，即 $a$ 与直和交换。
\end{proof}

\begin{lemma}
\label{duality-lemma-proper-flat-noetherian-relative}
设 $f : X \to Y$ 为 Noetherian 概形之间的完美真态射。
设 $a$ 为引理 \ref{duality-lemma-twisted-inverse-image} 中
$Rf_* : D_\QCoh(\mathcal{O}_X) \to D_\QCoh(\mathcal{O}_Y)$ 的右伴随。则
\begin{enumerate}
\item 对每个闭集 $T \subset Y$，若 $Q \in D_\QCoh(Y)$ 支撑于 $T$，
则 $a(Q)$ 支撑于 $f^{-1}(T)$；
\item 对每个开集 $V \subset Y$ 及任意 $K \in D_\QCoh(\mathcal{O}_Y)$，
映射 (\ref{duality-equation-sheafy}) 是同构；并且
\end{enumerate}
\end{lemma}

\begin{proof}
这由引理 \ref{duality-lemma-when-sheafy}、\ref{duality-lemma-proper-noetherian}
和 \ref{duality-lemma-proper-flat-noetherian} 得到。
\end{proof}

\begin{lemma}
\label{duality-lemma-compare-with-pullback-flat-proper-noetherian}
设 $f : X \to Y$ 为 Noetherian 概形之间的完美真态射。则对
$D_\QCoh(\mathcal{O}_Y)$ 的每个对象 $K$，映射
(\ref{duality-equation-compare-with-pullback}) 都是同构。
\end{lemma}

\begin{proof}
由引理 \ref{duality-lemma-proper-flat-noetherian}，$a$ 与直和交换。因此，
$D_\QCoh(\mathcal{O}_Y)$ 中使 (\ref{duality-equation-compare-with-pullback})
成为同构的对象构成 $D_\QCoh(\mathcal{O}_Y)$ 的严格满、饱和三角子范畴，
并且该子范畴在取直和下封闭。由于 $D_\QCoh(\mathcal{O}_Y)$ 是由单个
完美对象生成的模范畴（分别见《概形的导出范畴》，定理
\ref{perfect-theorem-DQCoh-is-Ddga} 和定理
\ref{perfect-theorem-bondal-van-den-Bergh}），可如《代数进阶》，注
\ref{more-algebra-remark-P-resolution} 那样论证，说明只需对一个完美对象证明
(\ref{duality-equation-compare-with-pullback}) 是同构。但该结论对完美对象成立；
见引理 \ref{duality-lemma-compare-with-pullback-perfect}。
\end{proof}

\begin{lemma}
\label{duality-lemma-proper-perfect-base-change}
设 $f : X \to Y$ 为 Noetherian 概形之间的完美真态射。
设 $g : Y' \to Y$ 为态射，且 $Y'$ Noetherian。若 $X$ 与
$Y'$ 在 $Y$ 上 Tor 独立，则对所有 $K \in D_\QCoh(\mathcal{O}_Y)$，
基变换映射 (\ref{duality-equation-base-change-map}) 是同构。
\end{lemma}

\begin{proof}
由引理 \ref{duality-lemma-proper-flat-noetherian-relative}，函子 $a$ 和 $a'$
的构造与到 $Y$ 和 $Y'$ 的开集的限制交换。因此可假设
$Y' \to Y$ 是仿射概形之间的态射；见注
\ref{duality-remark-check-over-affines}。此时结论由引理
\ref{duality-lemma-more-base-change} 得到。
\end{proof}



\section{有效 Cartier 除子的推前右伴随}
\label{duality-section-dualizing-Cartier}

\noindent
设 $X$ 为概形，$i : D \to X$ 为一个有效 Cartier 除子的嵌入。
记 $\mathcal{N} = i^*\mathcal{O}_X(D)$ 为 $i$ 的法层；见《态射》，节
\ref{morphisms-section-conormal-sheaf} 及《除子》，节
\ref{divisors-section-effective-Cartier-divisors}。回忆
$R\SheafHom(\mathcal{O}_D, -)$ 表示
$i_* : D(\mathcal{O}_D) \to D(\mathcal{O}_X)$ 的右伴随，并满足
$i_*R\SheafHom(\mathcal{O}_D, -) =
R\SheafHom_{\mathcal{O}_X}(i_*\mathcal{O}_D, -)$；见节
\ref{duality-section-sections-with-exact-support}。

\begin{lemma}
\label{duality-lemma-compute-for-effective-Cartier}
如上，设 $X$ 为概形，$D \subset X$ 为有效 Cartier 除子。
则在 $D(\mathcal{O}_D)$ 中存在典范同构
$R\SheafHom(\mathcal{O}_D, \mathcal{O}_X) = \mathcal{N}[-1]$。
\end{lemma}

\begin{proof}
等价地，这是说 $R\SheafHom(\mathcal{O}_D, \mathcal{O}_X)$
唯一非零的上同调层位于次数 $1$，且该层同构于 $\mathcal{N}$。
由于 $i_*$ 正合且全忠实，只需证明
$i_*R\SheafHom(\mathcal{O}_D, \mathcal{O}_X)$ 同构于
$i_*\mathcal{N}[-1]$。由引理
\ref{duality-lemma-sheaf-with-exact-support-ext}，有
$i_*R\SheafHom(\mathcal{O}_D, \mathcal{O}_X) =
R\SheafHom_{\mathcal{O}_X}(i_*\mathcal{O}_D, \mathcal{O}_X)$。
可使用分解
$$
0 \to \mathcal{I} \to \mathcal{O}_X \to i_*\mathcal{O}_D \to 0
$$
作计算，其中 $\mathcal{I}$ 是 $D$ 的理想层。由于局部计算给出
$R\SheafHom_{\mathcal{O}_X}(\mathcal{O}_X, \mathcal{O}_X) = \mathcal{O}_X$ 和
$R\SheafHom_{\mathcal{O}_X}(\mathcal{I}, \mathcal{O}_X) = \mathcal{O}_X(D)$，
可知
$$
R\SheafHom_{\mathcal{O}_X}(i_*\mathcal{O}_D, \mathcal{O}_X) =
(\mathcal{O}_X \to \mathcal{O}_X(D))
$$
其中右边的 $\mathcal{O}_X$ 位于次数 $0$，$\mathcal{O}_X(D)$ 位于次数 $1$。
结论由短正合列
$$
0 \to \mathcal{O}_X \to \mathcal{O}_X(D) \to i_*\mathcal{N} \to 0
$$
得到；该列来自 $D$ 是 $\mathcal{O}_X(D)$ 的典范截面的零点概形，
以及 $\mathcal{N} = i^*\mathcal{O}_X(D)$ 这一事实。
\end{proof}

\noindent
对 $D(\mathcal{O}_X)$ 的每个对象 $K$，存在典范映射
\begin{equation}
\label{duality-equation-map-effective-Cartier}
Li^*K
\otimes_{\mathcal{O}_D}^\mathbf{L}
R\SheafHom(\mathcal{O}_D, \mathcal{O}_X)
\longrightarrow
R\SheafHom(\mathcal{O}_D, K)
\end{equation}
它在 $D(\mathcal{O}_D)$ 中关于 $K$ 函子性，并与可区别三角相容。
具体地，此映射伴随于映射
$$
i_*(Li^*K \otimes^\mathbf{L}_{\mathcal{O}_D}
R\SheafHom(\mathcal{O}_D, \mathcal{O}_X)) =
K \otimes^\mathbf{L}_{\mathcal{O}_X}
R\SheafHom_{\mathcal{O}_X}(i_*\mathcal{O}_D, \mathcal{O}_X)
\longrightarrow K
$$
其中等式由《上同调》，引理
\ref{cohomology-lemma-projection-formula-closed-immersion} 给出，
而箭头来自映射 $\mathcal{O}_X \to i_*\mathcal{O}_D$ 所诱导的典范映射
$R\SheafHom_{\mathcal{O}_X}(i_*\mathcal{O}_D, \mathcal{O}_X) \to \mathcal{O}_X$。

\medskip\noindent
若 $K \in D_\QCoh(\mathcal{O}_X)$，则映射
(\ref{duality-equation-map-effective-Cartier}) 等于
(\ref{duality-equation-compare-with-pullback})；这里使用引理
\ref{duality-lemma-twisted-inverse-image-closed} 中的识别
$a(K) = R\SheafHom(\mathcal{O}_D, K)$。若
$K \in D_\QCoh(\mathcal{O}_X)$ 且 $X$ Noetherian，则下述引理是引理
\ref{duality-lemma-compare-with-pullback-flat-proper-noetherian} 的特殊情形。

\begin{lemma}
\label{duality-lemma-sheaf-with-exact-support-effective-Cartier}
如上，设 $X$ 为概形，$D \subset X$ 为有效 Cartier 除子。则
(\ref{duality-equation-map-effective-Cartier}) 与引理
\ref{duality-lemma-compute-for-effective-Cartier} 合在一起，定义了一个关于
$D(\mathcal{O}_X)$ 中的 $K$ 函子性的同构
$$
Li^*K \otimes_{\mathcal{O}_D}^\mathbf{L} \mathcal{N}[-1]
\longrightarrow
R\SheafHom(\mathcal{O}_D, K)
$$
\end{lemma}

\begin{proof}
由于 $i_*$ 在模上正合且全忠实，为证明该映射是同构，只需证明应用
$i_*$ 后它是同构。以下将不再说明地使用引理
\ref{duality-lemma-compute-for-effective-Cartier} 的证明中所用的短正合列
$0 \to \mathcal{I} \to \mathcal{O}_X \to i_*\mathcal{O}_D \to 0$
和
$0 \to \mathcal{O}_X \to \mathcal{O}_X(D) \to i_*\mathcal{N} \to 0$。
由《上同调》，引理
\ref{cohomology-lemma-projection-formula-closed-immersion}
（它用于定义映射 (\ref{duality-equation-map-effective-Cartier})），左边变为
$$
K \otimes_{\mathcal{O}_X}^\mathbf{L} i_*\mathcal{N}[-1] =
K \otimes_{\mathcal{O}_X}^\mathbf{L} (\mathcal{O}_X \to \mathcal{O}_X(D))
$$
右边变为
\begin{align*}
R\SheafHom_{\mathcal{O}_X}(i_*\mathcal{O}_D, K) & =
R\SheafHom_{\mathcal{O}_X}((\mathcal{I} \to \mathcal{O}_X), K) \\
& =
R\SheafHom_{\mathcal{O}_X}((\mathcal{I} \to \mathcal{O}_X), \mathcal{O}_X)
\otimes_{\mathcal{O}_X}^\mathbf{L} K
\end{align*}
最后一个等式由《上同调》，引理
\ref{cohomology-lemma-dual-perfect-complex} 得到。由于该映射来自同构
$$
R\SheafHom_{\mathcal{O}_X}((\mathcal{I} \to \mathcal{O}_X), \mathcal{O}_X)
= (\mathcal{O}_X \to \mathcal{O}_X(D))
$$
结论显然。
\end{proof}





\section{例中的推前右伴随}
\label{duality-section-examples}

\noindent
本节在若干例子中计算推前的右伴随。这些同构是典范的，但仅在最弱的
可能意义下如此；也就是说，我们既不证明也不声称这些同构与基变换、
态射复合等各种运算相容。关于这类问题有大量文献。读者可从
\cite{RD}、\cite{Conrad-GD} 中的材料开始（这些文献采用不同的对偶理论
出发点，但讨论了为相对对偶化复形构造典范代表的问题），然后继续查阅
Joseph Lipman 及其合作者的工作。

\begin{lemma}
\label{duality-lemma-upper-shriek-P1}
设 $Y$ 为拟紧拟分离概形。设 $\mathcal{E}$ 为秩为 $n + 1$ 的有限局部自由
$\mathcal{O}_Y$-模，其行列式为
$\mathcal{L} = \wedge^{n + 1}(\mathcal{E})$。设投影为
$f : X = \mathbf{P}(\mathcal{E}) \to Y$。设 $a$ 为引理
\ref{duality-lemma-twisted-inverse-image} 中
$Rf_* : D_\QCoh(\mathcal{O}_X) \to D_\QCoh(\mathcal{O}_Y)$ 的右伴随。
则存在同构
$$
c : f^*\mathcal{L}(-n - 1)[n] \longrightarrow a(\mathcal{O}_Y)
$$
特别地，若 $\mathcal{E} = \mathcal{O}_Y^{\oplus n + 1}$，则
$X = \mathbf{P}^n_Y$，并得到
$a(\mathcal{O}_Y) = \mathcal{O}_X(-n - 1)[n]$。
\end{lemma}

\begin{proof}
在《概形的上同调》，引理
\ref{coherent-lemma-cohomology-projective-bundle} 的证明中，我们构造了典范同构
$$
R^nf_*(f^*\mathcal{L}(-n - 1)) \longrightarrow \mathcal{O}_Y
$$
此外，$Rf_*(f^*\mathcal{L}(-n - 1))[n] = R^nf_*(f^*\mathcal{L}(-n - 1))$，
即其他高阶正像均为零。因此得到同构
$$
Rf_*(f^*\mathcal{L}(-n - 1)[n]) \longrightarrow \mathcal{O}_Y
$$
由于 $a$ 是 $Rf_*$ 的右伴随，此同构确定了引理陈述中的 $c$。
由引理 \ref{duality-lemma-proper-noetherian}，$a$ 的构造关于基底是局部的。
特别地，为检验 $c$ 是同构，可以在 $Y$ 上局部工作。换言之，可假设
$Y$ 仿射且 $\mathcal{E} = \mathcal{O}_Y^{\oplus n + 1}$。
此时层 $\mathcal{O}_X, \mathcal{O}_X(-1), \ldots,
\mathcal{O}_X(-n)$ 生成 $D_\QCoh(X)$；见《概形的导出范畴》，引理
\ref{perfect-lemma-generator-P1}。因此，只需证明映射
$c : \mathcal{O}_X(-n - 1)[n] \to a(\mathcal{O}_Y)$ 在函子
$$
F_{i, p}(-) = \Hom_{D(\mathcal{O}_X)}(\mathcal{O}_X(i), (-)[p])
$$
下变为同构，其中 $i \in \{-n, \ldots, 0\}$ 且
$p \in \mathbf{Z}$。对 $F_{0, p}$，这由箭头 $c$ 的构造成立！
对 $i \in \{-n, \ldots, -1\}$，有
$$
\Hom_{D(\mathcal{O}_X)}(\mathcal{O}_X(i), \mathcal{O}_X(-n - 1)[n + p]) =
H^p(X, \mathcal{O}_X(-n - 1 - i)) = 0
$$
这是射影空间上同调计算的结果（《概形的上同调》，引理
\ref{coherent-lemma-cohomology-projective-space-over-ring}）；并且有
$$
\Hom_{D(\mathcal{O}_X)}(\mathcal{O}_X(i), a(\mathcal{O}_Y)[p]) =
\Hom_{D(\mathcal{O}_Y)}(Rf_*\mathcal{O}_X(i), \mathcal{O}_Y[p]) = 0
$$
因为同一引理给出 $Rf_*\mathcal{O}_X(i) = 0$。因此
$F_{i, p}(c)$ 的源和目标均为零，故 $F_{i, p}(c)$ 必为同构。
证明完成。
\end{proof}

\begin{example}
\label{duality-example-base-change-wrong}
若 $f$ 为完美真态射而 $g$ 完美，则基变换映射
(\ref{duality-equation-base-change-map}) 不一定是同构。设 $k$ 为域。
令 $Y = \mathbf{A}^2_k$，并令 $f : X \to Y$ 为 $Y$ 在原点处的吹起。
记 $E \subset X$ 为例外除子。则 $f$ 可分解为
$$
X \xrightarrow{i} \mathbf{P}^1_Y \xrightarrow{p} Y
$$
由此得到分解 $a = c \circ b$，其中 $a$、$b$、$c$ 分别是引理
\ref{duality-lemma-twisted-inverse-image} 中 $Rf_*$、$Rp_*$、$Ri_*$ 的右伴随。
记 $\mathcal{O}(n)$ 为 $\mathbf{P}^1_Y$ 上结构层的 Serre 扭曲，
记 $\mathcal{O}_X(n)$ 为其到 $X$ 的限制。注意
$X \subset \mathbf{P}^1_Y$ 由一次方程截出，故
$\mathcal{O}(X) = \mathcal{O}(1)$。由引理
\ref{duality-lemma-upper-shriek-P1}，有
$b(\mathcal{O}_Y) = \mathcal{O}(-2)[1]$。由引理
\ref{duality-lemma-twisted-inverse-image-closed}，有
$$
a(\mathcal{O}_Y) = c(b(\mathcal{O}_Y)) =
c(\mathcal{O}(-2)[1]) =
R\SheafHom(\mathcal{O}_X, \mathcal{O}(-2)[1]) =
\mathcal{O}_X(-1)
$$
最后一个等式由引理
\ref{duality-lemma-sheaf-with-exact-support-effective-Cartier}。令
$Y' = \Spec(k)$ 为 $Y$ 中的原点。$a(\mathcal{O}_Y)$ 到
$X' = E = \mathbf{P}^1_k$ 的限制是次数为 $-1$、置于上同调次数
$0$ 的可逆层。另一方面，
$a'(\mathcal{O}_{\Spec(k)}) = \mathcal{O}_E(-2)[1]$
是次数为 $-2$、置于上同调次数 $-1$ 的可逆层，因而二者不同。
在此例中，引理 \ref{duality-lemma-more-base-change} 的 Tor 独立性假设不成立。
\end{example}

\begin{lemma}
\label{duality-lemma-ext}
设 $Y$ 为赋环空间，$\mathcal{I} \subset \mathcal{O}_Y$ 为理想层。
令 $\mathcal{O}_X = \mathcal{O}_Y/\mathcal{I}$ 及
$\mathcal{N} =
\SheafHom_{\mathcal{O}_Y}(\mathcal{I}/\mathcal{I}^2, \mathcal{O}_X)$。
则存在典范同构
$c : \mathcal{N} \to
\SheafExt^1_{\mathcal{O}_Y}(\mathcal{O}_X, \mathcal{O}_X)
$。
\end{lemma}

\begin{proof}
考虑典范短正合列
\begin{equation}
\label{duality-equation-second-order-thickening}
0 \to \mathcal{I}/\mathcal{I}^2 \to \mathcal{O}_Y/\mathcal{I}^2 \to
\mathcal{O}_X \to 0
\end{equation}
设 $U \subset X$ 为开集，$s \in \mathcal{N}(U)$。沿 $s$ 对
(\ref{duality-equation-second-order-thickening}) 作推出，可得到
$\mathcal{O}_X|_U$ 被 $\mathcal{O}_X|_U$ 的扩张 $E_s$。
这又定义了 $\SheafExt^1_{\mathcal{O}_Y}(\mathcal{O}_X, \mathcal{O}_X)$
在 $U$ 上的截面 $c(s)$。见《上同调》，引理
\ref{cohomology-lemma-section-RHom-over-U} 和《导出范畴》，引理
\ref{derived-lemma-ext-1}。反过来，给定 $\mathcal{O}_U$-模扩张
$$
0 \to \mathcal{O}_X|_U \to \mathcal{E} \to \mathcal{O}_X|_U \to 0
$$
可找到开覆盖 $U = \bigcup U_i$ 及截面
$e_i \in \mathcal{E}(U_i)$，它们映到 $1 \in \mathcal{O}_X(U_i)$。
于是 $e_i$ 定义映射
$\mathcal{O}_Y|_{U_i} \to \mathcal{E}|_{U_i}$，其核包含
$\mathcal{I}^2$。由此可见，$\mathcal{E}|_{U_i}$ 来自上述推出。
这说明 $c$ 满射。单射性的证明从略。
\end{proof}

\begin{lemma}
\label{duality-lemma-regular-ideal-ext}
设 $Y$ 为赋环空间，$\mathcal{I} \subset \mathcal{O}_Y$ 为理想层，
令 $\mathcal{O}_X = \mathcal{O}_Y/\mathcal{I}$。若 $\mathcal{I}$
是 Koszul-正则的（《除子》，定义
\ref{divisors-definition-regular-ideal-sheaf}），则
$R\SheafHom_{\mathcal{O}_Y}(\mathcal{O}_X, \mathcal{O}_X)$ 上的复合定义同构
$$
\wedge^i(\SheafExt^1_{\mathcal{O}_Y}(\mathcal{O}_X, \mathcal{O}_X))
\longrightarrow
\SheafExt^i_{\mathcal{O}_Y}(\mathcal{O}_X, \mathcal{O}_X)
$$
对所有 $i$ 均成立。
\end{lemma}

\begin{proof}
这里的复合是指《上同调》，引理
\ref{cohomology-lemma-internal-hom-composition} 中的映射
$$
R\SheafHom_{\mathcal{O}_Y}(\mathcal{O}_X, \mathcal{O}_X)
\otimes_{\mathcal{O}_Y}^\mathbf{L}
R\SheafHom_{\mathcal{O}_Y}(\mathcal{O}_X, \mathcal{O}_X)
\longrightarrow
R\SheafHom_{\mathcal{O}_Y}(\mathcal{O}_X, \mathcal{O}_X)
$$
它诱导乘法映射
$$
\SheafExt^a_{\mathcal{O}_Y}(\mathcal{O}_X, \mathcal{O}_X)
\otimes_{\mathcal{O}_Y}
\SheafExt^b_{\mathcal{O}_Y}(\mathcal{O}_X, \mathcal{O}_X)
\longrightarrow
\SheafExt^{a + b}_{\mathcal{O}_Y}(\mathcal{O}_X, \mathcal{O}_X)
$$
可与《代数进阶》，方程
(\ref{more-algebra-equation-simple-tor-product}) 比较。引理陈述的含义是诱导映射
$$
\SheafExt^1_{\mathcal{O}_Y}(\mathcal{O}_X, \mathcal{O}_X)
\otimes \ldots \otimes
\SheafExt^1_{\mathcal{O}_Y}(\mathcal{O}_X, \mathcal{O}_X)
\longrightarrow
\SheafExt^i_{\mathcal{O}_Y}(\mathcal{O}_X, \mathcal{O}_X)
$$
经由外积因子分解，并由此诱导同构。为证明这一点，可以在 $Y$ 上局部工作。
因此可假设存在 $\mathcal{O}_Y$ 的整体截面
$f_1, \ldots, f_r$，它们生成 $\mathcal{I}$ 并构成 Koszul 正则序列。记
$$
\mathcal{A} = \mathcal{O}_Y\langle \xi_1, \ldots, \xi_r\rangle
$$
为严格交换微分分次 $\mathcal{O}_Y$-代数层；它是在
$\mathcal{O}_Y$ 上由次数 $-1$ 的 $\xi_1, \ldots, \xi_r$
生成的（除幂）多项式代数，微分 $\text{d}$ 由规则
$\text{d}\xi_i = f_i$ 给出。记 $\mathcal{A}^\bullet$ 为其底层
$\mathcal{O}_Y$-模复形，即上述 Koszul 复形。于是典范映射
$\mathcal{A}^\bullet \to \mathcal{O}_X$ 是拟同构。由《上同调》，引理
\ref{cohomology-lemma-Rhom-strictly-perfect}，得到拟同构
$$
R\SheafHom_{\mathcal{O}_Y}(\mathcal{O}_X, \mathcal{O}_X) \to
\SheafHom^\bullet(\mathcal{A}^\bullet, \mathcal{A}^\bullet) \to
\SheafHom^\bullet(\mathcal{A}^\bullet, \mathcal{O}_X)
$$
后一个复形的微分为零，因而
$$
\SheafExt^i_{\mathcal{O}_Y}(\mathcal{O}_X, \mathcal{O}_X)
\cong \SheafHom_{\mathcal{O}_Y}(\mathcal{A}^{-i}, \mathcal{O}_X)
$$
对 $j \in \{1, \ldots, r\}$，令
$\delta_j : \mathcal{A} \to \mathcal{A}$ 为次数 $1$ 的导子，满足
$\delta_j(\xi_i) = \delta_{ij}$（Kronecker 符号）。计算表明
$\delta_j \circ \text{d} = - \text{d} \circ \delta_j$，故得到复形态射
$$
\delta_j : \mathcal{A}^\bullet \to \mathcal{A}^\bullet[1].
$$
因此 $\delta_j$ 定义了相应 $\SheafExt$-层的一个截面。另一次计算表明，
$\delta_1, \ldots, \delta_r$ 映到
$\SheafHom_{\mathcal{O}_Y}(\mathcal{A}^{-1}, \mathcal{O}_X)$
在 $\mathcal{O}_X$ 上的一组基。显然，作为 $\mathcal{A}$ 的自同态，
$\delta_j \circ \delta_j = 0$，且
$\delta_j \circ \delta_{j'} = - \delta_{j'} \circ \delta_j$；
故在 $\SheafExt$-层中亦如此。这说明上述映射经由外幂因子分解。
为看出所得映射是所需同构，读者只需验证元素
$$
\delta_{j_1} \circ \ldots \circ \delta_{j_i}
$$
对 $j_1 < \ldots < j_i$ 映到层
$\SheafHom_{\mathcal{O}_Y}(\mathcal{A}^{-i}, \mathcal{O}_X)$
在 $\mathcal{O}_X$ 上的一组基。
\end{proof}

\begin{lemma}
\label{duality-lemma-regular-immersion-ext}
设 $Y$ 为赋环空间，$\mathcal{I} \subset \mathcal{O}_Y$ 为理想层。
令 $\mathcal{O}_X = \mathcal{O}_Y/\mathcal{I}$ 及
$\mathcal{N} =
\SheafHom_{\mathcal{O}_Y}(\mathcal{I}/\mathcal{I}^2, \mathcal{O}_X)$。
若 $\mathcal{I}$ 是 Koszul-正则的（《除子》，定义
\ref{divisors-definition-regular-ideal-sheaf}），则
$$
R\SheafHom_{\mathcal{O}_Y}(\mathcal{O}_X, \mathcal{O}_Y) =
\wedge^r \mathcal{N}[-r]
$$
其中 $r : Y \to \{1, 2, 3, \ldots \}$ 把 $y$ 映到 $y$ 的某个邻域中
生成 $\mathcal{I}$ 所需的最少生成元数。
\end{lemma}

\begin{proof}
由引理 \ref{duality-lemma-ext} 和 \ref{duality-lemma-regular-ideal-ext}，可知对
$i \geq 0$ 有同构
$\wedge^i\mathcal{N} \to
\SheafExt^i_{\mathcal{O}_Y}(\mathcal{O}_X, \mathcal{O}_X)$。
因此，只需证明映射 $\mathcal{O}_Y \to \mathcal{O}_X$ 诱导同构
$$
\SheafExt^r_{\mathcal{O}_Y}(\mathcal{O}_X, \mathcal{O}_Y)
\longrightarrow
\SheafExt^r_{\mathcal{O}_Y}(\mathcal{O}_X, \mathcal{O}_X)
$$
并证明当 $i \not = r$ 时，
$\SheafExt^i_{\mathcal{O}_Y}(\mathcal{O}_X, \mathcal{O}_Y)$ 为零。
这些陈述关于 $Y$ 是局部的。因此可假设有 $\mathcal{O}_Y$ 的整体截面
$f_1, \ldots, f_r$，它们生成 $\mathcal{I}$ 并构成 Koszul 正则序列。
令 $\mathcal{A}^\bullet$ 为引理
\ref{duality-lemma-regular-ideal-ext} 的证明中引入的关于
$f_1, \ldots, f_r$ 的 Koszul 复形。则由《上同调》，引理
\ref{cohomology-lemma-Rhom-strictly-perfect}，有
$$
R\SheafHom_{\mathcal{O}_Y}(\mathcal{O}_X, \mathcal{O}_Y) =
\SheafHom^\bullet(\mathcal{A}^\bullet, \mathcal{O}_Y)
$$
记 $1 : \mathcal{A}^\bullet \to \mathcal{O}_Y$ 为微分分次
$\mathcal{O}_Y$-代数的映射，它在次数 $0$ 上由恒等映射
$\mathcal{A}^0 = \mathcal{O}_Y \to \mathcal{O}_Y$ 给出。
取引理 \ref{duality-lemma-regular-ideal-ext} 证明中的 $\delta_j$，通过把
$\delta_{j_1} \ldots \delta_{j_i}$ 映到次数 $i$ 中的
$1 \circ \delta_{j_1} \circ \ldots \circ \delta_{j_i}$，得到分次
$\mathcal{O}_Y$-模同构
$$
\mathcal{O}_Y\langle \delta_1, \ldots, \delta_r\rangle
\longrightarrow
\SheafHom^\bullet(\mathcal{A}^\bullet, \mathcal{O}_Y)
$$
在此同构下，右边的微分在左边诱导微分 $\text{d}$。由符号规则，
$\text{d}(1) = - \sum f_j \delta_j$。由于
$\delta_j : \mathcal{A}^\bullet \to \mathcal{A}^\bullet[1]$
是复形态射，故
$$
\text{d}(\delta_{j_1} \ldots \delta_{j_i}) =
(- \sum f_j \delta_j )\delta_{j_1} \ldots \delta_{j_i}
$$
注意在微分分次代数 $\mathcal{A}$ 上有
$\text{d} = \sum f_j \delta_j$。因此，由规则
$$
1 \circ \delta_{j_1} \ldots \delta_{j_i} \longmapsto
(\delta_{j_1} \circ \ldots \circ \delta_{j_i})(\xi_1 \ldots \xi_r)
$$
定义的映射在 $r$ 为奇数时给出复形同构
$$
\SheafHom^\bullet(\mathcal{A}^\bullet, \mathcal{O}_Y)
\longrightarrow \mathcal{A}^\bullet[-r]
$$
而在 $r$ 为偶数时与微分相差一个符号而交换。无论如何，这些复形有同构的
上同调，从而得到所需消没。在映射
$$
\SheafHom^\bullet(\mathcal{A}^\bullet, \mathcal{O}_Y)
\longrightarrow
\SheafHom^\bullet(\mathcal{A}^\bullet, \mathcal{O}_X)
$$
下次数 $r$ 的上同调同构也由此直接得到。
（注意，我们对 Koszul 复形的约定确实参与了同构
$R\SheafHom_{\mathcal{O}_X}(\mathcal{O}_X, \mathcal{O}_Y) =
\wedge^r \mathcal{N}[-r]$ 的定义。）
\end{proof}

\begin{lemma}
\label{duality-lemma-regular-immersion}
设 $Y$ 为拟紧拟分离概形，$i : X \to Y$ 为 Koszul-正则闭浸入。
设 $a$ 为引理 \ref{duality-lemma-twisted-inverse-image} 中
$Ri_* : D_\QCoh(\mathcal{O}_X) \to D_\QCoh(\mathcal{O}_Y)$ 的右伴随。
则存在同构
$$
\wedge^r\mathcal{N}[-r] \longrightarrow a(\mathcal{O}_Y)
$$
其中
$\mathcal{N} = \SheafHom_{\mathcal{O}_X}(\mathcal{C}_{X/Y}, \mathcal{O}_X)$
是 $i$ 的法层（《态射》，节
\ref{morphisms-section-conormal-sheaf}），而 $r$ 是将其秩视为
$X$ 上局部常值函数所得的函数。
\end{lemma}

\begin{proof}
由引理 \ref{duality-lemma-twisted-inverse-image-closed} 和
\ref{duality-lemma-sheaf-with-exact-support-ext}，回忆 $a(\mathcal{O}_Y)$ 是
$D_\QCoh(\mathcal{O}_X)$ 的对象，其到 $Y$ 的推前为
$R\SheafHom_{\mathcal{O}_Y}(i_*\mathcal{O}_X, \mathcal{O}_Y)$。
因此结论由引理 \ref{duality-lemma-regular-immersion-ext} 得到。
\end{proof}

\begin{lemma}
\label{duality-lemma-smooth-proper}
设 $S$ 为 Noetherian 概形，$f : X \to S$ 为相对维数 $d$ 的光滑真态射。
设 $a$ 为引理 \ref{duality-lemma-twisted-inverse-image} 中
$Rf_* : D_\QCoh(\mathcal{O}_X) \to D_\QCoh(\mathcal{O}_S)$ 的右伴随。
则在 $D(\mathcal{O}_X)$ 中存在同构
$$
\wedge^d \Omega_{X/S}[d] \longrightarrow a(\mathcal{O}_S)
$$
\end{lemma}

\begin{proof}
如注 \ref{duality-remark-relative-dualizing-complex}，令
$\omega_{X/S}^\bullet = a(\mathcal{O}_S)$。令 $c$ 为引理
\ref{duality-lemma-twisted-inverse-image} 中对应于
$\Delta : X \to X \times_S X$ 的右伴随。由于 $\Delta$ 是光滑态射的对角，
它是 Koszul-正则浸入；见《除子》，引理
\ref{divisors-lemma-immersion-smooth-into-smooth-regular-immersion}。
特别地，$\Delta$ 是完美真态射（《态射进阶》，引理
\ref{more-morphisms-lemma-regular-immersion-perfect}），从而得到
\begin{align*}
\mathcal{O}_X
& =
c(L\text{pr}_1^*\omega_{X/S}^\bullet) \\
& =
L\Delta^*(L\text{pr}_1^*\omega_{X/S}^\bullet)
\otimes_{\mathcal{O}_X}^\mathbf{L}
c(\mathcal{O}_{X \times_S X}) \\
& =
\omega_{X/S}^\bullet \otimes_{\mathcal{O}_X}^\mathbf{L}
c(\mathcal{O}_{X \times_S X}) \\
& =
\omega_{X/S}^\bullet
\otimes_{\mathcal{O}_X}^\mathbf{L}
\wedge^d(\mathcal{N}_\Delta)[-d]
\end{align*}
第一个等式是 (\ref{duality-equation-pre-rigid})，因为
$\omega_{X/S}^\bullet = a(\mathcal{O}_S)$。第二个等式由引理
\ref{duality-lemma-compare-with-pullback-flat-proper-noetherian}。第三个等式成立，因为
$\text{pr}_1 \circ \Delta = \text{id}_X$。第四个等式由引理
\ref{duality-lemma-regular-immersion}。注意 $\wedge^d(\mathcal{N}_\Delta)$
是可逆 $\mathcal{O}_X$-模。因此
$\wedge^d(\mathcal{N}_\Delta)[-d]$ 是 $D(\mathcal{O}_X)$ 中的可逆对象，
从而 $a(\mathcal{O}_S) = \omega_{X/S}^\bullet = \wedge^d(\mathcal{C}_\Delta)[d]$。
由于 $\Delta$ 的余法层 $\mathcal{C}_\Delta$ 由《态射》，引理
\ref{morphisms-lemma-differentials-diagonal} 等于 $\Omega_{X/S}$，证明完成。
\end{proof}







\section{上叹号函子}
\label{duality-section-upper-shriek}

\noindent
本节利用紧化，为在一个固定 Noetherian 基底上有限型且分离的概形之间的
态射构造函子 $f^!$。如凝聚对偶理论中通常的情形，需要证明若干图交换。
建议读者在读完构造后跳过这些引理的验证，直接进入下一节对上叹号函子性质的讨论。

\begin{situation}
\label{duality-situation-shriek}
这里 $S$ 是 Noetherian 概形，而 $\textit{FTS}_S$ 是如下范畴：
\begin{enumerate}
\item 对象是 $S$ 上的概形 $X$，其结构态射 $X \to S$ 既分离又有限型；
\item 对象之间的态射 $f : X \to Y$ 是 $S$ 上的概形态射。
\end{enumerate}
\end{situation}

\noindent
在情形 \ref{duality-situation-shriek} 中，给定 $\textit{FTS}_S$ 中的态射
$f : X \to Y$，我们将定义三角范畴间的正合函子
$$
f^! : D_\QCoh^+(\mathcal{O}_Y) \to D_\QCoh^+(\mathcal{O}_X)
$$
具体地，选取在 $Y$ 上的紧化 $X \to \overline{X}$；其存在性由
《平坦性进阶》，定理 \ref{flat-theorem-nagata} 和引理
\ref{flat-lemma-compactifyable}。记结构态射为
$\overline{f} : \overline{X} \to Y$。令
$\overline{a} : D_\QCoh(\mathcal{O}_Y) \to D_\QCoh(\mathcal{O}_{\overline{X}})$
为引理 \ref{duality-lemma-twisted-inverse-image} 中所构造的
$R\overline{f}_*$ 的右伴随。对 $K \in D_\QCoh^+(\mathcal{O}_Y)$，令
$$
f^!K  = \overline{a}(K)|_X
$$
由引理 \ref{duality-lemma-twisted-inverse-image-bounded-below}，所得对象属于
$D_\QCoh^+(\mathcal{O}_X)$。

\begin{lemma}
\label{duality-lemma-shriek-well-defined}
在情形 \ref{duality-situation-shriek} 中，设 $f : X \to Y$ 为
$\textit{FTS}_S$ 的态射。函子 $f^!$ 在典范同构意义下与紧化的选择无关。
\end{lemma}

\begin{proof}
$X$ 在 $Y$ 上的紧化范畴定义于《平坦性进阶》，节
\ref{flat-section-compactify}。由《平坦性进阶》，定理
\ref{flat-theorem-nagata} 和引理 \ref{flat-lemma-compactifyable}，
该范畴非空。对紧化的每个选择
$$
j : X \to \overline{X},\quad \overline{f} : \overline{X} \to Y
$$
上述构造配以函子
$j^* \circ \overline{a} :
D_\QCoh^+(\mathcal{O}_Y) \to D_\QCoh^+(\mathcal{O}_X)$，
其中 $\overline{a}$ 是引理 \ref{duality-lemma-twisted-inverse-image} 中构造的
$R\overline{f}_*$ 的右伴随。

\medskip\noindent
假设给定 $Y$ 上紧化 $j_i : X \to \overline{X}_i$ 之间的态射
$g : \overline{X}_1 \to \overline{X}_2$，满足
$g^{-1}(j_2(X)) = j_1(X)$\footnote{按我们对紧化的定义，这一点可能不成立。
见《平坦性进阶》，节 \ref{flat-section-compactify}。}。
令 $\overline{c}$ 为引理 \ref{duality-lemma-twisted-inverse-image} 中对应于 $g$ 的右伴随。
则 $\overline{c} \circ \overline{a}_2 = \overline{a}_1$，因为这些函子伴随于
$R\overline{f}_{2, *} \circ Rg_* = R(\overline{f}_2 \circ g)_*$。
由 (\ref{duality-equation-sheafy})，存在函子的典范变换
$$
j_1^* \circ \overline{c} \longrightarrow j_2^*
$$
其定义域和值域为
$D^+_\QCoh(\mathcal{O}_{\overline{X}_2}) \to D^+_\QCoh(\mathcal{O}_X)$；
由引理 \ref{duality-lemma-proper-noetherian}，它是同构。复合
$$
j_1^* \circ \overline{a}_1 \longrightarrow
j_1^* \circ \overline{c} \circ \overline{a}_2 \longrightarrow
j_2^* \circ \overline{a}_2
$$
是函子同构，记作 $\alpha_g$。

\medskip\noindent
考虑 $X$ 在 $Y$ 上的两个紧化
$j_i : X \to \overline{X}_i$，$i = 1, 2$。由《平坦性进阶》，引理
\ref{flat-lemma-compactifications-cofiltered} 第 (b) 部分，可找到具有稠密像的紧化
$j : X \to \overline{X}$ 以及紧化态射
$g_i : \overline{X} \to \overline{X}_i$。由《平坦性进阶》，引理
\ref{flat-lemma-compactifications-cofiltered} 第 (c) 部分，有
$g_i^{-1}(j_i(X)) = j(X)$。因此由上一段得到同构
$$
\alpha_{g_i} :
j^* \circ \overline{a}
\longrightarrow
j_i^* \circ \overline{a}_i
$$
从而得到同构
$$
\alpha_{g_2} \circ \alpha_{g_1}^{-1} :
j_1^* \circ \overline{a}_1 \to j_2^* \circ \overline{a}_2
$$
为完成证明，必须说明这些同构定义良好。我们断言，只需证明上一段中所构造的
同构的复合仍是这种同构（精确陈述见下一段）。建议读者在纸片上检验这一点；
不过我们也在本段余下部分完整写出细节。具体地，再选取一个具有稠密像的紧化
$j' : X \to \overline{X}'$ 以及紧化态射
$g'_i : \overline{X}' \to \overline{X}_i$。由《平坦性进阶》，引理
\ref{flat-lemma-compactifications-cofiltered}，可找到具有稠密像的紧化
$j'' : X \to \overline{X}''$ 以及紧化态射
$h : \overline{X}'' \to \overline{X}$ 和
$h' : \overline{X}'' \to \overline{X}'$。甚至可以假设
$g_1 \circ h = g'_1 \circ h'$ 且 $g_2 \circ h = g'_2 \circ h'$。
下一段的结果给出
$$
\alpha_{g_i} \circ \alpha_h = \alpha_{g_i \circ h} =
\alpha_{g'_i \circ h'} = \alpha_{g'_i} \circ \alpha_{h'}
$$
其中 $i = 1, 2$。由于这些全都是函子同构，可知所需等式
$\alpha_{g_2} \circ \alpha_{g_1}^{-1} =
\alpha_{g'_2} \circ \alpha_{g'_1}^{-1}$ 成立。

\medskip\noindent
假设给定紧化 $j_i : X \to \overline{X}_i$，其中
$i = 1, 2, 3$，以及紧化态射
$g : \overline{X}_1 \to \overline{X}_2$ 和
$h : \overline{X}_2 \to \overline{X}_3$，满足
$g^{-1}(j_2(X)) = j_1(X)$ 且 $h^{-1}(j_2(X)) = j_3(X)$。
令 $\overline{a}_i$ 如上。上述断言是指
$$
\alpha_g \circ \alpha_h = \alpha_{g \circ h} :
j_1^* \circ \overline{a}_1 \to j_3^* \circ \overline{a}_3
$$
令 $\overline{c}$（相应地\ $\overline{d}$）为引理
\ref{duality-lemma-twisted-inverse-image} 中对应于 $g$（相应地\ $h$）的右伴随。
则 $\overline{c} \circ \overline{a}_2 = \overline{a}_1$ 且
$\overline{d} \circ \overline{a}_3 = \overline{a}_2$，并存在典范变换
$$
j_1^* \circ \overline{c} \longrightarrow j_2^*
\quad\text{且}\quad
j_2^* \circ \overline{d} \longrightarrow j_3^*
$$
其函子分别为
$D^+_\QCoh(\mathcal{O}_{\overline{X}_2}) \to D^+_\QCoh(\mathcal{O}_X)$
和
$D^+_\QCoh(\mathcal{O}_{\overline{X}_3}) \to D^+_\QCoh(\mathcal{O}_X)$；
理由同上。记 $\overline{e}$ 为引理
\ref{duality-lemma-twisted-inverse-image} 中对应于 $h \circ g$ 的右伴随。
由 (\ref{duality-equation-sheafy})，存在函子的典范变换
$$
j_1^* \circ \overline{e} \longrightarrow j_3^*
$$
其定义域和值域为
$D^+_\QCoh(\mathcal{O}_{\overline{X}_3}) \to D^+_\QCoh(\mathcal{O}_X)$。
展开来说，必须证明复合
$$
\alpha_h \circ \alpha_g :
j_1^* \circ \overline{a}_1 \to
j_1^* \circ \overline{c} \circ \overline{a}_2 \to
j_2^* \circ \overline{a}_2 \to
j_2^* \circ \overline{d} \circ \overline{a}_3 \to
j_3^* \circ \overline{a}_3
$$
与复合
$$
\alpha_{h \circ g} :
j_1^* \circ \overline{a}_1 \to
j_1^* \circ \overline{e} \circ \overline{a}_3 \to
j_3^* \circ \overline{a}_3
$$
相同。将其分成两部分。第一部分是证明图
$$
\xymatrix{
\overline{a}_1 \ar[r] \ar[d] & \overline{c} \circ \overline{a}_2 \ar[d] \\
\overline{e} \circ \overline{a}_3 \ar[r] &
\overline{c} \circ \overline{d} \circ \overline{a}_3
}
$$
交换，其中下方水平箭头来自识别
$\overline{e} = \overline{c} \circ \overline{d}$。这是因为相应的总导出正像函子图
$$
\xymatrix{
R\overline{f}_{1, *} \ar[r] \ar[d] & Rg_* \circ R\overline{f}_{2, *} \ar[d] \\
R(h \circ g)_* \circ R\overline{f}_{3, *} \ar[r] &
Rg_* \circ Rh_* \circ R\overline{f}_{3, *}
}
$$
交换（此处插入将来引用）。第二部分是证明复合
$$
j_1^* \circ \overline{c} \circ \overline{d} \to
j_2^* \circ \overline{d} \to j_3^*
$$
通过识别 $\overline{e} = \overline{c} \circ \overline{d}$ 等于映射
$$
j_1^* \circ \overline{e} \to j_3^*
$$
这已在引理 \ref{duality-lemma-compose-base-change-maps} 中证明
（注意，本情形中该引理的态射 $f', g'$ 等于 $\text{id}_X$）。
\end{proof}

\begin{lemma}
\label{duality-lemma-upper-shriek-composition}
在情形 \ref{duality-situation-shriek} 中，设 $f : X \to Y$ 和
$g : Y \to Z$ 为 $\textit{FTS}_S$ 中可复合的态射。则存在典范同构
$(g \circ f)^! \to f^! \circ g^!$。
\end{lemma}

\begin{proof}
选取 $Y$ 在 $Z$ 上的紧化 $i : Y \to \overline{Y}$，再选取
$X$ 在 $\overline{Y}$ 上的紧化 $X \to \overline{X}$。这里两次使用了
《平坦性进阶》，定理 \ref{flat-theorem-nagata} 和引理
\ref{flat-lemma-compactifyable}。令 $\overline{a}$ 为引理
\ref{duality-lemma-twisted-inverse-image} 中对应于
$\overline{X} \to \overline{Y}$ 的右伴随，令 $\overline{b}$ 为引理
\ref{duality-lemma-twisted-inverse-image} 中对应于 $\overline{Y} \to Z$ 的右伴随。
则 $\overline{a} \circ \overline{b}$ 是引理
\ref{duality-lemma-twisted-inverse-image} 中对应于复合
$\overline{X} \to Z$ 的右伴随。因此
$g^! = i^* \circ \overline{b}$，且
$(g \circ f)^! = (X \to \overline{X})^* \circ \overline{a} \circ \overline{b}$。
令 $U$ 为 $Y$ 在 $\overline{X}$ 中的逆像，从而得到交换图
$$
\xymatrix{
X \ar[r]_j \ar[d] & U \ar[dl] \ar[r]_{j'} & \overline{X} \ar[dl] \\
Y \ar[r]_i \ar[d] & \overline{Y} \ar[dl] \\
Z
}
$$
令 $\overline{a}'$ 为引理 \ref{duality-lemma-twisted-inverse-image} 中对应于
$U \to Y$ 的右伴随。则 $f^! = j^* \circ \overline{a}'$。
由 (\ref{duality-equation-sheafy}) 得到
$$
\gamma : (j')^* \circ \overline{a} \to \overline{a}' \circ i^*
$$
并用它定义
$$
(g \circ f)^! =
(j' \circ j)^* \circ \overline{a} \circ \overline{b} =
j^* \circ (j')^* \circ \overline{a} \circ \overline{b}
\to
j^* \circ \overline{a}' \circ i^* \circ \overline{b} =
f^! \circ g^!
$$
由引理 \ref{duality-lemma-proper-noetherian}，它在
$D_\QCoh^+(\mathcal{O}_Z)$ 的对象上是同构。为完成证明，需说明此同构与所作选择无关。

\medskip\noindent
假设有两个图
$$
\vcenter{
\xymatrix{
X \ar[r]_{j_1} \ar[d] & U_1 \ar[dl] \ar[r]_{j'_1} & \overline{X}_1 \ar[dl] \\
Y \ar[r]_{i_1} \ar[d] & \overline{Y}_1 \ar[dl] \\
Z
}
}
\quad\text{且}\quad
\vcenter{
\xymatrix{
X \ar[r]_{j_2} \ar[d] & U_2 \ar[dl] \ar[r]_{j'_2} & \overline{X}_2 \ar[dl] \\
Y \ar[r]_{i_2} \ar[d] & \overline{Y}_2 \ar[dl] \\
Z
}
}
$$
首先可选取 $Y$ 在 $Z$ 上、具有稠密像且同时支配
$\overline{Y}_1$ 和 $\overline{Y}_2$ 的紧化
$i : Y \to \overline{Y}$；见《平坦性进阶》，引理
\ref{flat-lemma-compactifications-cofiltered}。由《平坦性进阶》，引理
\ref{flat-lemma-right-multiplicative-system} 及《范畴》，引理
\ref{categories-lemma-morphisms-right-localization} 和
\ref{categories-lemma-equality-morphisms-right-localization}，可选取
$X$ 在 $\overline{Y}$ 上、具有稠密像的紧化
$X \to \overline{X}$，以及态射
$\overline{X} \to \overline{X}_1$ 和
$\overline{X} \to \overline{X}_2$，使得复合
$\overline{X} \to \overline{Y} \to \overline{Y}_1$ 等于复合
$\overline{X} \to \overline{X}_1 \to \overline{Y}_1$，并使复合
$\overline{X} \to \overline{Y} \to \overline{Y}_2$ 等于复合
$\overline{X} \to \overline{X}_2 \to \overline{Y}_2$。
因此，只需在有如下交换图时比较这些图所确定的映射：
$$
\xymatrix{
X \ar[rr]_{j_1} \ar@{=}[d] & &
U_1 \ar[d] \ar[ddll] \ar[rr]_{j'_1} & &
\overline{X}_1 \ar[d] \ar[ddll] \\
X \ar'[r][rr]^-{j_2} \ar[d] & &
U_2 \ar'[dl][ddll] \ar'[r][rr]^-{j'_2} & &
\overline{X}_2 \ar[ddll] \\
Y \ar[rr]^{i_1} \ar@{=}[d] & & \overline{Y}_1 \ar[d] \\
Y \ar[rr]^{i_2} \ar[d] & & \overline{Y}_2 \ar[dll] \\
Z
}
$$
并且紧化 $X \to \overline{X}_1$ 和
$Y \to \overline{Y}_2$ 具有稠密像。分别用
$\overline{a}_i$、$\overline{a}'_i$、$\overline{c}$ 和
$\overline{c}'$ 表示引理 \ref{duality-lemma-twisted-inverse-image} 中对应于
$\overline{X}_i \to \overline{Y}_i$、$U_i \to Y$、
$\overline{X}_1 \to \overline{X}_2$ 和 $U_1 \to U_2$ 的右伴随。
以下每个方块
$$
\xymatrix{
X \ar[r] \ar[d] \ar@{}[dr]|A & U_1 \ar[d] \\
X \ar[r] & U_2
}
\quad
\xymatrix{
U_2 \ar[r] \ar[d] \ar@{}[dr]|B & \overline{X}_2 \ar[d] \\
Y \ar[r] & \overline{Y}_2
}
\quad
\xymatrix{
U_1 \ar[r] \ar[d] \ar@{}[dr]|C & \overline{X}_1 \ar[d] \\
Y \ar[r] & \overline{Y}_1
}
\quad
\xymatrix{
Y \ar[r] \ar[d] \ar@{}[dr]|D & \overline{Y}_1 \ar[d] \\
Y \ar[r] & \overline{Y}_2
}
\quad
\xymatrix{
X \ar[r] \ar[d] \ar@{}[dr]|E & \overline{X}_1 \ar[d] \\
X \ar[r] & \overline{X}_2
}
$$
都是笛卡尔的（对 A、D、E，见《平坦性进阶》，引理
\ref{flat-lemma-compactifications-cofiltered} 第 (c) 部分；对 B、C，回忆
$U_i$ 是 $Y$ 沿 $\overline{X}_i \to \overline{Y}_i$ 的逆像），
因而分别产生如下基变换映射 (\ref{duality-equation-sheafy})：
$$
\begin{matrix}
\gamma_A : j_1^* \circ \overline{c}' \to j_2^* &
\gamma_B : (j_2')^* \circ \overline{a}_2 \to \overline{a}'_2 \circ i_2^* &
\gamma_C : (j_1')^* \circ \overline{a}_1 \to \overline{a}'_1 \circ i_1^* \\
\gamma_D : i_1^* \circ \overline{d} \to i_2^* &
\gamma_E : (j'_1 \circ j_1)^* \circ \overline{c} \to (j'_2 \circ j_2)^*
\end{matrix}
$$
记 $f_1^! = j_1^* \circ \overline{a}'_1$、
$f_2^! = j_2^* \circ \overline{a}'_2$、
$g_1^! = i_1^* \circ \overline{b}_1$、
$g_2^! = i_2^* \circ \overline{b}_2$、
$(g \circ f)_1^! =
(j_1' \circ j_1)^* \circ \overline{a}_1 \circ \overline{b}_1$，以及
$(g \circ f)^!_2 =
(j_2' \circ j_2)^* \circ \overline{a}_2 \circ \overline{b}_2$。
证明第一段及引理 \ref{duality-lemma-shriek-well-defined} 中给出的构造使用
\begin{enumerate}
\item $\gamma_C$ 作为映射 $(g \circ f)^!_1 \to f_1^! \circ g_1^!$；
\item $\gamma_B$ 作为映射 $(g \circ f)^!_2 \to f_2^! \circ g_2^!$；
\item $\gamma_A$ 作为映射 $f_1^! \to f_2^!$；
\item $\gamma_D$ 作为映射 $g_1^! \to g_2^!$；并且
\item $\gamma_E$ 作为映射 $(g \circ f)^!_1 \to (g \circ f)^!_2$。
\end{enumerate}
必须证明图
$$
\xymatrix{
(g \circ f)^!_1 \ar[r]_{\gamma_E} \ar[d]_{\gamma_C} &
(g \circ f)^!_2 \ar[d]_{\gamma_B} \\
f_1^! \circ g_1^! \ar[r]^{\gamma_A \circ \gamma_D} & f_2^! \circ g_2^!
}
$$
交换。 将使用引理 \ref{duality-lemma-compose-base-change-maps}、
\ref{duality-lemma-compose-base-change-maps-horizontal}，以及注
\ref{duality-remark-going-around} 中的记号滥用（特别地，从记号中省略与恒等变换的
$\star$ 积）。可写成 $\gamma_E = \gamma_A \circ \gamma_F$，其中
$$
\xymatrix{
U_1 \ar[r] \ar[d] \ar@{}[rd]|F & \overline{X}_1 \ar[d] \\
U_2 \ar[r] & \overline{X}_2
}
$$
于是
$$
\gamma_B \circ \gamma_E = \gamma_B \circ \gamma_A  \circ \gamma_F
= \gamma_A \circ \gamma_B \circ \gamma_F
$$
最后一个等式成立，是因为方块 $A$ 与 $B$ 只在一个点相交
（类似于注 \ref{duality-remark-going-around} 中最后一个论证）。因此只需证明
$\gamma_D \circ \gamma_C = \gamma_B \circ \gamma_F$。由于二者都等于方块
$$
\xymatrix{
U_1 \ar[r] \ar[d] & \overline{X}_1 \ar[d] \\
Y \ar[r] & \overline{Y}_2
}
$$
的映射 (\ref{duality-equation-sheafy})，结论成立。
\end{proof}

\begin{lemma}
\label{duality-lemma-pseudo-functor}
在情形 \ref{duality-situation-shriek} 中，引理
\ref{duality-lemma-shriek-well-defined} 和
\ref{duality-lemma-upper-shriek-composition} 的构造定义了从范畴
$\textit{FTS}_S$ 到范畴的 $2$-范畴的伪函子（见《范畴》，定义
\ref{categories-definition-functor-into-2-category}）。
\end{lemma}

\begin{proof}
为证明这一点，必须证明对态射
$f : X \to Y$、$g : Y \to Z$、$h : Z \to T$，图
$$
\xymatrix{
(h \circ g \circ f)^! \ar[r]_{\gamma_{A + B}} \ar[d]_{\gamma_{B + C}} &
f^! \circ (h \circ g)^! \ar[d]^{\gamma_C} \\
(g \circ f)^! \circ h^! \ar[r]^{\gamma_A} & f^! \circ g^! \circ h^!
}
$$
交换（$\gamma$ 的含义见下文）。为此，先选取 $Z$ 在 $T$ 上的紧化
$\overline{Z}$，再选取 $Y$ 在 $\overline{Z}$ 上的紧化
$\overline{Y}$，最后选取 $X$ 在 $\overline{Y}$ 上的紧化
$\overline{X}$。这里使用《平坦性进阶》，定理
\ref{flat-theorem-nagata} 和引理 \ref{flat-lemma-compactifyable}。
令 $W \subset \overline{Y}$ 为 $Z$ 沿
$\overline{Y} \to \overline{Z}$ 的逆像，并令
$U \subset V \subset \overline{X}$ 为 $Y \subset W$ 沿
$\overline{X} \to \overline{Y}$ 的逆像。得到图
$$
\xymatrix{
X \ar[d]_f \ar[r] & U \ar[r] \ar[d] \ar@{}[dr]|A &
V \ar[d] \ar[r] \ar@{}[rd]|B & \overline{X} \ar[d] \\
Y \ar[d]_g \ar[r] & Y \ar[r] \ar[d] & W \ar[r] \ar[d] \ar@{}[rd]|C &
\overline{Y} \ar[d] \\
Z \ar[d]_h \ar[r] & Z \ar[d] \ar[r] & Z \ar[d] \ar[r] & \overline{Z} \ar[d] \\
T \ar[r] & T \ar[r] & T \ar[r] & T
}
$$
不引入大量记号，而完全像引理
\ref{duality-lemma-upper-shriek-composition} 的证明那样论证，可知第一个图中的映射
按所标示方式使用矩形 $A + B$、$B + C$、$A$、$C$ 的映射
(\ref{duality-equation-sheafy})。由引理
\ref{duality-lemma-compose-base-change-maps} 和
\ref{duality-lemma-compose-base-change-maps-horizontal}，有
$\gamma_{A + B} = \gamma_A \circ \gamma_B$ 且
$\gamma_{B + C} = \gamma_C \circ \gamma_B$，故只要
$\gamma_A \circ \gamma_C = \gamma_C \circ \gamma_A$，所需等式即成立。
这确实成立，因为方块 $A$ 与 $C$ 只在一个点相交
（类似于注 \ref{duality-remark-going-around} 中最后一个论证）。
\end{proof}

\begin{lemma}
\label{duality-lemma-map-pullback-to-shriek-well-defined}
在情形 \ref{duality-situation-shriek} 中，设
$f : X \to Y$ 是 $\textit{FTS}_S$ 的态射。存在关于
$D^+_\QCoh(\mathcal{O}_Y)$ 中的 $K$ 具有函子性的典范映射
$$
\mu_{f, K} :
Lf^*K \otimes_{\mathcal{O}_X}^\mathbf{L} f^!\mathcal{O}_Y
\longrightarrow
f^!K
$$
。若 $g : Y \to Z$ 是 $\textit{FTS}_S$ 的另一个态射，则对所有
$K \in D^+_\QCoh(\mathcal{O}_Z)$，图
$$
\xymatrix{
Lf^*(Lg^*K \otimes_{\mathcal{O}_Y}^\mathbf{L} g^!\mathcal{O}_Z)
\otimes_{\mathcal{O}_X}^\mathbf{L} f^!\mathcal{O}_Y
\ar@{=}[d] \ar[r]_-{\mu_f} &
f^!(Lg^*K \otimes_{\mathcal{O}_Y}^\mathbf{L} g^!\mathcal{O}_Z)
\ar[r]_-{f^!\mu_g} &
f^!g^!K \ar@{=}[d] \\
Lf^*Lg^*K \otimes_{\mathcal{O}_X}^\mathbf{L} Lf^* g^!\mathcal{O}_Z
\otimes_{\mathcal{O}_X}^\mathbf{L} f^!\mathcal{O}_Y \ar[r]^-{\mu_f} &
Lf^*Lg^*K \otimes_{\mathcal{O}_X}^\mathbf{L} f^!g^!\mathcal{O}_Z
\ar[r]^-{\mu_{g \circ f}} & f^!g^!K
}
$$
交换。
\end{lemma}

\begin{proof}
若 $f$ 是真态射，则 $f^! = a$，可使用
(\ref{duality-equation-compare-with-pullback})；若 $g$ 也是真态射，则引理
\ref{duality-lemma-transitivity-compare-with-pullback}（在更一般的情形下）证明该图交换。

\medskip\noindent
先定义映射 $\mu_{f, K}$。选取 $X$ 在 $Y$ 上的一个紧化
$j : X \to \overline{X}$。由于 $f^!$ 定义为
$j^* \circ \overline{a}$，将映射 (\ref{duality-equation-compare-with-pullback})
$$
L\overline{f}^*K \otimes_{\mathcal{O}_{\overline{X}}}^\mathbf{L}
\overline{a}(\mathcal{O}_Y)
\longrightarrow
\overline{a}(K)
$$
限制到 $X$，便得到 $\mu_{f, K}$。为说明它与紧化的选择无关，像引理
\ref{duality-lemma-shriek-well-defined} 的证明那样论证。建议读者先阅读该引理的证明。

\medskip\noindent
设给定 $Y$ 上两个紧化 $j_i : X \to \overline{X}_i$ 之间的态射
$g : \overline{X}_1 \to \overline{X}_2$，且
$g^{-1}(j_2(X)) = j_1(X)$。记 $\overline{c}$ 为引理
\ref{duality-lemma-twisted-inverse-image} 对态射 $g$ 给出的推前的右伴随。映射
$$
L\overline{f}_1^*K \otimes_{\mathcal{O}_{\overline{X}}}^\mathbf{L}
\overline{a}_1(\mathcal{O}_Y)
\longrightarrow
\overline{a}_1(K)
\quad\text{且}\quad
L\overline{f}_2^*K \otimes_{\mathcal{O}_{\overline{X}}}^\mathbf{L}
\overline{a}_2(\mathcal{O}_Y)
\longrightarrow
\overline{a}_2(K)
$$
嵌入交换图
$$
\xymatrix{
Lg^*(L\overline{f}_2^*K \otimes^\mathbf{L}
\overline{a}_2(\mathcal{O}_Y))
\otimes^\mathbf{L} \overline{c}(\mathcal{O}_{\overline{X}_2})
\ar@{=}[d] \ar[r]_-\sigma &
\overline{c}(L\overline{f}_2^*K \otimes^\mathbf{L}
\overline{a}_2(\mathcal{O}_Y)) \ar[r] &
\overline{c}(\overline{a}_2(K)) \ar@{=}[d] \\
L\overline{f}_1^*K \otimes^\mathbf{L} Lg^*\overline{a}_2(\mathcal{O}_Y)
\otimes^\mathbf{L} \overline{c}(\mathcal{O}_{\overline{X}_2})
\ar[r]^-{1 \otimes \tau} &
L\overline{f}_1^*K \otimes^\mathbf{L} \overline{a}_1(\mathcal{O}_Y) \ar[r] &
\overline{a}_1(K)
}
$$
，这由引理 \ref{duality-lemma-transitivity-compare-with-pullback} 得到。由引理
\ref{duality-lemma-compare-on-open}，映射 $\sigma$ 和 $\tau$ 限制到 $X$ 后是同构。
事实上还可说得更多。回忆在引理 \ref{duality-lemma-shriek-well-defined} 的证明中，
用映射 (\ref{duality-equation-sheafy})
$\gamma : j_1^* \circ \overline{c} \to j_2^*$ 构造了同构
$\alpha_g : j_1^* \circ \overline{a}_1 \to j_2^* \circ \overline{a}_2$。
用 $j_1$ 拉回映射 $\sigma$；若通过
$j_1^* \circ \overline{c} \to j_2^*$ 将
$j_1^*\overline{c}(\mathcal{O}_{\overline{X}_2})$ 与 $\mathcal{O}_X$ 等同，
则所得映射是
$j_2^*\left(L\overline{f}_2^*K \otimes^\mathbf{L}
\overline{a}_2(\mathcal{O}_Y)\right)$ 上的恒等映射；见引理
\ref{duality-lemma-restriction-compare-with-pullback}。类似地，映射
$\tau : Lg^*\overline{a}_2(\mathcal{O}_Y)
\otimes^\mathbf{L} \overline{c}(\mathcal{O}_{\overline{X}_2}) \to
\overline{a}_1(\mathcal{O}_Y) = \overline{c}(\overline{a}_2(\mathcal{O}_Y))$
拉回后，是 $j_2^*\overline{a}_2(\mathcal{O}_Y)$ 上的恒等映射。
所以用 $j_1$ 拉回，并在可以之处应用 $\gamma$，得到交换图
$$
\xymatrix{
j_2^*\left(L\overline{f}_2^*K \otimes^\mathbf{L}
\overline{a}_2(\mathcal{O}_Y)\right) \ar[r] \ar[d] &
j_2^*\overline{a}_2(K) \\
j_1^*L\overline{f}_1^*K \otimes^\mathbf{L} j_2^*\overline{a}_2(\mathcal{O}_Y) &
j_1^*(L\overline{f}_1^*K \otimes^\mathbf{L} \overline{a}_1(\mathcal{O}_Y))
\ar[r] \ar[l]_{1 \otimes \alpha_g} &
j_1^* \overline{a}_1(K) \ar[lu]_{\alpha_g}
}
$$
该图的交换性恰好说明：用紧化 $\overline{X}_1$ 构造的
$\mu_{f, K}$，经由引理 \ref{duality-lemma-shriek-well-defined} 的证明所用的等同
$\alpha_g$，与用紧化 $\overline{X}_2$ 构造的 $\mu_{f, K}$ 相同。
再完全像引理 \ref{duality-lemma-shriek-well-defined} 的证明那样使用若干范畴论论证，
便知 $\mu_{f, K}$ 定义良好（略去一个小细节）。

\medskip\noindent
由此，引理陈述中图的交换性来自同构
$(g \circ f)^! \to f^! \circ g^!$ 的构造（即引理
\ref{duality-lemma-upper-shriek-composition} 证明的第一部分，其中使用
$\overline{X} \to \overline{Y} \to Z$），以及引理
\ref{duality-lemma-transitivity-compare-with-pullback} 对
$\overline{X} \to \overline{Y} \to Z$ 的结论。
\end{proof}



\section{上叹号函子的性质}
\label{duality-section-upper-shriek-properties}

\noindent
下面给出上叹号函子的一些性质。

\begin{lemma}
\label{duality-lemma-shriek-open-immersion}
在情形 \ref{duality-situation-shriek} 中，设 $Y$ 是
$\textit{FTS}_S$ 的对象，且 $j : X \to Y$ 是开浸入。
则存在函子的典范同构 $j^! = j^*$。
\end{lemma}

\noindent
对 $\textit{FTS}_S$ 的 \'etale 态射 $f : X \to Y$，
也有 $f^* \cong f^!$；见引理 \ref{duality-lemma-shriek-etale}。

\begin{proof}
在此情形下，可选取 $\overline{X} = Y$ 作为紧化。
引理 \ref{duality-lemma-twisted-inverse-image} 对
$\text{id} : Y \to Y$ 给出的右伴随是恒等函子，故由定义
$j^! = j^*$。
\end{proof}

\begin{lemma}
\label{duality-lemma-restrict-before-or-after}
在情形 \ref{duality-situation-shriek} 中，设
$$
\xymatrix{
U \ar[r]_j \ar[d]_g & X \ar[d]^f \\
V \ar[r]^{j'} & Y
}
$$
是 $\textit{FTS}_S$ 中的交换图，其中 $j$ 和 $j'$ 是开浸入。
则作为函子
$D^+_\QCoh(\mathcal{O}_Y) \to D^+(\mathcal{O}_U)$，有
$j^* \circ f^! = g^! \circ (j')^*$。
\end{lemma}

\begin{proof}
令 $h = f \circ j = j' \circ g$。由引理
\ref{duality-lemma-upper-shriek-composition}，有
$h^! = j^! \circ f^! = g^! \circ (j')^!$。由引理
\ref{duality-lemma-shriek-open-immersion}，有
$j^! = j^*$ 且 $(j')^! = (j')^*$。
\end{proof}

\begin{lemma}
\label{duality-lemma-shriek-affine-line}
在情形 \ref{duality-situation-shriek} 中，设 $Y$ 是 $\textit{FTS}_S$ 的对象，
且 $f : X = \mathbf{A}^1_Y \to Y$ 是投影。则存在函子的一个
（非典范）同构 $f^!(-) \cong Lf^*(-) [1]$。
\end{lemma}

\begin{proof}
由于 $X = \mathbf{A}^1_Y \subset \mathbf{P}^1_Y$，且
$\mathcal{O}_{\mathbf{P}^1_Y}(-2)|_X \cong \mathcal{O}_X$，
结论由引理 \ref{duality-lemma-upper-shriek-P1} 和
\ref{duality-lemma-compare-with-pullback-flat-proper-noetherian} 得到。
\end{proof}

\begin{lemma}
\label{duality-lemma-shriek-closed-immersion}
在情形 \ref{duality-situation-shriek} 中，设 $Y$ 是
$\textit{FTS}_S$ 的对象，且 $i : X \to Y$ 是闭浸入。
则存在函子的典范同构
$i^!(-) = R\SheafHom(\mathcal{O}_X, -)$。
\end{lemma}

\begin{proof}
这只是引理 \ref{duality-lemma-twisted-inverse-image-closed} 的改述。
\end{proof}

\begin{remark}[上叹号的局部描述]
\label{duality-remark-local-calculation-shriek}
在情形 \ref{duality-situation-shriek} 中，设 $f : X \to Y$ 是
$\textit{FTS}_S$ 的态射。利用上述引理，可以如下局部计算 $f^!$。
设给定仿射开子集
$$
\xymatrix{
U \ar[r]_j \ar[d]_g & X \ar[d]^f \\
V \ar[r]^i & Y
}
$$
由于 $j^! \circ f^! = g^! \circ i^!$
（引理 \ref{duality-lemma-upper-shriek-composition}），且 $j^!$ 和 $i^!$
由限制给出（引理 \ref{duality-lemma-shriek-open-immersion}），可知对任意
$E \in D^+_\QCoh(\mathcal{O}_X)$，有
$$
(f^!E)|_U = g^!(E|_V)
$$
。写成 $U = \Spec(A)$、$V = \Spec(R)$，并令
$\varphi : R \to A$ 为对应于 $g$ 的有限型环同态。选取一个表示
$A = P/I$，其中 $P = R[x_1, \ldots, x_n]$ 是 $R$ 上 $n$ 个变量的
多项式代数。选取与 $E|_V$ 对应的对象 $K \in D^+(R)$
（《概形的导出范畴》，引理
\ref{perfect-lemma-affine-compare-bounded}）。那么断言 $f^!E|_U$ 对应于
$$
\varphi^!(K) = R\Hom(A, K \otimes_R^\mathbf{L} P)[n]
$$
，其中 $R\Hom(A, -) : D(P) \to D(A)$ 是《对偶复形》，第
\ref{dualizing-section-trivial} 节的函子，而
$\varphi^! : D(R) \to D(A)$ 是《对偶复形》，第
\ref{dualizing-section-relative-dualizing-complex-algebraic} 节的函子。
事实上，所选表示给出分解
$$
U \rightarrow \mathbf{A}^n_V \to \mathbf{A}^{n - 1}_V \to \ldots \to
\mathbf{A}^1_V \to V
$$
。恰好应用引理 \ref{duality-lemma-shriek-affine-line} $n$ 次，可知
$(\mathbf{A}^n_V \to V)^!(E|_V)$ 对应于
$K \otimes_R^\mathbf{L} P[n]$。由引理
\ref{duality-lemma-sheaf-with-exact-support-quasi-coherent} 和
\ref{duality-lemma-shriek-closed-immersion}，最后一步对应于应用
$R\Hom(A, -)$。
\end{remark}

\begin{lemma}
\label{duality-lemma-shriek-coherent}
在情形 \ref{duality-situation-shriek} 中，设 $f : X \to Y$ 是
$\textit{FTS}_S$ 的态射。则 $f^!$ 将
$D_{\textit{Coh}}^+(\mathcal{O}_Y)$ 映入
$D_{\textit{Coh}}^+(\mathcal{O}_X)$。
\end{lemma}

\begin{proof}
这个问题在 $X$ 上是局部的，故可设 $X$ 和 $Y$ 是仿射概形。
此时可将 $f : X \to Y$ 分解为
$$
X \xrightarrow{i} \mathbf{A}^n_Y \to \mathbf{A}^{n - 1}_Y \to \ldots \to
\mathbf{A}^1_Y \to Y
$$
，其中 $i$ 是闭浸入。该引理由以下各项推出：由引理
\ref{duality-lemma-shriek-affine-line}、
\ref{duality-lemma-sheaf-with-exact-support-coherent}、《对偶复形》，引理
\ref{dualizing-lemma-dualizing-polynomial-ring}，以及归纳法。
\end{proof}

\begin{lemma}
\label{duality-lemma-shriek-dualizing}
在情形 \ref{duality-situation-shriek} 中，设 $f : X \to Y$ 是
$\textit{FTS}_S$ 的态射。若 $K$ 是 $Y$ 的对偶复形，
则 $f^!K$ 是 $X$ 的对偶复形。
\end{lemma}

\begin{proof}
这个问题在 $X$ 上是局部的，故可设 $X$ 和 $Y$ 是仿射概形。
此时可将 $f : X \to Y$ 分解为
$$
X \xrightarrow{i} \mathbf{A}^n_Y \to \mathbf{A}^{n - 1}_Y \to \ldots \to
\mathbf{A}^1_Y \to Y
$$
，其中 $i$ 是闭浸入。由引理 \ref{duality-lemma-shriek-affine-line}、
《对偶复形》，引理 \ref{dualizing-lemma-dualizing-polynomial-ring}
以及归纳法，可知 $p^!K$ 是 $\mathbf{A}^n_Y$ 上的对偶复形，
其中 $p : \mathbf{A}^n_Y \to Y$ 是投影。类似地，由《对偶复形》，引理
\ref{dualizing-lemma-dualizing-quotient} 以及引理
\ref{duality-lemma-sheaf-with-exact-support-quasi-coherent} 和
\ref{duality-lemma-shriek-closed-immersion}，可知 $i^!$ 将对偶复形变为对偶复形。
\end{proof}

\begin{lemma}
\label{duality-lemma-shriek-via-duality}
在情形 \ref{duality-situation-shriek} 中，设 $f : X \to Y$ 是
$\textit{FTS}_S$ 的态射。设 $K$ 是 $Y$ 上的对偶复形。
对 $M \in D_{\textit{Coh}}(\mathcal{O}_Y)$，令
$D_Y(M) = R\SheafHom_{\mathcal{O}_Y}(M, K)$；对
$E \in D_{\textit{Coh}}(\mathcal{O}_X)$，令
$D_X(E) = R\SheafHom_{\mathcal{O}_X}(E, f^!K)$。
则对 $M \in D_{\textit{Coh}}^+(\mathcal{O}_Y)$，存在典范同构
$$
f^!M \longrightarrow D_X(Lf^*D_Y(M))
$$
。
\end{lemma}

\begin{proof}
选取 $X$ 在 $Y$ 上的紧化 $j : X \subset \overline{X}$
（《平坦性进阶》，定理 \ref{flat-theorem-nagata} 和引理
\ref{flat-lemma-compactifyable}）。令 $a$ 为引理
\ref{duality-lemma-twisted-inverse-image} 对 $\overline{X} \to Y$ 给出的右伴随。
对 $E \in D_{\textit{Coh}}(\mathcal{O}_{\overline{X}})$，令
$D_{\overline{X}}(E) = R\SheafHom_{\mathcal{O}_{\overline{X}}}(E, a(K))$。
由于 $R\SheafHom$ 的形成与限制到开集可交换，且
$f^! = j^* \circ a$，故只需证明：对
$M \in D_{\textit{Coh}}(\mathcal{O}_Y)$，存在典范同构
$$
a(M) \longrightarrow D_{\overline{X}}(L\overline{f}^*D_Y(M))
$$
。对 $F \in D_\QCoh(\mathcal{O}_X)$，有
\begin{align*}
\Hom_{\overline{X}}(
F, D_{\overline{X}}(L\overline{f}^*D_Y(M)))
& =
\Hom_{\overline{X}}(
F \otimes_{\mathcal{O}_X}^\mathbf{L} L\overline{f}^*D_Y(M), a(K)) \\
& =
\Hom_Y(
R\overline{f}_*(F \otimes_{\mathcal{O}_X}^\mathbf{L} L\overline{f}^*D_Y(M)),
K) \\
& =
\Hom_Y(
R\overline{f}_*(F) \otimes_{\mathcal{O}_Y}^\mathbf{L} D_Y(M),
K) \\
& =
\Hom_Y(
R\overline{f}_*(F), D_Y(D_Y(M))) \\
& =
\Hom_Y(R\overline{f}_*(F), M) \\
& = \Hom_{\overline{X}}(F, a(M))
\end{align*}
第一个等式由《上同调》，引理 \ref{cohomology-lemma-internal-hom} 得到。
第二个等式来自 $a$ 的定义。第三个等式由《概形的导出范畴》，引理
\ref{perfect-lemma-cohomology-base-change} 得到。第四个等式由《上同调》，
引理 \ref{cohomology-lemma-internal-hom} 和 $D_Y$ 的定义得到。
第五个等式由引理 \ref{duality-lemma-dualizing-schemes} 得到。
最后一个等式来自 $a$ 的定义。因此由 Yoneda 引理可知
$a(M) = D_{\overline{X}}(L\overline{f}^*D_Y(M))$。
\end{proof}

\begin{lemma}
\label{duality-lemma-perfect-comparison-shriek}
在情形 \ref{duality-situation-shriek} 中，设 $f : X \to Y$ 是
$\textit{FTS}_S$ 的态射。假设 $f$ 是完美态射（例如平坦态射）。则
\begin{enumerate}
\item[(a)] $f^!\mathcal{O}_Y$ 属于 $D_{\textit{Coh}}^b(\mathcal{O}_X)$；
\item[(b)] $f^!\mathcal{O}_Y$ 在 $D(f^{-1}\mathcal{O}_Y)$ 中具有有限 Tor 维数；
\item[(c)] $\mathcal{O}_X \to
R\SheafHom_{\mathcal{O}_X}(f^!\mathcal{O}_Y, f^!\mathcal{O}_Y)$
是同构；
\item[(d)] $f^!$ 将 $D_{\textit{Coh}}^b(\mathcal{O}_Y)$ 映入
$D_{\textit{Coh}}^b(\mathcal{O}_X)$；
\item[(e)] 引理 \ref{duality-lemma-map-pullback-to-shriek-well-defined} 的映射
$\mu_{f,  K} :
Lf^*K \otimes_{\mathcal{O}_X}^\mathbf{L} f^!\mathcal{O}_Y
\to
f^!K$
对所有 $K \in D_\QCoh^+(\mathcal{O}_Y)$ 都是同构。
\end{enumerate}
\end{lemma}

\begin{proof}
（有限表示的平坦态射是完美态射；见《态射进阶》，引理
\ref{more-morphisms-lemma-flat-finite-presentation-perfect}。）
断言 (a)、(b)、(c) 在 $X$ 上是局部的，故可设 $X$ 和 $Y$ 是仿射的。
此时注 \ref{duality-remark-local-calculation-shriek} 将 (a)、(b)、(c) 分别化为
《对偶复形》，引理 \ref{dualizing-lemma-relative-dualizing-algebraic}
的 (1)(a)、(1)(b)、(1)(c)。（使用《概形的导出范畴》，引理
\ref{perfect-lemma-tor-dimension-rel-affine}、
\ref{perfect-lemma-pseudo-coherent}、
\ref{perfect-lemma-identify-pseudo-coherent-noetherian}、
\ref{perfect-lemma-quasi-coherence-internal-hom} 以使各断言的表述相符。）

\medskip\noindent
为证明 (d) 和 (e)，先作一系列预备说明。
\begin{enumerate}
\item 已知 $f^!$ 将 $D_{\textit{Coh}}^+(\mathcal{O}_Y)$ 映入
$D_{\textit{Coh}}^+(\mathcal{O}_X)$；见引理
\ref{duality-lemma-shriek-coherent}。
\item 若 $f$ 是开浸入，则 $f^! = f^*$，且 (d)、(e) 成立，
因为在 $f^!$ 和 $\mu_f$ 的构造中可取 $\overline{X} = Y$；见引理
\ref{duality-lemma-shriek-open-immersion}。
\item 若 $f$ 是完美真态射，则由引理
\ref{duality-lemma-compare-with-pullback-flat-proper-noetherian}，(e) 成立。
\item 若存在开覆盖 $X = \bigcup U_i$，且 (d) 对 $U_i \to Y$ 成立，
则 (d) 对 $X \to Y$ 成立；(e) 也一样。这是因为 $f^!$ 和 $\mu_f$
的构造与转到开子概形可交换。
\item 若 $g : Y \to Z$ 是 $\textit{FTS}_S$ 中的第二个完美态射，
且 (e) 对 $f$ 和 $g$ 成立，则
$f^!g^!\mathcal{O}_Z =
Lf^*g^!\mathcal{O}_Z \otimes_{\mathcal{O}_X}^\mathbf{L} f^!\mathcal{O}_Y$，
并且由引理 \ref{duality-lemma-map-pullback-to-shriek-well-defined} 的交换图，
(e) 对 $g \circ f$ 成立。
\item 若 (d)、(e) 对 $f$ 和 $g$ 都成立，则它们对 $g \circ f$ 也成立。
事实上，由上一点，$f^!g^!\mathcal{O}_Z$ 上有界，而
$L(g \circ f)^*$ 具有有限上同调维数；于是 (d) 由上面已经看到的 (e) 推出。
\end{enumerate}
由此可知，只需在 $X$ 仿射时证明 (d)、(e)。选取浸入
$X \to \mathbf{A}^n_Y$（《态射》，引理
\ref{morphisms-lemma-quasi-affine-finite-type-over-S}），并将其分解为
$X \to U \to \mathbf{A}^n_Y \to Y$，其中 $X \to U$ 是闭浸入，
$U \subset \mathbf{A}^n_Y$ 是开集。注意由《态射进阶》，引理
\ref{more-morphisms-lemma-perfect-permanence}，$X \to U$ 是完美闭浸入。
所以只需对完美闭浸入以及投影 $\mathbf{A}^n_Y \to Y$ 证明该引理。

\medskip\noindent
设 $f : X \to Y$ 是完美闭浸入。已知 (d) 成立。设
$K \in D^b_{\textit{Coh}}(\mathcal{O}_Y)$。则
$f^!K = R\SheafHom(\mathcal{O}_X, K)$
（引理 \ref{duality-lemma-shriek-closed-immersion}），且
$f_*f^!K = R\SheafHom_{\mathcal{O}_Y}(f_*\mathcal{O}_X, K)$。
因为 $f$ 是完美的，复形 $f_*\mathcal{O}_X$ 是完美的，故
$R\SheafHom_{\mathcal{O}_Y}(f_*\mathcal{O}_X, K)$ 上有界。
这证明 (d) 成立。略去若干细节。

\medskip\noindent
设 $f : \mathbf{A}^n_Y \to Y$ 是投影。反复应用引理
\ref{duality-lemma-shriek-affine-line}，可知 (d) 成立。最后，(e) 成立，
因为它对 $\mathbf{P}^n_Y \to Y$ 成立（该态射平坦且真），而
$\mathbf{A}^n_Y \subset \mathbf{P}^n_Y$ 是开集。
\end{proof}

\begin{lemma}
\label{duality-lemma-flat-shriek-relatively-perfect}
在情形 \ref{duality-situation-shriek} 中，设 $f : X \to Y$ 是
$\textit{FTS}_S$ 的态射。若 $f$ 平坦，则
$f^!\mathcal{O}_Y$ 是 $D(\mathcal{O}_X)$ 中的 $Y$-完美对象，且
$\mathcal{O}_X \to
R\SheafHom_{\mathcal{O}_X}(f^!\mathcal{O}_Y, f^!\mathcal{O}_Y)$
是同构。
\end{lemma}

\begin{proof}
两个断言在 $X$ 上都是局部的，故可设 $X$ 和 $Y$ 是仿射的。
此时注 \ref{duality-remark-local-calculation-shriek} 将该引理化为一个代数引理，
即《对偶复形》，引理 \ref{dualizing-lemma-relative-dualizing-algebraic}。
（使用《概形的导出范畴》，引理
\ref{perfect-lemma-affine-locally-rel-perfect} 以使表述相符。）
\end{proof}

\begin{lemma}
\label{duality-lemma-lci-shriek}
在情形 \ref{duality-situation-shriek} 中，设 $f : X \to Y$ 是
$\textit{FTS}_S$ 的态射。假设 $f : X \to Y$ 是局部完全交态射。则
\begin{enumerate}
\item $f^!\mathcal{O}_Y$ 是 $D(\mathcal{O}_X)$ 中的可逆对象；并且
\item $f^!$ 将完美复形映为完美复形。
\end{enumerate}
\end{lemma}

\begin{proof}
回忆局部完全交态射是完美态射；见《态射进阶》，引理
\ref{more-morphisms-lemma-lci-properties}。由引理
\ref{duality-lemma-perfect-comparison-shriek}，只需证明 $f^!\mathcal{O}_Y$
是 $D(\mathcal{O}_X)$ 中的可逆对象。这个问题在 $X$ 和 $Y$ 上是局部的，
故可设 $X \to Y$ 分解为 $X \to \mathbf{A}^n_Y \to Y$，
其中第一个箭头是 Koszul 正则浸入；见《态射进阶》，第
\ref{more-morphisms-section-lci} 节。由引理
\ref{duality-lemma-shriek-affine-line}，结论对 $\mathbf{A}^n_Y \to Y$ 成立。
因此只需在 $f$ 是 Koszul 正则浸入时证明该引理。
再次局部地工作，归约到 $X = \Spec(A)$、$Y = \Spec(B)$ 的情形，
其中 $A = B/(f_1, \ldots, f_r)$，且
$f_1, \ldots, f_r \in B$ 是正则序列（这里使用：对 Noetherian 局部环，
Koszul 正则与正则这两个概念相同；见《代数进阶》，引理
\ref{more-algebra-lemma-noetherian-finite-all-equivalent}）。
所以 $X \to Y$ 是复合
$$
X = X_r \to X_{r - 1} \to \ldots \to X_1 \to X_0 = Y
$$
，其中每个箭头都是一个有效 Cartier 除子的包含。这样便归约到有效
Cartier 除子 $i : D \to X$ 的包含情形。此时由引理
\ref{duality-lemma-compute-for-effective-Cartier}，有
$i^!\mathcal{O}_X = \mathcal{N}[1]$，证明完成。
\end{proof}







\section{上叹号的基变换}
\label{duality-section-base-change-shriek}

\noindent
在情形 \ref{duality-situation-shriek} 中，设
$$
\xymatrix{
X' \ar[r]_{g'} \ar[d]_{f'} & X \ar[d]^f \\
Y' \ar[r]^g & Y
}
$$
是 $\textit{FTS}_S$ 中的笛卡尔图，且 $X$ 和 $Y'$ 在 $Y$ 上
Tor 独立。目前的设置不足以在这种一般性下构造基变换映射
$L(g')^* \circ f^! \to (f')^! \circ Lg^*$。原因是一般无法选取
$Y$ 上的紧化 $j : X \to \overline{X}$，使得
$\overline{X}$ 和 $Y'$ 在 $Y$ 上 Tor 独立；因此第
\ref{duality-section-base-change-map} 节对基变换映射的构造不适用\footnote{
熟悉导出代数几何的读者会认识到，这并不是一个“真正的”问题。
事实上，取 $\overline{X}'$ 为 $\overline{X}$ 和 $Y'$ 在 $Y$ 上的
导出纤维积，就可以完全像引理 \ref{duality-lemma-base-change-shriek-flat}
的证明那样论证并定义该映射。毕竟，$X$ 与 $Y'$ 的 Tor 独立性保证
$X'$ 是导出概形 $\overline{X}'$ 的开子概形。}。

\medskip\noindent
第 \ref{duality-section-relative-dualizing-complexes} 节将给出一种部分补救。
也就是说，若态射 $f$ 平坦，就有合适的相对对偶复形概念，并可利用引理
\ref{duality-lemma-compactifyable-relative-dualizing}、
\ref{duality-lemma-base-change-relative-dualizing} 和
\ref{duality-lemma-perfect-comparison-shriek} 构造典范基变换同构。
如果确实需要使用这一点，将在本章后文补充精确陈述与证明。

\begin{lemma}
\label{duality-lemma-base-change-shriek-flat}
在情形 \ref{duality-situation-shriek} 中，设
$$
\xymatrix{
X' \ar[r]_{g'} \ar[d]_{f'} & X \ar[d]^f \\
Y' \ar[r]^g & Y
}
$$
是 $\textit{FTS}_S$ 中的笛卡尔图，且 $g$ 平坦。则在
$D_\QCoh^+(\mathcal{O}_Y)$ 上存在同构
$L(g')^* \circ f^! \to (f')^! \circ Lg^*$。
\end{lemma}

\begin{proof}
事实上，因为 $g$ 平坦，所以对 $X$ 在 $Y$ 上的紧化
$j : X \to \overline{X}$ 的每个选择，概形 $\overline{X}$ 都与
$Y'$ Tor 独立。记 $j' : X' \to \overline{X}'$ 为 $j$ 的基变换，
并记 $\overline{g}' : \overline{X}' \to \overline{X}$ 为投影。
将基变换映射定义为复合
$$
L(g')^* \circ f^! = L(g')^* \circ j^* \circ a =
(j')^* \circ L(\overline{g}')^* \circ a \longrightarrow
(j')^* \circ a' \circ Lg^* = (f')^! \circ Lg^*
$$
，其中中间箭头是基变换映射 (\ref{duality-equation-base-change-map})，
而 $a$ 和 $a'$ 是引理 \ref{duality-lemma-twisted-inverse-image}
分别对 $\overline{X} \to Y$ 和 $\overline{X}' \to Y'$ 给出的推前的右伴随。
此构造与紧化的选择无关（如果确实需要这一结果，将陈述一个精确引理并予以证明）。

\medskip\noindent
为完成证明，只需说明基变换映射
$L(g')^* \circ a \to a' \circ Lg^*$ 在
$D_\QCoh^+(\mathcal{O}_Y)$ 上是同构。由引理
\ref{duality-lemma-proper-noetherian}，$a$ 和 $a'$ 的形成与限制到
$Y$ 和 $Y'$ 的仿射开集可交换。因此由注
\ref{duality-remark-check-over-affines}，可设 $Y$ 和 $Y'$ 是仿射的。
因此结论由引理 \ref{duality-lemma-more-base-change} 得到。
\end{proof}

\begin{lemma}
\label{duality-lemma-shriek-etale}
在情形 \ref{duality-situation-shriek} 中，设 $f : X \to Y$ 是
$\textit{FTS}_S$ 的 \'etale 态射。则在
$D^+_\QCoh(\mathcal{O}_Y)$ 上，函子满足 $f^! \cong f^*$。
\end{lemma}

\begin{proof}
将使用如下事实：\'etale 态射是平坦的、syntomic 的，并且是局部完全交态射
（《态射》，引理 \ref{morphisms-lemma-etale-syntomic} 和
\ref{morphisms-lemma-etale-flat}，以及《态射进阶》，引理
\ref{more-morphisms-lemma-flat-lci}）。由引理
\ref{duality-lemma-perfect-comparison-shriek}，只需证明
$f^!\mathcal{O}_Y = \mathcal{O}_X$。由引理
\ref{duality-lemma-lci-shriek}，已知 $f^!\mathcal{O}_Y$ 是可逆模。考虑交换图
$$
\xymatrix{
X \times_Y X \ar[r]_{p_2} \ar[d]_{p_1} & X \ar[d]^f \\
X \ar[r]^f & Y
}
$$
和对角态射 $\Delta : X \to X \times_Y X$。由于 $\Delta$ 是开浸入
（由《态射》，引理
\ref{morphisms-lemma-diagonal-unramified-morphism} 和
\ref{morphisms-lemma-etale-smooth-unramified}），故由引理
\ref{duality-lemma-shriek-open-immersion}，有 $\Delta^! = \Delta^*$。
由引理 \ref{duality-lemma-upper-shriek-composition}，有
$\Delta^! \circ p_1^! \circ f^! = f^!$。将引理
\ref{duality-lemma-base-change-shriek-flat} 应用于该图，得到
$p_1^!\mathcal{O}_X = p_2^*f^!\mathcal{O}_Y$。因此
$$
f^!\mathcal{O}_Y = \Delta^!p_1^!f^!\mathcal{O}_Y =
\Delta^*(p_1^*f^!\mathcal{O}_Y \otimes p_1^!\mathcal{O}_X) =
\Delta^*(p_2^*f^!\mathcal{O}_Y \otimes p_1^*f^!\mathcal{O}_Y) =
(f^!\mathcal{O}_Y)^{\otimes 2}
$$
，其中第二步再次使用了引理
\ref{duality-lemma-perfect-comparison-shriek}。故如所需，
$f^!\mathcal{O}_Y = \mathcal{O}_X$。
\end{proof}


\noindent
在本节余下部分，将陈述若干容易证明的结果；它们应当是良好的基变换映射理论的推论。

\begin{lemma}[权宜基变换]
\label{duality-lemma-base-change-locally}
在情形 \ref{duality-situation-shriek} 中，设
$$
\xymatrix{
X' \ar[r]_{g'} \ar[d]_{f'} & X \ar[d]^f \\
Y' \ar[r]^g & Y
}
$$
是 $\textit{FTS}_S$ 中的笛卡尔图。设
$E \in D^+_\QCoh(\mathcal{O}_Y)$ 是对象，且
$Lg^*E$ 属于 $D^+(\mathcal{O}_Y)$。若 $f$ 平坦，则对映入
$Y$、$Y'$ 和 $X$ 的仿射开集的仿射开集 $U' \subset X'$，
$L(g')^*f^!E$ 与 $(f')^!Lg^*E$ 限制后是
$D(\mathcal{O}_{U'})$ 中的同构对象。
\end{lemma}

\begin{proof}
由假设，立即归约到 $X$、$Y$、$Y'$、$X'$ 都仿射的情形。
设 $Y = \Spec(R)$、$Y' = \Spec(R')$、$X = \Spec(A)$、
$X' = \Spec(A')$。则 $A' = A \otimes_R R'$。设 $E$ 对应于
$K \in D^+(R)$。记给定映射为 $\varphi : R \to A$ 和
$\varphi' : R' \to A'$。由注
\ref{duality-remark-local-calculation-shriek}，可知 $L(g')^*f^!E$ 和
$(f')^!Lg^*E$ 分别对应于
$\varphi^!(K) \otimes_A^\mathbf{L} A'$ 和
$(\varphi')^!(K \otimes_R^\mathbf{L} R')$，其中
$\varphi^!$ 和 $(\varphi')^!$ 是《对偶复形》，第
\ref{dualizing-section-relative-dualizing-complex-algebraic} 节的函子。
结论由《对偶复形》，引理 \ref{dualizing-lemma-bc-flat} 得到。
\end{proof}

\begin{lemma}
\label{duality-lemma-relative-dualizing-fibres}
在情形 \ref{duality-situation-shriek} 中，设 $f : X \to Y$ 是
$\textit{FTS}_S$ 的态射。假设 $f$ 平坦。在
$D^b_{\textit{Coh}}(X)$ 中令
$\omega_{X/Y}^\bullet = f^!\mathcal{O}_Y$。设 $y \in Y$，且
$h : X_y \to X$ 是投影。则 $Lh^*\omega_{X/Y}^\bullet$
是 $X_y$ 上的对偶复形。
\end{lemma}

\begin{proof}
由引理 \ref{duality-lemma-perfect-comparison-shriek}，复形
$\omega_{X/Y}^\bullet$ 属于 $D^b_{\textit{Coh}}$。
成为对偶复形是局部性质。因此由引理
\ref{duality-lemma-base-change-locally}，只需证明
$(X_y \to y)^!\mathcal{O}_y$ 是 $X_y$ 上的对偶复形。
这由引理 \ref{duality-lemma-shriek-dualizing} 得到。
\end{proof}








\section{一种对偶性理论}
\label{duality-section-duality}

\noindent
本节具体说明：上述非常一般的结果对固定 Noetherian 基概形上的有限型分离概形
给出何种对偶性理论。

\medskip\noindent
回忆 Noetherian 概形 $X$ 上的对偶复形，是 $D(\mathcal{O}_X)$
中的对象，并在每个仿射局部给出相应环的对偶复形；见定义
\ref{duality-definition-dualizing-scheme}。

\medskip\noindent
给定 Noetherian 概形 $S$，用 $\textit{FTS}_S$ 表示 $S$ 上有限型且分离的
概形所成的范畴。则：
\begin{enumerate}
\item 函子 $f^!$ 使 $D_\QCoh^+$ 成为 $\textit{FTS}_S$ 上的伪函子；
\item 若 $f : X \to Y$ 是 $\textit{FTS}_S$ 中的真态射，则 $f^!$ 是
$Rf_* : D_\QCoh(\mathcal{O}_X) \to D_\QCoh(\mathcal{O}_Y)$
的右伴随在 $D_\QCoh^+(\mathcal{O}_Y)$ 上的限制，并且对所有
$K \in D_{\textit{Coh}}^-(\mathcal{O}_X)$ 和
$M \in D_\QCoh^+(\mathcal{O}_Y)$，存在典范同构
$$
Rf_*R\SheafHom_{\mathcal{O}_X}(K, f^!M)
\to
R\SheafHom_{\mathcal{O}_Y}(Rf_*K, M)
$$
；
\item 若 $\textit{FTS}_S$ 的对象 $X$ 有对偶复形
$\omega_X^\bullet$，则函子
$D_X = R\SheafHom_{\mathcal{O}_X}(-, \omega_X^\bullet)$
定义 $D_{\textit{Coh}}(\mathcal{O}_X)$ 的一个对合，它交换
$D_{\textit{Coh}}^+(\mathcal{O}_X)$ 与
$D_{\textit{Coh}}^-(\mathcal{O}_X)$，并固定
$D_{\textit{Coh}}^b(\mathcal{O}_X)$；
\item 若 $f : X \to Y$ 是 $\textit{FTS}_S$ 的态射，且
$\omega_Y^\bullet$ 是 $Y$ 上的对偶复形，则
\begin{enumerate}
\item $\omega_X^\bullet = f^!\omega_Y^\bullet$ 是 $X$ 的对偶复形；
\item 对 $M \in D_{\textit{Coh}}^+(\mathcal{O}_Y)$，典范地有
$f^!M = D_X(Lf^*D_Y(M))$；并且
\item 若另外 $f$ 是真态射，则对 $K$，若其属于
$D^-_{\textit{Coh}}(\mathcal{O}_X)$，有
$$
Rf_*R\SheafHom_{\mathcal{O}_X}(K, \omega_X^\bullet) =
R\SheafHom_{\mathcal{O}_Y}(Rf_*K, \omega_Y^\bullet)
$$
；
\end{enumerate}
\item 若 $f : X \to Y$ 是 $\textit{FTS}_S$ 中的闭浸入，则
$f^!(-) = R\SheafHom(\mathcal{O}_X, -)$；
\item 若 $f : Y \to X$ 是 $\textit{FTS}_S$ 中的有限态射，则
$f_*f^!(-) = R\SheafHom_{\mathcal{O}_X}(f_*\mathcal{O}_Y, -)$；
\item 若 $f : X \to Y$ 是一个有效 Cartier 除子到
$\textit{FTS}_S$ 某对象中的包含，则
$f^!(-) = Lf^*(-) \otimes_{\mathcal{O}_X} f^*\mathcal{O}_Y(X)[-1]$；
\item 若 $f : X \to Y$ 是余维数为 $c$ 的 Koszul 正则浸入，
其目标是 $\textit{FTS}_S$ 的对象，则
$f^!(-) \cong Lf^*(-) \otimes_{\mathcal{O}_X} \wedge^c\mathcal{N}[-c]$；并且
\item 若 $f : X \to Y$ 是 $\textit{FTS}_S$ 中相对维数为 $d$ 的光滑真态射，则
$f^!(-) \cong Lf^*(-)  \otimes_{\mathcal{O}_X} \Omega^d_{X/Y}[d]$。
\end{enumerate}
这些结论由引理
\ref{duality-lemma-dualizing-schemes}、
\ref{duality-lemma-iso-on-RSheafHom}、
\ref{duality-lemma-twisted-inverse-image-closed}、
\ref{duality-lemma-finite-twisted}、
\ref{duality-lemma-sheaf-with-exact-support-effective-Cartier}、
\ref{duality-lemma-regular-immersion}、
\ref{duality-lemma-smooth-proper}、
\ref{duality-lemma-upper-shriek-composition}、
\ref{duality-lemma-pseudo-functor}、
\ref{duality-lemma-shriek-closed-immersion}、
\ref{duality-lemma-shriek-dualizing}、
\ref{duality-lemma-shriek-via-duality}、
\ref{duality-lemma-perfect-comparison-shriek}，以及例
\ref{duality-example-iso-on-RSheafHom-noetherian} 得到。
这些函子是通过一个非常抽象的程序得到的，而该程序最终依赖于调用一个存在性定理
（《导出范畴》，命题 \ref{derived-proposition-brown}）。这意味着一般而言，
没有函子 $f^!$ 的显式描述。这有时会造成问题。但事实上，知道对偶复形的存在性和
对偶同构，往往已经足以唯一确定 $f^!$。






\section{对偶复形的黏合}
\label{duality-section-glue}

\noindent
现在利用对偶复形的黏合，为 $S$ 上所有有限型概形得到一种理论，
其中给定情形 \ref{duality-situation-dualizing} 中那样的偶对
$(S, \omega_S^\bullet)$。这类似于 \cite[Remark on page 310]{RD}。

\begin{situation}
\label{duality-situation-dualizing}
这里 $S$ 是 Noetherian 概形，$\omega_S^\bullet$ 是对偶复形。
\end{situation}

\noindent
在情形 \ref{duality-situation-dualizing} 中，设 $X$ 是 $S$ 上的有限型概形。
设 $\mathcal{U} : X = \bigcup_{i = 1, \ldots, n} U_i$ 是由
$\textit{FTS}_S$ 的对象组成的 $X$ 的有限开覆盖；见情形
\ref{duality-situation-shriek}。这仅仅意味着态射 $U_i \to S$ 是分离的
（它们已经是有限型的）。由《概形》，引理
\ref{schemes-lemma-compose-after-separated}，$S$ 上每个有限型仿射概形
都是 $\textit{FTS}_S$ 的对象，故这样的开覆盖当然存在。
于是对每个 $i, j, k \in \{1, \ldots, n\}$，态射
$p_i : U_i \to S$、$p_{ij} : U_i \cap U_j \to S$、
$p_{ijk} : U_i \cap U_j \cap U_k \to S$ 都是分离的，
而且这些概形各自都是 $\textit{FTS}_S$ 的对象。由这样的开覆盖得到
\begin{enumerate}
\item $\omega_i^\bullet = p_i^!\omega_S^\bullet$，它是 $U_i$ 上的对偶复形；
见第 \ref{duality-section-duality} 节；
\item 对每个 $i, j$，有典范同构
$\varphi_{ij} :
\omega_i^\bullet|_{U_i \cap U_j} \to \omega_j^\bullet|_{U_i \cap U_j}$；并且
\item
\label{duality-item-cocycle-glueing}
对每个 $i, j, k$，有
$$
\varphi_{ik}|_{U_i \cap U_j \cap U_k} =
\varphi_{jk}|_{U_i \cap U_j \cap U_k} \circ
\varphi_{ij}|_{U_i \cap U_j \cap U_k}
$$
，该等式位于 $D(\mathcal{O}_{U_i \cap U_j \cap U_k})$ 中。
\end{enumerate}
这里，在 (2) 中使用了 $(U_i \cap U_j \to U_i)^!$ 由限制给出
（引理 \ref{duality-lemma-shriek-open-immersion}），以及由引理
\ref{duality-lemma-upper-shriek-composition} 得到的典范同构
$$
(U_i \cap U_j \to U_i)^! \circ p_i^! = p_{ij}^! =
(U_i \cap U_j \to U_j)^! \circ p_j^!
$$
；为得到 (3)，使用了引理 \ref{duality-lemma-pseudo-functor} 所说明的上叹号函子构成伪函子。

\medskip\noindent
在刚刚描述的情形中，一个{\it 相对于 $\omega_S^\bullet$ 和
$\mathcal{U}$ 规范化的对偶复形}是偶对 $(K, \alpha_i)$，其中
$K \in D(\mathcal{O}_X)$，而
$\alpha_i : K|_{U_i} \to \omega_i^\bullet$ 是同构，并且
$\varphi_{ij}$ 由
$\alpha_j|_{U_i \cap U_j} \circ \alpha_i^{-1}|_{U_i \cap U_j}$ 给出。
由于成为概形上的对偶复形是局部性质，可知相对于
$\omega_S^\bullet$ 和 $\mathcal{U}$ 规范化的对偶复形确实是对偶复形。

\begin{lemma}
\label{duality-lemma-good-dualizing-unique}
在情形 \ref{duality-situation-dualizing} 中，设 $X$ 是 $S$ 上的有限型概形，
而 $\mathcal{U}$ 是由 $S$ 上分离概形组成的 $X$ 的有限开覆盖。
若存在相对于 $\omega_S^\bullet$ 和 $\mathcal{U}$ 规范化的对偶复形，
则它在唯一同构意义下唯一。
\end{lemma}

\begin{proof}
若 $(K, \alpha_i)$ 和 $(K', \alpha_i')$ 是两个这样的复形，则考虑
$L = R\SheafHom_{\mathcal{O}_X}(K, K')$。由引理
\ref{duality-lemma-dualizing-unique-schemes} 及其证明，这是
$D(\mathcal{O}_X)$ 中的可逆对象。利用 $\alpha_i$ 和 $\alpha'_i$，得到同构
$$
\alpha_i^t \otimes \alpha'_i :
L|_{U_i} \longrightarrow
R\SheafHom_{\mathcal{O}_X}(\omega_i^\bullet, \omega_i^\bullet) =
\mathcal{O}_{U_i}[0]
$$
。这已经蕴含在 $D(\mathcal{O}_X)$ 中有 $L = H^0(L)[0]$。
此外，$H^0(L)$ 是可逆层，并在 $X$ 的开集 $U_i$ 上带有给定的平凡化。
最后，
$\alpha_j|_{U_i \cap U_j} \circ \alpha_i^{-1}|_{U_i \cap U_j}$
和
$\alpha'_j|_{U_i \cap U_j} \circ (\alpha'_i)^{-1}|_{U_i \cap U_j}$
都给出 $\varphi_{ij}$ 这一条件，蕴含转移映射均为 $1$，因而得到同构
$H^0(L) = \mathcal{O}_X$。
\end{proof}

\begin{lemma}
\label{duality-lemma-good-dualizing-independence-covering}
在情形 \ref{duality-situation-dualizing} 中，设 $X$ 是 $S$ 上的有限型概形，
而 $\mathcal{U}$、$\mathcal{V}$ 是由 $S$ 上分离概形组成的
$X$ 的两个有限开覆盖。若存在相对于 $\omega_S^\bullet$ 和
$\mathcal{U}$ 规范化的对偶复形，则存在相对于
$\omega_S^\bullet$ 和 $\mathcal{V}$ 规范化的对偶复形，
并且这些复形典范同构。
\end{lemma}

\begin{proof}
只需在 $\mathcal{U}$ 由开集 $U_1, \ldots, U_n$ 给出、而
$\mathcal{V}$ 由开集 $U_1, \ldots, U_{n + m}$ 给出时证明。
事实上，可以并且确实假设 $m = 1$。从相对于
$\omega_S^\bullet$ 和 $\mathcal{V}$ 规范化的对偶复形
$(K, \alpha_i)$ 得到相对于 $\omega_S^\bullet$ 和 $\mathcal{U}$
规范化的对偶复形，只需忘掉 $i = n + 1$ 时的 $\alpha_i$。
反过来，设 $(K, \alpha_i)$ 是相对于 $\omega_S^\bullet$ 和
$\mathcal{U}$ 规范化的对偶复形。为完成证明，需要构造满足所需条件的映射
$\alpha_{n + 1} : K|_{U_{n + 1}} \to \omega_{n + 1}^\bullet$。
为此注意到 $U_{n + 1} = \bigcup U_i \cap U_{n + 1}$ 是开覆盖。
显然，$(K|_{U_{n + 1}}, \alpha_i|_{U_i \cap U_{n + 1}})$ 是相对于
$\omega_S^\bullet$ 和覆盖
$U_{n + 1} = \bigcup U_i \cap U_{n + 1}$ 规范化的对偶复形。
另一方面，由条件 (\ref{duality-item-cocycle-glueing})，偶对
$(\omega_{n + 1}^\bullet|_{U_{n + 1}}, \varphi_{n + 1i})$
是另一个相对于 $\omega_S^\bullet$ 和覆盖
$U_{n + 1} = \bigcup U_i \cap U_{n + 1}$ 规范化的对偶复形。
由引理 \ref{duality-lemma-good-dualizing-unique}，得到与给定局部同构相容的唯一同构
$$
\alpha_{n + 1} : K|_{U_{n + 1}} \longrightarrow \omega_{n + 1}^\bullet
$$
。证明这意味着它满足所需性质，是一个令人愉快的练习。
\end{proof}

\begin{lemma}
\label{duality-lemma-existence-good-dualizing}
在情形 \ref{duality-situation-dualizing} 中，设 $X$ 是 $S$ 上的有限型概形，
而 $\mathcal{U}$ 是由 $S$ 上分离概形组成的 $X$ 的有限开覆盖。
则存在相对于 $\omega_S^\bullet$ 和 $\mathcal{U}$ 规范化的对偶复形。
\end{lemma}

\begin{proof}
设 $\mathcal{U} : X = \bigcup_{i = 1, \ldots, n} U_i$。
对 $n$ 作归纳来证明该引理。基础情形 $n = 1$ 是立即的。
假设 $n > 1$。令 $X' = U_1 \cup \ldots \cup U_{n - 1}$，
并令 $(K', \{\alpha'_i\}_{i = 1, \ldots, n - 1})$ 是相对于
$\omega_S^\bullet$ 和
$\mathcal{U}' : X' = \bigcup_{i = 1, \ldots, n - 1} U_i$
规范化的对偶复形。显然，
$(K'|_{X' \cap U_n}, \alpha'_i|_{U_i \cap U_n})$ 是相对于
$\omega_S^\bullet$ 和覆盖
$X' \cap U_n = \bigcup_{i = 1, \ldots, n - 1} U_i \cap U_n$
规范化的对偶复形。另一方面，由条件
(\ref{duality-item-cocycle-glueing})，偶对
$(\omega_n^\bullet|_{X' \cap U_n}, \varphi_{ni})$ 是另一个相对于
$\omega_S^\bullet$ 和覆盖
$X' \cap U_n = \bigcup_{i = 1, \ldots, n - 1} U_i \cap U_n$
规范化的对偶复形。由引理 \ref{duality-lemma-good-dualizing-unique}，
得到与给定局部同构相容的唯一同构
$$
\epsilon : K'|_{X' \cap U_n} \longrightarrow \omega_i^\bullet|_{X' \cap U_n}
$$
。由《上同调》，引理 \ref{cohomology-lemma-glue}，得到
$K \in D(\mathcal{O}_X)$，以及同构
$\beta : K|_{X'} \to K'$ 和
$\gamma : K|_{U_n} \to \omega_n^\bullet$，使得
$\epsilon = \gamma|_{X'\cap U_n} \circ \beta|_{X' \cap U_n}^{-1}$。
然后定义
$$
\alpha_i = \alpha'_i \circ \beta|_{U_i}, i = 1, \ldots, n - 1,
\text{ 且 }
\alpha_n = \gamma
$$
。还需验证 $\varphi_{ij}$ 由
$\alpha_j|_{U_i \cap U_j} \circ \alpha_i^{-1}|_{U_i \cap U_j}$ 给出。
当 $i, j \leq n - 1$ 时，这由 $\alpha_i'$ 的相应条件得到；
当 $i = j = n$ 时也很清楚。若 $i < j = n$，则
$$
\alpha_n|_{U_i \cap U_n} \circ \alpha_i^{-1}|_{U_i \cap U_n} =
\gamma|_{U_i \cap U_n} \circ \beta^{-1}|_{U_i \cap U_n}
\circ (\alpha'_i)^{-1}|_{U_i \cap U_n} =
\epsilon|_{U_i \cap U_n} \circ (\alpha'_i)^{-1}|_{U_i \cap U_n}
$$
。这恰好等于 $\alpha_{in}$，因为 $\epsilon$ 是与映射
$\alpha_i'$ 和 $\alpha_{ni}$ 相容的唯一映射。
\end{proof}

\noindent
设 $(S, \omega_S^\bullet)$ 如情形 \ref{duality-situation-dualizing} 所述。
上述引理的要点是：给定任意 $S$ 上的有限型概形 $X$，存在在唯一同构意义下
给定的偶对 $(K, \alpha_U)$；它由对象 $K \in D(\mathcal{O}_X)$，
以及对每个在 $S$ 上分离的开子概形 $U \subset X$ 的同构
$\alpha_U : K|_U \to \omega_U^\bullet$ 组成。这里
$\omega_U^\bullet = (U \to S)^!\omega_S^\bullet$ 是 $U$ 上的对偶复形；
见第 \ref{duality-section-duality} 节。此外，若
$\mathcal{U} : X = \bigcup U_i$ 是由在 $S$ 上分离的开集组成的有限开覆盖，
则 $(K, \alpha_{U_i})$ 是相对于 $\omega_S^\bullet$ 和
$\mathcal{U}$ 规范化的对偶复形。事实上，在唯一同构意义下的唯一性由引理
\ref{duality-lemma-good-dualizing-unique} 给出；对一个开覆盖的存在性由引理
\ref{duality-lemma-existence-good-dualizing} 给出；而 $K$ 随后适用于所有开覆盖这一事实
由引理 \ref{duality-lemma-good-dualizing-independence-covering} 给出。

\begin{definition}
\label{duality-definition-good-dualizing}
设 $S$ 是 Noetherian 概形，$\omega_S^\bullet$ 是 $S$ 上的对偶复形。
设 $X$ 是 $S$ 上的有限型概形。上面构造的复形 $K$ 称为
{\it 相对于 $\omega_S^\bullet$ 规范化的对偶复形}，并记作
$\omega_X^\bullet$。
\end{definition}

\noindent
正如术语所示，相对于 $\omega_S^\bullet$ 规范化的对偶复形不仅是
$X$ 的导出范畴中的对象，还配备了上面描述的局部同构。若 $X$ 在 $S$ 上分离，
这与令 $\omega_X^\bullet = p^!\omega_S^\bullet$ 并不冲突，
其中 $p : X \to S$ 是结构态射。更一般地，有下面的合理性检验。

\begin{lemma}
\label{duality-lemma-good-over-both}
设 $(S, \omega_S^\bullet)$ 如情形 \ref{duality-situation-dualizing} 所述。
设 $f : X \to Y$ 是 $S$ 上有限型概形之间的态射。
设 $\omega_X^\bullet$ 和 $\omega_Y^\bullet$ 是相对于
$\omega_S^\bullet$ 规范化的对偶复形。则 $\omega_X^\bullet$
是相对于 $\omega_Y^\bullet$ 规范化的对偶复形。
\end{lemma}

\begin{proof}
这只是一个记账问题。选取有限仿射开覆盖
$\mathcal{V} : Y = \bigcup V_j$。对每个 $j$，选取有限仿射开覆盖
$f^{-1}(V_j) = U_{ji}$。令 $\mathcal{U} : X = \bigcup U_{ji}$。
概形 $V_j$ 和 $U_{ji}$ 在 $S$ 上分离，故对
$q_j : V_j \to S$、$p_{ji} : U_{ji} \to S$、
$f_{ji} : U_{ji} \to V_j$、$f_{ji}' : U_{ji} \to Y$
都有上叹号函子。令 $(L, \beta_j)$ 是相对于
$\omega_S^\bullet$ 和 $\mathcal{V}$ 规范化的对偶复形。
令 $(K, \gamma_{ji})$ 是相对于 $\omega_S^\bullet$ 和
$\mathcal{U}$ 规范化的对偶复形。（换言之，
$L = \omega_Y^\bullet$ 且 $K = \omega_X^\bullet$。）可定义
$$
\alpha_{ji} :
K|_{U_{ji}} \xrightarrow{\gamma_{ji}}
p_{ji}^!\omega_S^\bullet = f_{ji}^!q_j^!\omega_S^\bullet
\xrightarrow{f_{ji}^!\beta_j^{-1}} f_{ji}^!(L|_{V_j}) =
(f_{ji}')^!(L)
$$
。为完成证明，必须说明
$\alpha_{ji}|_{U_{ji} \cap U_{j'i'}}
\circ \alpha_{j'i'}^{-1}|_{U_{ji} \cap U_{j'i'}}$
是典范同构
$(f_{ji}')^!(L)|_{U_{ji} \cap U_{j'i'}} \to
(f_{j'i'}')^!(L)|_{U_{ji} \cap U_{j'i'}}$。这是形式的，略去细节。
\end{proof}

\begin{lemma}
\label{duality-lemma-open-immersion-good-dualizing-complex}
设 $(S, \omega_S^\bullet)$ 如情形 \ref{duality-situation-dualizing} 所述。
设 $j : X \to Y$ 是 $S$ 上有限型概形之间的开浸入。
设 $\omega_X^\bullet$ 和 $\omega_Y^\bullet$ 是相对于
$\omega_S^\bullet$ 规范化的对偶复形。则存在典范同构
$\omega_X^\bullet = \omega_Y^\bullet|_X$。
\end{lemma}

\begin{proof}
这由定义 \ref{duality-definition-good-dualizing} 紧前面给出的规范化对偶复形构造立即得到。
\end{proof}

\begin{lemma}
\label{duality-lemma-proper-map-good-dualizing-complex}
设 $(S, \omega_S^\bullet)$ 如情形 \ref{duality-situation-dualizing} 所述。
设 $f : X \to Y$ 是 $S$ 上有限型概形之间的真态射。
设 $\omega_X^\bullet$ 和 $\omega_Y^\bullet$ 是相对于
$\omega_S^\bullet$ 规范化的对偶复形。令 $a$ 为引理
\ref{duality-lemma-twisted-inverse-image} 对 $f$ 给出的右伴随。
则存在典范同构 $a(\omega_Y^\bullet) = \omega_X^\bullet$。
\end{lemma}

\begin{proof}
设 $p : X \to S$ 和 $q : Y \to S$ 是结构态射。若 $X$ 和 $Y$
在 $S$ 上分离，则结论来自
$\omega_X^\bullet = p^!\omega_S^\bullet$、
$\omega_Y^\bullet = q^!\omega_S^\bullet$、$f^! = a$，以及
$f^! \circ q^! = p^!$（引理 \ref{duality-lemma-upper-shriek-composition}）。
在一般情形下，先使用引理 \ref{duality-lemma-good-over-both} 归约到
$Y = S$。此时 $X$ 和 $Y$ 在 $S$ 上分离，而刚才已经看到了结论。
\end{proof}

\noindent
设 $(S, \omega_S^\bullet)$ 如情形 \ref{duality-situation-dualizing} 所述。
对 $S$ 上的有限型概形 $X$，用 $\omega_X^\bullet$ 表示 $X$ 的
相对于 $\omega_S^\bullet$ 规范化的对偶复形。像引理
\ref{duality-lemma-dualizing-schemes} 中那样定义
$D_X(-) = R\SheafHom_{\mathcal{O}_X}(-, \omega_X^\bullet)$。
设 $f : X \to Y$ 是 $S$ 上有限型概形之间的态射。定义
$$
f_{new}^! = D_X \circ Lf^* \circ D_Y :
D_{\textit{Coh}}^+(\mathcal{O}_Y)
\to
D_{\textit{Coh}}^+(\mathcal{O}_X)
$$
。若 $f : X \to Y$ 和 $g : Y \to Z$ 是 $S$ 上有限型概形之间
可复合的态射，则定义
\begin{align*}
(g \circ f)^!_{new} & = D_X \circ L(g \circ f)^* \circ D_Z \\
& = D_X \circ Lf^* \circ Lg^* \circ D_Z \\
& \to D_X \circ Lf^* \circ D_Y \circ D_Y \circ Lg^* \circ D_Z \\
& = f^!_{new} \circ g^!_{new}
\end{align*}
，其中箭头在引理 \ref{duality-lemma-dualizing-schemes} 中定义。
下面的引理汇集这些结果。

\begin{lemma}
\label{duality-lemma-duality-bootstrap}
设 $(S, \omega_S^\bullet)$ 如情形 \ref{duality-situation-dualizing} 所述。
按上面的方式，对 $S$ 上所有有限型概形（的态射）定义
$f^!_{new}$ 和 $\omega_X^\bullet$。则：
\begin{enumerate}
\item 函子 $f^!_{new}$ 和箭头
$(g \circ f)^!_{new} \to f^!_{new} \circ g^!_{new}$
使 $D_{\textit{Coh}}^+$ 成为从 $S$ 上有限型概形的范畴到范畴的
$2$-范畴的伪函子；
\item $\omega_X^\bullet = (X \to S)^!_{new} \omega_S^\bullet$；
\item 函子 $D_X$ 定义 $D_{\textit{Coh}}(\mathcal{O}_X)$ 的一个对合，
它交换 $D_{\textit{Coh}}^+(\mathcal{O}_X)$ 与
$D_{\textit{Coh}}^-(\mathcal{O}_X)$，并固定
$D_{\textit{Coh}}^b(\mathcal{O}_X)$；
\item 对 $S$ 上有限型概形之间的态射 $f : X \to Y$，有
$\omega_X^\bullet = f^!_{new}\omega_Y^\bullet$；
\item 对 $M \in D_{\textit{Coh}}^+(\mathcal{O}_Y)$，有
$f^!_{new}M = D_X(Lf^*D_Y(M))$；并且
\item 若另外 $f$ 是真态射，则 $f^!_{new}$ 同构于
$Rf_* : D_\QCoh(\mathcal{O}_X) \to D_\QCoh(\mathcal{O}_Y)$
的右伴随在 $D_{\textit{Coh}}^+(\mathcal{O}_Y)$ 上的限制，
而且对 $K \in D^-_{\textit{Coh}}(\mathcal{O}_X)$ 和
$M \in D_{\textit{Coh}}^+(\mathcal{O}_Y)$，存在典范同构
$$
Rf_*R\SheafHom_{\mathcal{O}_X}(K, f^!_{new}M)
\to
R\SheafHom_{\mathcal{O}_Y}(Rf_*K, M)
$$
，并且对 $K \in D^-_{\textit{Coh}}(\mathcal{O}_X)$，有
$$
Rf_*R\SheafHom_{\mathcal{O}_X}(K, \omega_X^\bullet) =
R\SheafHom_{\mathcal{O}_Y}(Rf_*K, \omega_Y^\bullet)
$$
，并且
\end{enumerate}
若 $X$ 在 $S$ 上分离，则 $\omega_X^\bullet$ 典范同构于
$(X \to S)^!\omega_S^\bullet$；若 $f$ 是 $S$ 上分离概形之间的态射，
则存在典范同构\footnote{尚未检查这些同构与
$(g \circ f)^! \to f^! \circ g^!$ 和
$(g \circ f)^!_{new} \to f^!_{new} \circ g^!_{new}$ 相容。
若后文需要，将在这里完成检查。}
$f_{new}^!K = f^!K$，其中 $K$ 属于 $D_{\textit{Coh}}^+$。
\end{lemma}

\begin{proof}
设 $f : X \to Y$、$g : Y \to Z$、$h : Z \to T$ 是
$S$ 上有限型概形之间的态射。必须证明图
$$
\xymatrix{
(h \circ g \circ f)^!_{new} \ar[r] \ar[d] &
f^!_{new} \circ (h \circ g)^!_{new} \ar[d] \\
(g \circ f)^!_{new} \circ h^!_{new} \ar[r] &
f^!_{new} \circ g^!_{new} \circ h^!_{new}
}
$$
交换。令 $\eta_Y : \text{id} \to D_Y^2$ 和
$\eta_Z : \text{id} \to D_Z^2$ 为引理
\ref{duality-lemma-dualizing-schemes} 的典范同构。然后利用《范畴》，引理
\ref{categories-lemma-properties-2-cat-cats}，一个略去的计算表明，
两个箭头
$(h \circ g \circ f)^!_{new} \to f^!_{new} \circ g^!_{new} \circ h^!_{new}$
都由
$$
1 \star \eta_Y \star 1 \star \eta_Z \star 1 :
D_X \circ Lf^* \circ Lg^* \circ Lh^* \circ D_T
\longrightarrow
D_X \circ Lf^* \circ D_Y^2 \circ Lg^* \circ D_Z^2 \circ Lh^* \circ D_T
$$
给出。这证明了 (1)。(2) 由 $(X \to S)^!_{new}$ 的定义以及
$D_S(\omega_S^\bullet) = \mathcal{O}_S$ 这一事实立即得到。
(3) 是引理 \ref{duality-lemma-dualizing-schemes}。(4) 由与 (2) 相同的论证得到。
(5) 是 $f^!_{new}$ 的定义。

\medskip\noindent
证明 (6)。令 $a$ 为引理 \ref{duality-lemma-twisted-inverse-image}
对 $S$ 上有限型概形之间的真态射 $f : X \to Y$ 给出的右伴随。
问题在于不知道 $X$ 或 $Y$ 在 $S$ 上是否分离（一般确实不分离），
所以不能立即把引理 \ref{duality-lemma-shriek-via-duality} 应用于基概形
$S$ 上的 $f$。为绕过这一点，使用引理
\ref{duality-lemma-proper-map-good-dualizing-complex} 的典范等同
$\omega_X^\bullet = a(\omega_Y^\bullet)$。因此，将引理
\ref{duality-lemma-shriek-via-duality} 应用于以 $Y$ 为基概形的
$f : X \to Y$，可知 $f^!_{new}$ 是 $a$ 在
$D_{\textit{Coh}}^+(\mathcal{O}_Y)$ 上的限制！所显示的等式由例
\ref{duality-example-iso-on-RSheafHom-noetherian} 得到。

\medskip\noindent
最后的断言由规范化对偶复形的构造以及已经使用过的引理
\ref{duality-lemma-shriek-via-duality} 得到。
\end{proof}

\begin{remark}
\label{duality-remark-independent-omega-S}
设 $S$ 是具有对偶复形的 Noetherian 概形。设 $f : X \to Y$ 是
$S$ 上有限型概形之间的态射。则函子
$$
f_{new}^! : D^+_{Coh}(\mathcal{O}_Y) \to D^+_{Coh}(\mathcal{O}_X)
$$
在典范同构意义下与对偶复形 $\omega_S^\bullet$ 的选择无关。
下面略述证明。任意第二个对偶复形都形如
$\omega_S^\bullet \otimes_{\mathcal{O}_S}^\mathbf{L} \mathcal{L}$，
其中 $\mathcal{L}$ 是 $D(\mathcal{O}_S)$ 中的可逆对象；见引理
\ref{duality-lemma-dualizing-unique-schemes}。对任意有限型分离态射
$p : U \to S$，由引理 \ref{duality-lemma-compare-with-pullback-perfect}，有
$p^!(\omega_S^\bullet \otimes^\mathbf{L}_{\mathcal{O}_S} \mathcal{L}) =
p^!(\omega_S^\bullet) \otimes^\mathbf{L}_{\mathcal{O}_U} Lp^*\mathcal{L}$。
因此，若 $\omega_X^\bullet$ 和 $\omega_Y^\bullet$ 是相对于
$\omega_S^\bullet$ 规范化的对偶复形，则
$\omega_X^\bullet \otimes_{\mathcal{O}_X}^\mathbf{L} La^*\mathcal{L}$ 和
$\omega_Y^\bullet \otimes_{\mathcal{O}_Y}^\mathbf{L} Lb^*\mathcal{L}$
是相对于
$\omega_S^\bullet \otimes_{\mathcal{O}_S}^\mathbf{L} \mathcal{L}$
规范化的对偶复形（其中 $a : X \to S$ 和 $b : Y \to S$ 是结构态射）。
于是结论由下列等式得到：
\begin{align*}
& R\SheafHom_{\mathcal{O}_X}(Lf^*R\SheafHom_{\mathcal{O}_Y}(K,
\omega_Y^\bullet \otimes_{\mathcal{O}_Y}^\mathbf{L} Lb^*\mathcal{L}),
\omega_X^\bullet \otimes_{\mathcal{O}_X}^\mathbf{L} La^*\mathcal{L}) \\
& = R\SheafHom_{\mathcal{O}_X}(Lf^*R(\SheafHom_{\mathcal{O}_Y}(K,
\omega_Y^\bullet) \otimes_{\mathcal{O}_Y}^\mathbf{L} Lb^*\mathcal{L}),
\omega_X^\bullet \otimes_{\mathcal{O}_X}^\mathbf{L} La^*\mathcal{L}) \\
& = R\SheafHom_{\mathcal{O}_X}(Lf^*R\SheafHom_{\mathcal{O}_Y}(K,
\omega_Y^\bullet) \otimes_{\mathcal{O}_X}^\mathbf{L} La^*\mathcal{L},
\omega_X^\bullet \otimes_{\mathcal{O}_X}^\mathbf{L} La^*\mathcal{L}) \\
& = R\SheafHom_{\mathcal{O}_X}(Lf^*R\SheafHom_{\mathcal{O}_Y}(K,
\omega_Y^\bullet), \omega_X^\bullet)
\end{align*}
，其中 $K \in D^+_{Coh}(\mathcal{O}_Y)$。
最后一个等式成立，是因为 $La^*\mathcal{L}$ 在
$D(\mathcal{O}_X)$ 中可逆。
\end{remark}


\begin{example}
\label{duality-example-trace-proper}
设 $S$ 是 Noetherian 概形，$\omega_S^\bullet$ 是对偶复形。
设 $f : X \to Y$ 是 $S$ 上有限型概形之间的真态射。
设 $\omega_X^\bullet$ 和 $\omega_Y^\bullet$ 是相对于
$\omega_S^\bullet$ 规范化的对偶复形。在此情形下，有
$a(\omega_Y^\bullet) = \omega_X^\bullet$
（引理 \ref{duality-lemma-proper-map-good-dualizing-complex}），故迹映射
（第 \ref{duality-section-trace} 节）是典范箭头
$$
\text{Tr}_f : Rf_*\omega_X^\bullet \longrightarrow \omega_Y^\bullet
$$
，它产生同构（引理 \ref{duality-lemma-duality-bootstrap}）
$$
\Hom_X(L, \omega_X^\bullet) = \Hom_Y(Rf_*L, \omega_Y^\bullet)
$$
和
$$
Rf_*R\SheafHom_{\mathcal{O}_X}(L, \omega_X^\bullet) =
R\SheafHom_{\mathcal{O}_Y}(Rf_*L, \omega_Y^\bullet)
$$
，其中 $L$ 属于 $D_\QCoh(\mathcal{O}_X)$。
\end{example}

\begin{remark}
\label{duality-remark-dualizing-finite}
设 $S$ 是 Noetherian 概形，$\omega_S^\bullet$ 是对偶复形。
设 $f : X \to Y$ 是 $S$ 上有限型概形之间的有限态射。
设 $\omega_X^\bullet$ 和 $\omega_Y^\bullet$ 是相对于
$\omega_S^\bullet$ 规范化的对偶复形。由引理
\ref{duality-lemma-finite-twisted} 和
\ref{duality-lemma-proper-map-good-dualizing-complex}，在
$D_\QCoh^+(f_*\mathcal{O}_X)$ 中有
$$
f_*\omega_X^\bullet = R\SheafHom(f_*\mathcal{O}_X, \omega_Y^\bullet)
$$
，而例 \ref{duality-example-trace-proper} 的迹映射是映射
$$
\text{Tr}_f : Rf_*\omega_X^\bullet = f_*\omega_X^\bullet =
R\SheafHom(f_*\mathcal{O}_X, \omega_Y^\bullet) \longrightarrow
\omega_Y^\bullet
$$
，它常称为“在 $1$ 处求值”。
\end{remark}

\begin{remark}
\label{duality-remark-relative-dualizing-complex-shriek}
设 $f : X \to Y$ 是情形 \ref{duality-situation-dualizing} 中那样的偶对
$(S, \omega_S^\bullet)$ 上有限型概形之间的平坦真态射。
相对对偶复形（注 \ref{duality-remark-relative-dualizing-complex}）是
$\omega_{X/Y}^\bullet = a(\mathcal{O}_Y)$。由引理
\ref{duality-lemma-proper-map-good-dualizing-complex}，在
$$
\omega_X^\bullet = a(\omega_Y^\bullet) =
Lf^*\omega_Y^\bullet \otimes_{\mathcal{O}_X}^\mathbf{L} \omega_{X/Y}^\bullet
$$
中有第一个典范同构。该等式位于 $D(\mathcal{O}_X)$ 中。
第二个典范同构由注 \ref{duality-remark-relative-dualizing-complex} 中的讨论得到。
\end{remark}




\section{维数函数}
\label{duality-section-dimension-functions}

\noindent
需要进一步了解转到另一个概形上的有限型概形时，维数函数如何变化。

\begin{lemma}
\label{duality-lemma-good-dualizing-normalized}
设 $S$ 是 Noetherian 概形，$\omega_S^\bullet$ 是对偶复形。
设 $X$ 是 $S$ 上的有限型概形，$\omega_X^\bullet$ 是相对于
$\omega_S^\bullet$ 规范化的对偶复形。若 $x \in X$ 是位于
$S$ 的闭点 $s$ 之上的闭点，且 $\omega_{S, s}^\bullet$ 是
$\mathcal{O}_{S, s}$ 上的规范化对偶复形，则
$\omega_{X, x}^\bullet$ 是 $\mathcal{O}_{X, x}$ 上的规范化对偶复形。
\end{lemma}

\begin{proof}
可用 $x$ 的一个仿射邻域替换 $X$，故可以并且确实假设
$f : X \to S$ 分离。于是 $\omega_X^\bullet = f^!\omega_S^\bullet$。
必须说明 $R\Hom_{\mathcal{O}_{X, x}}(\kappa(x), \omega_{X, x}^\bullet)$
集中于次数 $0$。令 $i_x : x \to X$ 表示包含态射；由于 $x$ 是闭点，
它是闭浸入。因此由引理 \ref{duality-lemma-shriek-closed-immersion}，
$R\Hom_{\mathcal{O}_{X, x}}(\kappa(x), \omega_{X, x}^\bullet)$
表示 $i_x^!\omega_X^\bullet$。考虑交换图
$$
\xymatrix{
x \ar[r]_{i_x} \ar[d]_\pi & X \ar[d]^f \\
s \ar[r]^{i_s} & S
}
$$
。由《态射》，引理
\ref{morphisms-lemma-closed-point-fibre-locally-finite-type}，
扩张 $\kappa(x)/\kappa(s)$ 是有限的，因而 $\pi$ 是有限态射。所以
$$
i_x^!\omega_X^\bullet = i_x^! f^! \omega_S^\bullet =
\pi^! i_s^! \omega_S^\bullet
$$
。于是若 $\omega_{S, s}^\bullet$ 是 $\mathcal{O}_{S, s}$ 上的
规范化对偶复形，则由上面的同一论证，
$i_s^!\omega_S^\bullet = \kappa(s)[0]$。有
$$
R\pi_*(\pi^!(\kappa(s)[0])) =
R\SheafHom_{\mathcal{O}_s}(R\pi_*(\kappa(x)[0]), \kappa(s)[0]) =
\widetilde{\Hom_{\kappa(s)}(\kappa(x), \kappa(s))}
$$
。第一个等式来自例 \ref{duality-example-iso-on-RSheafHom-noetherian}，
其中取 $L = \kappa(x)[0]$。第二个等式成立，是因为 $\pi_*$ 正合。
因此 $\pi^!(\kappa(s)[0])$ 支撑在次数 $0$，结论成立。
\end{proof}

\begin{lemma}
\label{duality-lemma-good-dualizing-dimension-function}
设 $S$ 是 Noetherian 概形，$\omega_S^\bullet$ 是对偶复形。
设 $f : X \to S$ 是有限型态射，$\omega_X^\bullet$ 是相对于
$\omega_S^\bullet$ 规范化的对偶复形。对所有 $x \in X$，有
$$
\delta_X(x) - \delta_S(f(x)) = \text{trdeg}_{\kappa(f(x))}(\kappa(x))
$$
，其中 $\delta_S$\ 和 $\delta_X$\ 分别是
$\omega_S^\bullet$ 和 $\omega_X^\bullet$ 的维数函数；见引理
\ref{duality-lemma-dimension-function-scheme}。
\end{lemma}

\begin{proof}
可用 $x$ 的一个仿射邻域替换 $X$。故可以并且确实假设存在
$S$ 上的紧化 $X \subset \overline{X}$。再用
$\overline{X}$ 替换 $X$，可设 $X$ 在 $S$ 上为真。还可用含 $x$ 的
$X$ 的连通分支替换 $X$，从而假设 $X$ 连通。其次，回忆
$\delta_X$ 和函数
$x \mapsto \delta_S(f(x)) + \text{trdeg}_{\kappa(f(x))}(\kappa(x))$
都是 $X$ 上的维数函数；见《态射》，引理
\ref{morphisms-lemma-dimension-function-propagates}
（并使用由引理 \ref{duality-lemma-dimension-function-scheme} 得到的
$S$ 是普遍链状的事实）。由《拓扑》，引理
\ref{topology-lemma-dimension-function-unique}，二者之差局部常值；
由于 $X$ 连通，它是常值的。因此只需在 $X$ 的任一点证明等式。
由《性质》，引理 \ref{properties-lemma-locally-Noetherian-closed-point}，
概形 $X$ 有闭点 $x$。由于 $X \to S$ 是真态射，$x$ 的像 $s$
在 $S$ 中闭。于是应用引理 \ref{duality-lemma-good-dualizing-normalized} 即得结论。
\end{proof}

\begin{lemma}
\label{duality-lemma-shriek}
在情形 \ref{duality-situation-shriek} 中，设 $f : X \to Y$ 是
$\textit{FTS}_S$ 的态射。设 $x \in X$ 的像为 $y \in Y$。则
$$
H^i(f^!\mathcal{O}_Y)_x \not = 0
\Rightarrow - \dim_x(X_y) \leq i.
$$
\end{lemma}

\begin{proof}
因为该陈述在 $X$ 上是局部的，可设 $X$ 和 $Y$ 是仿射概形。
写成 $X = \Spec(A)$ 和 $Y = \Spec(R)$。
$f^!\mathcal{O}_Y$ 对应于《对偶复形》，第
\ref{dualizing-section-relative-dualizing-complexes-Noetherian} 节的
相对对偶复形 $\omega_{A/R}^\bullet$；这由注
\ref{duality-remark-local-calculation-shriek} 得到。因此该引理由《对偶复形》，引理
\ref{dualizing-lemma-relative-dualizing-trivial-vanishing} 得到。
\end{proof}

\begin{lemma}
\label{duality-lemma-flat-shriek}
在情形 \ref{duality-situation-shriek} 中，设 $f : X \to Y$ 是
$\textit{FTS}_S$ 的态射。设 $x \in X$ 的像为 $y \in Y$。
若 $f$ 平坦，则
$$
H^i(f^!\mathcal{O}_Y)_x \not = 0
\Rightarrow - \dim_x(X_y) \leq i \leq 0.
$$
事实上，若 $f$ 的所有纤维维数都 $\leq d$，则作为
$D(X, f^{-1}\mathcal{O}_Y)$ 的对象，$f^!\mathcal{O}_Y$
的 Tor 振幅位于 $[-d, 0]$。
\end{lemma}

\begin{proof}
完全像引理 \ref{duality-lemma-shriek} 的证明中那样论证，结论由《对偶复形》，引理
\ref{dualizing-lemma-relative-dualizing-flat-vanishing} 得到。
\end{proof}

\begin{lemma}
\label{duality-lemma-shriek-over-CM}
在情形 \ref{duality-situation-shriek} 中，设 $f : X \to Y$ 是
$\textit{FTS}_S$ 的态射。设 $x \in X$ 的像为 $y \in Y$。假设
\begin{enumerate}
\item $\mathcal{O}_{Y, y}$ 是 Cohen-Macaulay 的；并且
\item 对含 $x$ 的 $X$ 的任意不可约分支的泛点 $\xi$，有
$\text{trdeg}_{\kappa(f(\xi))}(\kappa(\xi)) \leq r$。
\end{enumerate}
则
$$
H^i(f^!\mathcal{O}_Y)_x \not = 0
\Rightarrow - r \leq i
$$
，而且作为 $\mathcal{O}_{X, x}$-模，茎
$H^{-r}(f^!\mathcal{O}_Y)_x$ 满足 $(S_2)$。
\end{lemma}

\begin{proof}
用 $x$ 的一个开邻域替换 $X$ 后，可设 $X$ 的每个不可约分支都经过 $x$。
于是完全像引理 \ref{duality-lemma-shriek} 的证明中那样论证，结论由《对偶复形》，引理
\ref{dualizing-lemma-relative-dualizing-CM-vanishing} 得到。
\end{proof}

\begin{lemma}
\label{duality-lemma-flat-quasi-finite-shriek}
在情形 \ref{duality-situation-shriek} 中，设 $f : X \to Y$ 是
$\textit{FTS}_S$ 的态射。若 $f$ 平坦且拟有限，则
$$
f^!\mathcal{O}_Y = \omega_{X/Y}[0]
$$
，其中 $\omega_{X/Y}$ 是某个在 $Y$ 上平坦的凝聚
$\mathcal{O}_X$-模。
\end{lemma}

\begin{proof}
这是引理 \ref{duality-lemma-flat-shriek} 以及如下事实的推论：由引理
\ref{duality-lemma-shriek-coherent}，$f^!\mathcal{O}_Y$ 的上同调层是凝聚的。
\end{proof}

\begin{lemma}
\label{duality-lemma-CM-shriek}
在情形 \ref{duality-situation-shriek} 中，设 $f : X \to Y$ 是
$\textit{FTS}_S$ 的态射。若 $f$ 是 Cohen-Macaulay 态射
（《态射进阶》，定义 \ref{more-morphisms-definition-CM}），则
$$
f^!\mathcal{O}_Y = \omega_{X/Y}[d]
$$
，其中 $\omega_{X/Y}$ 是某个在 $Y$ 上平坦的凝聚
$\mathcal{O}_X$-模，而 $d$ 是 $X$ 上给出 $X$ 在 $Y$ 上相对维数的
局部常值函数。
\end{lemma}

\begin{proof}
由《态射》，引理
\ref{morphisms-lemma-flat-finite-presentation-CM-fibres-relative-dimension}，
相对维数 $d$ 定义良好且局部常值。由引理
\ref{duality-lemma-shriek-coherent}，$f^!\mathcal{O}_Y$ 的上同调层是凝聚的。
若能说明 $f^!\mathcal{O}_Y$ 的其他上同调层均为零，则由引理
\ref{duality-lemma-flat-shriek} 得到 $\omega_{X/Y}$ 的平坦性。

\medskip\noindent
这个问题在 $X$ 上是局部的，故可设 $X$ 和 $Y$ 是仿射的，且态射的相对维数为
$d$。若 $d = 0$，结论直接由引理
\ref{duality-lemma-flat-quasi-finite-shriek} 得到。若 $d > 0$，则可设存在分解
$$
X \xrightarrow{g} \mathbf{A}^d_Y \xrightarrow{p} Y
$$
，其中 $g$ 拟有限且平坦；见《态射进阶》，引理
\ref{more-morphisms-lemma-flat-finite-presentation-characterize-CM}。
于是 $f^! = g^! \circ p^!$。由引理
\ref{duality-lemma-shriek-affine-line}，有
$p^!\mathcal{O}_Y \cong \mathcal{O}_{\mathbf{A}^d_Y}[-d]$。
最后用 $d = 0$ 的情形即得结论。
\end{proof}

\begin{remark}
\label{duality-remark-the-same-is-true}
设 $S$ 是配备对偶复形 $\omega_S^\bullet$ 的 Noetherian 概形。
在此情形下，引理 \ref{duality-lemma-shriek}、
\ref{duality-lemma-flat-shriek}、
\ref{duality-lemma-flat-quasi-finite-shriek} 和 \ref{duality-lemma-CM-shriek}
对 $S$ 上有限型概形之间的任意态射 $f : X \to Y$ 都成立，
但要用 $f_{new}^!$ 替换 $f^!$。这是清楚的，因为每种情形的证明都立即归约到
仿射情形，而此时由引理 \ref{duality-lemma-duality-bootstrap}，有
$f^! = f_{new}^!$。
\end{remark}




\section{对偶模}
\label{duality-section-dualizing-module}


\noindent
本节是《对偶复形》，第 \ref{dualizing-section-dualizing-module} 节的延续。

\medskip\noindent
设 $X$ 是 Noetherian 概形，$\omega_X^\bullet$ 是对偶复形。
设 $n \in \mathbf{Z}$ 是使 $H^n(\omega_X^\bullet)$ 非零的最小整数。
换言之，$-n$ 是与 $\omega_X^\bullet$ 相伴的维数函数的最大值
（引理 \ref{duality-lemma-dimension-function-scheme}）。有时
$H^n(\omega_X^\bullet)$ 称为 $X$ 的{\it 对偶模}或{\it 对偶层}，
并常记作 $\omega_X$。下文说“设 $\omega_X$ 为一个对偶模”时，
意指上述情形。

\medskip\noindent
在 Noetherian 概形 $X$ 上使用对偶模 $\omega_X$ 时必须谨慎：
\begin{enumerate}
\item 从 $X$ 转到 $X$ 的开集 $U$ 时，整数 $n$ 可能改变；
此时 $\omega_X|_U = \omega_U$ 不成立；
\item 对偶复形不唯一；对偶模仅在张量一个可逆模的意义下唯一。
\end{enumerate}
第二个问题往往无关紧要，因为将研究基变换 $S$ 上的有限型概形
$X$，其中该基变换配备固定的对偶复形 $\omega_S^\bullet$，而
$\omega_X^\bullet$ 将取为相对于 $\omega_S^\bullet$ 规范化的对偶复形。
若 $X$ 等维，则第一个问题不会发生；更精确地说，只要与
$\omega_X^\bullet$ 相伴的维数函数（引理
\ref{duality-lemma-dimension-function-scheme}）将 $X$ 的每个泛点映到同一个整数即可。

\begin{example}
\label{duality-example-proper-over-local}
设 $S = \Spec(A)$，其中 $(A, \mathfrak m, \kappa)$ 是 Noetherian 局部环，
而 $\omega_S^\bullet$ 对应于规范化对偶复形
$\omega_A^\bullet$。若 $f : X \to S$ 在 $S$ 上为真，且
$\omega_X^\bullet = f^!\omega_S^\bullet$，则凝聚层
$$
\omega_X = H^{-\dim(X)}(\omega_X^\bullet)
$$
是对偶模，并常称为 $X$ 的对偶模（这里默认已指定 $S$ 和
$\omega_S^\bullet$）。将会看到它具有良好性质。
\end{example}

\begin{example}
\label{duality-example-equidimensional-over-field}
设 $X$ 是域 $k$ 上的等维有限型概形。通常取
$\omega_X^\bullet$ 为相对于 $k[0]$ 规范化的对偶复形，并把
$$
\omega_X = H^{-\dim(X)}(\omega_X^\bullet)
$$
称为 $X$ 的对偶模。若 $X$ 在 $k$ 上分离，则
$\omega_X^\bullet = f^!\mathcal{O}_{\Spec(k)}$，其中
$f : X \to \Spec(k)$ 是结构态射；这由引理
\ref{duality-lemma-duality-bootstrap} 得到。若 $X$ 在 $k$ 上为真，
则这是例 \ref{duality-example-proper-over-local} 的特殊情形。
\end{example}

\begin{lemma}
\label{duality-lemma-dualizing-module}
设 $X$ 是连通 Noetherian 概形，$\omega_X$ 是 $X$ 上的对偶模。
$\omega_X$ 的支集是相对于任意维数函数具有最大维数的不可约分支之并，
而且 $\omega_X$ 是满足性质 $(S_2)$ 的凝聚 $\mathcal{O}_X$-模。
\end{lemma}

\begin{proof}
按上面讨论的约定，存在对偶复形 $\omega_X^\bullet$，使得
$\omega_X$ 是其最左边的非零上同调层。由于 $X$ 连通，任意两个维数函数
相差一个常数（《拓扑》，引理
\ref{topology-lemma-dimension-function-unique}）。因此可使用与
$\omega_X^\bullet$ 相伴的维数函数（引理
\ref{duality-lemma-dimension-function-scheme}）。有了这些说明，该引理由
《对偶复形》，引理 \ref{dualizing-lemma-depth-dualizing-module}
和各定义得到（特别是《概形的上同调》，定义
\ref{coherent-definition-depth}）。
\end{proof}

\begin{lemma}
\label{duality-lemma-vanishing-good-dualizing}
设 $X/A$、$\omega_X^\bullet$、$\omega_X$ 如例
\ref{duality-example-proper-over-local} 所述。则
\begin{enumerate}
\item $H^i(\omega_X^\bullet) \not = 0 \Rightarrow
i \in \{-\dim(X), \ldots, 0\}$；
\item $H^i(\omega_X^\bullet)$ 的支集维数至多为 $-i$；
\item $\text{Supp}(\omega_X)$ 是维数为 $\dim(X)$ 的分支之并；并且
\item $\omega_X$ 满足性质 $(S_2)$。
\end{enumerate}
\end{lemma}

\begin{proof}
令 $\delta_X$ 和 $\delta_S$ 为与 $\omega_X^\bullet$ 和
$\omega_S^\bullet$ 相伴的维数函数，如引理
\ref{duality-lemma-good-dualizing-dimension-function} 所述。由于 $X$ 在 $A$ 上为真，
$X$ 的每个闭子概形都含有一个映到闭点 $s \in S$ 的闭点 $x$，并且
$\delta_X(x) = \delta_S(s) = 0$。因此对任意点 $\xi \in X$，有
$\delta_X(\xi) = \dim(\overline{\{\xi\}})$。所以可通过考察
$\Spec(\mathcal{O}_{X, x})$ 上的情形来检验引理的每个陈述；此时结论由
《对偶复形》，引理 \ref{dualizing-lemma-sitting-in-degrees} 和
\ref{dualizing-lemma-depth-dualizing-module} 得到。略去若干细节。
最后两个陈述也可由引理 \ref{duality-lemma-dualizing-module} 推出。
\end{proof}

\begin{lemma}
\label{duality-lemma-dualizing-module-proper-over-A}
设带对偶模 $\omega_X$ 的 $X/A$ 如例
\ref{duality-example-proper-over-local} 所述。令 $d = \dim(X_s)$ 为闭纤维的维数。
若 $\dim(X) = d + \dim(A)$，则对偶模 $\omega_X$ 表示凝聚
$\mathcal{O}_X$-模范畴上的函子
$$
\mathcal{F} \longmapsto \Hom_A(H^d(X, \mathcal{F}), \omega_A)
$$
。
\end{lemma}

\begin{proof}
有
\begin{align*}
\Hom_X(\mathcal{F}, \omega_X)
& =
\Ext^{-\dim(X)}_X(\mathcal{F}, \omega_X^\bullet) \\
& =
\Hom_X(\mathcal{F}[\dim(X)], \omega_X^\bullet) \\
& =
\Hom_X(\mathcal{F}[\dim(X)], f^!(\omega_A^\bullet)) \\
& =
\Hom_S(Rf_*\mathcal{F}[\dim(X)], \omega_A^\bullet) \\
& =
\Hom_A(H^d(X, \mathcal{F}), \omega_A)
\end{align*}
第一个等式成立，是因为当 $i < -\dim(X)$ 时
$H^i(\omega_X^\bullet) = 0$；见引理
\ref{duality-lemma-vanishing-good-dualizing} 和《导出范畴》，引理
\ref{derived-lemma-negative-exts}。第二个等式由 Ext 群的定义得到。
第三个等式来自对 $\omega_X^\bullet$ 的选择。第四个等式成立，
是因为 $f^!$ 是引理 \ref{duality-lemma-twisted-inverse-image} 对 $f$ 给出的右伴随；
见第 \ref{duality-section-duality} 节。最后一个等式成立，是因为当 $i > d$ 时
$R^if_*\mathcal{F}$ 为零（《概形的上同调》，引理
\ref{coherent-lemma-higher-direct-images-zero-above-dimension-fibre}），
并且当 $j < -\dim(A)$ 时 $H^j(\omega_A^\bullet)$ 为零。
\end{proof}








\section{Cohen-Macaulay 概形}
\label{duality-section-CM}

\noindent
本节是《对偶复形》，第 \ref{dualizing-section-CM} 节的延续。
对 Cohen-Macaulay 概形，对偶性具有特别简单的形式。

\begin{lemma}
\label{duality-lemma-dualizing-module-CM-scheme}
设 $X$ 是带对偶复形 $\omega_X^\bullet$ 的局部 Noetherian 概形。
\begin{enumerate}
\item $X$ 是 Cohen-Macaulay 的 $\Leftrightarrow$
$\omega_X^\bullet$ 局部仅有一个非零上同调层；
\item $\mathcal{O}_{X, x}$ 是 Cohen-Macaulay 的 $\Leftrightarrow$
$\omega_{X, x}^\bullet$ 仅有一个非零上同调；
\item $U = \{x \in X \mid \mathcal{O}_{X, x}\text{ 为 Cohen--Macaulay}\}$
是开集且为 Cohen-Macaulay 概形。
\end{enumerate}
若 $X$ 连通且为 Cohen-Macaulay 概形，则存在整数 $n$ 和凝聚
Cohen-Macaulay $\mathcal{O}_X$-模 $\omega_X$，使得
$\omega_X^\bullet = \omega_X[-n]$。
\end{lemma}

\begin{proof}
由定义和《对偶复形》，引理
\ref{dualizing-lemma-dualizing-localize}，对每个 $x \in X$，
复形 $\omega_{X, x}^\bullet$ 是 $\mathcal{O}_{X, x}$ 上的对偶复形。
由《对偶复形》，引理 \ref{dualizing-lemma-apply-CM}，可知 (2) 成立。

\medskip\noindent
为证明 (3)，假设 $\mathcal{O}_{X, x}$ 是 Cohen-Macaulay 的。
令 $n_x$ 为使 $H^{n_{x}}(\omega_{X, x}^\bullet)$ 非零的唯一整数。
对 $x$ 的仿射邻域 $V \subset X$，有
$\omega_X^\bullet|_V$ 属于 $D^b_{\textit{Coh}}(\mathcal{O}_V)$，
故只有有限多个凝聚模 $H^i(\omega_X^\bullet)|_V$ 非零。
因此缩小 $V$ 后，可设只有 $H^{n_x}$ 非零；见《模》，引理
\ref{modules-lemma-finite-type-stalk-zero}。于是可知对每个
$v \in V$，$\mathcal{O}_{X, v}$ 都是 Cohen-Macaulay 的。
这同时证明 $U$ 是开集和 Cohen-Macaulay 概形。

\medskip\noindent
证明 (1)。蕴含 $\Leftarrow$ 由 (2) 得到。蕴含 $\Rightarrow$
由上一段的讨论得到：那里已经说明，若 $\mathcal{O}_{X, x}$
是 Cohen-Macaulay 的，则在 $x$ 的一个邻域中，复形
$\omega_X^\bullet$ 只有一个非零上同调层。

\medskip\noindent
假设 $X$ 连通且为 Cohen-Macaulay 概形。上面说明映射
$x \mapsto n_x$ 局部常值。由于 $X$ 连通，它是常值的，设其值为 $n$。
令 $\omega_X = H^n(\omega_X^\bullet)$。由《对偶复形》，引理
\ref{dualizing-lemma-apply-CM}（以及《概形的上同调》，定义
\ref{coherent-definition-Cohen-Macaulay}），$\omega_X$ 是
Cohen-Macaulay 的，故引理成立。
\end{proof}

\begin{lemma}
\label{duality-lemma-has-dualizing-module-CM-scheme}
设 $X$ 是局部 Noetherian 概形。若存在凝聚层 $\omega_X$，使得
$\omega_X[0]$ 是 $X$ 上的对偶复形，则 $X$ 是 Cohen-Macaulay 概形。
\end{lemma}

\begin{proof}
这由《对偶复形》，引理
\ref{dualizing-lemma-has-dualizing-module-CM} 和我们的定义立即得到。
\end{proof}

\begin{lemma}
\label{duality-lemma-affine-flat-Noetherian-CM}
在情形 \ref{duality-situation-shriek} 中，设 $f : X \to Y$ 是
$\textit{FTS}_S$ 的态射，$x \in X$。若 $f$ 平坦，则下列条件等价：
\begin{enumerate}
\item $f$ 在 $x$ 处是 Cohen-Macaulay 的；
\item $f^!\mathcal{O}_Y$ 在 $x$ 的一个邻域中只有一个非零上同调层。
\end{enumerate}
\end{lemma}

\begin{proof}
引理的一个方向由引理 \ref{duality-lemma-CM-shriek} 得到。为证明逆向，
可设 $f^!\mathcal{O}_Y$ 只有一个非零上同调层。令 $y = f(x)$。
设 $\xi_1, \ldots, \xi_n \in X_y$ 是纤维 $X_y$ 中特化到 $x$ 的泛点。
令 $d_1, \ldots, d_n$ 为 $X_y$ 相应不可约分支的维数。
由《态射进阶》，引理
\ref{more-morphisms-lemma-flat-finite-presentation-CM-open}，
态射 $f : X \to Y$ 在 $\eta_i$ 处是 Cohen-Macaulay 的。
所以由引理 \ref{duality-lemma-CM-shriek}，有 $d_1 = \ldots = d_n$。
若 $d$ 表示其共同值，则 $d = \dim_x(X_y)$。缩小 $X$ 后，可设所有纤维
维数至多为 $d$（《态射》，引理
\ref{morphisms-lemma-openness-bounded-dimension-fibres}）。于是由引理
\ref{duality-lemma-flat-shriek}，唯一的非零上同调层
$\omega = H^{-d}(f^!\mathcal{O}_Y)$ 在 $Y$ 上平坦。因此若
$h : X_y \to X$ 表示典范态射，则由《概形的导出范畴》，引理
\ref{perfect-lemma-tor-independence-and-tor-amplitude}，有
$Lh^*(f^!\mathcal{O}_Y) = Lh^*(\omega[d]) = (h^*\omega)[d]$。
所以由引理 \ref{duality-lemma-relative-dualizing-fibres}，
$h^*\omega[d]$ 是 $X_y$ 的对偶复形。进而由引理
\ref{duality-lemma-dualizing-module-CM-scheme}，$X_y$ 是 Cohen-Macaulay 的。
这如所需地证明 $f$ 在 $x$ 处是 Cohen-Macaulay 的。
\end{proof}

\begin{remark}
\label{duality-remark-CM-morphism-compare-dualizing}
在情形 \ref{duality-situation-shriek} 中，设 $f : X \to Y$ 是
$\textit{FTS}_S$ 的态射。假设 $f$ 是相对维数为 $d$ 的
Cohen-Macaulay 态射。令 $\omega_{X/Y} = H^{-d}(f^!\mathcal{O}_Y)$
为 $f^!\mathcal{O}_Y$ 的唯一非零上同调层；见引理
\ref{duality-lemma-CM-shriek}。则存在典范同构
$$
f^!K = Lf^*K \otimes_{\mathcal{O}_X}^\mathbf{L} \omega_{X/Y}[d]
$$
，其中 $K \in D^+_\QCoh(\mathcal{O}_Y)$；见引理
\ref{duality-lemma-perfect-comparison-shriek}。特别地，若 $S$
有对偶复形 $\omega_S^\bullet$，
$\omega_Y^\bullet = (Y \to S)^!\omega_S^\bullet$，且
$\omega_X^\bullet = (X \to S)^!\omega_S^\bullet$，则
$$
\omega_X^\bullet =
Lf^*\omega_Y^\bullet \otimes_{\mathcal{O}_X}^\mathbf{L} \omega_{X/Y}[d]
$$
。进一步，若 $X$ 和 $Y$ 连通且为 Cohen-Macaulay 概形，
而 $\omega_Y$ 和 $\omega_X$ 分别表示 $\omega_Y^\bullet$ 和
$\omega_X^\bullet$ 的唯一非零上同调层，则
$$
\omega_X = f^*\omega_Y \otimes_{\mathcal{O}_X} \omega_{X/Y}.
$$
类似结论对 $S$ 上任意有限型概形 $X$ 和 $Y$ 也成立
（即不必在 $S$ 上分离），其中对偶复形像第
\ref{duality-section-glue} 节那样相对于 $\omega_S^\bullet$ 规范化。
\end{remark}





\section{Gorenstein 概形}
\label{duality-section-gorenstein}

\noindent
本节是《对偶复形》，第 \ref{dualizing-section-gorenstein} 节的延续。

\begin{definition}
\label{duality-definition-gorenstein}
设 $X$ 是概形。若 $X$ 局部 Noetherian，且对所有 $x \in X$，
$\mathcal{O}_{X, x}$ 都是 Gorenstein 环，则称 $X$ 为{\it Gorenstein} 概形。
\end{definition}

\noindent
这个定义是合理的，因为一个 Noetherian 环称为 Gorenstein 环，当且仅当
其所有局部环都是 Gorenstein 环；见《对偶复形》，定义
\ref{dualizing-definition-gorenstein}。

\begin{lemma}
\label{duality-lemma-gorenstein-CM}
Gorenstein 概形是 Cohen-Macaulay 的。
\end{lemma}

\begin{proof}
在仿射局部考察，结论由代数中的相应结果得到，即《对偶复形》，引理
\ref{dualizing-lemma-gorenstein-CM}。
\end{proof}

\begin{lemma}
\label{duality-lemma-regular-gorenstein}
正则概形是 Gorenstein 的。
\end{lemma}

\begin{proof}
在仿射局部考察，结论由代数中的相应结果得到，即《对偶复形》，引理
\ref{dualizing-lemma-regular-gorenstein}。
\end{proof}

\begin{lemma}
\label{duality-lemma-gorenstein}
设 $X$ 是局部 Noetherian 概形。
\begin{enumerate}
\item 若 $X$ 有对偶复形 $\omega_X^\bullet$，则
\begin{enumerate}
\item $X$ 是 Gorenstein 的 $\Leftrightarrow$
$\omega_X^\bullet$ 是 $D(\mathcal{O}_X)$ 中的可逆对象；
\item $\mathcal{O}_{X, x}$ 是 Gorenstein 的 $\Leftrightarrow$
$\omega_{X, x}^\bullet$ 是 $D(\mathcal{O}_{X, x})$ 中的可逆对象；
\item $U = \{x \in X \mid \mathcal{O}_{X, x}\text{ is Gorenstein}\}$
是 Gorenstein 开子概形。
\end{enumerate}
\item 若 $X$ 是 Gorenstein 的，则 $X$ 有对偶复形，当且仅当
$\mathcal{O}_X[0]$ 是对偶复形。
\end{enumerate}
\end{lemma}

\begin{proof}
在仿射局部考察，结论由代数中的相应结果得到，即《对偶复形》，引理
\ref{dualizing-lemma-gorenstein}。
\end{proof}

\begin{lemma}
\label{duality-lemma-gorenstein-lci}
若 $f : Y \to X$ 是局部完全交态射，且 $X$ 是 Gorenstein 概形，
则 $Y$ 是 Gorenstein 概形。
\end{lemma}

\begin{proof}
由《态射进阶》，引理 \ref{more-morphisms-lemma-affine-lci}，
只需证明关于环同态的相应陈述。这就是《对偶复形》，引理
\ref{dualizing-lemma-gorenstein-lci}。
\end{proof}

\begin{lemma}
\label{duality-lemma-gorenstein-local-syntomic}
性质 $\mathcal{P}(S) =$``$S$ is Gorenstein'' 在 syntomic 拓扑中是局部的。
\end{lemma}

\begin{proof}
设 $\{S_i \to S\}$ 是 syntomic 覆盖。概形 $S$ 局部 Noetherian，
当且仅当每个 $S_i$ 都是 Noetherian 的；见《下降》，引理
\ref{descent-lemma-Noetherian-local-fppf}。因此现在可设 $S$ 和 $S_i$
都局部 Noetherian。若 $S$ 是 Gorenstein 的，则由引理
\ref{duality-lemma-gorenstein-lci}，每个 $S_i$ 都是 Gorenstein 的。
反过来，若每个 $S_i$ 都是 Gorenstein 的，则对每个点
$s \in S$，可选取 $i$ 和映到 $s$ 的 $t \in S_i$。
于是 $\mathcal{O}_{S, s} \to \mathcal{O}_{S_i, t}$ 是平坦局部环同态，
且 $\mathcal{O}_{S_i, t}$ 是 Gorenstein 的。因此由《对偶复形》，引理
\ref{dualizing-lemma-flat-under-gorenstein}，$\mathcal{O}_{S, s}$
是 Gorenstein 的。
\end{proof}






\section{Gorenstein 态射}
\label{duality-section-gorenstein-morphisms}

\noindent
本节是一个系列的一部分。关于正规态射、正则态射和 Cohen-Macaulay 态射的
相应各节见《态射进阶》，第
\ref{more-morphisms-section-normal}、
\ref{more-morphisms-section-regular} 和
\ref{more-morphisms-section-CM} 节。

\medskip\noindent
下面的引理说明，定义几何 Gorenstein 概形没有意义，因为它们与
Gorenstein 概形相同。

\begin{lemma}
\label{duality-lemma-gorenstein-base-change}
设 $X$ 是域 $k$ 上的局部 Noetherian 概形，$k'/k$ 是有限生成域扩张。
设 $x \in X$ 是点，$x' \in X_{k'}$ 是位于 $x$ 上方的点。则
$$
\mathcal{O}_{X, x}\text{ is Gorenstein}
\Leftrightarrow
\mathcal{O}_{X_{k'}, x'}\text{ is Gorenstein}
$$
。若 $X$ 在 $k$ 上局部有限型，则对任意域扩张 $k'/k$ 也有同样结论。
\end{lemma}

\begin{proof}
在两种情形下，环同态
$\mathcal{O}_{X, x} \to \mathcal{O}_{X_{k'}, x'}$ 都是 Noetherian
局部环之间的忠实平坦局部同态。因此若
$\mathcal{O}_{X_{k'}, x'}$ 是 Gorenstein 的，则由《对偶复形》，引理
\ref{dualizing-lemma-flat-under-gorenstein}，$\mathcal{O}_{X, x}$
也是 Gorenstein 的。反向上升也使用《对偶复形》，引理
\ref{dualizing-lemma-flat-under-gorenstein}。所以必须证明
$$
\mathcal{O}_{X_{k'}, x'}/\mathfrak m_x \mathcal{O}_{X_{k'}, x'} =
\kappa(x) \otimes_k k'
$$
是 Gorenstein 的。注意在第一种情形下 $k \to k'$ 有限生成，
在第二种情形下 $k \to \kappa(x)$ 有限生成。因此结论来自
Gorenstein 满足性质 (A)；见《对偶复形》，引理
\ref{dualizing-lemma-formal-fibres-gorenstein}。
\end{proof}

\noindent
上述引理保证下面确实是 Gorenstein 态射的正确定义。

\begin{definition}
\label{duality-definition-gorenstein-morphism}
设 $f : X \to Y$ 是概形态射。假设所有纤维 $X_y$ 都是局部
Noetherian 概形。
\begin{enumerate}
\item 设 $x \in X$ 且 $y = f(x)$。若 $f$ 在 $x$ 处平坦，
而概形 $X_y$ 在 $x$ 处的局部环是 Gorenstein 的，则称 $f$
{\it 在 $x$ 处为 Gorenstein}；
\item 若 $f$ 在 $X$ 的每一点都是 Gorenstein 的，则称 $f$
为{\it Gorenstein 态射}。
\end{enumerate}
\end{definition}

\noindent
下面给出一个等价表述。

\begin{lemma}
\label{duality-lemma-gorenstein-morphism}
设 $f : X \to Y$ 是概形态射。假设 $f$ 的所有纤维都局部 Noetherian。
则下列条件等价：
\begin{enumerate}
\item $f$ 是 Gorenstein 的；并且
\item $f$ 平坦，且其纤维都是 Gorenstein 概形。
\end{enumerate}
\end{lemma}

\begin{proof}
这直接由定义得到。
\end{proof}

\begin{lemma}
\label{duality-lemma-gorenstein-CM-morphism}
Gorenstein 态射是 Cohen-Macaulay 的。
\end{lemma}

\begin{proof}
结论由引理 \ref{duality-lemma-gorenstein-CM} 和定义得到。
\end{proof}

\begin{lemma}
\label{duality-lemma-lci-gorenstein}
syntomic 态射是 Gorenstein 的。等价地，平坦局部完全交态射是
Gorenstein 的。
\end{lemma}

\begin{proof}
回忆 syntomic 态射平坦，且其纤维是域上的局部完全交；见《态射》，引理
\ref{morphisms-lemma-syntomic-flat-fibres}。由引理
\ref{duality-lemma-gorenstein-lci}，域上的局部完全交是 Gorenstein 概形，
故结论成立。性质“syntomic”和“平坦且为局部完全交态射”由《态射进阶》，引理
\ref{more-morphisms-lemma-flat-lci} 可知等价。
\end{proof}

\begin{lemma}
\label{duality-lemma-composition-gorenstein}
设 $f : X \to Y$ 和 $g : Y \to Z$ 是态射。假设 $f$、$g$ 和
$g \circ f$ 的纤维 $X_y$、$Y_z$、$X_z$ 都局部 Noetherian。
\begin{enumerate}
\item 若 $f$ 在 $x$ 处为 Gorenstein，且 $g$ 在 $f(x)$ 处为
Gorenstein，则 $g \circ f$ 在 $x$ 处为 Gorenstein；
\item 若 $f$ 和 $g$ 都是 Gorenstein 的，则 $g \circ f$ 是 Gorenstein 的；
\item 若 $g \circ f$ 在 $x$ 处为 Gorenstein，且 $f$ 在 $x$ 处平坦，
则 $f$ 在 $x$ 处为 Gorenstein，而 $g$ 在 $f(x)$ 处为 Gorenstein；
\item 若 $g \circ f$ 是 Gorenstein 的且 $f$ 平坦，则 $f$ 是
Gorenstein 的，而 $g$ 在 $f$ 的像中的每一点都是 Gorenstein 的。
\end{enumerate}
\end{lemma}

\begin{proof}
翻译为代数后，结论由《对偶复形》，引理
\ref{dualizing-lemma-flat-under-gorenstein} 得到。
\end{proof}

\begin{lemma}
\label{duality-lemma-flat-morphism-from-gorenstein-scheme}
\begin{slogan}
平坦纤维化的总空间具有 Gorenstein 性，蕴含基与纤维也具有 Gorenstein 性
\end{slogan}
设 $f : X \to Y$ 是局部 Noetherian 概形之间的平坦态射。
若 $X$ 是 Gorenstein 的，则 $f$ 是 Gorenstein 的，并且对所有
$x \in X$，$\mathcal{O}_{Y, f(x)}$ 都是 Gorenstein 的。
\end{lemma}

\begin{proof}
翻译为代数后，结论由《对偶复形》，引理
\ref{dualizing-lemma-flat-under-gorenstein} 得到。
\end{proof}

\begin{lemma}
\label{duality-lemma-base-change-gorenstein}
设 $f : X \to Y$ 是概形态射。假设所有纤维 $X_y$ 都是局部
Noetherian 概形。设 $Y' \to Y$ 局部有限型，并设
$f' : X' \to Y'$ 是 $f$ 的基变换。设 $x' \in X'$ 的像为 $x \in X$。
\begin{enumerate}
\item 若 $f$ 在 $x$ 处为 Gorenstein，则
$f' : X' \to Y'$ 在 $x'$ 处为 Gorenstein；
\item 若 $f$ 在 $x$ 处平坦，且 $f'$ 在 $x'$ 处为 Gorenstein，
则 $f$ 在 $x$ 处为 Gorenstein；
\item 若 $Y' \to Y$ 在 $f'(x')$ 处平坦，且 $f'$ 在 $x'$ 处为
Gorenstein，则 $f$ 在 $x$ 处为 Gorenstein。
\end{enumerate}
\end{lemma}

\begin{proof}
注意，对 $Y' \to Y$ 的假设蕴含：若 $y' \in Y'$ 映到 $y \in Y$，
则域扩张 $\kappa(y')/\kappa(y)$ 有限生成。因此所有纤维
$X'_{y'} = (X_y)_{\kappa(y')}$ 也都局部 Noetherian；见《簇》，引理
\ref{varieties-lemma-locally-Noetherian-base-change}。所以该引理有意义。
令 $y' = f'(x')$ 和 $y = f(x)$。于是得到局部环的交换图
$$
\xymatrix{
\mathcal{O}_{X', x'} & \mathcal{O}_{X, x} \ar[l] \\
\mathcal{O}_{Y', y'} \ar[u] & \mathcal{O}_{Y, y} \ar[l] \ar[u]
}
$$
，其中左上角是右上角与左下角在右下角上的张量积的一个局部化。

\medskip\noindent
假设 $f$ 在 $x$ 处为 Gorenstein。
$\mathcal{O}_{Y, y} \to \mathcal{O}_{X, x}$ 的平坦性蕴含
$\mathcal{O}_{Y', y'} \to \mathcal{O}_{X', x'}$ 的平坦性；见《代数》，引理
\ref{algebra-lemma-base-change-flat-up-down}。
$\mathcal{O}_{X, x}/\mathfrak m_y\mathcal{O}_{X, x}$ 是 Gorenstein 的
这一事实蕴含
$\mathcal{O}_{X', x'}/\mathfrak m_{y'}\mathcal{O}_{X', x'}$
是 Gorenstein 的；见引理 \ref{duality-lemma-gorenstein-base-change}。
所以 $f'$ 在 $x'$ 处为 Gorenstein。

\medskip\noindent
假设 $f$ 在 $x$ 处平坦，且 $f'$ 在 $x'$ 处为 Gorenstein。
$\mathcal{O}_{X', x'}/\mathfrak m_{y'}\mathcal{O}_{X', x'}$
是 Gorenstein 的这一事实蕴含
$\mathcal{O}_{X, x}/\mathfrak m_y\mathcal{O}_{X, x}$ 是 Gorenstein 的；
见引理 \ref{duality-lemma-gorenstein-base-change}。所以 $f$ 在 $x$ 处为 Gorenstein。

\medskip\noindent
假设 $Y' \to Y$ 在 $y'$ 处平坦，且 $f'$ 在 $x'$ 处为 Gorenstein。
$\mathcal{O}_{Y', y'} \to \mathcal{O}_{X', x'}$ 和
$\mathcal{O}_{Y, y} \to \mathcal{O}_{Y', y'}$ 的平坦性蕴含
$\mathcal{O}_{Y, y} \to \mathcal{O}_{X, x}$ 的平坦性；见《代数》，引理
\ref{algebra-lemma-base-change-flat-up-down}。
$\mathcal{O}_{X', x'}/\mathfrak m_{y'}\mathcal{O}_{X', x'}$
是 Gorenstein 的这一事实蕴含
$\mathcal{O}_{X, x}/\mathfrak m_y\mathcal{O}_{X, x}$ 是 Gorenstein 的；
见引理 \ref{duality-lemma-gorenstein-base-change}。所以 $f$ 在 $x$ 处为 Gorenstein。
\end{proof}

\begin{lemma}
\label{duality-lemma-flat-lft-base-change-gorenstein}
设 $f : X \to Y$ 是平坦且局部有限型的概形态射。则集合
$\{x \in X \mid f\text{ is Gorenstein at }x\}$
的形成与任意基变换交换。
\end{lemma}

\begin{proof}
该假设蕴含 $f$ 的任意纤维都在某个域上局部有限型，因而局部
Noetherian；任意基变换也同样如此。因此该陈述有意义。考察纤维后，
问题归结为：设 $X$ 是域 $k$ 上局部有限型的概形，设 $K/k$ 是域扩张，
并设点 $x_K \in X_K$ 的像为 $x \in X$。要证明：
$\mathcal{O}_{X_K, x_K}$ 是 Gorenstein 的，当且仅当
$\mathcal{O}_{X, x}$ 是 Gorenstein 的。

\medskip\noindent
可用少量代数如下解决该问题。选取包含 $x$ 的仿射开集
$\Spec(A) \subset X$。设 $x$ 对应于 $\mathfrak p \subset A$。
令 $A_K = A \otimes_k K$，则包含 $x_K$ 的仿射开集为
$\Spec(A_K) \subset X_K$。设 $x_K$ 对应于 $\mathfrak p_K \subset A_K$。
令 $\omega_A^\bullet$ 为 $A$ 的一个对偶复形。由《对偶复形》，引理
\ref{dualizing-lemma-base-change-dualizing-over-field}，
$\omega_A^\bullet \otimes_A A_K$ 是 $A_K$ 的一个对偶复形。
现在结论成立，因为
$A_\mathfrak p \to (A_K)_{\mathfrak p_K}$ 是 Noetherian 环的平坦局部同态，
所以 $(\omega_A^\bullet)_\mathfrak p$ 是 $D(A_\mathfrak p)$ 中的可逆对象，
当且仅当
$(\omega_A^\bullet)_\mathfrak p \otimes_{A_\mathfrak p} (A_K)_{\mathfrak p_K}$
是 $D((A_K)_{\mathfrak p_K})$ 中的可逆对象。
略去了一些细节；提示：考察上同调模。
\end{proof}

\begin{lemma}
\label{duality-lemma-affine-flat-Noetherian-gorenstein}
在情形 \ref{duality-situation-shriek} 中，设 $f : X \to Y$ 是
$\textit{FTS}_S$ 的态射，并设 $x \in X$。若 $f$ 平坦，则下列条件等价：
\begin{enumerate}
\item $f$ 在 $x$ 处为 Gorenstein；
\item $f^!\mathcal{O}_Y$ 在 $x$ 的某个邻域内与一个可逆对象同构。
\end{enumerate}
特别地，$f$ 为 Gorenstein 的点组成 $X$ 中的开集。
\end{lemma}

\begin{proof}
令 $\omega^\bullet = f^!\mathcal{O}_Y$。由引理
\ref{duality-lemma-relative-dualizing-fibres}，这是一个上同调层相干的有界复形，
且它在纤维 $X_y$ 上的导出限制 $Lh^*\omega^\bullet$ 是 $X_y$ 上的对偶复形。
以 $i : x \to X_y$ 表示该点的嵌入。则下列条件等价：
\begin{enumerate}
\item $f$ 在 $x$ 处为 Gorenstein；
\item $\mathcal{O}_{X_y, x}$ 是 Gorenstein 的；
\item $Lh^*\omega^\bullet$ 在 $x$ 的某个邻域内可逆；
\item $Li^* Lh^* \omega^\bullet$ 恰有一个非零上同调，且它在
$\kappa(x)$ 上的维数为 $1$；
\item $L(h \circ i)^* \omega^\bullet$ 恰有一个非零上同调，且它在
$\kappa(x)$ 上的维数为 $1$；
\item $\omega^\bullet$ 在 $x$ 的某个邻域内可逆。
\end{enumerate}
(1) 与 (2) 的等价性由定义得到（因为 $f$ 平坦）。
(2) 与 (3) 的等价性来自引理 \ref{duality-lemma-gorenstein}。
(3) 与 (4) 的等价性来自《代数进阶》，引理
\ref{more-algebra-lemma-lift-bounded-pseudo-coherent-to-perfect}。
(4) 与 (5) 等价，是因为 $Li^* Lh^* = L(h \circ i)^*$。
(5) 与 (6) 的等价性来自《代数进阶》，引理
\ref{more-algebra-lemma-lift-bounded-pseudo-coherent-to-perfect}。
由此引理得证。
\end{proof}

\begin{lemma}
\label{duality-lemma-flat-finite-presentation-characterize-gorenstein}
设 $f : X \to S$ 是平坦且局部有限表示的概形态射。设 $x \in X$ 的像为
$s \in S$，并令 $d = \dim_x(X_s)$。则下列条件等价：
\begin{enumerate}
\item $f$ 在 $x$ 处为 Gorenstein；
\item 存在 $x$ 的开邻域 $U \subset X$ 以及 $S$ 上的局部拟有限态射
$U \to \mathbf{A}^d_S$，它在 $x$ 处为 Gorenstein；
\item 存在 $x$ 的开邻域 $U \subset X$ 以及 $S$ 上的局部拟有限
Gorenstein 态射 $U \to \mathbf{A}^d_S$；
\item 对 $x$ 的任意开邻域 $U \subset X$ 及任意 $S$-态射
$g : U \to \mathbf{A}^d_S$，都有：
$g$ 在 $x$ 处拟有限 $\Rightarrow$ $g$ 在 $x$ 处为 Gorenstein。
\end{enumerate}
特别地，$f$ 为 Gorenstein 的点组成 $X$ 中的开集。
\end{lemma}

\begin{proof}
选取满足 $x \in U$ 的仿射开集 $U = \Spec(A) \subset X$，以及满足
$f(U) \subset V$ 的仿射开集 $V = \Spec(R) \subset S$。则 $R \to A$
是有限表示的平坦环映射。设 $\mathfrak p \subset A$ 是对应于 $x$ 的素理想。
用一个主局部化替换 $A$ 后，可以假设存在拟有限映射
$R[x_1, \ldots, x_d]  \to A$；见《代数》，引理
\ref{algebra-lemma-quasi-finite-over-polynomial-algebra}。
因此至少存在一对 $(U, g)$，其中 $U \subset X$ 是 $x$ 的开邻域，且
$g : U \to \mathbf{A}^d_S$ 是局部\footnote{若 $S$ 拟分离，则 $g$
是拟有限的。}拟有限态射。

\medskip\noindent
由此，该引理转化为下面的代数问题（略去转化过程）。给定有限表示的平坦映射
$R \to A$、素理想 $\mathfrak p \subset A$，以及在 $\mathfrak p$ 处拟有限的
$\varphi : R[x_1, \ldots, x_d] \to A$，则下列条件等价：
\begin{enumerate}
\item[(a)] $\Spec(\varphi)$ 在 $\mathfrak p$ 处为 Gorenstein；并且
\item[(b)] $\Spec(A) \to \Spec(R)$ 在 $\mathfrak p$ 处为 Gorenstein；
\item[(c)] $\Spec(A) \to \Spec(R)$ 在 $\mathfrak p$ 的某个开邻域内为
Gorenstein。
\end{enumerate}
在每种情形下，$R[x_1, \ldots, x_n] \to A$ 在 $\mathfrak p$ 处平坦，
故由平坦性的开性（《代数》，定理
\ref{algebra-theorem-openness-flatness}），可以假设
$R[x_1, \ldots, x_n] \to A$ 平坦（用适当的主局部化替换 $A$）。
由《代数》，引理 \ref{algebra-lemma-flat-finite-presentation-limit-flat}，
存在 $R_0 \subset R$ 和 $R_0[x_1, \ldots, x_n] \to A_0$，使得
$R_0$ 在 $\mathbf{Z}$ 上有限型，$R_0 \to A_0$ 有限型，并且
$R_0[x_1, \ldots, x_n] \to A_0$ 平坦。注意，由引理
\ref{duality-lemma-base-change-gorenstein}，有限型平坦态射为 Gorenstein 的点集
与基变换交换。如此便归结到 $R$ 为 Noetherian 的情形。

\medskip\noindent
于是可以假设 $X$ 和 $S$ 都是仿射的，且 $f$ 有如下分解：
$$
X \xrightarrow{g} \mathbf{A}^n_S \xrightarrow{p} S
$$
，其中 $g$ 平坦且拟有限，而 $S$ 是 Noetherian 的。于是 $X$ 和
$\mathbf{A}^n_S$ 在 $S$ 上分离，并且
$$
f^!\mathcal{O}_S = g^!p^!\mathcal{O}_S = g^!\mathcal{O}_{\mathbf{A}^n_S}[n]
$$
；这是上叹号函子的已知性质（引理
\ref{duality-lemma-upper-shriek-composition} 和
\ref{duality-lemma-shriek-affine-line}）。因此由引理
\ref{duality-lemma-affine-flat-Noetherian-gorenstein}，(a)、(b)、(c) 等价。
\end{proof}

\begin{lemma}
\label{duality-lemma-gorenstein-local-source-and-target}
性质 $\mathcal{P}(f)=$“$f$ 的纤维局部 Noetherian，且 $f$ 为
Gorenstein”在目标的 fppf 拓扑下是局部的，在源的 syntomic 拓扑下也是局部的。
\end{lemma}

\begin{proof}
有
$\mathcal{P}(f) =
\mathcal{P}_1(f) \wedge \mathcal{P}_2(f)$，
其中 $\mathcal{P}_1(f)=$“$f$ 平坦”，而
$\mathcal{P}_2(f)=$“$f$ 的纤维局部 Noetherian 且为 Gorenstein”。
已知 $\mathcal{P}_1$ 在源和目标的 fppf 拓扑下都是局部的；见《下降》，引理
\ref{descent-lemma-descending-property-flat} 和
\ref{descent-lemma-flat-fpqc-local-source}。所以只须处理 $\mathcal{P}_2$。

\medskip\noindent
设 $f : X \to Y$ 是概形态射，且
$\{\varphi_i : Y_i \to Y\}_{i \in I}$ 是 $Y$ 的一个 fppf 覆盖。
以 $f_i : X_i \to Y_i$ 表示 $f$ 沿 $\varphi_i$ 的基变换。
取 $i \in I$ 和一点 $y_i \in Y_i$，令 $y = \varphi_i(y_i)$。注意
$$
X_{i, y_i} = \Spec(\kappa(y_i)) \times_{\Spec(\kappa(y))} X_y.
$$
并且 $\kappa(y_i)/\kappa(y)$ 是有限生成域扩张。因此若 $X_y$ 局部
Noetherian，则 $X_{i, y_i}$ 也局部 Noetherian；见《簇》，引理
\ref{varieties-lemma-locally-Noetherian-base-change}。若进一步 $X_y$ 为
Gorenstein，则 $X_{i, y_i}$ 也为 Gorenstein；见引理
\ref{duality-lemma-gorenstein-base-change}。因此 $\mathcal{P}_2$ 在目标上是 fppf 局部的。

\medskip\noindent
设 $\{X_i \to X\}$ 是 $X$ 的一个 syntomic 覆盖，并设 $y \in Y$。
此时 $\{X_{i, y} \to X_y\}$ 是该纤维的一个 syntomic 覆盖。因此，由引理
\ref{duality-lemma-gorenstein-local-syntomic}，$\mathcal{P}_2$ 在源的 syntomic
拓扑下是局部的。
\end{proof}




\section{对偶复形进阶}
\label{duality-section-more-dualizing}

\noindent
下面是一些不太适合放在别处的引理。

\begin{lemma}
\label{duality-lemma-descent-ascent}
设 $f : X \to Y$ 是局部 Noetherian 概形的态射。假设下列条件之一成立：
\begin{enumerate}
\item $f$ 是满的 syntomic 态射；或
\item $f$ 是满的平坦局部完全交态射；或
\item $f$ 是有限型的满 Gorenstein 态射。
\end{enumerate}
则 $K \in D_\QCoh(\mathcal{O}_Y)$ 是 $Y$ 上的对偶复形，当且仅当
$Lf^*K$ 是 $X$ 上的对偶复形。
\end{lemma}

\begin{proof}
取仿射开集并使用《概形的导出范畴》，引理
\ref{perfect-lemma-affine-compare-bounded}，该结论转化为《对偶复形》，引理
\ref{dualizing-lemma-descent-ascent}。
\end{proof}







\section{域上固有概形的对偶性}
\label{duality-section-duality-proper-over-field}

\noindent
本节阐明上面关于对偶复形的一般材料对域上固有概形的对偶性所产生的结论。

\begin{lemma}
\label{duality-lemma-duality-proper-over-field}
设 $X$ 是域 $k$ 上的固有概形。存在一个具有下列性质的对偶复形
$\omega_X^\bullet$：
\begin{enumerate}
\item 仅当 $i \in [-\dim(X), 0]$ 时，$H^i(\omega_X^\bullet)$ 才可能非零；
\item $\omega_X = H^{-\dim(X)}(\omega_X^\bullet)$ 是相干
$(S_2)$-模，其支撑由维数为 $\dim(X)$ 的不可约分支组成；
\item $H^i(\omega_X^\bullet)$ 的支撑维数至多为 $-i$；
\item 对闭点 $x \in X$，模
$H^i(\omega_{X, x}^\bullet) \oplus \ldots \oplus H^0(\omega_{X, x}^\bullet)$
非零，当且仅当 $\text{depth}(\mathcal{O}_{X, x}) \leq -i$；
\item 对 $K \in D_\QCoh(\mathcal{O}_X)$，存在与平移和可区别三角相容的函子性同构
\footnote{由 Yoneda 引理，该性质在 $D_\QCoh(\mathcal{O}_X)$ 中将
$\omega_X^\bullet$ 刻画到唯一同构。由于 $\omega_X^\bullet$ 属于
$D^b_{\textit{Coh}}(\mathcal{O}_X)$，事实上只须考虑
$K \in D^b_{\textit{Coh}}(\mathcal{O}_X)$。}
$$
\Ext^i_X(K, \omega_X^\bullet) = \Hom_k(H^{-i}(X, K), k)
$$
；
\item 对 $X$ 上的拟相干层 $\mathcal{F}$，存在函子性同构
$\Hom(\mathcal{F}, \omega_X) = \Hom_k(H^{\dim(X)}(X, \mathcal{F}), k)$；并且
\item 若 $X \to \Spec(k)$ 光滑且相对维数为 $d$，则
$\omega_X^\bullet \cong \wedge^d\Omega_{X/k}[d]$ 且
$\omega_X \cong \wedge^d\Omega_{X/k}$。
\end{enumerate}
\end{lemma}

\begin{proof}
以 $f : X \to \Spec(k)$ 表示结构态射。设 $a$ 为该态射的推前的右伴随；
见引理 \ref{duality-lemma-twisted-inverse-image}。考虑相对对偶复形
$$
\omega_X^\bullet = a(\mathcal{O}_{\Spec(k)})
$$
。参见注 \ref{duality-remark-relative-dualizing-complex}。由于 $f$ 固有，按定义有
$f^!(\mathcal{O}_{\Spec(k)}) = a(\mathcal{O}_{\Spec(k)})$；见第
\ref{duality-section-upper-shriek} 节。应用引理 \ref{duality-lemma-shriek-dualizing}，
可知 $\omega_X^\bullet$ 是对偶复形。此外，$\omega_X^\bullet$ 和
$\omega_X$ 分别如例 \ref{duality-example-proper-over-local} 和例
\ref{duality-example-equidimensional-over-field} 中所述。

\medskip\noindent
(1)、(2)、(3) 来自引理 \ref{duality-lemma-vanishing-good-dualizing}。

\medskip\noindent
对闭点 $x \in X$，$\omega_{X, x}^\bullet$ 是
$\mathcal{O}_{X, x}$ 上的规范化对偶复形；见引理
\ref{duality-lemma-good-dualizing-normalized}。于是 (4) 来自《对偶复形》，引理
\ref{dualizing-lemma-depth-in-terms-dualizing-complex}。

\medskip\noindent
(5) 由构造成立，因为 $a$ 是
$Rf_* : D_\QCoh(\mathcal{O}_X) \to D(\mathcal{O}_{\Spec(k)}) = D(k)$
的右伴随，而后者可与 $K \mapsto R\Gamma(X, K)$ 等同。这里还使用了：
$k$-模的导出范畴 $D(k)$ 等同于分次 $k$-向量空间范畴。

\medskip\noindent
当 $\mathcal{F}$ 相干时，(6) 来自引理
\ref{duality-lemma-dualizing-module-proper-over-A}；一般情形则把 (5) 用于
$K = \mathcal{F}[0]$ 和 $i = -\dim(X)$ 并展开定义即可。

\medskip\noindent
(7) 来自引理 \ref{duality-lemma-smooth-proper}。
\end{proof}

\begin{remark}
\label{duality-remark-duality-proper-over-field}
设 $k$、$X$、$\omega_X^\bullet$ 如引理
\ref{duality-lemma-duality-proper-over-field} 中所述。经由 (5) 中的函子性同构，
复形 $\omega_X^\bullet$ 上的恒等映射对应于一个映射
$$
t : H^0(X, \omega_X^\bullet) \longrightarrow k
$$
。对 $D_\QCoh(\mathcal{O}_X)$ 中的任意 $K$，(5) 中
$\Hom(K, \omega_X^\bullet)$ 与 $H^0(X, K)^\vee$ 的等同对应于配对
$$
\Hom_X(K, \omega_X^\bullet) \times H^0(X, K) \longrightarrow k,\quad
(\alpha, \beta) \longmapsto t(\alpha(\beta))
$$
。这来自 (5) 中同构的函子性。类似地，对任意 $i \in \mathbf{Z}$ 得到配对
$$
\Ext^i_X(K, \omega_X^\bullet) \times H^{-i}(X, K) \longrightarrow k,\quad
(\alpha, \beta) \longmapsto t(\alpha(\beta))
$$
。为定义 $\alpha(\beta)$，这里把 $\alpha$ 看作态射
$K[-i] \to \omega_X^\bullet$，把 $\beta$ 看作 $H^0(X, K[-i])$ 的元素。
注意，当 $K$ 任意时，只知道该配对在一侧非退化：该配对分别给出
$\Hom_X(K, \omega_X^\bullet)$ 和\ $\Ext^i_X(K, \omega_X^\bullet)$，
与 $H^0(X, K)$ 的 $k$-线性对偶及\ $H^{-i}(X, K)$ 的线性对偶之间的同构；
反向一般不成立。若 $K$
属于 $D^b_{\textit{Coh}}(\mathcal{O}_X)$，则
$\Hom_X(K, \omega_X^\bullet)$、$\Ext_X(K, \omega_X^\bullet)$、
$H^0(X, K)$ 和 $H^i(X, K)$ 都是有限维 $k$-向量空间（由《概形的导出范畴》，引理
\ref{perfect-lemma-coherent-internal-hom} 和
\ref{perfect-lemma-direct-image-coherent-bdd-below}），
并且这些配对在通常意义下是完美的。
\end{remark}

\begin{remark}
\label{duality-remark-coherent-duality-proper-over-field}
继续注 \ref{duality-remark-duality-proper-over-field} 中的讨论，并沿用记号
$k$、$X$、$\omega_X^\bullet$、$t$。若 $\mathcal{F}$ 是相干
$\mathcal{O}_X$-模，则得到有限维 $k$-向量空间之间的完美配对
$$
\langle -, - \rangle :
\Ext^i_X(\mathcal{F}, \omega_X^\bullet) \times H^{-i}(X,\mathcal{F})
\longrightarrow k,\quad
(\alpha, \beta) \longmapsto t(\alpha(\beta))
$$
。这些配对满足下列（显然的）函子性：若
$\varphi : \mathcal{F} \to \mathcal{G}$ 是相干 $\mathcal{O}_X$-模的同态，则
$$
\langle \alpha \circ \varphi, \beta \rangle =
\langle \alpha, \varphi(\beta) \rangle
$$
对 $\alpha \in \Ext^i_X(\mathcal{G}, \omega_X^\bullet)$ 和
$\beta \in H^{-i}(X, \mathcal{F})$ 成立。换言之，经由这些配对，由
$\varphi$ 诱导的 $k$-线性映射
$\Ext^i_X(\mathcal{G}, \omega_X^\bullet) \to
\Ext^i_X(\mathcal{F}, \omega_X^\bullet)$，是由 $\varphi$ 诱导的
$k$-线性映射 $H^{-i}(X, \mathcal{F}) \to H^{-i}(X, \mathcal{G})$
的 $k$-线性对偶。按这种方式表述后，当 $\varphi$ 是拟相干
$\mathcal{O}_X$-模的同态时，结论仍然成立。
\end{remark}

\begin{lemma}
\label{duality-lemma-duality-proper-over-field-perfect}
设 $k$、$X$、$\omega_X^\bullet$ 如引理
\ref{duality-lemma-duality-proper-over-field} 中所述，并设
$t : H^0(X, \omega_X^\bullet) \to k$ 如注
\ref{duality-remark-duality-proper-over-field} 中所述。设
$E \in D(\mathcal{O}_X)$ 是完美的。则对所有 $i$，配对
$$
H^i(X, \omega_X^\bullet \otimes_{\mathcal{O}_X}^\mathbf{L} E^\vee)
\times
H^{-i}(X, E) \longrightarrow k, \quad
(\xi, \eta) \longmapsto
t((1_{\omega_X^\bullet} \otimes \epsilon)(\xi \cup \eta))
$$
都是完美的。这里 $\cup$ 表示《上同调》，第
\ref{cohomology-section-cup-product} 节的杯积，而
$\epsilon : E^\vee \otimes_{\mathcal{O}_X}^\mathbf{L} E \to \mathcal{O}_X$
如《上同调》，例 \ref{cohomology-example-dual-derived} 中所述。
\end{lemma}

\begin{proof}
用 $E[-i]$ 替换 $E$，便归结到 $i = 0$ 的情形。由《上同调》，引理
\ref{cohomology-lemma-ext-composition-is-cup}，该配对就是注
\ref{duality-remark-duality-proper-over-field} 中讨论的配对，故结论来自该注中的讨论。
\end{proof}

\begin{lemma}
\label{duality-lemma-duality-proper-over-field-CM}
设 $X$ 是域 $k$ 上的固有概形，它是 Cohen-Macaulay 的、等维的，且维数为
$d$。则引理 \ref{duality-lemma-duality-proper-over-field} 中的模 $\omega_X$
具有下列性质：
\begin{enumerate}
\item $\omega_X$ 是 $X$ 上的对偶模（第 \ref{duality-section-dualizing-module} 节）；
\item $\omega_X$ 是支撑为 $X$ 的相干 Cohen-Macaulay 模；
\item 对 $K \in D_\QCoh(X)$，存在与平移和可区别三角相容的函子性同构
$\Ext^i_X(K, \omega_X[d]) = \Hom_k(H^{-i}(X, K), k)$；
\item 对 $X$ 上的拟相干层 $\mathcal{F}$，存在函子性同构
$\Ext^{d - i}(\mathcal{F}, \omega_X) = \Hom_k(H^i(X, \mathcal{F}), k)$。
\end{enumerate}
\end{lemma}

\begin{proof}
由引理 \ref{duality-lemma-duality-proper-over-field}，$\omega_X$ 显然是对偶模
（因为它是对偶复形最左侧的非零上同调层）。有
$\omega_X^\bullet = \omega_X[d]$；并且因为 $X$ 是 Cohen-Macaulay 的，
$\omega_X$ 也是 Cohen-Macaulay 的；见引理
\ref{duality-lemma-dualizing-module-CM-scheme}。其余陈述由此与引理
\ref{duality-lemma-duality-proper-over-field} 中的相应陈述合并得到。
\end{proof}

\begin{remark}
\label{duality-remark-rework-duality-locally-free-CM}
设 $X$ 是域 $k$ 上的固有 Cohen-Macaulay 概形，它等维且维数为 $d$。
设 $\omega_X^\bullet$ 和 $\omega_X$ 如引理
\ref{duality-lemma-duality-proper-over-field} 中所述。由引理
\ref{duality-lemma-duality-proper-over-field-CM}，有
$\omega_X^\bullet = \omega_X[d]$。设
$t : H^d(X, \omega_X) \to k$ 是注
\ref{duality-remark-duality-proper-over-field} 中的映射。设 $\mathcal{E}$ 是有限局部自由
$\mathcal{O}_X$-模，其对偶为 $\mathcal{E}^\vee$。则有完美配对
$$
H^i(X, \omega_X \otimes_{\mathcal{O}_X} \mathcal{E}^\vee)
\times
H^{d - i}(X, \mathcal{E})
\longrightarrow
k,\quad
(\xi, \eta) \longmapsto t(1 \otimes \epsilon)(\xi \cup \eta))
$$
，其中 $\cup$ 是杯积，而
$\epsilon : \mathcal{E}^\vee \otimes_{\mathcal{O}_X} \mathcal{E}
\to \mathcal{O}_X$ 是求值映射。这是引理
\ref{duality-lemma-duality-proper-over-field-perfect} 的特殊情形。
\end{remark}

\noindent
下面对该对偶复形作一个相容性检验。

\begin{lemma}
\label{duality-lemma-sanity-check-duality}
设 $X$ 是域 $k$ 上的固有概形，且 $\omega_X^\bullet$、$\omega_X$
如引理 \ref{duality-lemma-duality-proper-over-field} 中所述。
\begin{enumerate}
\item 若对某个域 $k'$，$X \to \Spec(k)$ 分解为
$X \to \Spec(k') \to \Spec(k)$，则对态射 $X \to \Spec(k')$，
$\omega_X^\bullet$ 和 $\omega_X$ 仍如引理
\ref{duality-lemma-duality-proper-over-field} 中所述；
\item 若 $K/k$ 是域扩张，则 $\omega_X^\bullet$ 和 $\omega_X$
到基变换 $X_K$ 上的拉回，对态射 $X_K \to \Spec(K)$ 而言仍如引理
\ref{duality-lemma-duality-proper-over-field} 中所述。
\end{enumerate}
\end{lemma}

\begin{proof}
以 $f : X \to \Spec(k)$ 表示结构态射，并以
$f' : X \to \Spec(k')$ 表示给定的分解。在引理
\ref{duality-lemma-duality-proper-over-field} 的证明中，取了
$\omega_X^\bullet = a(\mathcal{O}_{\Spec(k)})$，其中 $a$ 是引理
\ref{duality-lemma-twisted-inverse-image} 对 $f$ 给出的右伴随。因此须证明
$a(\mathcal{O}_{\Spec(k)}) \cong a'(\mathcal{O}_{\Spec(k)})$，
其中 $a'$ 是引理 \ref{duality-lemma-twisted-inverse-image} 对 $f'$ 给出的右伴随。
由于 $k' \subset H^0(X, \mathcal{O}_X)$，可知 $k'/k$ 是有限扩张
（《概形的上同调》，引理
\ref{coherent-lemma-proper-over-affine-cohomology-finite}）。
由伴随的唯一性，有 $a = a' \circ b$，其中 $b$ 是引理
\ref{duality-lemma-twisted-inverse-image} 对 $g : \Spec(k') \to \Spec(k)$ 给出的右伴随。
换言之，有 $f^! = (f')^! \circ g^!$。所以只须证明作为 $k'$-模有
$\Hom_k(k', k) \cong k'$；见例
\ref{duality-example-affine-twisted-inverse-image}。这是因为两者都是相同维数的
$k'$-向量空间（即维数为 $1$）。

\medskip\noindent
证明 (2)。结论成立，因为由引理 \ref{duality-lemma-more-base-change}，$a$
满足基变换。参见注 \ref{duality-remark-relative-dualizing-complex} 中的讨论。
\end{proof}








\section{相对对偶复形}
\label{duality-section-relative-dualizing-complexes}

\noindent
对有限表示的固有平坦态射，已有刚性相对对偶复形；见注
\ref{duality-remark-relative-dualizing-complex} 和引理 \ref{duality-lemma-van-den-bergh}。
对 Noetherian 概形之间的分离有限型态射 $f : X \to Y$，可以考虑
$f^!\mathcal{O}_Y$。本节为概形（不必局部 Noetherian）之间的平坦且局部
有限表示态射（不必拟分离或拟紧）定义相对对偶复形。我们将证明这类复形存在、
在唯一同构意义下唯一，并与上述情形一致。阅读本节前，请先阅读《对偶复形》，
第 \ref{dualizing-section-relative-dualizing-complexes} 节。

\begin{definition}
\label{duality-definition-relative-dualizing-complex}
设 $X \to S$ 是平坦且局部有限表示的概形态射。任取开集
$W \subset X \times_S X$，使得对角态射
$\Delta_{X/S} : X \to X \times_S X$ 经过闭浸入
$\Delta : X \to W$ 分解。所谓一个{\it 相对对偶复形}，是一个偶对
$(K, \xi)$，其中 $K \in D(\mathcal{O}_X)$ 是对象，且有
$D(\mathcal{O}_W)$ 中的映射
$$
\xi : \Delta_*\mathcal{O}_X \longrightarrow L\text{pr}_1^*K|_W
$$
，满足：
\begin{enumerate}
\item $K$ 是 $S$-完美的（《概形的导出范畴》，定义
\ref{perfect-definition-relatively-perfect}）；并且
\item $\xi$ 给出 $\Delta_*\mathcal{O}_X$ 与
$R\SheafHom_{\mathcal{O}_W}(
\Delta_*\mathcal{O}_X, L\text{pr}_1^*K|_W)$ 之间的同构。
\end{enumerate}
\end{definition}

\noindent
由引理 \ref{duality-lemma-sheaf-with-exact-support-ext}，条件 (2) 等价于存在
$D(\mathcal{O}_X)$ 中的同构
$$
\mathcal{O}_X \longrightarrow
R\SheafHom(\mathcal{O}_X, L\text{pr}_1^*K|_W)
$$
，其沿 $\Delta$ 的推前等于 $\xi$。由于
$R\SheafHom(\mathcal{O}_X, L\text{pr}_1^*K|_W)$ 与开集 $W$ 的选择无关，
偶对 $(K, \xi)$ 所成的范畴也与此选择无关。若 $X \to S$ 分离，则可取
$W = X \times_S X$。下面的引理将把许多论证归结到环的情形。

\begin{lemma}
\label{duality-lemma-relative-dualizing-complex-algebra}
设 $X \to S$ 是平坦且局部有限表示的概形态射，并设 $(K, \xi)$
是一个相对对偶复形。则对任意交换图
$$
\xymatrix{
\Spec(A) \ar[d] \ar[r] & X \ar[d] \\
\Spec(R) \ar[r] & S
}
$$
，若水平箭头都是开浸入，则 $K$ 在 $\Spec(A)$ 上的限制，经由《概形的导出范畴》，
引理 \ref{perfect-lemma-affine-compare-bounded}，对应于《对偶复形》，定义
\ref{dualizing-definition-relative-dualizing-complex} 意义下 $R \to A$
的一个相对对偶复形。
\end{lemma}

\begin{proof}
由于 $R\SheafHom$ 的形成与限制到开集交换，可以直接假设
$X = \Spec(A)$ 且 $S = \Spec(R)$。注意，$D(\mathcal{O}_X)$ 中的相对完美对象
是伪凝聚的，因而属于 $D_\QCoh(\mathcal{O}_X)$（《概形的导出范畴》，引理
\ref{perfect-lemma-pseudo-coherent}）。所以该陈述有意义。
由《概形的导出范畴》，引理
\ref{perfect-lemma-quasi-coherence-pushforward}、
\ref{perfect-lemma-quasi-coherence-pullback}、
\ref{perfect-lemma-quasi-coherence-internal-hom}，取 $\Delta_*$、
$L\text{pr}_1^*$、$R\SheafHom$ 与代数侧的相应运算相容。对于最后一项，注意
$L\text{pr}_1^*K$ 是 $S$-完美的（因而下有界），而
$\Delta_*\mathcal{O}_X$ 是 $D(\mathcal{O}_W)$ 中的伪凝聚对象；翻译为代数，
这表示 $A$ 作为 $A \otimes_R A$-模是伪凝聚的。这由《代数进阶》，引理
\ref{more-algebra-lemma-more-relative-pseudo-coherent-is-moot}
应用于 $R \to A \otimes_R A \to A$ 得到。因此恰好恢复《对偶复形》，定义
\ref{dualizing-definition-relative-dualizing-complex} 中的条件。
\end{proof}

\begin{lemma}
\label{duality-lemma-relative-dualizing-RHom}
设 $X \to S$ 是平坦且局部有限表示的概形态射，并设 $(K, \xi)$
是一个相对对偶复形。则
$\mathcal{O}_X \to R\SheafHom_{\mathcal{O}_X}(K, K)$
是同构。
\end{lemma}

\begin{proof}
在仿射局部考察，并使用引理 \ref{duality-lemma-relative-dualizing-complex-algebra}，
便归结到《对偶复形》，引理
\ref{dualizing-lemma-relative-dualizing-RHom} 的代数情形。
\end{proof}

\begin{lemma}
\label{duality-lemma-uniqueness-relative-dualizing}
设 $X \to S$ 是平坦且局部有限表示的概形态射。若 $(K, \xi)$ 和
$(L, \eta)$ 是 $X/S$ 上的两个相对对偶复形，则存在唯一的同构
$K \to L$，它把 $\xi$ 送到 $\eta$。
\end{lemma}

\begin{proof}
设 $U \subset X$ 是映入 $S$ 中某个仿射开集的仿射开集。由引理
\ref{duality-lemma-relative-dualizing-complex-algebra} 及《对偶复形》，引理
\ref{dualizing-lemma-uniqueness-relative-dualizing}，存在同构
$K|_U \to L|_U$。读者可以在概形情形中复用该引理的论证来证明本情形。
这里改用粘合论证。

\medskip\noindent
假设已有同构 $\alpha : K \to L$。则对某个可逆截面
$u \in H^0(W, \Delta_*\mathcal{O}_X) = H^0(X, \mathcal{O}_X)$，有
$\alpha(\xi) = u \eta$（因为按假设，$\eta$ 和 $\alpha(\xi)$
都是某个可逆 $\Delta_*\mathcal{O}_X$-模的生成元）。因此以
$u^{-1}\alpha$ 替换 $\alpha$ 后，有 $\alpha(\xi) = \eta$。
由引理 \ref{duality-lemma-relative-dualizing-RHom}，$K$ 的自同构群是
$H^0(X, \mathcal{O}_X^*)$，所以这样的 $\alpha$ 至多有一个。

\medskip\noindent
令 $\mathcal{B}$ 为 $X$ 中映入 $S$ 的某个仿射开集的全体仿射开集。
由前两段的讨论，对每个 $U \in \mathcal{B}$，都有唯一同构
$\alpha_U : K|_U \to L|_U$ 把 $\xi$ 送到 $\eta$。由引理
\ref{duality-lemma-relative-dualizing-RHom}，对 $X$ 的任意开集 $U$ 及
$i < 0$，有 $\text{Ext}^i(K|_U, K|_U) = 0$。把《上同调》，引理
\ref{cohomology-lemma-vanishing-and-glueing} 应用于
$\text{id} : X \to X$，得到唯一态射 $\alpha : K \to L$，
它对所有 $U \in \mathcal{B}$ 都与 $\alpha_U$ 相符。由于局部如此，
$\alpha$ 把 $\xi$ 送到 $\eta$。
\end{proof}

\begin{lemma}
\label{duality-lemma-existence-relative-dualizing}
设 $X \to S$ 是平坦且局部有限表示的概形态射。则存在相对对偶复形
$(K, \xi)$。
\end{lemma}

\begin{proof}
令 $\mathcal{B}$ 为 $X$ 中映入 $S$ 的某个仿射开集的全体仿射开集。
对每个 $U$，都有 $U$ 在 $S$ 上的相对对偶复形 $(K_U, \xi_U)$。
具体地，选取仿射开集 $V \subset S$，使得 $U \to X \to S$ 经过 $V$ 分解。
写成 $U = \Spec(A)$ 和 $V = \Spec(R)$。由《对偶复形》，引理
\ref{dualizing-lemma-base-change-relative-dualizing}，存在 $R \to A$
的相对对偶复形 $K_A \in D(A)$。反向应用引理
\ref{duality-lemma-relative-dualizing-complex-algebra} 的证明，这决定了一个
$V$-完美对象 $K_U \in D(\mathcal{O}_U)$，以及
$D(\mathcal{O}_{U \times_V U})$ 中的映射
$$
\xi : \Delta_*\mathcal{O}_U \to L\text{pr}_1^*K_U
$$
。由于 $V$-完美等同于 $S$-完美，且
$U \times_V U = U \times_S U$，可知 $(K_U, \xi_U)$ 满足要求。

\medskip\noindent
若 $U' \subset U \subset X$ 且 $U', U \in \mathcal{B}$，则由引理
\ref{duality-lemma-uniqueness-relative-dualizing}，有唯一同构
$\rho_{U'}^U : K_U|_{U'} \to K_{U'}$ 属于 $D(\mathcal{O}_{U'})$，并把
$\xi_U|_{U' \times_S U'}$ 送到 $\xi_{U'}$（显然，相对对偶复形到开集的限制
仍是相对对偶复形）。唯一性保证，对 $\mathcal{B}$ 中的
$U'' \subset U' \subset U$，有
$\rho^U_{U''} = \rho^V_{U''} \circ \rho ^U_{U'}|_{U''}$。
把引理 \ref{duality-lemma-relative-dualizing-RHom} 应用于 $U/S$ 和 $K_U$，可知对
$U \in \mathcal{B}$ 及 $i < 0$，有 $\text{Ext}^i(K_U, K_U) = 0$。
因此 BBD 粘合引理（《上同调》，定理
\ref{cohomology-theorem-glueing-bbd-general}）说明存在唯一的解，即对象
$K \in D(\mathcal{O}_X)$ 和同构 $\rho_U : K|_U \to K_U$，使得对所有
$U' \subset U$、$U, U' \in \mathcal{B}$，都有
$\rho^U_{U'} \circ \rho_U|_{U'} = \rho_{U'}$。

\medskip\noindent
为完成证明，还须构造 $D(\mathcal{O}_W)$ 中的映射
$$
\xi : \Delta_*\mathcal{O}_X \longrightarrow L\text{pr}_1^*K|_W
$$
，它诱导从 $\Delta_*\mathcal{O}_X$ 到
$R\SheafHom_{\mathcal{O}_W}(\Delta_*\mathcal{O}_X, L\text{pr}_1^*K|_W)$
的同构。由于可以改变 $W$，取
$W = \bigcup_{U \in \mathcal{B}} U \times_S U$。利用 $\rho_U$ 得到同构
$$
R\SheafHom_{\mathcal{O}_W}(
\Delta_*\mathcal{O}_X, L\text{pr}_1^*K|_W)|_{U \times_S U}
\xrightarrow{\rho_U}
R\SheafHom_{\mathcal{O}_{U \times_S U}}(
\Delta_*\mathcal{O}_U, L\text{pr}_1^*K_U)
$$
。由于这些开集 $U \times_S U$ 覆盖 $W$，可知
$R\SheafHom_{\mathcal{O}_W}(\Delta_*\mathcal{O}_X, L\text{pr}_1^*K|_W)$
的上同调层除次数 $0$ 外均为零。此外，得到同构
$$
H^0\left(U \times_S U, R\SheafHom_{\mathcal{O}_W}(\Delta_*\mathcal{O}_X,
L\text{pr}_1^*K|_W)\right)
\xrightarrow{\rho_U}
H^0\left((R\SheafHom_{\mathcal{O}_{U \times_S U}}(
\Delta_*\mathcal{O}_U, L\text{pr}_1^*K_U)\right)
$$
。令左端的元素 $\tau_U$ 在此映射下映到 $\xi_U$。当
$U' \subset U \subset X$ 且开集 $U', U \in \mathcal{B}$ 时，
$\rho^U_{U'}$、$\xi_U$、$\xi_{U'}$、$\rho_U$、$\rho_{U'}$ 之间的相容性
蕴含 $\tau_U|_{U' \times_S U'} = \tau_{U'}$。因此得到第 $0$ 上同调层
$H^0(R\SheafHom_{\mathcal{O}_W}(\Delta_*\mathcal{O}_X, L\text{pr}_1^*K|_W))$
的整体截面 $\tau$。由于
$R\SheafHom_{\mathcal{O}_W}(\Delta_*\mathcal{O}_X, L\text{pr}_1^*K|_W)$
的其余上同调层均为零，该整体截面 $\tau$ 决定所需的态射 $\xi$。
由于 $\xi$ 在 $U \times_S U$ 上的限制给出 $\xi_U$，可知它满足定义
\ref{duality-definition-relative-dualizing-complex} 的最后一个条件。
\end{proof}

\begin{lemma}
\label{duality-lemma-base-change-relative-dualizing}
考虑概形的笛卡尔方块
$$
\xymatrix{
X' \ar[d]_{f'} \ar[r]_{g'} & X \ar[d]^f \\
S' \ar[r]^g & S
}
$$
。假设 $X \to S$ 平坦且局部有限表示。设 $(K, \xi)$ 是 $f$
的一个相对对偶复形。令 $K' = L(g')^*K$，并令 $\xi'$ 为 $\xi$
的导出基变换（见证明）。则 $(K', \xi')$ 是 $f'$ 的相对对偶复形。
\end{lemma}

\begin{proof}
考虑笛卡尔方块
$$
\xymatrix{
X' \ar[d]_{\Delta_{X'/S'}} \ar[r] & X \ar[d]^{\Delta_{X/S}} \\
X' \times_{S'} X' \ar[r]^{g' \times g'} & X \times_S X
}
$$
。选取开集 $W \subset X \times_S X$，使 $\Delta_{X/S}$ 经过闭浸入
$\Delta : X \to W$ 分解。选取开集 $W' \subset X' \times_{S'} X'$，
使 $\Delta_{X'/S'}$ 经过闭浸入 $\Delta' : X \to W'$ 分解，且
$(g' \times g')(W') \subset W$。仍以
$g' \times g' : W' \to W$ 表示所诱导的态射。于是
$$
L(g' \times g')^*\Delta_*\mathcal{O}_X =
\Delta'_*\mathcal{O}_{X'}
\quad\text{且}\quad
L(g' \times g')^*L\text{pr}_1^*K|_W =
L\text{pr}_1^*K'|_{W'}
$$
。第一个等式成立，是因为 $X$ 与 $X' \times_{S'} X'$ 在
$X \times_S X$ 上 Tor 独立（例如见《态射进阶》，引理
\ref{more-morphisms-lemma-case-of-tor-independence}）。第二个等式来自导出拉回的传递性
（《上同调》，引理 \ref{cohomology-lemma-derived-pullback-composition}）。
因此 $\xi' = L(g' \times g')^*\xi$ 可视为映射
$$
\xi' : \Delta'_*\mathcal{O}_{X'} \longrightarrow L\text{pr}_1^*K'|_{W'}
$$
。至此，该引理的证明是直接的。首先，由《概形的导出范畴》，引理
\ref{perfect-lemma-base-change-relatively-perfect}，$K'$ 是 $S'$-完美的。
为检验 $\xi'$ 诱导从 $\Delta'_*\mathcal{O}_{X'}$ 到
$R\SheafHom_{\mathcal{O}_{W'}}(
\Delta'_*\mathcal{O}_{X'}, L\text{pr}_1^*K'|_{W'})$ 的同构，
可以在仿射局部工作。由引理
\ref{duality-lemma-relative-dualizing-complex-algebra}，归结到《对偶复形》，引理
\ref{dualizing-lemma-base-change-relative-dualizing} 中已证的相应代数陈述。
\end{proof}

\begin{lemma}
\label{duality-lemma-flat-proper-relative-dualizing}
设 $S$ 是拟紧拟分离概形，$f : X \to S$ 是有限表示的固有平坦态射。
注 \ref{duality-remark-relative-dualizing-complex} 中的相对对偶复形
$\omega_{X/S}^\bullet$，连同 (\ref{duality-equation-pre-rigid})，是定义
\ref{duality-definition-relative-dualizing-complex} 意义下的相对对偶复形。
\end{lemma}

\begin{proof}
引理 \ref{duality-lemma-properties-relative-dualizing} 已证明
$\omega_{X/S}^\bullet$ 是 $S$-完美的。令 $c$ 为引理
\ref{duality-lemma-twisted-inverse-image} 对对角态射
$\Delta : X \to X \times_S X$ 给出的右伴随。把 $\Delta_*$ 应用于
(\ref{duality-equation-pre-rigid})，得到同构
$$
\Delta_*\mathcal{O}_X \to
\Delta_*(c(L\text{pr}_1^*\omega_{X/S}^\bullet)) =
R\SheafHom_{\mathcal{O}_{X \times_S X}}(
\Delta_*\mathcal{O}_X, L\text{pr}_1^*\omega_{X/S}^\bullet)
$$
。该等式由引理 \ref{duality-lemma-twisted-inverse-image-closed} 和
\ref{duality-lemma-sheaf-with-exact-support-ext} 得到。证明完成。
\end{proof}

\begin{remark}
\label{duality-remark-relative-dualizing-complex-bis}
设 $X \to S$ 是有限表示的平坦固有概形态射。由引理
\ref{duality-lemma-existence-relative-dualizing}，存在定义
\ref{duality-definition-relative-dualizing-complex} 意义下的相对对偶复形
$(\omega_{X/S}^\bullet, \xi)$。考虑任意态射 $g : S' \to S$，其中
$S'$ 拟紧拟分离（例如 $S$ 的一个仿射开集）。由引理
\ref{duality-lemma-base-change-relative-dualizing}，
$(L(g')^*\omega_{X/S}^\bullet, L(g')^*\xi)$ 是基变换
$f' : X' \to S'$ 的相对对偶复形，取定义
\ref{duality-definition-relative-dualizing-complex} 的意义。设
$\omega_{X'/S'}^\bullet$ 是注 \ref{duality-remark-relative-dualizing-complex}
意义下 $X' \to S'$ 的相对对偶复形。结合引理
\ref{duality-lemma-flat-proper-relative-dualizing} 和
\ref{duality-lemma-uniqueness-relative-dualizing}，可知存在唯一同构
$$
\omega_{X'/S'}^\bullet \longrightarrow L(g')^*\omega_{X/S}^\bullet
$$
，它与 (\ref{duality-equation-pre-rigid}) 和 $L(g')^*\xi$ 相容。这些同构与
$S$ 上拟紧拟分离概形之间的态射，以及引理
\ref{duality-lemma-proper-flat-base-change} 的基变换同构相容（如果以后需要这种相容性，
我们将在这里仔细陈述并证明它）。
\end{remark}

\begin{lemma}
\label{duality-lemma-compactifyable-relative-dualizing}
在情形 \ref{duality-situation-shriek} 中，设 $f : X \to Y$ 是
$\textit{FTS}_S$ 的态射。若 $f$ 平坦，则 $f^!\mathcal{O}_Y$
是定义 \ref{duality-definition-relative-dualizing-complex} 意义下 $X$ 在 $Y$
上的某个相对对偶复形的第一分量。
\end{lemma}

\begin{proof}
由引理 \ref{duality-lemma-flat-shriek-relatively-perfect}，
$f^!\mathcal{O}_Y$ 是 $Y$-完美的。由于 $f$ 分离，对角态射
$\Delta : X \to X \times_Y X$ 是闭浸入，并且
$\Delta_*\Delta^!(-) =
R\SheafHom_{\mathcal{O}_{X \times_Y X}}(\mathcal{O}_X, -)$；见引理
\ref{duality-lemma-twisted-inverse-image-closed} 和
\ref{duality-lemma-sheaf-with-exact-support-ext}。因此，为完成证明，只须证明
$\Delta^!(L\text{pr}_1^*f^!(\mathcal{O}_Y)) \cong \mathcal{O}_X$，
其中 $\text{pr}_1 : X \times_Y X \to X$ 是第一投影。有
$$
\mathcal{O}_X = \Delta^! \text{pr}_1^!\mathcal{O}_X =
\Delta^! \text{pr}_1^! L\text{pr}_2^*\mathcal{O}_Y =
\Delta^!(L\text{pr}_1^* f^!\mathcal{O}_Y)
$$
，其中 $\text{pr}_2 : X \times_Y X \to X$ 是第二投影，并使用了引理
\ref{duality-lemma-base-change-shriek-flat} 的基变换同构
$\text{pr}_1^! \circ L\text{pr}_2^* = L\text{pr}_1^* \circ f^!$。
\end{proof}

\begin{lemma}
\label{duality-lemma-relative-dualizing-composition}
设 $f : Y \to X$ 和 $X \to S$ 是有限表示的平坦概形态射。设
$(K, \xi)$ 和 $(M, \eta)$ 分别是 $X \to S$ 和 $Y \to X$
的相对对偶复形。令 $E = M \otimes_{\mathcal{O}_Y}^\mathbf{L} Lf^*K$。
则对适当的 $\zeta$，$(E, \zeta)$ 是 $Y \to S$ 的相对对偶复形。
\end{lemma}

\begin{proof}
使用引理 \ref{duality-lemma-relative-dualizing-complex-algebra} 及该引理的代数版本
（《对偶复形》，引理
\ref{dualizing-lemma-relative-dualizing-composition}），可知在仿射局部，
$E$ 是某个相对对偶复形的第一分量。特别地，$E$ 是 $S$-完美的，
因为这可在仿射局部检验；见《概形的导出范畴》，引理
\ref{perfect-lemma-affine-locally-rel-perfect}。

\medskip\noindent
先在态射 $X \to S$ 和 $Y \to X$ 都分离的情形中证明 $\zeta$ 存在；
此时 $\Delta_{X/S}$、$\Delta_{Y/X}$、$\Delta_{Y/S}$ 都是闭浸入。
考虑下图
$$
\xymatrix{
& & Y \ar@{=}[r] & Y \ar[d]^f \\
Y \ar[r]_{\Delta_{Y/X}} &
Y \times_X Y \ar[d]_m \ar[r]_\delta \ar[ru]_q &
Y \times_S Y \ar[d]^{f \times f} \ar[ru]_p & X\\
& X \ar[r]^{\Delta_{X/S}} & X \times_S X \ar[ru]_r
}
$$
，其中 $p$、$q$、$r$ 是第一投影。由引理
\ref{duality-lemma-sheaf-with-exact-support-internal-home}，有
$$
R\SheafHom_{\mathcal{O}_{Y \times_S Y}}(
\Delta_{Y/S, *}\mathcal{O}_Y, Lp^*E) =
R\delta_*\left(R\SheafHom_{\mathcal{O}_{Y \times_X Y}}(
\Delta_{Y/X, *}\mathcal{O}_Y,
R\SheafHom(\mathcal{O}_{Y \times_X Y}, Lp^*E))\right)
$$
。由引理 \ref{duality-lemma-sheaf-with-exact-support-tensor}，有
$$
R\SheafHom(\mathcal{O}_{Y \times_X Y}, Lp^*E) =
R\SheafHom(\mathcal{O}_{Y \times_X Y}, L(f \times f)^*Lr^*K)
\otimes_{\mathcal{O}_{Y \times_S Y}}^\mathbf{L} Lq^*M
$$
。由引理 \ref{duality-lemma-flat-bc-sheaf-with-exact-support}，有
$$
R\SheafHom(\mathcal{O}_{Y \times_X Y}, L(f \times f)^*Lr^*K) =
Lm^*R\SheafHom(\mathcal{O}_X, Lr^*K)
$$
。最后一个表达式经 $\xi$ 同构于
$Lm^*\mathcal{O}_X = \mathcal{O}_{Y \times_X Y}$，故前一个表达式同构于
$Lq^*M$。于是
$$
R\SheafHom_{\mathcal{O}_{Y \times_S Y}}(
\Delta_{Y/S, *}\mathcal{O}_Y, Lp^*E) =
R\delta_*\left(R\SheafHom_{\mathcal{O}_{Y \times_X Y}}(
\Delta_{Y/X, *}\mathcal{O}_Y, Lq^*M)\right)
$$
。括号内的对象经 $\eta$ 同构于
$\Delta_{Y/X, *}*\mathcal{O}_X$。这就完成了分离情形的证明。

\medskip\noindent
在一般情形中，选取开集 $W \subset X \times_S X$，使
$\Delta_{X/S}$ 经过闭浸入 $\Delta : X \to W$ 分解；并选取开集
$V \subset Y \times_X Y$，使 $\Delta_{Y/X}$ 经过闭浸入
$\Delta' : Y \to V$ 分解。最后，选取开集 $W' \subset Y \times_S Y$，
使它与 $Y \times_X Y$ 的交为 $V$，并映入 $W$。于是考虑图
$$
\xymatrix{
& & Y \ar@{=}[r] & Y \ar[d]^f \\
Y \ar[r]_{\Delta'} &
V \ar[d]_m \ar[r]_\delta \ar[ru]_q &
W' \ar[d]^{f \times f} \ar[ru]_p & X\\
& X \ar[r]^\Delta & W \ar[ru]_r
}
$$
，并使用与前面完全相同的论证。
\end{proof}






\section{lci 态射的基本类}
\label{duality-section-fundamental-class}

\noindent
本节将使用第 \ref{duality-section-examples} 节中的计算。因此，本节的结果也具有
与该节相同类型的非唯一性。

\begin{lemma}
\label{duality-lemma-determinant}
设 $X$ 是局部赋环空间，并设
$$
\mathcal{E}_1 \xrightarrow{\alpha} \mathcal{E}_0 \to \mathcal{F} \to 0
$$
是 $\mathcal{O}_X$-模的短正合序列。假设 $\mathcal{E}_1$ 和
$\mathcal{E}_0$ 分别是秩 $r_1, r_0$ 的局部自由模。则存在典范映射
$$
\wedge^{r_0 - r_1}\mathcal{F}
\longrightarrow
\wedge^{r_1}(\mathcal{E}_1^\vee) \otimes \wedge^{r_0}\mathcal{E}_0
$$
；它在 $x \in X$ 处的茎上是同构，当且仅当 $\mathcal{F}$ 在 $x$
的某个开邻域内局部自由且秩为 $r_0 - r_1$。
\end{lemma}

\begin{proof}
若 $r_1 > r_0$，则按约定 $\wedge^{r_0 - r_1}\mathcal{F} = 0$，
唯一的映射不可能是同构。因此可以假设 $r = r_0 - r_1 \geq 0$。
用下式定义该映射：
$$
s_1 \wedge \ldots \wedge s_r \mapsto
t_1^\vee \wedge \ldots \wedge t_{r_1}^\vee \otimes
\alpha(t_1) \wedge \ldots \wedge \alpha(t_{r_1}) \wedge
\tilde s_1 \wedge \ldots \wedge \tilde s_r
$$
，其中 $t_1, \ldots, t_{r_1}$ 是 $\mathcal{E}_1$ 的一个局部基，
相应地 $t_1^\vee, \ldots, t_{r_1}^\vee$ 是 $\mathcal{E}_1^\vee$ 的对偶基，
而 $s'_i$ 是 $s_i$ 到 $\mathcal{E}_0$ 的一个局部截面的提升。
略去该定义为良定义的证明。

\medskip\noindent
若 $\mathcal{F}$ 局部自由且秩为 $r$，直接可验该映射是同构。反之，
假设该映射在 $x$ 处的茎上是同构。则 $\wedge^r\mathcal{F}_x$ 可逆，
因而 $\mathcal{F}_x$ 可由至多 $r$ 个元素生成。只有当 $\alpha$
模 $\mathfrak m_x$ 的秩为 $r$ 时才可能如此；换言之，$\alpha$
模 $\mathfrak m_x$ 具有最大秩。这蕴含 $\alpha$ 在 $x$ 的某个邻域内
具有最大秩，故 $\mathcal{F}$ 在该邻域内局部自由且秩为 $r$，正合所求。
\end{proof}

\begin{lemma}
\label{duality-lemma-fundamental-class-lci}
设 $Y$ 是 Noetherian 概形，$f : X \to Y$ 是局部完全交态射，且它分解为
浸入 $X \to P$ 后接固有光滑态射 $P \to Y$。设 $r$ 是 $X$
上的局部常值函数，使得
$\omega_{X/Y} = H^{-r}(f^!\mathcal{O}_Y)$ 是 $f^!\mathcal{O}_Y$
唯一的非零上同调层；见引理 \ref{duality-lemma-lci-shriek}。则存在映射
$$
\wedge^r\Omega_{X/Y} \longrightarrow \omega_{X/Y}
$$
；它在点 $x$ 处的茎上是同构，当且仅当 $f$ 在 $x$ 处光滑。
\end{lemma}

\begin{proof}
该假设蕴含 $X$ 在 $Y$ 上可紧化，故 $f^!$ 有定义；见第
\ref{duality-section-upper-shriek} 节。设 $j : W \to P$ 是开子概形，使
$X \to P$ 经过闭浸入 $i : X \to W$ 分解。此外，若
$g : P \to Y$ 为给定态射，则 $f^! = i^! \circ j^! \circ g^!$。
由引理 \ref{duality-lemma-smooth-proper}，有
$g^!\mathcal{O}_Y = \wedge^d\Omega_{P/Y}[d]$，其中 $d$ 是给出 $P$
在 $Y$ 上相对维数的局部常值函数。还有 $j^! = j^*$，并且
$i^!\mathcal{O}_W = \wedge^c\mathcal{N}[-c]$，其中 $c$ 是 $X$ 在 $W$
中的余维（$X$ 上的局部常值函数），而 $\mathcal{N}$ 是 Koszul 正则浸入
$i$ 的法层；见引理 \ref{duality-lemma-regular-immersion}。合并以上各式，得到
$$
f^!\mathcal{O}_Y =
\left(\wedge^c\mathcal{N} \otimes_{\mathcal{O}_X}
\wedge^d\Omega_{P/Y}|_X\right)[d - c]
$$
，这里还使用了引理 \ref{duality-lemma-perfect-comparison-shriek}。
因此作为 $X$ 上的局部常值函数，有 $r = d|_X - c$。态射
$X \to P$ 的余法层是模 $\mathcal{I}/\mathcal{I}^2$，其中
$\mathcal{I} \subset \mathcal{O}_W$ 是 $i$ 的理想层；见《态射》，第
\ref{morphisms-section-conormal-sheaf} 节。考虑《态射》，引理
\ref{morphisms-lemma-differentials-relative-immersion} 的典范正合序列
$$
\mathcal{I}/\mathcal{I}^2 \to
\Omega_{P/Y}|_X \to \Omega_{X/Y} \to 0
$$
。应用引理 \ref{duality-lemma-determinant} 即得到所需映射。

\medskip\noindent
若 $f$ 在 $x$ 处光滑，则由引理 \ref{duality-lemma-determinant} 以及
$\Omega_{X/Y}$ 在 $x$ 处局部自由且秩为 $r$，该映射是同构。反之，
假设所构造的映射在 $x$ 处的茎上是同构。该引理说明，用 $x$
的一个开邻域替换 $X$ 后，$\Omega_{X/Y}$ 自由且秩为 $r$。另一方面，
还可假设 $X = \Spec(A)$、$Y = \Spec(R)$，其中
$A = R[x_1, \ldots, x_n]/(f_1, \ldots, f_m)$，而
$f_1, \ldots, f_m$ 是 Koszul 正则序列（这来自局部完全交态射的定义）。
显然这蕴含 $r = n - m$。于是偏导数矩阵
$\partial f_j/\partial x_i$ 在 $x$ 处剩余域上的秩为 $m$。重排变量后，
可假设 $(\partial f_j/\partial x_i)_{1 \leq i, j \leq m}$ 的行列式
在 $x$ 的某个开邻域内可逆。因此 $R \to A$ 在该点光滑；例如见《代数》，
例 \ref{algebra-example-make-standard-smooth}。
\end{proof}

\begin{lemma}
\label{duality-lemma-fundamental-class-almost-lci}
设 $f : X \to Y$ 是概形态射，并设 $r \geq 0$。假设：
\begin{enumerate}
\item $Y$ 是 Cohen-Macaulay 的（《性质》，定义
\ref{properties-definition-Cohen-Macaulay}）；
\item $f$ 分解为 $X \to P \to Y$，其中第一个态射是浸入，第二个态射光滑且固有；
\item 若 $x \in X$ 且 $\dim(\mathcal{O}_{X, x}) \leq 1$，则 $f$
在 $x$ 处为 Koszul（《态射进阶》，定义
\ref{more-morphisms-definition-lci}）；并且
\item 若 $\xi$ 是 $X$ 的某个不可约分支的泛点，则
$\text{trdeg}_{\kappa(f(\xi))} \kappa(\xi) = r$。
\end{enumerate}
那么，令 $\omega_{X/Y} = H^{-r}(f^!\mathcal{O}_Y)$，存在映射
$$
\wedge^r\Omega_{X/Y} \longrightarrow \omega_{X/Y}
$$
，它在 $f$ 光滑的轨迹上是同构。
\end{lemma}

\begin{proof}
设 $U \subset X$ 是 $f$ 为局部完全交态射的开子概形。由假设 (4)，
$f$ 在所有泛点处的相对维数都是 $r$，故引理
\ref{duality-lemma-fundamental-class-lci} 中的局部常值函数恒等于 $r$，并得到映射
$$
\wedge^r\Omega_{X/Y}|_U = \wedge^r \Omega_{U/Y}
\longrightarrow
\omega_{U/Y} = \omega_{X/Y}|_U
$$
；它在 $f$ 的光滑点处是同构（该轨迹包含于 $U$，因为光滑态射是局部完全交态射）。
由引理 \ref{duality-lemma-shriek-over-CM} 及 $Y$ 为 Cohen-Macaulay 的假设，
模 $\omega_{X/Y}$ 满足 $(S_2)$。由条件 (3)，$U$ 包含所有余维为 $1$
的点；再用《除子》，引理 \ref{divisors-lemma-depth-2-hartog}，得到
$j_*\omega_{U/Y} = \omega_{X/Y}$。因此 $U$ 上的映射延拓到 $X$，证明完成。
\end{proof}






\section{相干模的零延拓}
\label{duality-section-extension-by-zero}

\noindent
本节及随后几节的材料可见 \cite{RD} 中 Deligne 所写的附录。

\medskip\noindent
本节中，$j : U \to X$ 是 Noetherian 概形的开浸入。考虑如下构造的
$D^b_{\textit{Coh}}(\mathcal{O}_X)$ 中的逆系统 $(K_n)$。设
$\mathcal{F}^\bullet$ 是相干 $\mathcal{O}_X$-模的有界复形，并设
$\mathcal{I} \subset \mathcal{O}_X$ 是拟相干理想层，满足
$V(\mathcal{I}) = X \setminus U$。可令
$$
K_n = \mathcal{I}^n\mathcal{F}^\bullet
$$
。更精确地，$K_n$ 是 $D^b_{\textit{Coh}}(\mathcal{O}_X)$ 中由如下复形
表示的对象：它的次数 $q$ 项是 $\mathcal{F}^q$ 的相干子模
$\mathcal{I}^n\mathcal{F}^q$。注意，映射
$\ldots \to K_3 \to K_2 \to K_1$ 限制到 $U$ 后诱导同构。称这样的系统为
{\it Deligne 系统}。

\begin{lemma}
\label{duality-lemma-lift-map}
设 $j : U \to X$ 是 Noetherian 概形的开浸入。设 $(K_n)$ 是 Deligne 系统，
并以 $K \in D^b_{\textit{Coh}}(\mathcal{O}_U)$ 表示常值系统
$(K_n|_U)$ 的值。设 $L$ 是 $D^b_{\textit{Coh}}(\mathcal{O}_X)$ 中的对象。
则 $\colim \Hom_X(K_n, L) = \Hom_U(K, L|_U)$。
\end{lemma}

\begin{proof}
设 $L \to M \to N \to L[1]$ 是
$D^b_{\textit{Coh}}(\mathcal{O}_X)$ 中的可区别三角。得到交换图
$$
\xymatrix{
\ldots \ar[r] &
\colim \Hom_X(K_n, L) \ar[r] \ar[d] &
\colim \Hom_X(K_n, M) \ar[r] \ar[d] &
\colim \Hom_X(K_n, N) \ar[r] \ar[d] &
\ldots \\
\ldots \ar[r] &
\Hom_U(K, L|_U) \ar[r] &
\Hom_U(K, M|_U) \ar[r] &
\Hom_U(K, N|_U) \ar[r] &
\ldots
}
$$
。由《导出范畴》，引理 \ref{derived-lemma-representable-homological} 和
《代数》，引理 \ref{algebra-lemma-directed-colimit-exact}，两行都正合。
因此，若引理的陈述对 $N[-1]$、$L$、$N$、$L[1]$ 成立，则由五引理，
它对 $M$ 也成立。由此，利用 $L$ 的典范截断的可区别三角（见《导出范畴》，
注 \ref{derived-remark-truncation-distinguished-triangle}），可归结到 $L$
仅有一个非零上同调层的情形。

\medskip\noindent
选取相干 $\mathcal{O}_X$-模的有界复形 $\mathcal{F}^\bullet$，以及定义
$X \setminus U$ 的拟相干理想 $\mathcal{I} \subset \mathcal{O}_X$，使
$K_n$ 由 $\mathcal{I}^n\mathcal{F}^\bullet$ 表示。利用“朴素”截断，得到逐项
分裂且彼此相容的复形短正合序列
$$
0 \to \sigma_{\geq a + 1} \mathcal{I}^n\mathcal{F}^\bullet \to
\mathcal{I}^n\mathcal{F}^\bullet \to
\sigma_{\leq a} \mathcal{I}^n\mathcal{F}^\bullet \to 0
$$
，它们进而对应于 $D^b_{\textit{Coh}}(\mathcal{O}_X)$ 中彼此相容的
可区别三角系统。重复上述论证，归结到 $\mathcal{F}^\bullet$
仅有一个非零项的情形，也就是下一段所讨论的情形。

\medskip\noindent
给定相干 $\mathcal{O}_X$-模 $\mathcal{F}$ 和相干 $\mathcal{O}_X$-模
$\mathcal{G}$，须证明对所有 $i \geq 0$，典范映射
$$
\colim \Ext^i_X(\mathcal{I}^n\mathcal{F}, \mathcal{G})
\longrightarrow
\Ext^i_U(\mathcal{F}|_U, \mathcal{G}|_U)
$$
都是同构。当 $i = 0$ 时，这就是《概形的上同调》，引理
\ref{coherent-lemma-homs-over-open}。以下假设 $i > 0$。

\medskip\noindent
单射性。设 $\xi \in \Ext^i_X(\mathcal{I}^n\mathcal{F}, \mathcal{G})$
的限制在 $U$ 上为零。须证明存在 $m \geq n$，使 $\xi$ 到
$\mathcal{I}^m\mathcal{F} = \mathcal{I}^{m - n}\mathcal{I}^n\mathcal{F}$
上的限制为零。以 $\mathcal{I}^n\mathcal{F}$ 替换 $\mathcal{F}$ 后，
可以假设 $n = 0$；即有
$\xi \in \Ext^i_X(\mathcal{F}, \mathcal{G})$，其限制在 $U$ 上为零。
由《概形的导出范畴》，命题 \ref{perfect-proposition-DCoh}，有
$D^b_{\textit{Coh}}(\mathcal{O}_X) = D^b(\textit{Coh}(\mathcal{O}_X))$。
因此可以在相干 $\mathcal{O}_X$-模的 Abel 范畴中计算 $\Ext$ 群。
这说明存在满射 $\alpha : \mathcal{F}'' \to \mathcal{F}$，使得
$\xi \circ \alpha = 0$（这里使用了 $i > 0$）。令
$\mathcal{F}' = \Ker(\alpha)$，于是有短正合序列
$$
0 \to \mathcal{F}' \to \mathcal{F}'' \to \mathcal{F} \to 0
$$
。所以 $\xi$ 是某个元素
$\xi' \in \Ext^{i - 1}_X(\mathcal{F}', \mathcal{G})$ 的像，且该元素
到 $U$ 上的限制属于映射
$\Ext^{i - 1}_U(\mathcal{F}''|_U, \mathcal{G}|_U) \to
\Ext^{i - 1}_U(\mathcal{F}'|_U, \mathcal{G}|_U)$ 的像。
由 Artin--Rees，逆系统 $(\mathcal{I}^n\mathcal{F}')$ 和
$(\mathcal{I}^n \mathcal{F}'' \cap \mathcal{F}')$ pro-同构；见《概形的上同调》，
引理 \ref{coherent-lemma-Artin-Rees}。由于有相容的短正合序列系统
$$
0 \to
\mathcal{F}' \cap \mathcal{I}^n\mathcal{F}'' \to
\mathcal{I}^n\mathcal{F}'' \to
\mathcal{I}^n\mathcal{F} \to 0
$$
，得到交换图
$$
\xymatrix{
\colim \Ext^{i - 1}_X(\mathcal{I}^n\mathcal{F}'', \mathcal{G})
\ar[r] \ar[d] &
\colim \Ext^{i - 1}_X(\mathcal{F}' \cap \mathcal{I}^n\mathcal{F}'', \mathcal{G})
\ar[r] \ar[d] &
\colim \Ext^i_X(\mathcal{I}^n\mathcal{F}, \mathcal{G})
\ar[d] \\
\Ext^{i - 1}_U(\mathcal{F}''|_U, \mathcal{G}|_U) \ar[r] &
\Ext^{i - 1}_U(\mathcal{F}'|_U, \mathcal{G}|_U) \ar[r] &
\Ext^{i - 1}_U(\mathcal{F}|_U, \mathcal{G}|_U)
}
$$
，两行都正合。对 $i$ 作归纳并使用上面对逆系统的说明，可知左边两个竖直箭头
都是同构。现在 $\xi$ 给出右上群中的元素，它是中上群中 $\xi'$ 的像；
后者又映到中下群中的一个元素，而该元素来自左下群中的某个元素。因此，对某个
$n$，$\xi$ 在 $\Ext^i_X(\mathcal{I}^n\mathcal{F}, \mathcal{G})$
中的像为零，正合所求。

\medskip\noindent
满射性。设 $\xi \in \Ext^i_U(\mathcal{F}|_U, \mathcal{G}|_U)$。
像上面一样论证并使用 $i > 0$，可找到相干 $\mathcal{O}_U$-模的满射
$\mathcal{H} \to \mathcal{F}|_U$，使 $\xi$ 在
$\Ext^i_U(\mathcal{H}, \mathcal{G}|_U)$ 中的像为零。然后可找到相干
$\mathcal{O}_X$-模的映射 $\varphi : \mathcal{F}'' \to \mathcal{F}$，
它限制到 $U$ 后就是 $\mathcal{H} \to \mathcal{F}|_U$；见《性质》，引理
\ref{properties-lemma-extend-finite-presentation}。注意，该引理并不保证
$\varphi$ 是满射，但这不影响论证（多做一点工作也可以选到满的 $\varphi$）。
记 $\mathcal{F}' = \Ker(\varphi)$。短正合序列
$$
0 \to \mathcal{F}'|_U \to \mathcal{F}''|_U \to \mathcal{F}|_U \to 0
$$
说明 $\xi$ 是
$\Ext^{i - 1}_U(\mathcal{F}'|_U, \mathcal{G}|_U)$ 中 $\xi'$ 的像。
对 $i$ 作归纳，可找到 $n$，使 $\xi'$ 是
$\Ext^{i - 1}_X(\mathcal{I}^n\mathcal{F}', \mathcal{G})$ 中某个
$\xi'_n$ 的像。由 Artin--Rees，可找到 $m \geq n$，使
$\mathcal{F}' \cap \mathcal{I}^m\mathcal{F}'' \subset
\mathcal{I}^n\mathcal{F}'$。利用短正合序列
$$
0 \to \mathcal{F}' \cap \mathcal{I}^m\mathcal{F}'' \to
\mathcal{I}^m\mathcal{F}'' \to \mathcal{I}^m\Im(\varphi) \to 0
$$
，$\xi'_n$ 在
$\Ext^{i - 1}_X(\mathcal{F}' \cap \mathcal{I}^m\mathcal{F}'', \mathcal{G})$
中的像，经边界映射映到
$\Ext^i_X(\mathcal{I}^m\Im(\varphi), \mathcal{G})$ 中的元素 $\xi_m$，
后者映到 $\xi$。由于 $\Im(\varphi)$ 与 $\mathcal{F}$ 在 $U$ 上相同，
$\mathcal{F}/\mathcal{I}^m\Im(\varphi)$ 支撑于 $X \setminus U$。
因此存在 $l \geq m$，使
$\mathcal{I}^l\mathcal{F} \subset \mathcal{I}^m\Im(\varphi)$；见《概形的上同调》，
引理 \ref{coherent-lemma-power-ideal-kills-sheaf}。取 $\xi_m$ 在
$\Ext^i_X(\mathcal{I}^l\mathcal{F}, \mathcal{G})$ 中的像即得结论。
\end{proof}

\begin{lemma}
\label{duality-lemma-lift-map-plus}
引理 \ref{duality-lemma-lift-map} 的结论对
$L \in D^+_{\textit{Coh}}(\mathcal{O}_X)$ 仍然成立。
\end{lemma}

\begin{proof}
事实上，若 $(K_n)$ 是 Deligne 系统，则存在 $b \in \mathbf{Z}$，使得当
$i > b$ 时 $H^i(K_n) = 0$。于是
$\Hom(K_n, L) = \Hom(K_n, \tau_{\leq b}L)$ 且
$\Hom(K, L) = \Hom(K, \tau_{\leq b}L)$。把该引理的结论用于
$\tau_{\leq b}L$ 即得。
\end{proof}

\begin{lemma}
\label{duality-lemma-extension-by-zero}
设 $j : U \to X$ 是 Noetherian 概形的开浸入。
\begin{enumerate}
\item 设 $(K_n)$ 和 $(L_n)$ 是 Deligne 系统，并设 $K$ 和 $L$ 分别是
常值系统 $(K_n|_U)$ 和 $(L_n|_U)$ 的值。给定 $D(\mathcal{O}_U)$ 中的态射
$\alpha : K \to L$，存在唯一的 $D^b_{\textit{Coh}}(\mathcal{O}_X)$
的 pro-系统态射 $(K_n) \to (L_n)$，其限制到 $U$ 后为 $\alpha$；
\item 给定 $K \in D^b_{\textit{Coh}}(\mathcal{O}_U)$，存在 Deligne 系统
$(K_n)$，使 $(K_n|_U)$ 是值为 $K$ 的常值系统；
\item (2) 中 $D^b_{\textit{Coh}}(\mathcal{O}_X)$ 的 pro-对象 $(K_n)$
在唯一同构意义下唯一（作为 pro-对象）。
\end{enumerate}
\end{lemma}

\begin{proof}
(1) 直接来自引理 \ref{duality-lemma-lift-map}，以及 pro-系统之间的态射等同于它们
所共表示的函子之间的态射这一事实；见《范畴》，注
\ref{categories-remark-pro-category-copresheaves}。

\medskip\noindent
设 $K$ 如 (2) 所述。可以选取 $K' \in D^b_{\textit{Coh}}(\mathcal{O}_X)$，
使其限制到 $U$ 后同构于 $K$；见《概形的导出范畴》，引理
\ref{perfect-lemma-lift-coherent}。由《概形的导出范畴》，命题
\ref{perfect-proposition-DCoh}，可用相干 $\mathcal{O}_X$-模的有界复形
$\mathcal{F}^\bullet$ 表示 $K'$。选取拟相干理想层
$\mathcal{I} \subset \mathcal{O}_X$，其零化轨迹为 $X \setminus U$
（例如令 $\mathcal{I}$ 对应于 $X \setminus U$ 上的约化诱导子概形结构）。
于是可像本节开头那样，令 $K_n$ 为复形
$\mathcal{I}^n\mathcal{F}^\bullet$ 所表示的对象。

\medskip\noindent
(3) 直接由 (1) 和 (2) 得到。
\end{proof}

\begin{lemma}
\label{duality-lemma-extension-by-zero-triangle}
设 $j : U \to X$ 是 Noetherian 概形的开浸入，并设
$$
K \to L \to M \to K[1]
$$
是 $D^b_{\textit{Coh}}(\mathcal{O}_U)$ 中的可区别三角。则存在
$D^b_{\textit{Coh}}(\mathcal{O}_X)$ 中的可区别三角的逆系统
$$
K_n \to L_n \to M_n \to K_n[1]
$$
，使 $(K_n)$、$(L_n)$、$(M_n)$ 都是 Deligne 系统，并且这些可区别三角
限制到 $U$ 后同构于原来的可区别三角。
\end{lemma}

\begin{proof}
取引理 \ref{duality-lemma-extension-by-zero} 的 (2) 中那样的 $(K_n)$。
选取 $D^b_{\textit{Coh}}(\mathcal{O}_X)$ 的对象 $L'$，使其限制到 $U$
后为 $L$（该引理说明这是可以做到的）。由引理 \ref{duality-lemma-lift-map}，
可找到 $n$ 以及 $X$ 上的态射 $K_n \to L'$，其限制到 $U$ 后为给定箭头
$K \to L$。因此存在 $D^b_{\textit{Coh}}(\mathcal{O}_X)$ 中的态射
$K' \to L'$，其限制到 $U$ 后为给定箭头 $K \to L$。

\medskip\noindent
由《概形的导出范畴》，命题 \ref{perfect-proposition-DCoh}，可找到相干
$\mathcal{O}_X$-模的有界复形之间的态射
$\alpha^\bullet : \mathcal{F}^\bullet \to
\mathcal{G}^\bullet$，它表示 $K' \to L'$。选取拟相干理想层
$\mathcal{I} \subset \mathcal{O}_X$，其零化轨迹为 $X \setminus U$。
令 $K_n = \mathcal{I}^n\mathcal{F}^\bullet$，且
$L_n = \mathcal{I}^n\mathcal{G}^\bullet$。注意，$\alpha^\bullet$ 诱导复形态射
$\alpha_n^\bullet : \mathcal{I}^n\mathcal{F}^\bullet \to
\mathcal{I}^n\mathcal{G}^\bullet$。由《导出范畴》，第
\ref{derived-section-cones} 节中的锥构造，显然有
$$
C(\alpha_n)^\bullet = \mathcal{I}^nC(\alpha^\bullet)
$$
，故可令 $M_n = C(\alpha_n)^\bullet$。也就是说，得到相容的可区别三角系统
（见《导出范畴》，第 \ref{derived-section-canonical-delta-functor} 节中的讨论）
$$
K_n \to L_n \to M_n \to K_n[1]
$$
；由公理 TR3 及《导出范畴》，引理
\ref{derived-lemma-third-isomorphism-triangle}，它限制到 $U$ 后同构于原来的
可区别三角。
\end{proof}

\begin{remark}
\label{duality-remark-extension-by-zero}
设 $j : U \to X$ 是 Noetherian 概形的开浸入。把
$K \in D^b_{\textit{Coh}}(\mathcal{O}_U)$ 送到其限制在 $U$ 上为 $K$
的 Deligne 系统，这决定了函子
$$
Rj_! :
D^b_{\textit{Coh}}(\mathcal{O}_U)
\longrightarrow
\text{Pro-}D^b_{\textit{Coh}}(\mathcal{O}_X)
$$
。由引理 \ref{duality-lemma-extension-by-zero-triangle}，该函子是“正合的”；
由引理 \ref{duality-lemma-lift-map}，它是函子
$j^* : D^b_{\textit{Coh}}(\mathcal{O}_X) \to
D^b_{\textit{Coh}}(\mathcal{O}_U)$ 的“左伴随”。
\end{remark}

\begin{remark}
\label{duality-remark-extension-by-zero-linear-pro-system}
设 $(A_n)$ 和 $(B_n)$ 是范畴 $\mathcal{C}$ 的逆系统。所谓从 $(A_n)$ 到
$(B_n)$ 的线性 pro-态射，是指对某些固定整数 $c, d \geq 1$，给定对所有
$n \geq 1$ 都彼此相容的态射族 $\varphi_n : A_{cn + d} \to B_n$。
若存在 $c'' \geq \max(c, c')$ 和 $d'' \geq \max(d, d')$，使得对所有 $n$，
所诱导的两个态射 $A_{c'' n + d''} \to B_n$ 都相同，则称
$(\varphi_n : A_{cn + d} \to B_n)$ 与
$(\psi_n : A_{c'n + d'} \to B_n)$ 决定同一个态射。给定 $U$ 上的值后，
Deligne 系统 $(K_n)$ 很可能在线性 pro-同构意义下是良定义的。如果以后需要，
我们将在这里仔细陈述并证明这一点。
\end{remark}

\begin{lemma}
\label{duality-lemma-deligne-system-2-out-of-3}
设 $j : U \to X$ 是 Noetherian 概形的开浸入。设
$$
K_n \to L_n \to M_n \to K_n[1]
$$
是 $D^b_{\textit{Coh}}(\mathcal{O}_X)$ 中可区别三角的逆系统。若
$(K_n)$ 和 $(M_n)$ 都 pro-同构于 Deligne 系统，则 $(L_n)$ 也是如此。
\end{lemma}

\begin{proof}
由于 $(K_n|_U)$ 和 $(M_n|_U)$ 都 pro-同构于常值系统，它们本质常值。
以 $K$ 和 $M$ 表示其值。由《导出范畴》，引理
\ref{derived-lemma-essentially-constant-2-out-of-3}，逆系统 $L_n|_U$
也本质常值；以 $L$ 表示其值。取
$N \in D^b_{\textit{Coh}}(\mathcal{O}_X)$，考虑交换图
$$
\xymatrix{
\ldots \ar[r] &
\colim \Hom_X(M_n, N) \ar[r] \ar[d] &
\colim \Hom_X(L_n, N) \ar[r] \ar[d] &
\colim \Hom_X(K_n, N) \ar[r] \ar[d] &
\ldots \\
\ldots \ar[r] &
\Hom_U(M, N|_U) \ar[r] &
\Hom_U(L, N|_U) \ar[r] &
\Hom_U(K, N|_U) \ar[r] &
\ldots
}
$$
。由引理 \ref{duality-lemma-lift-map} 以及同构的 ind-系统有相同余极限这一事实，
中间箭头左右各两个竖直箭头都是同构。由五引理，中间的竖直箭头也是同构。
现在，设 $(L'_n)$ 是限制到 $U$ 后常值为 $L$ 的 Deligne 系统
（由引理 \ref{duality-lemma-extension-by-zero}，它存在），则同样有
$\colim \Hom_X(L'_n, N) = \Hom_U(L, N|_U)$。因此由《范畴》，注
\ref{categories-remark-pro-category-copresheaves}，pro-系统 $(L_n)$ 和
$(L'_n)$ pro-同构。
\end{proof}

\begin{lemma}
\label{duality-lemma-consequence-Artin-Rees-bis}
设 $X$ 是 Noetherian 概形，$\mathcal{I} \subset \mathcal{O}_X$
是拟相干理想层，$\mathcal{F}^\bullet$ 是相干 $\mathcal{O}_X$-模的复形，
且 $p \in \mathbf{Z}$。令 $\mathcal{H} = H^p(\mathcal{F}^\bullet)$ 和
$\mathcal{H}_n = H^p(\mathcal{I}^n\mathcal{F}^\bullet)$。则有典范的
$\mathcal{O}_X$-模映射
$$
\ldots \to \mathcal{H}_3 \to \mathcal{H}_2 \to \mathcal{H}_1 \to \mathcal{H}
$$
。存在 $c > 0$，使得对 $n \geq c$，$\mathcal{H}_n \to \mathcal{H}$
的像包含于 $\mathcal{I}^{n - c}\mathcal{H}$，并有典范的
$\mathcal{O}_X$-模映射
$\mathcal{I}^n\mathcal{H} \to \mathcal{H}_{n - c}$，使复合
$$
\mathcal{I}^n \mathcal{H} \to \mathcal{H}_{n - c} \to
\mathcal{I}^{n - 2c}\mathcal{H}
\quad\text{且}\quad
\mathcal{H}_n \to \mathcal{I}^{n - c}\mathcal{H} \to \mathcal{H}_{n - 2c}
$$
都是典范映射。特别地，逆系统 $(\mathcal{H}_n)$ 和
$(\mathcal{I}^n\mathcal{H})$ 作为 $\textit{Mod}(\mathcal{O}_X)$
的 pro-对象同构。
\end{lemma}

\begin{proof}
若 $X$ 仿射，翻译为代数后，这就是《代数进阶》，引理
\ref{more-algebra-lemma-consequence-Artin-Rees-bis}。一般情形中，完全照搬
该引理的证明，只把代数中的 Artin--Rees 引用替换为《概形的上同调》，引理
\ref{coherent-lemma-Artin-Rees}。略去细节。
\end{proof}

\begin{lemma}
\label{duality-lemma-characterize-extension-by-zero-algebra}
设 $j : U \to X$ 是 Noetherian 概形的开浸入，$a \leq b$ 是整数，
且 $(K_n)$ 是 $D^b_{\textit{Coh}}(\mathcal{O}_X)$ 的逆系统，满足当
$i \not \in [a, b]$ 时 $H^i(K_n) = 0$。则下列条件等价：
\begin{enumerate}
\item $(K_n)$ pro-同构于一个 Deligne 系统；
\item 对每个 $p \in \mathbf{Z}$，存在相干 $\mathcal{O}_X$-模
$\mathcal{F}$，使 pro-系统 $(H^p(K_n))$ 和
$(\mathcal{I}^n\mathcal{F})$ pro-同构。
\end{enumerate}
\end{lemma}

\begin{proof}
假设 (1)。为证明 (2)，可以假设 $(K_n)$ 本身是 Deligne 系统。
按定义，可以选取相干 $\mathcal{O}_X$-模的有界复形
$\mathcal{F}^\bullet$，以及定义 $X \setminus U$ 的拟相干理想层，使
$K_n$ 由 $\mathcal{I}^n\mathcal{F}^\bullet$ 表示。结论于是来自引理
\ref{duality-lemma-consequence-Artin-Rees-bis}。

\medskip\noindent
假设 (2)。对 $b - a$ 作归纳，证明 $(K_n)$ 满足 (1)。若 $a = b$，
则 (1) 本质上就是假设。若 $a < b$，考虑相容的可区别三角系统
$$
\tau_{\leq a}K_n \to K_n \to \tau_{\geq a + 1}K_n \to (\tau_{\leq a}K_n)[1]
$$
；见《导出范畴》，注
\ref{derived-remark-truncation-distinguished-triangle}。对 $b - a$ 作归纳，
可知 $\tau_{\leq a}K_n$ 和 $\tau_{\geq a + 1}K_n$ 都 pro-同构于
Deligne 系统。再由引理 \ref{duality-lemma-deligne-system-2-out-of-3} 即得结论。
\end{proof}

\begin{lemma}
\label{duality-lemma-characterize-extension-by-zero}
设 $j : U \to X$ 是 Noetherian 概形的开浸入，$(K_n)$ 是
$D^b_{\textit{Coh}}(\mathcal{O}_X)$ 中的逆系统，并设
$X = W_1 \cup \ldots \cup W_r$ 是开覆盖。则下列条件等价：
\begin{enumerate}
\item $(K_n)$ pro-同构于一个 Deligne 系统；
\item 对每个 $i$，限制 $(K_n|_{W_i})$ 都 pro-同构于开浸入
$U \cap W_i \to W_i$ 的一个 Deligne 系统。
\end{enumerate}
\end{lemma}

\begin{proof}
对 $r$ 作归纳。若 $r = 1$，结论显然。假设 $r > 1$，令
$V = W_1 \cup \ldots \cup W_{r - 1}$。由归纳假设，$(K_n|_V)$
是 Deligne 系统。于是归结到下一段的讨论。

\medskip\noindent
假设 $X = V \cup W$ 是开覆盖，且 $(K_n|_W)$、$(K_n|_V)$ 都
pro-同构于 Deligne 系统。须证明 $(K_n)$ pro-同构于一个 Deligne 系统。
注意，$(K_n|_{V \cap W})$ pro-同构于一个 Deligne 系统
（由 Deligne 系统的构造，限制到开集显然保持这项性质）。特别地，pro-系统
$(K_n|_{U \cap V})$、$(K_n|_{U \cap W})$、
$(K_n|_{U \cap V \cap W})$ 都本质常值。由《上同调》，引理
\ref{cohomology-lemma-exact-sequence-j-star} 中的可区别三角和《导出范畴》，
引理 \ref{derived-lemma-essentially-constant-2-out-of-3}，可知
$(K_n|_U)$ 本质常值。以
$K \in D^b_{\textit{Coh}}(\mathcal{O}_U)$ 表示该系统的值。设 $L$
是 $D^b_{\textit{Coh}}(\mathcal{O}_X)$ 中的对象，考虑图
$$
\xymatrix{
\colim \Ext^{-1}(K_n|_V, L|_V) \oplus
\colim \Ext^{-1}(K_n|_W, L|_W) \ar[r] \ar[d] &
\Ext^{-1}(K|_{U \cap V}, L|_{U \cap V}) \oplus
\Ext^{-1}(K|_{U \cap W}, L|_{U \cap W}) \ar[d] \\
\colim \Ext^{-1}(K_n|_{V \cap W}, L|_{V \cap W}) \ar[r] \ar[d] &
\Ext^{-1}(K|_{U \cap V \cap W}, L|_{U \cap V \cap W}) \ar[d] \\
\colim \Hom(K_n, L) \ar[d] \ar[r] &
\Hom(K|_U, L|_U) \ar[d] \\
\colim \Hom(K_n|_V, L|_V) \oplus \colim \Hom(K_n|_W, L|_W) \ar[r] \ar[d] &
\Hom(K|_{U \cap V}, L|_{U \cap V}) \oplus
\Hom(K|_{U \cap W}, L|_{U \cap W}) \ar[d] \\
\colim \Hom(K_n|_{V \cap W}, L|_{V \cap W}) \ar[r] &
\Hom(K|_{U \cap V \cap W}, L|_{U \cap V \cap W})
}
$$
。由《上同调》，引理 \ref{cohomology-lemma-mayer-vietoris-hom}，
以及滤过余极限正合这一事实，两个竖直序列都正合。由引理
\ref{duality-lemma-lift-map} 以及 pro-同构的系统有相同余极限这一事实，除中间箭头外
的所有水平箭头都是同构。由五引理，中间箭头也是同构。因此 $(K_n)$
pro-同构于 $K$ 的一个 Deligne 系统。即若 $(K'_n)$ 是限制到 $U$
后常值为 $K$ 的 Deligne 系统（由引理 \ref{duality-lemma-extension-by-zero}，
它存在），则同样有 $\colim \Hom_X(K'_n, L) = \Hom_U(K, L|_U)$。
故由《范畴》，注 \ref{categories-remark-pro-category-copresheaves}，
pro-系统 $(K_n)$ 和 $(K'_n)$ pro-同构。
\end{proof}

\begin{lemma}
\label{duality-lemma-tensoring-Deligne-system}
设 $j : U \to X$ 是 Noetherian 概形的开浸入，并设拟相干理想层
$\mathcal{I} \subset \mathcal{O}_X$ 满足 $V(\mathcal{I}) = X \setminus U$。
设 $K$ 属于 $D^b_{\textit{Coh}}(\mathcal{O}_X)$。则
$$
K \otimes_{\mathcal{O}_X}^\mathbf{L} \mathcal{I}^n
$$
pro-同构于一个 Deligne 系统，它在 $U$ 上的常值为 $K|_U$。
\end{lemma}

\begin{proof}
由引理 \ref{duality-lemma-characterize-extension-by-zero}，该问题在 $X$ 上是局部的。
因此可以假设 $X$ 是某个 Noetherian 环的谱。此时结论来自代数版本，
即《代数进阶》，引理 \ref{more-algebra-lemma-tensoring-Deligne-system}。
\end{proof}







\section{紧支撑上同调的预备知识}
\label{duality-section-preliminaries-compactly-supported}

\noindent
在情形 \ref{duality-situation-shriek} 中，设 $f : X \to Y$ 是范畴
$\textit{FTS}_S$ 中的态射。利用上一节的构造，将构造函子
$$
Rf_! :
D^b_{\textit{Coh}}(\mathcal{O}_X)
\longrightarrow
\text{Pro-}D^b_{\textit{Coh}}(\mathcal{O}_Y)
$$
。若 $f$ 是开浸入，它退化为注 \ref{duality-remark-extension-by-zero} 中的函子；
一般地，它使用 $f$ 的紧化来构造。在此之前，需要下面的引理来证明该构造良定义。

\begin{lemma}
\label{duality-lemma-well-defined-pre}
设 $f : X \to Y$ 是 Noetherian 概形的固有态射，$V \subset Y$
是开子概形，并令 $U = f^{-1}(V)$。图示为
$$
\xymatrix{
U \ar[r]_j \ar[d]_g & X \ar[d]^f \\
V \ar[r]^{j'} & Y
}
$$
。则有函子的典范同构 $Rj'_! \circ Rg_* \to Rf_* \circ Rj_!$，其定义域和值域为
$D^b_{\textit{Coh}}(\mathcal{O}_U) \to
\text{Pro-}D^b_{\textit{Coh}}(\mathcal{O}_Y)$；其中 $Rj_!$ 和
$Rj'_!$ 如注 \ref{duality-remark-extension-by-zero} 中所述。
\end{lemma}

\begin{proof}[第一个证明]
设 $K$ 是 $D^b_{\textit{Coh}}(\mathcal{O}_U)$ 中的对象。设 $(K_n)$
是 $U \to X$ 的 Deligne 系统，其限制到 $U$ 后常值为 $K$。
这当然意味着 $(K_n)$ 在 $\text{Pro-}D^b_{\textit{Coh}}(\mathcal{O}_X)$
中表示 $Rj_!K$。注意，$Rj'_!Rg_*K$ 和 $Rf_*Rj_!K$ 限制到 $V$
后都是值为 $Rg_*K$ 的常值 pro-对象。对前者这是直接的；对后者，它来自
$(Rf_*K_n)|_V = Rg_*(K_n|_U) = Rg_*K$。由引理
\ref{duality-lemma-extension-by-zero} 中 Deligne 系统的唯一性，只须证明
$(Rf_*K_n)$ pro-同构于一个 Deligne 系统。所引引理还会说明所得同构具有函子性。

\medskip\noindent
证明 $(Rf_*K_n)$ pro-同构于一个 Deligne 系统。首先，由上述唯一性，
该问题与对应于 $K$ 的 Deligne 系统 $(K_n)$ 的选择无关。由引理
\ref{duality-lemma-extension-by-zero-triangle} 和
\ref{duality-lemma-deligne-system-2-out-of-3}，若
$D^b_{\textit{Coh}}(\mathcal{O}_U)$ 中有可区别三角
$$
K \to L \to M \to K[1]
$$
，且结论对 $K$ 和 $M$ 成立，则它对 $L$ 也成立。利用典范截断的可区别三角
（《导出范畴》，注 \ref{derived-remark-truncation-distinguished-triangle}），
归结到下一段研究的问题。

\medskip\noindent
设 $\mathcal{F}$ 是相干 $\mathcal{O}_X$-模，并设
$\mathcal{J} \subset \mathcal{O}_Y$ 是定义 $Y \setminus V$ 的拟相干理想层。
以 $\mathcal{J}^n\mathcal{F}$ 表示
$f^*\mathcal{J}^n \otimes \mathcal{F} \to \mathcal{F}$ 的像。
须证明 $(Rf_*(\mathcal{J}^n\mathcal{F}))$ 是 Deligne 系统。由引理
\ref{duality-lemma-characterize-extension-by-zero}，问题在 $Y$ 上是局部的。
因此可以假设 $Y = \Spec(A)$ 仿射，且 $\mathcal{J}$ 对应于理想
$I \subset A$。由引理 \ref{duality-lemma-characterize-extension-by-zero-algebra}，
只须证明对某个有限 $A$-模 $M$，上同调模的逆系统
$(H^p(X, I^n\mathcal{F}))$ pro-同构于逆系统 $(I^n M)$。
这已在《概形的上同调》，引理
\ref{coherent-lemma-cohomology-powers-ideal-application} 中证明。
\end{proof}

\begin{proof}[第二个证明]
设 $K$ 是 $D^b_{\textit{Coh}}(\mathcal{O}_U)$ 中的对象，$L$ 是
$D^b_{\textit{Coh}}(\mathcal{O}_Y)$ 中的对象。将构造双射
$$
\Hom_{\text{Pro-}D^b_{\textit{Coh}}(\mathcal{O}_Y)}(Rj'_!Rg_*K, L)
\longrightarrow
\Hom_{\text{Pro-}D^b_{\textit{Coh}}(\mathcal{O}_Y)}(Rf_*Rj_!K, L)
$$
，它对 $K$ 和 $L$ 都具有函子性。固定 $K$，由《范畴》，注
\ref{categories-remark-pro-category-copresheaves}，这将决定 pro-对象的同构
$Rf_*Rj_!K \to Rj'_!Rg_*K$；让 $K$ 变化，便得到函子的同构。
为实际构造该同构，使用下列函子性等式链：
\begin{align*}
\Hom_{\text{Pro-}D^b_{\textit{Coh}}(\mathcal{O}_Y)}(Rj'_!Rg_*K, L)
& =
\Hom_V(Rg_*K, L|_V) \\
& =
\Hom_U(K, g^!(L|_V)) \\
& =
\Hom_U(K, f^!L|_U)) \\
& =
\Hom_{\text{Pro-}D^b_{\textit{Coh}}(\mathcal{O}_X)}(Rj_!K, f^!L) \\
& =
\Hom_{\text{Pro-}D^b_{\textit{Coh}}(\mathcal{O}_Y)}(Rf_*Rj_!K, L)
\end{align*}
。第一个等式来自引理 \ref{duality-lemma-lift-map}。第二个等式成立，是因为 $g$
作为 $f$ 到 $V$ 的基变换是固有的，故按构造 $g^!$ 是推前的右伴随；见第
\ref{duality-section-upper-shriek} 节。第三个等式成立，是因为由引理
\ref{duality-lemma-restrict-before-or-after}，有 $g^!(L|_V) = f^!L|_U$。
由引理 \ref{duality-lemma-shriek-coherent}，$f^!L$ 属于
$D^+_{\textit{Coh}}(\mathcal{O}_X)$，故第四个等式来自引理
\ref{duality-lemma-lift-map-plus}。第五个等式再次来自：$f$ 固有时，$f^!$
是 $Rf_*$ 的右伴随。
\end{proof}

\begin{lemma}
\label{duality-lemma-well-defined}
设 $j : U \to X$ 是 Noetherian 概形的开浸入，$j' : U \to X'$
是 $U$ 在 $X$ 上的一个紧化（见证明），并以 $f : X' \to X$
表示结构态射。则有函子的典范同构 $Rj_! \to Rf_* \circ R(j')_!$，
其定义域和值域为 $D^b_{\textit{Coh}}(\mathcal{O}_U) \to
\text{Pro-}D^b_{\textit{Coh}}(\mathcal{O}_X)$；其中 $Rj_!$ 和
$Rj'_!$ 如注 \ref{duality-remark-extension-by-zero} 中所述。
\end{lemma}

\begin{proof}
$X'$ 是 $U$ 在 $X$ 上的紧化，恰好意味着 $f : X' \to X$ 固有、
$j'$ 是开浸入且 $j = f \circ j'$；见《平坦性进阶》，第
\ref{flat-section-compactify} 节。若 $j'(U) = f^{-1}(j(U))$，
则结论直接来自引理 \ref{duality-lemma-well-defined-pre}。若
$j'(U) \not = f^{-1}(j(U))$，以 $X'' \subset X'$ 表示
$j' : U \to X'$ 的概形论闭包，并以 $j'' : U \to X''$ 表示相应的开浸入。
图示为
$$
\xymatrix{
& & X'' \ar[d]^{f'} \\
& & X' \ar[d]^f \\
U \ar[rr]^j \ar[rru]^{j'} \ar[rruu]^{j''} & & X
}
$$
。由《平坦性进阶》，引理 \ref{flat-lemma-compactifications-cofiltered}
的 (c) 及上述讨论，有同构
$Rf'_* \circ Rj''_! = Rj'_!$ 和 $R(f \circ f')_* \circ Rj''_! = Rj_!$。
由于 $R(f \circ f')_* = Rf_* \circ Rf'_*$，结论成立。
\end{proof}

\begin{remark}
\label{duality-remark-covariance-open-j-lower-shriek}
设 $X \supset U \supset U'$ 是 Noetherian 概形 $X$ 的开子概形，
并以 $j : U \to X$ 和 $j' : U' \to X$ 表示包含态射。断言存在典范映射
$$
Rj'_!(K|_{U'}) \longrightarrow Rj_!K
$$
，它对 $D^b_{\textit{Coh}}(\mathcal{O}_U)$ 中的 $K$ 具有函子性。
事实上，由引理 \ref{duality-lemma-lift-map}，对
$D^b_{\textit{Coh}}(\mathcal{O}_X)$ 中任意 $L$，有映射
\begin{align*}
\Hom_{\text{Pro-}D^b_{\textit{Coh}}(\mathcal{O}_X)}(Rj_!K, L)
& =
\Hom_U(K, L|_U) \\
& \to
\Hom_{U'}(K|_{U'}, L|_{U'}) \\
& =
\Hom_{\text{Pro-}D^b_{\textit{Coh}}(\mathcal{O}_X)}(Rj'_!(K|_{U'}), L)
\end{align*}
，它对 $L$ 和 $K'$ 具有函子性。由对 $L$ 的函子性和《范畴》，注
\ref{categories-remark-pro-category-copresheaves}，得到典范映射
$Rj'_!(K|_{U'}) \to Rj_!K$；上述箭头对 $K$ 的函子性又说明该映射对 $K$
具有函子性。

\medskip\noindent
下面显式构造该箭头。设 $\mathcal{F}^\bullet$ 是相干
$\mathcal{O}_X$-模的有界复形，其限制到 $U$ 后在导出范畴中表示 $K$。
引理 \ref{duality-lemma-extension-by-zero} 的证明已说明这样的复形总是存在。
令 $\mathcal{I}$ 和\ $\mathcal{I}'$ 是 $X$ 上的拟相干理想层，分别满足
$V(\mathcal{I}) = X \setminus U$ 和\ $V(\mathcal{I}') = X \setminus U'$。
以 $\mathcal{I} + \mathcal{I}'$ 替换 $\mathcal{I}$ 后，可以假设
$\mathcal{I}' \subset \mathcal{I}$。按构造，$Rj_!K$ 和\ $Rj'_!(K|_{U'})$
分别由 $D^b_{\textit{Coh}}(\mathcal{O}_X)$ 的逆系统
$(K_n)$ 和\ $(K'_n)$ 表示，其中
$$
K_n = \mathcal{I}^n\mathcal{F}^\bullet
\quad\text{分别}\quad
K'_n = (\mathcal{I}')^n\mathcal{F}^\bullet
$$
。显然，上面构造的映射由包含
$(\mathcal{I}')^n \subset \mathcal{I}^n$ 所给出的映射 $K'_n \to K_n$ 给出。
\end{remark}






\section{相干模的紧支撑上同调}
\label{duality-section-compactly-supported}

\noindent
在情形 \ref{duality-situation-shriek} 中，给定 $\textit{FTS}_S$ 中的态射
$f : X \to Y$，定义函子
$$
Rf_! : D^b_{\textit{Coh}}(\mathcal{O}_X)
\longrightarrow
\text{Pro-}D^b_{\textit{Coh}}(\mathcal{O}_Y)
$$
。具体地，选取在 $Y$ 上的紧化 $j : X \to \overline{X}$；
由《平坦性进阶》，定理 \ref{flat-theorem-nagata} 和引理
\ref{flat-lemma-compactifyable}，这是可以做到的。以
$\overline{f} : \overline{X} \to Y$ 表示结构态射。对
$K \in D^b_{\textit{Coh}}(\mathcal{O}_X)$，令
$$
Rf_!K  = R\overline{f}_* Rj_! K
$$
，其中 $Rj_!$ 如注 \ref{duality-remark-extension-by-zero} 中所述。

\begin{lemma}
\label{duality-lemma-lower-shriek-well-defined}
函子 $Rf_!$ 在同构意义下与紧化的选择无关。
\end{lemma}

\noindent
事实上，函子 $Rf_!$ 将被刻画为 $f^!$ 的“左伴随”，从而在唯一同构意义下确定。

\begin{proof}
考虑 $X$ 在 $Y$ 上的紧化所成的范畴。由《平坦性进阶》，定理
\ref{flat-theorem-nagata} 及引理
\ref{flat-lemma-compactifications-cofiltered}、
\ref{flat-lemma-compactifyable}，该范畴余滤过。对紧化的每个选择
$$
j : X \to \overline{X},\quad \overline{f} : \overline{X} \to Y
$$
，上述构造都给出函子 $R\overline{f}_* \circ Rj_!$。设
$g : \overline{X}_1 \to \overline{X}_2$ 是 $Y$ 上两个紧化
$j_i : X \to \overline{X}_i$ 之间的态射。则得到同构
$$
R\overline{f}_{2, *} \circ Rj_{2, !} =
R\overline{f}_{2, *} \circ Rg_* \circ j_{1, !} =
R\overline{f}_{1, *} \circ Rj_{1, !}
$$
，其中第一个等式使用了引理 \ref{duality-lemma-well-defined}。由此可知，
该函子在同构意义下与紧化的选择无关。
\end{proof}

\begin{proposition}
\label{duality-proposition-duality-compactly-supported}
在情形 \ref{duality-situation-shriek} 中，设 $f : X \to Y$ 是
$\textit{FTS}_S$ 的态射。则函子 $Rf_!$ 与 $f^!$ 在下述意义下伴随：
对所有 $K \in D^b_{\textit{Coh}}(\mathcal{O}_X)$ 和
$L \in D^+_{\textit{Coh}}(\mathcal{O}_Y)$，有
$$
\Hom_X(K, f^!L) =
\Hom_{\text{Pro-}D^+_{\textit{Coh}}(\mathcal{O}_Y)}(Rf_!K, L)
$$
，并且它对 $K$ 和 $L$ 双函子性地成立。
\end{proposition}

\begin{proof}
选取在 $Y$ 上的紧化 $j : X \to \overline{X}$，并以
$\overline{f} : \overline{X} \to Y$ 表示结构态射。于是
\begin{align*}
\Hom_X(K, f^!L)
& =
\Hom_X(K, j^*\overline{f}{}^!L) \\
& =
\Hom_{\text{Pro-}D^+_{\textit{Coh}}(\mathcal{O}_{\overline{X}})}
(Rj_!K, \overline{f}{}^!L) \\
& =
\Hom_{\text{Pro-}D^+_{\textit{Coh}}(\mathcal{O}_Y)}(Rf_*Rj_!K, L) \\
& =
\Hom_{\text{Pro-}D^+_{\textit{Coh}}(\mathcal{O}_Y)}(Rf_!K, L)
\end{align*}
。第一个等式直接来自第 \ref{duality-section-upper-shriek} 节中 $f^!$ 的构造。
由引理 \ref{duality-lemma-shriek-coherent}，$\overline{f}{}^!L$ 属于
$D^+_{\textit{Coh}}(\mathcal{O}_{\overline{X}})$，故第二个等式来自引理
\ref{duality-lemma-lift-map-plus}。由于 $\overline{f}$ 固有，按构造，函子
$\overline{f}{}^!$ 是推前的右伴随；这给出第三个等式。第四个等式成立，
是因为 $Rf_! = Rf_* Rj_!$。
\end{proof}

\begin{lemma}
\label{duality-lemma-compactly-supported-triangle}
在情形 \ref{duality-situation-shriek} 中，设 $f : X \to Y$ 是
$\textit{FTS}_S$ 的态射，并设
$$
K \to L \to M \to K[1]
$$
是 $D^b_{\textit{Coh}}(\mathcal{O}_X)$ 中的可区别三角。则存在
$D^b_{\textit{Coh}}(\mathcal{O}_Y)$ 中的可区别三角的逆系统
$$
K_n \to L_n \to M_n \to K_n[1]
$$
，使 pro-系统 $(K_n)$、$(L_n)$、$(M_n)$ 分别给出
$Rf_!K$、$Rf_!L$、$Rf_!M$。
\end{lemma}

\begin{proof}
选取在 $Y$ 上的紧化 $j : X \to \overline{X}$，并以
$\overline{f} : \overline{X} \to Y$ 表示结构态射。按照引理
\ref{duality-lemma-extension-by-zero-triangle}，选取对应于开浸入 $j$
和给定可区别三角的逆系统
$$
\overline{K}_n \to \overline{L}_n \to \overline{M}_n \to \overline{K}_n[1]
$$
，它位于 $D^b_{\textit{Coh}}(\mathcal{O}_{\overline{X}})$ 中。
取 $K_n = R\overline{f}_*\overline{K}_n$，并类似地定义 $L_n$ 和 $M_n$。
由 $Rf_!$ 的定义，这满足要求。
\end{proof}

\begin{remark}
\label{duality-remark-compose-inverse-systems}
设 $\mathcal{C}$ 是一个范畴。假设给定 $\mathcal{C}$ 的 pro-对象范畴中的
逆系统的逆系统
$$
\ldots \xrightarrow{\alpha_4} (M_{3, n}) \xrightarrow{\alpha_3} (M_{2, n})
\xrightarrow{\alpha_2} (M_{1, n})
$$
。换言之，箭头 $\alpha_i$ 是 pro-对象的态射。由《范畴》，例
\ref{categories-example-pro-morphism-inverse-systems}，每个 $\alpha_i$
都可由偶对 $(m_i, a_i)$ 表示，其中
$m_i : \mathbf{N} \to \mathbf{N}$ 是递增函数，而
$a_{i, n} : M_{i, m_i(n)} \to M_{i - 1, n}$ 是 $\mathcal{C}$ 中的态射，
使图
$$
\xymatrix{
\ldots \ar[r] &
M_{i, m_i(3)} \ar[d]^{a_{i, 3}} \ar[r] &
M_{i, m_i(2)} \ar[d]^{a_{i, 2}} \ar[r] &
M_{i, m_i(1)} \ar[d]^{a_{i, 1}} \\
\ldots \ar[r] &
M_{i - 1, 3} \ar[r] &
M_{i - 1, 2} \ar[r] &
M_{i - 1, 1}
}
$$
交换。用 $\max(n, m_i(n))$ 替换 $m_i(n)$，并相应调整态射 $a_i(n)$
（如所引例中那样），可以假设 $m_i(n) \geq n$。此时考虑逆系统
$$
\ldots \to
M_{4, m_4(m_3(m_2(4)))} \to
M_{3, m_3(m_2(3))} \to
M_{2, m_2(2)} \to
M_{1, 1}
$$
，其一般项为
$$
M_k = M_{k, m_k(m_{k - 1}(\ldots (m_2(k))\ldots))}
$$
。对 $\mathcal{C}$ 的任意对象 $N$，有
$$
\colim_i \colim_n \Mor_\mathcal{C}(M_{i, n}, N) =
\colim_k \Mor_\mathcal{C}(M_k, N) 
$$
。略去细节。换言之，逆系统 $(M_k)$ 具有性质
$$
\colim_i \Mor_{\text{Pro-}\mathcal{C}}((M_{i, n}), N) =
\Mor_{\text{Pro-}\mathcal{C}}((M_k), N)
$$
。由《范畴》，注 \ref{categories-remark-pro-category-copresheaves} 中的讨论，
该性质在 pro-同构意义下确定逆系统 $(M_k)$。这样，便可把
$\text{Pro-}\mathcal{C}$ 中的某些逆系统化为具有可数指标范畴的 pro-对象。
\end{remark}

\begin{remark}
\label{duality-remark-composition-lower-shriek}
在情形 \ref{duality-situation-shriek} 中，设 $f : X \to Y$ 和 $g : Y \to Z$
是 $\textit{FTS}_S$ 的可复合态射。定义复合函子
$$
Rg_! \circ Rf_! :
D^b_{\textit{Coh}}(\mathcal{O}_X)
\longrightarrow
\text{Pro-}D^b_{\textit{Coh}}(\mathcal{O}_Z)
$$
。具体地，由 $Rf_!$ 的构造，对
$D^b_{\textit{Coh}}(\mathcal{O}_X)$ 中的 $K$，输出 $Rf_!K$ 是
$D^b_{\textit{Coh}}(\mathcal{O}_Y)$ 中某个逆系统 $(M_n)$ 的
pro-同构类。由于 $Rg_!$ 以相同方式构造，
$$
\ldots \to Rg_!M_3 \to Rg_!M_2 \to Rg_!M_1
$$
是 $\text{Pro-}D^b_{\textit{Coh}}(\mathcal{O}_Y)$ 的逆系统。
由注 \ref{duality-remark-compose-inverse-systems} 中的讨论，存在唯一的
$D^b_{\textit{Coh}}(\mathcal{O}_Z)$ 中逆系统的 pro-同构类，记为
$Rg_! Rf_! K$，使得
$$
\Hom_{\text{Pro-}D^b_{\textit{Coh}}(\mathcal{O}_Z)}(Rg_!Rf_!K, L) =
\colim_n \Hom_{\text{Pro-}D^b_{\textit{Coh}}(\mathcal{O}_Z)}(Rg_!M_n, L)
$$
。略去验证该构造对 $K$ 具有函子性所需的讨论，因为这会直接由下一个引理得到。
\end{remark}

\begin{lemma}
\label{duality-lemma-composition-lower-shriek}
在情形 \ref{duality-situation-shriek} 中，设 $f : X \to Y$ 和 $g : Y \to Z$
是 $\textit{FTS}_S$ 的可复合态射。沿用注
\ref{duality-remark-composition-lower-shriek} 的记号，有
$Rg_! \circ Rf_! = R(g \circ f)_!$。
\end{lemma}

\begin{proof}
由《范畴》，注 \ref{categories-remark-pro-category-copresheaves} 中的讨论，
只须证明对 $D^b_{\textit{Coh}}(\mathcal{O}_Z)$ 中的 $L$ 计算 $\Hom$
时得到相同答案。沿用注 \ref{duality-remark-composition-lower-shriek} 中的记号，计算如下：
\begin{align*}
\Hom_Z(Rg_!Rf_!K, L)
& =
\colim_n \Hom_Z(Rg_!M_n, L) \\
& =
\colim_n \Hom_Y(M_n, g^!L) \\
& =
\Hom_Y(Rf_!K, g^!L) \\
& =
\Hom_X(K, f^!g^!L) \\
& =
\Hom_X(K, (g \circ f)^!L) \\
& =
\Hom_Z(R(g \circ f)_!K, L)
\end{align*}
。第一个等式是 $Rg_!Rf_!K$ 的定义。第二个等式是命题
\ref{duality-proposition-duality-compactly-supported} 对 $g$ 的应用。
第三个等式使用 $Rf_!K$ 由 $(M_n)$ 给出这一事实。第四个等式是命题
\ref{duality-proposition-duality-compactly-supported} 对 $f$ 的应用。
第五个等式是引理 \ref{duality-lemma-upper-shriek-composition}。第六个等式是命题
\ref{duality-proposition-duality-compactly-supported} 对 $g \circ f$ 的应用。
\end{proof}

\begin{remark}
\label{duality-remark-covariance-open-lower-shriek}
在情形 \ref{duality-situation-shriek} 中，设 $f : X \to Y$ 是
$\textit{FTS}_S$ 的态射，且 $U \subset X$ 是开集。令
$g = f|_U : U \to Y$。则存在典范态射
$$
Rg_!(K|_U) \longrightarrow Rf_!K
$$
，它对 $D^b_{\textit{Coh}}(\mathcal{O}_X)$ 中的 $K$ 具有函子性，
并且至少可以用三种方式定义。
\begin{enumerate}
\item 以 $i : U \to X$ 表示包含态射。由引理
\ref{duality-lemma-composition-lower-shriek}，有 $Rg_! = Rf_! \circ Ri_!$。
把 $Rf_!$ 应用于映射 $Ri_!(K|_U) \to K$ 即可；这是注
\ref{duality-remark-covariance-open-j-lower-shriek} 的特殊情形。
\item 选取 $X$ 在 $Y$ 上的紧化 $j : X \to \overline{X}$，其结构态射为
$\overline{f} : \overline{X} \to Y$。令
$j' = j \circ i : U \to \overline{X}$。利用
$Rf_! = R\overline{f}_* \circ Rj_!$ 和
$Rg_! = R\overline{f}_* \circ Rj'_!$，并把 $R\overline{f}_*$ 应用于注
\ref{duality-remark-covariance-open-j-lower-shriek} 中的映射
$Rj'_!(K|_U) \to Rj_!K$。
\item 使用
\begin{align*}
\Hom_{\text{Pro-}D^b_{\textit{Coh}}(\mathcal{O}_Y)}(Rf_!K, L)
& =
\Hom_X(K, f^!L) \\
& \to
\Hom_U(K|_U, f^!L|_U) \\
& =
\Hom_U(K|_U, g^!L) \\
& =
\Hom_{\text{Pro-}D^b_{\textit{Coh}}(\mathcal{O}_Y)}(Rg_!(K|_U), L)
\end{align*}
。它对 $L$ 和 $K$ 具有函子性。这里两次使用了命题
\ref{duality-proposition-duality-compactly-supported}，还使用了上叹号函子的构造所给出的
$g^! = i^* \circ f^!$。由对 $L$ 的函子性及《范畴》，注
\ref{categories-remark-pro-category-copresheaves}，得到
$\text{Pro-}D^b_{\textit{Coh}}(\mathcal{O}_Y)$ 中的典范映射
$Rg_!(K|_U) \to Rf_!K$；上述箭头对 $K$ 的函子性说明该映射也对 $K$
具有函子性。
\end{enumerate}
这三种构造给出同一个箭头；略去细节。
\end{remark}

\begin{remark}
\label{duality-remark-covariance-etale-lower-shriek}
把注 \ref{duality-remark-covariance-open-lower-shriek} 中紧支撑上同调的协变性推广到
\'etale 态射。具体地，在情形 \ref{duality-situation-shriek} 中，假设给定
$\textit{FTS}_S$ 中的交换图
$$
\xymatrix{
U \ar[rr]_h \ar[rd]_g & & X \ar[ld]^f \\
& Y
}
$$
，其中 $h$ 为 \'etale。则存在典范态射
$$
Rg_!(h^*K) \longrightarrow Rf_!K
$$
，它对 $D^b_{\textit{Coh}}(\mathcal{O}_X)$ 中的 $K$ 具有函子性。
用下列映射链定义该变换：
\begin{align*}
\Hom_{\text{Pro-}D^b_{\textit{Coh}}(\mathcal{O}_Y)}(Rf_!K, L)
& =
\Hom_X(K, f^!L) \\
& \to
\Hom_U(h^*K, h^*(f^!L)) \\
& =
\Hom_U(h^*K, h^!f^!L) \\
& =
\Hom_U(h^*K, g^!L) \\
& =
\Hom_{\text{Pro-}D^b_{\textit{Coh}}(\mathcal{O}_Y)}(Rg_!(h^*K), L)
\end{align*}
。它对 $L$ 和 $K$ 具有函子性。这里两次使用了命题
\ref{duality-proposition-duality-compactly-supported}，使用了引理
\ref{duality-lemma-shriek-etale} 的等式 $h^* = h^!$，还使用了引理
\ref{duality-lemma-upper-shriek-composition} 的等式 $h^! \circ f^! = g^!$。
由对 $L$ 的函子性及《范畴》，注
\ref{categories-remark-pro-category-copresheaves}，得到
$\text{Pro-}D^b_{\textit{Coh}}(\mathcal{O}_Y)$ 中的典范映射
$Rg_!(h^*K) \to Rf_!K$；上述箭头对 $K$ 的函子性说明该映射也对 $K$
具有函子性。
\end{remark}

\begin{remark}
\label{duality-remark-covariance-lower-shriek}
在注 \ref{duality-remark-covariance-open-lower-shriek} 和
\ref{duality-remark-covariance-etale-lower-shriek} 中已经看到，紧支撑上同调的构造
对开浸入和 \'etale 态射具有协变性。事实上，正确的一般性如下：给定
$\textit{FTS}_S$ 中的交换图
$$
\xymatrix{
U \ar[rr]_h \ar[rd]_g & & X \ar[ld]^f \\
& Y
}
$$
，其中 $h$ 平坦且拟有限，则存在典范变换
$$
Rg_! \circ h^* \longrightarrow Rf_!
$$
。像注 \ref{duality-remark-covariance-etale-lower-shriek} 中一样，可以使用
$D^+_{\textit{Coh}}(\mathcal{O}_X)$ 上的函子变换 $h^* \to h^!$
来构造该映射。回忆
$h^!K = h^*K \otimes \omega_{U/X}$，其中
$\omega_{U/X} = h^!\mathcal{O}_X$ 是平坦拟有限态射 $h$ 的相对对偶层
（见引理 \ref{duality-lemma-perfect-comparison-shriek} 和
\ref{duality-lemma-flat-quasi-finite-shriek}）。还记得，$\omega_{U/X}$
等于《判别式》，注
\ref{discriminant-remark-relative-dualizing-for-quasi-finite}
中由《判别式》，引理 \ref{discriminant-lemma-compare-dualizing}
构造的相对对偶模。因此，可以使用将在《判别式》，注
\ref{discriminant-remark-relative-dualizing-for-flat-quasi-finite}
中构造的迹元素 $\tau_{U/X} : \mathcal{O}_U \to \omega_{U/X}$
来定义该变换。如果以后需要，我们将在这里精确陈述并证明该结果。
\end{remark}







\section{紧支撑上同调的对偶性}
\label{duality-section-duality-compactly-supported}

\noindent
设 $k$ 是域，$U$ 是 $k$ 上有限型的分离概形，$K$ 是
$D^b_{\textit{Coh}}(\mathcal{O}_U)$ 中的对象。如下定义 $K$
的紧支撑上同调 $H^i_c(U, K)$。选取到 $k$ 上某个固有概形中的开浸入
$j : U \to X$，以及 $j : U \to X$ 的一个 Deligne 系统 $(K_n)$，
使其限制到 $U$ 后常值为 $K$。令
$$
H^i_c(U, K) = \lim H^i(X, K_n)
$$
。赋予它逆极限拓扑，并把它视为拓扑 $k$-向量空间（见《代数进阶》，第
\ref{more-algebra-section-topological-ring} 节）。这里有几点需要说明。

\medskip\noindent
第一，该定义与 $X$ 和 $(K_n)$ 的选择无关。事实上，若
$p : U \to \Spec(k)$ 表示结构态射，则已经知道
$Rp_!K = (R\Gamma(X, K_n))$ 在 $D(k)$ 中于 pro-同构意义下良定义，
所以定义 $H^i_c(U, K)$ 的逆极限也良定义。

\medskip\noindent
第二，使用表达式
$$
H^i(R\lim R\Gamma(X, K_n)) = R\Gamma(X, R\lim K_n)
$$
看起来可能更自然，但所得答案相同：由于 $k$-向量空间
$H^j(X, K_n)$ 都有限维，这些逆系统满足 Mittag--Leffler 条件，因而《上同调》，
引理 \ref{cohomology-lemma-RGamma-commutes-with-Rlim} 中的
$R^1\lim$ 项消失。

\medskip\noindent
若 $U' \subset U$ 是开子概形，则存在典范映射
$$
H^i_c(U', K|_{U'}) \longrightarrow H^i_c(U, K)
$$
，它对 $D^b_{\textit{Coh}}(\mathcal{O}_U)$ 中的 $K$ 具有函子性。
例如见注 \ref{duality-remark-covariance-open-lower-shriek}。事实上，由注
\ref{duality-remark-covariance-etale-lower-shriek}，更一般地，对 $k$
上有限型分离概形之间的 \'etale 态射 $U' \to U$，也存在这样的映射。

\medskip\noindent
若 $V$ 是 $k$-向量空间，如下在 $\Hom_k(V, k)$ 上赋予拓扑：把
$V = \bigcup V_i$ 写成其有限维 $k$-子向量空间的滤过并，并在
$\Hom_k(V, k) = \lim \Hom_k(V_i, k)$ 上使用逆极限拓扑。若
$\dim_k V < \infty$，则 $\Hom_k(V, k)$ 上的拓扑是离散的。更一般地，
若 $V = \colim_n V_n$ 写成有限维向量空间的有向余极限，则作为拓扑向量空间，
$\Hom_k(V, k) = \lim \Hom_k(V_n, k)$。

\begin{lemma}
\label{duality-lemma-duality-compact-support}
设 $p : U \to \Spec(k)$ 是有限型分离态射，其中 $k$ 是域。令
$\omega_{U/k}^\bullet = p^!\mathcal{O}_{\Spec(k)}$。则有拓扑
$k$-向量空间的典范同构
$$
\Hom_k(H^i(U, K), k) =
H^{-i}_c(U, R\SheafHom_{\mathcal{O}_U}(K, \omega_{U/k}^\bullet))
$$
，它对 $D^b_{\textit{Coh}}(\mathcal{O}_U)$ 中的 $K$ 具有函子性。
\end{lemma}

\begin{proof}
选取 $k$ 上的紧化 $j : U \to X$。设拟相干理想层
$\mathcal{I} \subset \mathcal{O}_X$ 满足 $V(\mathcal{I}) = X \setminus U$。
由《概形的导出范畴》，命题 \ref{perfect-proposition-DCoh}，可选取
$M \in D^b_{\textit{Coh}}(\mathcal{O}_X)$，使 $K = M|_U$。由引理
\ref{duality-lemma-lift-map}，有
$$
H^i(U, K) =
\Ext^i_U(\mathcal{O}_U, M|_U) =
\colim \Ext^i_X(\mathcal{I}^n, M) =
\colim H^i(X, R\SheafHom_{\mathcal{O}_X}(\mathcal{I}^n, M))
$$
。由于 $\mathcal{I}^n$ 是相干 $\mathcal{O}_X$-模，故 $\mathcal{I}^n$
属于 $D^-_{\textit{Coh}}(\mathcal{O}_X)$；因而由《概形的导出范畴》，引理
\ref{perfect-lemma-coherent-internal-hom}，
$R\SheafHom_{\mathcal{O}_X}(\mathcal{I}^n, M)$ 属于
$D^+_{\textit{Coh}}(\mathcal{O}_X)$。

\medskip\noindent
令 $\omega_{X/k}^\bullet = q^!\mathcal{O}_{\Spec(k)}$，其中
$q : X \to \Spec(k)$ 是结构态射；见第
\ref{duality-section-duality-proper-over-field} 节。于是
\begin{align*}
\Hom_k(
&
H^i(X, R\SheafHom_{\mathcal{O}_X}(\mathcal{I}^n, M)), k) \\
& =
\Ext^{-i}_X(R\SheafHom_{\mathcal{O}_X}(\mathcal{I}^n, M),
\omega_{X/k}^\bullet) \\
& =
H^{-i}(X, R\SheafHom_{\mathcal{O}_X}(R\SheafHom_{\mathcal{O}_X}(
\mathcal{I}^n, M), \omega_{X/k}^\bullet))
\end{align*}
；见引理 \ref{duality-lemma-duality-proper-over-field}。由引理
\ref{duality-lemma-internal-hom-evaluate-isom} 的 (1)，典范映射
$$
R\SheafHom_{\mathcal{O}_X}(M, \omega_{X/k}^\bullet)
\otimes_{\mathcal{O}_X}^\mathbf{L} \mathcal{I}^n
\longrightarrow
R\SheafHom_{\mathcal{O}_X}(R\SheafHom_{\mathcal{O}_X}(
\mathcal{I}^n, M), \omega_{X/k}^\bullet)
$$
是同构。注意，$\omega^\bullet_{U/k} = \omega^\bullet_{X/k}|_U$，
因为 $p^!$ 构造为 $q^!$ 后接限制到 $U$。所以
$R\SheafHom_{\mathcal{O}_X}(M, \omega_{X/k}^\bullet)$ 是
$D^b_{\textit{Coh}}(\mathcal{O}_X)$ 中的对象，且限制到 $U$ 后为
$R\SheafHom_{\mathcal{O}_U}(K, \omega_{U/k}^\bullet)$。由引理
\ref{duality-lemma-tensoring-Deligne-system}，可知
$$
\lim H^{-i}(X, R\SheafHom_{\mathcal{O}_X}(M, \omega_{X/k}^\bullet)
\otimes_{\mathcal{O}_X}^\mathbf{L} \mathcal{I}^n)
$$
是该引理等式右端的一种表示。合并以上所得全部同构，即得到该引理的同构。
\end{proof}

\begin{lemma}
\label{duality-lemma-duality-compact-support-restrict-open}
沿用引理 \ref{duality-lemma-duality-compact-support} 的记号，假设
$U' \subset U$ 是开子概形。则图
$$
\xymatrix{
\Hom_k(H^i(U, K), k) \ar[rr] & &
H^{-i}_c(U, R\SheafHom_{\mathcal{O}_U}(K, \omega_{U/k}^\bullet)) \\
\Hom_k(H^i(U', K|_{U'}), k) \ar[rr] \ar[u] & &
H^{-i}_c(U', R\SheafHom_{\mathcal{O}_{U'}}(K, \omega_{U'/k}^\bullet)) \ar[u]
}
$$
交换。这里，水平箭头是引理 \ref{duality-lemma-duality-compact-support} 的同构；
左侧竖直箭头是限制映射 $H^i(U, K) \to H^i(U', K|_{U'})$
的对偶映射；右侧竖直箭头来自注 \ref{duality-remark-covariance-open-lower-shriek}
（见该引理前面的讨论）。
\end{lemma}

\begin{proof}
强烈建议读者跳过这个证明。像引理 \ref{duality-lemma-duality-compact-support}
的证明那样选取 $X$ 和 $M$。以下从 $R\SheafHom$ 和
$\otimes^\mathbf{L}$ 中省略下标 $\mathcal{O}_X$。写成
$$
H^i(U, K) = \colim H^i(X, R\SheafHom(\mathcal{I}^n, M))
$$
以及
$$
H^i(U', K|_{U'}) = \colim H^i(X, R\SheafHom((\mathcal{I}')^n, M))
$$
，如引理 \ref{duality-lemma-duality-compact-support} 的证明中那样；这里按照注
\ref{duality-remark-covariance-open-j-lower-shriek} 中的讨论选取
$\mathcal{I}' \subset \mathcal{I}$，使映射
$H^i(U, K) \to H^i(U', K|_{U'})$ 由映射
$(\mathcal{I}')^n \to \mathcal{I}^n$ 诱导。类似地写成
$$
H^i_c(U, R\SheafHom(K, \omega_{U/k}^\bullet)) =
\lim H^i(X,  R\SheafHom(M, \omega_{X/k}^\bullet)
\otimes^\mathbf{L} \mathcal{I}^n)
$$
以及
$$
H^i_c(U', R\SheafHom(K|_{U'}, \omega_{U'/k}^\bullet)) =
\lim H^i(X,  R\SheafHom(M, \omega_{X/k}^\bullet)
\otimes^\mathbf{L} (\mathcal{I}')^n)
$$
，从而箭头 $H^i_c(U', R\SheafHom(K|_{U'}, \omega_{U'/k}^\bullet))
\to H^i_c(U, R\SheafHom(K, \omega_{U/k}^\bullet))$ 也由映射
$(\mathcal{I}')^n \to \mathcal{I}^n$ 导出。图
$$
\xymatrix{
R\SheafHom(M, \omega_{X/k}^\bullet)
\otimes^\mathbf{L} \mathcal{I}^n
\ar[rr] & &
R\SheafHom(R\SheafHom(\mathcal{I}^n, M), \omega_{X/k}^\bullet) \\
R\SheafHom(M, \omega_{X/k}^\bullet) \otimes^\mathbf{L} (\mathcal{I}')^n
\ar[rr] \ar[u] & &
R\SheafHom(R\SheafHom((\mathcal{I}')^n, M), \omega_{X/k}^\bullet) \ar[u]
}
$$
交换，因为《上同调》，引理
\ref{cohomology-lemma-internal-hom-evaluate} 中水平箭头的构造对三个变元
都具有函子性。因此最终归结为断言图
$$
\xymatrix{
\Hom_k(H^i(X, R\SheafHom(\mathcal{I}^n, M)), k) \ar[r] &
H^{-i}(X, R\SheafHom(R\SheafHom(
\mathcal{I}^n, M), \omega_{X/k}^\bullet)) \\
\Hom_k(H^i(X, R\SheafHom((\mathcal{I}')^n, M)), k) \ar[r] \ar[u] &
H^{-i}(X, R\SheafHom(R\SheafHom(
(\mathcal{I}')^n, M), \omega_{X/k}^\bullet)) \ar[u]
}
$$
交换。这成立是因为对偶同构
$$
\Hom_k(H^i(X, L), k) = \Ext^{-i}_X(L, \omega_{X/k}^\bullet) =
H^{-i}(X, R\SheafHom(L, \omega_{X/k}^\bullet))
$$
对 $D_\QCoh(\mathcal{O}_X)$ 中的 $L$ 具有函子性。
\end{proof}

\begin{lemma}
\label{duality-lemma-h0-compactly-supported}
设 $X$ 是域 $k$ 上的完备概形。设
$K \in D^b_{\textit{Coh}}(\mathcal{O}_X)$ 满足对 $i < 0$ 有
$H^i(K) = 0$。令 $\mathcal{F} = H^0(K)$。设 $Z \subset X$
是闭子集，其补集为 $U = X \setminus U$。则
$$
H^0_c(U, K|_U) \subset H^0(X, \mathcal{F})
$$
恰由 $\mathcal{F}$ 中在 $Z$ 的某个开邻域上消失的整体截面组成。
\end{lemma}

\begin{proof}
考虑注 \ref{duality-remark-covariance-open-lower-shriek} 中的映射
$H^0_c(U, K|_U) \to H^0_X(X, K) = H^0(X, K) = H^0(X, \mathcal{F})$。
为研究此映射，用一个有界复形 $\mathcal{F}^\bullet$ 表示 $K$，其中
对 $i < 0$ 有 $\mathcal{F}^i = 0$。根据定义，
$$
H^0_c(U, K|_U) = \lim H^0(X, \mathcal{I}^n\mathcal{F}^\bullet)
= \lim \Ker(
H^0(X, \mathcal{I}^n\mathcal{F}^0) \to
H^0(X, \mathcal{I}^n\mathcal{F}^1))
$$
由 Artin--Rees 引理（《概形上同调》，引理
\ref{coherent-lemma-Artin-Rees}），这等于
$\lim H^0(X, \mathcal{I}^n\mathcal{F})$。因此箭头
$H^0_c(U, K|_U) \to H^0(X, \mathcal{F})$ 是单射，其像由
$\mathcal{F}$ 中对每个 $n$ 都属于子层 $\mathcal{I}^n\mathcal{F}$
的整体截面组成。考察 $\mathcal{F}$ 的茎并应用 Krull 交定理
（《代数》，引理 \ref{algebra-lemma-intersect-powers-ideal-module-zero}），
即可把这些截面刻画为在 $Z$ 的某个邻域上消失的截面。函数情形见
《代数》，注 \ref{algebra-remark-intersection-powers-ideal} 中的讨论。
\end{proof}




\section{Lichtenbaum 定理}
\label{duality-section-lichtenbaum}

\noindent
以下定理由 Lichtenbaum 猜想，并由 Grothendieck 证明
（见 \cite{Hartshorne-local-cohomology}）。Kleiman 在
\cite{Kleiman-Lichtenbaum} 中给出了一个非常漂亮的证明。
该定理对带支集上同调情形的推广见
\cite{Lyubeznik-Lichtenbaum}。论证中最有趣的部分包含在下述引理的证明中。

\begin{lemma}
\label{duality-lemma-lichtenbaum}
设 $U$ 是代数簇，$\mathcal{F}$ 是凝聚 $\mathcal{O}_U$-模。
若 $H^d(U, \mathcal{F})$ 非零，则 $\dim(U) \geq d$；若等号成立，
则 $U$ 是完备的。
\end{lemma}

\begin{proof}
由《上同调》，命题 \ref{cohomology-proposition-vanishing-Noetherian}
中的 Grothendieck 消失结果，得到 $\dim(U) \geq d$。
假设 $\dim(U) = d$。选取一个使 $U$ 在 $X$ 中稠密的紧化
$U \to X$。（由《平坦性的进一步结果》，定理
\ref{flat-theorem-nagata} 和引理 \ref{flat-lemma-compactifyable}，
这样的紧化存在。）以其约化替换 $X$ 后，可知 $X$ 是维数为
$d$ 的完备代数簇，并且 $U$ 完备当且仅当 $U = X$。令
$Z = X \setminus U$。下面证明：若 $Z$ 非空，则
$H^d(U, \mathcal{F})$ 为零。

\medskip\noindent
选取一个凝聚 $\mathcal{O}_X$-模 $\mathcal{G}$，使其在 $U$ 上的限制为
$\mathcal{F}$；见《性质》，引理
\ref{properties-lemma-lift-finite-presentation}。令
$\omega_X^\bullet$ 表示第 \ref{duality-section-duality-proper-over-field} 节所述
$X$ 的对偶化复形，并令 $\omega_U^\bullet = \omega_X^\bullet|_U$。
由引理 \ref{duality-lemma-duality-compact-support}，$H^d(U, \mathcal{F})$
与
$$
H^{-d}_c(U, R\SheafHom_{\mathcal{O}_U}(\mathcal{F}, \omega_U^\bullet))
$$
互为对偶。由引理 \ref{duality-lemma-duality-proper-over-field}，
$\omega_X^\bullet$ 的上同调层在次数 $< -d$ 时消失，并且
$H^{-d}(\omega_X^\bullet) = \omega_X$ 是一个凝聚
$\mathcal{O}_X$-模，它满足 $(S_2)$，且支集为 $X$。特别地，
$\omega_X$ 是无挠的；见《除子》，引理
\ref{divisors-lemma-torsion-free-finite-noetherian-domain}。因此上同调层
$$
H^{-d}(R\SheafHom_{\mathcal{O}_X}(\mathcal{G}, \omega_X^\bullet)) =
\SheafHom(\mathcal{G}, \omega_X)
$$
是无挠的；见《除子》，引理
\ref{divisors-lemma-hom-into-torsion-free}。所以该层不存在在 $X$
的任何非空开集上消失的非零截面（这样的截面将是挠截面）。于是由引理
\ref{duality-lemma-h0-compactly-supported}，
$H^{-d}_c(U, R\SheafHom_{\mathcal{O}_U}(\mathcal{F}, \omega_U^\bullet))$
为零，因而如所欲，$H^d(U, \mathcal{F})$ 为零。
\end{proof}

\begin{theorem}
\label{duality-theorem-lichtenbaum}
设 $X$ 是域 $k$ 上的非空有限型分离概形。令 $d = \dim(X)$。
则下列条件等价：
\begin{enumerate}
\item 对 $X$ 上每个凝聚 $\mathcal{O}_X$-模 $\mathcal{F}$，都有
$H^d(X, \mathcal{F}) = 0$；
\item 对 $X$ 上每个拟凝聚 $\mathcal{O}_X$-模 $\mathcal{F}$，都有
$H^d(X, \mathcal{F}) = 0$；
\item 不存在维数为 $d$ 且在 $k$ 上完备的不可约分支
$X' \subset X$。
\end{enumerate}
\end{theorem}

\begin{proof}
假设存在不可约分支 $X' \subset X$，它是完备的且维数为 $d$；
这里把它视为整闭子概形。令 $\omega_{X'}$ 为 $X'$ 在 $k$ 上的一个
对偶化模；见引理 \ref{duality-lemma-duality-proper-over-field}。由该引理，
$H^d(X', \omega_{X'})$ 与 $H^0(X', \mathcal{O}_{X'})$ 互为对偶，
故前者非零。因此 $H^d(X, \omega_{X'}) = H^d(X', \omega_{X'})$
非零，从而条件 (1) 不成立。由此可见 (1) 蕴含 (3)。

\medskip\noindent
下面证明 (3) 蕴含 (1)。设 $\mathcal{F}$ 是凝聚
$\mathcal{O}_X$-模，且 $H^d(X, \mathcal{F})$ 非零。选取《概形上同调》，
引理 \ref{coherent-lemma-coherent-filter} 中的滤过
$$
0 = \mathcal{F}_0 \subset \mathcal{F}_1 \subset
\ldots \subset \mathcal{F}_m = \mathcal{F}
$$
得到正合序列
$$
H^d(X, \mathcal{F}_i) \to H^d(X, \mathcal{F}_{i + 1}) \to
H^d(X, \mathcal{F}_{i + 1}/\mathcal{F}_i)
$$
于是对某个 $i \in \{1, \ldots, m\}$，
$H^d(X, \mathcal{F}_{i + 1}/\mathcal{F}_i)$ 非零。由滤过的选取方式，
这意味着存在一个整闭子概形 $Z \subset X$ 以及一个非零凝聚理想层
$\mathcal{I} \subset \mathcal{O}_Z$，使 $H^d(Z, \mathcal{I})$ 非零。
由引理 \ref{duality-lemma-lichtenbaum}，得到 $\dim(Z) = d$ 且 $Z$ 在 $k$
上完备，这与 (3) 矛盾。因此 (3) 蕴含 (1)。

\medskip\noindent
最后证明：对任意诺特概形 $X$，(1) 与 (2) 等价。显然 (2) 蕴含
(1)。反过来，假设 (1) 成立，并设 $\mathcal{F}$ 是拟凝聚
$\mathcal{O}_X$-模。可将 $\mathcal{F} = \colim \mathcal{F}_i$
写成其凝聚子模的滤过余极限；见《性质》，引理
\ref{properties-lemma-quasi-coherent-colimit-finite-type}。于是由《上同调》，
引理 \ref{cohomology-lemma-quasi-separated-cohomology-colimit}，有
$H^d(X, \mathcal{F}) = \colim H^d(X, \mathcal{F}_i) = 0$。
所以 (2) 成立。
\end{proof}
